\providecommand{\toplevelprefix}{../..}  %
\documentclass[../../book-main_zh.tex]{subfiles}

\begin{document}

\chapter{真实世界数据与任务的表征学习}
\label{ch:applications}

\begin{quote}
    \hfill    ``{\em 最好的理论源于实践，而
    最好的实践源于理论}。''

    $~$ \hfill --- 唐纳德·克努斯 (Donald Knuth)
\end{quote}
\vspace{5mm}

\section{引言}
前面的章节系统地介绍了
数学问题、计算框架以及实用
算法，这些内容均与从高维数据中学习低维分布
息息相关。尽管这些方法的大多数理论证明
都是针对理想化的数据分布模型
（如子空间（混合）和/或高斯分布）建立的，但
这些计算方法背后的原理和思想
依然强大且通用，实际上它们旨在
应用于真实世界的数据集和任务。

为了帮助读者更好地理解本书的内容，并学习
如何将迄今为止学到的知识应用于真实世界数据，在
本书的最后，我们提供了一些代表性应用的
演示和案例。每个应用都提出了一个
针对真实世界任务和真实世界数据集（如
视觉数据、运动数据和文本数据）的解决方案，其中使用了我们
在本书中介绍的方法。本章展示的结果
旨在实现以下两个目的：
\begin{itemize}
    \item 首先，提供额外的实验细节和
        经验证据，以验证本书前面介绍的
        方法，并展示它们在真实世界环境中的
        巨大潜力；
    \item 其次，向读者介绍深度学习中某些现代的
        经验方法和任务，这些内容在研究或生产代码库之外
        通常缺乏详细的文档说明。
\end{itemize}
然而，坦率地说，这里展示的解决方案和结果
仅仅是为了验证该方法论的有效性。因此，
无论是在理论理解还是实际工程方面，未来都有很大的改进空间，
以进一步提升
最先进的性能。我们将通过本章，特别是最后一节 \ref{sec:scalable}，展示一些直接且
明显的改进。关于表征学习中可能更具变革性的想法，以及关于一般智能的遗留重大开放问题的讨论，我们将留待
最后的 \Cref{ch:future} 中作为未来研究的内容。

\subsection{本章概要}
\label{sec:application_outline}
对于自然界的智能而言，可以说我们大脑用来构建记忆和
知识的三种最重要的数据类型是
{\em 视觉数据}、{\em 身体运动}和{\em 自然语言}。因此，在本章中，我们将展示如何应用本书介绍的原理和方法来学习这些真实世界数据的良好表示，并展示这些表示如何促进与这些数据相关的重要实际任务。

首先，从 \ref{sec:contrastive_learning} 节到
\ref{sec:CUPID} 节，我们将从视觉数据开始，展示如何
学习从 2D 图像到 3D 对象的分布，因为视觉
记忆作为我们（以及动物）的低级但基础的模型，
存储了关于物理世界的知识，并有助于
在新环境中进行感知和预测。然后，在 \ref{sec:egoallo} 节中，
我们将展示如何学习身体运动的分布，
因为这对于我们（以及动物）在环境中迅速且准确地行动或交互至关重要。
最后，在 \ref{sec:clm_text} 节中，
我们将展示如何建模和学习
自然语言的分布。

显然，与自然语言相比，我们对
物理数据（视觉和运动）的覆盖将更为广泛。
尽管在过去几年中，人工智能行业在自然语言建模方面投入了大量资金，\footnote{包括其衍生形式，如编程语言、逻辑推演和数学证明。}人们已经开始意识到对
物理世界及我们在其中的行为的知识进行建模的重要性。近年来，人们对开发所谓的“世界模型”、“空间智能”、“物理智能”或“具身智能”的兴趣迅速增长，注意力也随之转移，因为它们对于需要与开放的物理世界交互的AI系统或智能体（例如智能机器人）至关重要。正如我们在 \ref{ch:intro} 章中所了解到的，这正是诺伯特·维纳（Norbert Wiener）在20世纪40年代发起的{\em 控制论}（Cybernetics）项目的目标和范围。  

在技术层面上，从 \Cref{sec:contrastive_learning} 到
\Cref{sec:image_classification}，我们将展示如何使用
本书介绍的原理来学习图像数据的良好表示，
分别在{\em 无监督、弱监督}和{\em 有监督}
的设置下进行。\footnote{这三种设置的目标都是学习数据分布的良好表示。它们的区别仅在于提供了多少有助于学习的关于分布的辅助信息。} 这既是对
图像数据处理、数据增强技术以及流行的深度架构 Transformer 的介绍，也为后续实现 CRATE 架构做铺垫。这些例子还
首次展示了利用本书介绍的原理可以实现的巨大简化。我们将继续
在 \Cref{sec:image_classification} 和 \Cref{sec:clm_text} 中分别针对图像和文本对 \textit{Transformer 网络架构}进行修改，
这展示了简化的白盒架构（包括 CRATE 及其变体）在图像和文本领域内进行 \textit{编码} 的能力。
我们还将在 \Cref{sec:image_completion} 中展示简化架构在 \textit{自编码} 方面的能力。 %

实际上，我们想要学习数据分布的良好
表示的主要原因是，它允许我们对该
分布进行采样（作为先验），以便在给定当前观察或信息为条件的情况下，重新生成、估计或预测
世界的状态。我们展示了学习到的表示的生成能力：
\begin{itemize}
\item 对于 2D 图像，在 \Cref{sec:auto-encode-img-gen} 和
\Cref{sec:cond_gen} 中；
\item 对于 3D 对象，在
\Cref{sec:Michelangelo} 和 \Cref{sec:shape-and-pose} 中； 
\item 对于人体运动，在 \Cref{sec:egoallo} 中；
\item 以及对于自然语言，在 \Cref{sec:clm_text} 中。
\end{itemize} 

\subsection{总体技术设置}\label{sec:experiment_setup}

在前面的章节中，我们提到了使用
表征学习技术大规模处理真实数据的不同设置。在
本章中，我们将非常详细地描述这些设置。
本节的目的是让读者能够
仅使用我们在本书中提供的描述、
前面章节介绍并在本章中扩展的原理，
以及从本章讨论其结果的少数论文中提取的超参数，
来复现本节（或整本书）中讨论的任何实验。
为此，我们将
使用可以直接在代码中实现的详细语言、
伪代码或数学符号，\textit{精确地}描述所有过程。
只要有可能，我们就会讨论具体的
实现如何与本书前面介绍的原理相联系。

\begin{figure}
    \includegraphics[width=\textwidth]{\toplevelprefix/chapters/chapter8/figs/encoder_pipeline.pdf}
    \caption{\small\textbf{编码器流水线示意图。} 数据
        \(\vX \in \cD\) 通过嵌入
        \(f_{\theta}^{\emb}\) 输入，得到 \((\R^{d})^{*}\) 中的一个序列。该
        嵌入通过主干网络 \(f_{\theta}^{\mathrm{bb}}\)
        以获取每个 token 的特征 \(\vZ_{\theta}(\vX)\)。我们可以
        使用提取映射 \(f_{\theta}^{\mathrm{ext}}\) 提取聚合特征
        \(\vz_{\theta}(\vX)\)。最后，为了在下游任务中使用
        该聚合特征，我们可以使用
    特定任务头 \(h_{\theta}\)。}
    \label{fig:overall_encoder_pipeline}
\end{figure}

\begin{figure}
    \includegraphics[width=\textwidth]{\toplevelprefix/chapters/chapter8/figs/autoencoder_pipeline.pdf}
    \caption{\small\textbf{自编码器流水线示意图。}
        数据 \(\vX \in \cD\) 通过嵌入
        \(f_{\theta}^{\emb}\) 输入，得到 \((\R^{d})^{*}\) 中的一个序列。该
        嵌入通过编码器主干网络
        \(f_{\theta}^{\mathrm{bb}}\) 以获取每个 token 的特征
        \(\vZ_{\theta}(\vX)\)。为了解码
        \(\vZ_{\theta}(\vX)\)，我们将其通过解码器主干网络
        \(g_{\eta}^{\mathrm{bb}}\)。为了将解码器主干网络的输出映射
        回数据空间 \(\cD\)，我们使用一个 \textit{逆嵌入}（unembedding）层
        \(g_{\eta}^{\mathrm{unemb}}\)，最终获得重建结果
        \(\hat{\vX}_{\theta, \eta}(\vX)\)（此处风格化为
    输入的像素化重建）。}
    \label{fig:overall_autoencoder_pipeline}
\end{figure}

%\paragraph{Technical Setup.}
让我们将可能的数据集定义为 \(\cD\)（例如，最终这
    将是图像集 \(\cI\)，或者是文本集
\(\cT\)），并将 \(\R^{d}\) 中有限 token 序列的集合
（即具有 \(d\) 行的矩阵集合）定义为 \((\R^{d})^{*} \doteq
\bigcup_{T = 1}^{\infty}\R^{d \times T}\)。为了讨论
本章介绍的丰富应用，我们首先回顾
在本书的其他部分中，我们讨论了两种不同类型的模型
架构。
\begin{itemize}
    \item 一种 \textit{编码器}（encoder）架构，由
        \(\theta\) 参数化，它由几个组件组成：
        \begin{itemize}
            \item 一个 \textit{嵌入}（embedding）\(f_{\theta}^{\emb} \colon
                \cD \to (\R^{d})^{*}\)，它将输入数据
                \(\cD\) 转换为一系列 \textit{token}，这些 token 被
                映射到或 \textit{嵌入} 到
                \(d\) 维空间中。\textit{在本章的其余部分，
                    我们通常会将 token 和嵌入
                等同起来。}
            \item 一个 \textit{编码器主干网络}（encoder backbone）
                \(f_{\theta}^{\backbone} \colon (\R^{d})^{*} \to
                (\R^{d})^{*}\)，它使用序列到序列的操作
                处理这一系列嵌入。
                该主干网络由前几章讨论的
                网络架构实现，但
                随着内容的深入，我们将给出更正式的描述。
            \item 一个 \textit{聚合特征提取器}（aggregate feature extractor）
                \(f_{\theta}^{\ext} \colon (\R^{d})^{*} \to \R^{d}\)，
                它提取整个序列的聚合表示。
                这用于为整个数据样本定义一个
                单一特征。
            \item 一个 \textit{特定任务头}（task-specific head）\(h_{\theta} \colon
                \R^{d} \to \R^{m}\)，它提取一个
                \(m\) 维输出用于预测。
        \end{itemize}
        我们还定义 \(f_{\theta} \doteq f_{\theta}^{\backbone}
        \circ f_{\theta}^{\emb} \colon \cD \to (\R^{d})^{*}\)。给定
        输入 \(\vX\)，我们记 \(\vZ_{\theta}(\vX) \doteq
        f_{\theta}(\vX)\) 且 \(\vz_{\theta}(\vX) \doteq
        f_{\theta}^{\ext}(\vZ_{\theta}(\vX))\)。整体流水线
        如 \Cref{fig:overall_encoder_pipeline} 所示。
    \item 一种 \textit{自编码器}（autoencoder）架构，它由
        几个组件组成：
        \begin{itemize}
            \item 一个 \textit{嵌入} \(f_{\theta}^{\emb} \colon
                \cD \to (\R^{d})^{*}\)，它将输入数据
                \(\cD\) 转换为一系列 \textit{token}，这些 token 被
                嵌入到 \(d\) 维空间中。
            \item 一个 \textit{编码器主干网络}
                \(f_{\theta}^{\backbone} \colon (\R^{d})^{*} \to
                (\R^{d})^{*}\)，它使用序列到序列的操作
                处理这一系列嵌入。
            \item 一个 \textit{解码器主干网络}（decoder backbone）\(g_{\eta}^{\backbone}
                \colon (\R^{d})^{*} \to (\R^{d})^{*}\)，它在
                概念上撤销了编码器主干网络的操作。
            \item 一个 \textit{逆嵌入} \(g_{\eta}^{\unemb} \colon
                (\R^{d})^{*} \to \cD\)，它作为嵌入的
                逆操作。
        \end{itemize}
        我们还定义 \(f_{\theta} \doteq f_{\theta}^{\backbone}
        \circ f_{\theta}^{\emb} \colon \cD \to (\R^{d})^{*}\) 且
        \(g_{\eta} \doteq g_{\eta}^{\unemb} \circ
        g_{\eta}^{\backbone} \colon (\R^{d})^{*} \to \cD\)。给定
        输入 \(\vX\)，我们记 \(\vZ_{\theta}(\vX) \doteq
        f_{\theta}(\vX)\) 且 \(\hat{\vX}_{\theta, \eta}(\vX) \doteq
        g_{\eta}(\vZ_{\theta}(\vX))\)。整体流水线
        如 \Cref{fig:overall_autoencoder_pipeline} 所示。
\end{itemize}

在本章中，我们将多次重复使用这种符号，因此
如果有不清楚的地方，请随时回头查阅。
我们对网络的这种分解也非常贴近大多数代码
实现，你可以通过定义这些网络来开始你的
编码项目。

\section{图像表示的无监督
学习}\label{sec:contrastive_learning}

学习图像数据分布的高质量且忠实的表示
是机器智能中的一个基本问题，正如我们
在 \ref{ch:lossy-compression} 章中详细讨论的那样。
在文献中，已经提出了许多针对此任务的流行方法，
其中许多方法似乎并没有使用本书中概述的技术和
原理。其中一种方法被称为
\textit{对比学习}（contrastive learning），之所以这样命名，是因为其学习
目标（粗略地说）是确保
“相似”数据的特征相似，而“不相似”数据的特征
相互远离。对比学习解决方案通常是
高度工程化、基于经验设计的方法。在本节中，
我们将描述一种名为 DINO \citep{caron2021emerging} 的此类方法，并使用本书中描述的原理来大幅简化其设计决策，同时改进学习到的表示。

\subsection{数据}\label{sub:contrastive_learning_data}
我们将用来探索和简化 DINO
方法的数据都是二维（2D）图像数据。 

对于 \textit{训练}，
我们将使用 ImageNet-1K 和 ImageNet-21K 数据集。数据集中的每个样本都是一张 RGB 图像（分辨率各异），以及一个指示图像包含的对象或场景的标签（即图像的 \textit{类别}）。ImageNet-1K 数据集包含 128 万张训练图像和 5 万张验证图像，分为 1000 个类别。ImageNet-21K 数据集包含 1420 万张训练图像和 2.18 万个类别，但这些类别并非互不相交（即某些类别是其他类别的子集）。由于我们进行的是自监督学习，标签在训练期间不会被使用，仅在评估期间使用。为了提高模型的鲁棒性，我们通常在处理过程中对每张图像应用 \textit{小幅度数据增强}，例如翻转、添加微小的随机噪声或随机大面积裁剪；我们没有将此包含在符号中，因为自然图像的每次增强本身（非常接近于）我们数据集中的自然图像。

对于 \textit{评估}，我们将使用一系列数据集。其中最常见的是 CIFAR-10。CIFAR-10 是一个包含 6 万张 32 \(\times\) 32 RGB 自然图像的数据集，分为 10 个类别，具有
预先建立的 5 万个样本的训练集和 1 万个样本的验证集。使用 CIFAR-10 的目的是确保在一种图像分布（ImageNet）上训练的模型能够泛化到另一种图像分布（CIFAR-10）。我们还参考了其他
类似的数据集，例如 CIFAR-100（与 CIFAR-10 不相交）、Oxford Flowers 和 Oxford Pets。ImageNet-1K 和 CIFAR-10 数据的示例见 \Cref{fig:in1k_cifar10_examples}。 

\begin{figure}
    \centering

    \hfill
    \begin{subfigure}{0.3\textwidth}
        \centering
        \includegraphics[width=\textwidth]{\toplevelprefix/chapters/chapter8/figs/imagenet.pdf}
        \caption{\small ImageNet-1K 样本。}
    \end{subfigure}
    \hfill
    \begin{subfigure}{0.26\textwidth}
        \centering
        \includegraphics[width=\textwidth]{\toplevelprefix/chapters/chapter8/figs/cifar10.pdf}
        \caption{\small CIFAR10 样本。}
    \end{subfigure}
    \hfill
    \phantom{}

    \caption{\small\textbf{来自 ImageNet-1K \textit{（左）}
        和 CIFAR-10 \textit{（右）}的图像。} 请注意，CIFAR-10 图像
        的分辨率要低得多，一般来说，这降低了
    学习该分布的复杂性。}
    \label{fig:in1k_cifar10_examples}
\end{figure}

在稍微更正式的层面上，我们的数据 \(\vX\) 将是图像；我们
设 \(\cI\) 为所有图像的集合。由于图像是一个矩形的
像素阵列，并且每个像素都有由 RGB、CMYK 或
其他颜色格式给出的颜色，我们说图像是 \(\R^{c
\times h \times w}\) 的一个元素——这里 \(c\) 是通道数（即
RGB 为 \(3\)，CMYK 为 \(4\)），\(h\) 是图像高度，而
\(w\) 是图像宽度。因此，所有图像的集合 \(\cI
\doteq \bigcup_{c, h, w = 1}^{\infty}\R^{c \times h \times w}\) 是
所有可能的此类数据的集合。同样，我们将反复使用这种符号。

\subsection{简化的 DINO 任务与目标}
\label{sub:contrastive_learning_objective}

我们的任务是学习数据的良好表示。总的来说，对比
学习通过以下方式做到这一点：定义我们希望特征反映输入图像的哪些属性，
构建共享这些属性但其他属性不同的图像，
并设置一个损失函数，
促使具有共享属性的图像的特征
相互接近，而具有不同属性的图像的特征相互远离。
这个学习问题的自然最优解是，
学习到的特征保留了输入的期望属性。
然而，在训练对比模型的过程中，
会出现许多实际和经验上的复杂情况。

在 DINO 的案例中，作者提出使用一种方法，
该方法为整张图像生成一个单一的特征向量，并希望
该特征向量包含“全局”（即图像级别）信息。
相应地，损失函数将促使具有相似全局
信息的图像具有相似的特征，而具有不同全局
信息的图像具有不同的特征。

这似乎很直观，但如前所述，即使在设置损失函数时，也有几个
经验上的考虑。首先也是最重要的是，我们应该如何促进相似性和差异性？
DINO \citep{caron2021emerging} 给出的答案是\footnote{在
作者看来，这令人费解……}将输出特征转换为
对应于某种概率分布的“逻辑值”（logits），并计算
它们的交叉熵。更具体地说，令 \(\Delta_{m} \doteq \{\vx
        \in \R^{m} \colon x_{i} \geq 0\ \forall i \in [m], \sum_{i =
1}^{m}x_{i} = 1\}\) 为 \(\R^{m}\) 中的概率向量空间，
并定义函数 \(d_{\CE} \colon \R^{m} \times \R^{m} \to \R\) 如下：
\begin{equation}\label{eq:cross_entropy_difference}
    d_{\CE}(\vp, \vq) \doteq \CE(\vp, \vq), \quad \forall \vp, \vq
    \in \Delta_{m}
\end{equation}
其中 \(\CE \colon \Delta_{m} \times \Delta_{m} \to \R\) 是
交叉熵，定义为
\begin{equation}\label{eq:expts_def_ce}
    \CE(\vp, \vq) \doteq -\sum_{i = 1}^{m}p_{i}\log q_{i}, \quad
    \forall \vp = (p_{1}, \dots, p_{m}), \vq = (q_{1}, \dots, q_{m})
    \in \Delta_{m}.
\end{equation}
在继续讨论之前，让我们建立一些关于
这个距离函数的直观理解。特别地，我们有
\begin{align}
    \CE(\vp, \vq)
    &= -\sum_{i = 1}^{m}p_{i}\log q_{i} = \sum_{i = 1}^{m}p_{i}\log(p_{i}/q_{i})
    - \sum_{i = 1}^{m}p_{i}\log p_{i} = \KL(\vp\mmid \vq) + H(\vp)
\end{align}
其中 \(\KL \colon \Delta_{m} \times \Delta_{m} \to \R\) 是 KL
散度，定义为
\begin{equation}
    \KL(\vp\mmid \vq) \doteq \sum_{i = 1}^{m}p_{i}\log(p_{i}/q_{i}),
\end{equation}
而 \(H \colon \Delta_{m} \to \R\) 是随机变量的熵。请注意，
当且仅当 \(\vp = \vq\) 时，\(\KL(\vp\mmid \vq)\) 取得最小值。
因此，最小化 \(d_{\CE}\) 会起到两个作用：它使得 \(\vp = \vq\)，
并且它使得 \(\vp\) 和 \(\vq\) 具有最小的熵（即，
    在一个分量上为 \(1\) 而在其他地方为 \(0\) 的向量——这些被称为
\textit{独热向量}（one-hot vectors））。总的来说，这个目标的意图不仅
\textit{仅仅}是匹配 \(\vp\) 和 \(\vq\)，而且还要以某种方式塑造它们，
使它们成为低熵的。在讨论公式时，
请牢记这一点。

下一个问题是，我们应该如何获取具有相似
全局信息的样本？DINO（以及几乎所有
对比学习）给出的答案是 \textit{数据增强}——从每个
样本中，制作几个共享所需属性的相关样本。
在 DINO 的案例中，我们使用输入图像的不同裁剪或
\textit{视图}（views）。回想一下，我们将图像建模为
集合 \(\cI\) 的一个元素。在这种符号下，视图是一个函数
\(v \colon \cI \to \cI\)。在 DINO 的案例中，视图是一个
\textit{随机调整大小的裁剪}：它对图像进行随机选择的矩形
裁剪（其占图像总面积的固定百分比 \(p_{v} \in [0, 1]\)），
按比例调整其大小，使得
\textit{较短的边}长为 \(S_{\rsz}\) 像素，然后将其调整
为固定形状 \((C, S_{v}, S_{v})\)，其中 \(S_{v} \geq 1\) 是
视图的大小，\(C\) 是原始图像中的通道数。

\begin{figure}
    \centering
    \includegraphics[width=0.7\textwidth]{\toplevelprefix/chapters/chapter8/figs/global_local_views.pdf}
    \caption{\textbf{DINO 中的局部和全局视图。} 局部视图和
        全局视图对输入图像进行矩形裁剪，并
        将其调整为正方形，然后输入到网络中
    进行处理。}
    \label{fig:dino_local_global_views}
\end{figure}

我们想要使用两种类型的视图，如
\Cref{fig:dino_local_global_views} 所示：
\begin{itemize}
    \item \textit{全局视图}（global views），即面积百分比参数为
        \(p_{\glo} \in [0, 1]\) 且输出
        形状为 \((C, S_{\glo}, S_{\glo})\) 的随机调整大小的裁剪；
    \item \textit{局部视图}（local views），即面积百分比参数为
        \(p_{\loc} \in [0, 1]\) 且输出
        形状为 \((C, S_{\loc}, S_{\loc})\) 的随机调整大小的裁剪。这里 \(p_{\loc} <
        p_{\glo}\) 且 \(S_{\loc} < S_{\glo}\)。
\end{itemize}

DINO 希望输入图像 \(\vX\) 的所有视图
\(\vX_{v} \doteq v(\vX)\) 的聚合特征 \(\vz_{\theta}(\vX_{v})
\doteq (f_{\theta}^{\ext} \circ f_{\theta})(\vX_{v})\) 彼此
保持一致。DINO 通过使用一个 \textit{“DINO
头”}\footnote{请注意，\(h_{\vW, \vmu}\) 是特定任务
    头，在 \Cref{sec:experiment_setup} 中它仅由
    \(\theta\) 参数化，而不是任何特定参数，但由于我们使用
    具有不同第二个参数值的两次 \(h\) 调用，
我们保留了指定的符号。} \(h_{\vW, \vmu}\)（由矩阵 \(\vW \in \R^{s \times d}\) 和向量
\(\vmu \in \R^{s}\) 参数化）来实现这一点，
使用以下简单的配方从聚合特征 \(\vz_{\theta}(\vX_{v})\) 中
提取概率向量 \(\vp_{\theta,
\vW, \vmu}(\vX_{v}) \doteq h_{\vW, \vmu}(\vz_{\theta}(\vX_{v}))\)：
\begin{equation}
    h_{\vW, \vmu}(\vz) \doteq \softmax([\vW\vz - \vmu]/\tau), \qquad
    \forall \vz \in \R^{d},
\end{equation}
其中 \(\softmax \colon \R^{s} \to \Delta_{s}\) 函数定义为
\begin{equation}
    \softmax\rp{\mat{x_{1} \\ \vdots \\ x_{s}}} \doteq
    \frac{1}{\sum_{i = 1}^{s}e^{x_{i}}}\mat{e^{x_{1}} \\ \vdots \\ e^{x_{s}}}
\end{equation}
且 \(\tau > 0\) 是一个“温度”参数，它控制着
softmax 输出的熵。

具体而言，DINO 最小化了每个全局视图 \(\vX_{g} \doteq
v_{g}(\vX)\) 的概率向量 \(\vp_{\theta, \vW, \vmu}(\vX_{g}) \doteq h_{\vW,
\vmu}(\vz_{\theta}(\vX_{g}))\) 与每个视图
\(\vX_{c} \doteq v_{c}(\vX)\) 的概率向量 \(\vp_{\theta, \vW}(\vX_{c})
\doteq h_{\vW, \vzero_{m}}(\vz_{\theta}(\vX_{c}))\) 之间的差异。
这里，\(v_{c}\) 可以是局部视图，\textit{也可以}是
全局视图。我们将在 \Cref{sub:contrastive_learning_architecture} 中
稍后讨论 \(f_{\theta}\) 和 \(f_{\theta}^{\ext}\) 的实现。
总体而言，DINO 解决以下问题：
\begin{eqnarray}
    &&\min_{\theta, \vW, \vmu}\cL_{\dino}(\theta, \vW, \vmu) \qquad
    \text{其中} \nonumber \\
    &&\cL_{\dino}(\theta, \vW, \vmu) \doteq \Ex[d_{\CE}(\vp_{\theta,
    \vW, \vmu}(\vX_{g}), \vp_{\theta, \vW}(\vX_{c}))],\label{eq:dino_loss}
\end{eqnarray}
其中期望是针对数据 \(\vX\)、全局视图 \(v_{g}\)
和其他视图 \(v_{c}\) 计算的。

然而，在这个特定情况下，如果你尝试实现
\eqref{eq:dino_loss} 并在真实网络上对其进行优化，你很
可能会遇到一个问题：在运行了几次
学习算法的迭代之后，特征映射
\(f_{\theta}^{\ext} \circ f_{\theta}\) 将\textit{变成
常数函数}！这当然优化了上述损失，因为它
最小化了同一图像不同视图的特征之间的距离。
但我们显然不想学习到这个解。

实际上，避免崩溃是对比学习中
一个非常常见的考虑因素。那么在这种情况下我们该怎么做呢？DINO 的解决方案
再次是基于经验设计的，它仔细调整了
参数 \(\vmu\)（使用批次中的所有样本进行更新）的优化，
以及一个“温度”超参数 \(\tau\)，
该超参数是 \(h_{\vW, \vmu}\) 实现的一部分，并
在 \Cref{sub:contrastive_learning_architecture} 中进行了讨论。给定
一组效果良好的特定超参数，这确实
足以确保表示不发生崩溃。然而，在
这种特殊配置之外，训练模型收敛是
困难的，并且训练过程高度不稳定。

为了改善这种状况，让我们讨论对公式的简化。
首先，我们不再使用聚合特征 \(\vz_{\theta}\) 的学习变换
\(h_{\vW, \vmu}\) 来计算概率向量，
而是可以\textit{直接使用聚合
表示}，忽略特定任务头（或者等效地，
将其设置为恒等映射）。但现在我们需要一种直接比较
向量的方法。使用我们在
\Cref{ch:deep-networks} 中的假设，即良好的表示应该具有
欧几里得（子空间）几何结构，一种更自然的差异度量是
\textit{平方 \(\ell^{2}\) 距离}
\(d_{\ell^{2}} \colon \R^{d} \times \R^{d} \to \R\)，定义为
\begin{equation}\label{eq:cosine_similarity}
    d_{\ell^{2}}(\vx, \vy) \doteq \frac{1}{2}\norm{\vx -
    \vy}_{2}^{2}, \qquad  \forall \vx, \vy \in \R^{d}.
\end{equation}
这种基于距离的分数计算起来比
交叉熵分数更高效。因此，在我们的简化中，
\(d_{\ell^{2}}\) 取代了 \(d_{\CE}\)。

以前，通过使用技巧更新 \(\vmu\) 和 \(\tau\) 来避免崩溃。
在我们的简化中，如果我们在表示空间内比较特征，
而不是将它们转换为概率，
我们就没有这两个参数，因此必须考虑一种
不同的方法来避免崩溃。为此，我们回归到
基本原理。避免崩溃的基本思想是，为了
确保所有样本不会返回完全相同的
特征，我们需要不同的样本具有不同的特征。换句话说，
我们希望特征的\textit{协方差}在某种意义上
是\textit{大}的。但从
\Cref{ch:lossy-compression,ch:deep-networks} 中，我们已经有了一个
衡量协方差矩阵大小的量。也就是说，我们使用
直观的（总体级别的）\textit{高斯编码率}
\(R\) 来确保具有不同全局信息的不同图像的全局视图特征
能够很好地分离且不崩溃（因此被扩展）。
整体修改后的损失
\(\cL_{\simdino}\) 变为：
\begin{equation}\label{eq:simdino_loss}
    \cL_{\simdino}(\theta) \doteq \Ex[d_{\ell^{2}}(\vz_{\theta}(\vX_{g}),
    \vz_{\theta}(\vX_{c}))] - \frac{\gamma}{2}\log\det\rp{\vI +
    \frac{d}{\eps^{2}}\Cov(\vz_{\theta}(\vX_{g}))},
\end{equation}
其中 \(\eps > 0\) 是固定的，并且适当的期望与
之前一样，是针对数据 \(\vX\)、全局视图 \(v_{g}\) 和其他
（局部或全局）视图 \(v_{c}\) 计算的。\eqref{eq:simdino_loss} 中的损失
本质上与 \ref{ch:lossy-compression} 章 \ref{sec:SimDINO} 节中为
简化的 DINO（“SimDINO”）引入的损失相同。
正如我们将看到的，如果实现得当，它的效果至少
与原始 DINO 一样好。

\subsection{架构：视觉
Transformer}\label{sub:contrastive_learning_architecture}

对于架构，我们使用标准的视觉 Transformer。以下是
这种架构在图像数据背景下的正式工作原理。
回顾 \Cref{sec:experiment_setup}，编码器架构有四个
组件，即嵌入、
主干网络、特征提取器和特定任务头。我们现在
讨论这四个部分。

\begin{figure}
    \centering
    \includegraphics[width=0.8\textwidth]{\toplevelprefix/chapters/chapter8/figs/patchify.pdf}
    \caption{\small\textbf{将图像转换为 \(5
        \times 5\) 个方形图块（patch）并按光栅顺序排列的示例。}
        每个图块的大小相同，图块网格的
        形状为 \((N_{H}, N_{W}) = (5, 5)\)。然后将图块网格
    按光栅顺序展开为长度为 \(5 \times 5 = 25\) 的序列。}
    \label{fig:patchify_rasterize}
\end{figure}

\begin{figure}
    \centering
    \includegraphics[width=\textwidth]{\toplevelprefix/chapters/chapter8/figs/transformer_embedding.pdf}
    \caption{\small\textbf{Transformer 嵌入流水线。} 给定
        按光栅顺序展开的图块序列 \(\vX^{\patch}\)，
        每个展开的图块被线性投影到特征空间，
        并配备一个（加性）位置编码和一个
        称为类别 token（class token）的附加 token。输出是
        第一层输入特征 \(\vZ_{\theta}^{1}(\vX) =
    f_{\theta}^{\emb}(\vX)\)。}
    \label{fig:transformer_embedding}
\end{figure}

\paragraph{嵌入。} 给定图像数据 \(\vX \in \cI\)，我们使用映射
\(f_{\theta}^{\emb}\) 将其嵌入为 \(\R^{d}\) 中的 token 序列，
具体如下。前两个步骤如
\Cref{fig:patchify_rasterize} 所示，后两个步骤如
\Cref{fig:transformer_embedding} 所示。
\begin{enumerate}
    \item 首先，我们将图像数据 \(\vX\) 转换为形状为
        \((C, P_{H}, P_{W})\) 的图块序列，其中 \(P_{H}\) 和
        \(P_{W}\) 是图块的尺寸。我们假设 \(P_{H}\)
        和 \(P_{W}\) 分别整除 \(\vX\) 的高度和宽度
        （在 \Cref{sub:contrastive_learning_objective} 的符号中，
            我们假设
            \(P_{H}\) 和 \(P_{W}\) 整除 \(S_{\loc}\) 和
        \(S_{\glo}\)）。设生成的图块网格具有
        \(N_{H}\) 行和 \(N_{W}\) 列。
    \item 我们将每个图块展开为长度为 \(D \doteq
        CP_{H}P_{W}\) 的向量。共有 \(N \doteq N_{H}N_{W}\) 个图块
        向量，我们将它们按“光栅顺序”（左上 \(\to\)
        右上 \(\to\) 左下 \(\to\) 右下）放入
        矩阵 \(\vX^{\patch} \in \R^{D \times N}\) 中，其中
        \(\vX^{\patch} \doteq f^{\patch}(\vX)\)。请注意，\(D\)
        仅取决于图块大小和通道数。由于
        后者在同一数据集的样本中通常是恒定的，
        因此 \(D\) 对于数据集中的所有图像都是相同的，
        而 \(N\) 对于较大和较小的图像则不同。
    \item 然后我们对 \(\vX^{\patch}
        \in \R^{D \times N}\) 执行以下操作，将其投影到 \(\R^{d \times n}\)
        其中 \(n \doteq N + 1\)：
        \begin{equation}
            \vX^{\patch} \mapsto [\vz_{\cls}^{1}, \vW^{\emb}\vX] + \vE^{\pos}.
        \end{equation}
        这里我们有三个可训练参数 \(\vW^{\emb}\)、
        \(\vz_{\cls}^{1}\) 和 \(\vE^{\pos}\)，它们的用途如下：
        \begin{itemize}
            \item \(\vW^{\emb} \in \R^{d \times D}\) 是一个矩阵，
                它将每个图块向量投影为一个 token 特征。
            \item \(\vz_{\cls}^{1} \in \R^{d}\) 是一个所谓的
                \textit{类别 token}（class token）或 \textit{寄存器 token}（register token）。
                类别 token 启发式地保存了整个数据的全局信息，
                并用于下游任务。在
                \Cref{ch:deep-networks} 的压缩深度网络框架中，
                我们期望在
                前向传播的过程中，类别 token 被投影到与
                显著或语义相关的 token 相同的子空间上。
            \item \(\vE^{\pos} \in \R^{d \times N}\) 是一个所谓的
                \textit{位置编码}（positional encoding），它将
                不同图块的 token 区分开来。也就是说，
                token 特征应该具有位置信息，以便
                整体映射 \(f^{\pre}\) 对图块的排列不是不变的，
                而 \(\vE^{\pos}\)
                插入了这种位置信息。
                \begin{itemize}
                    \item 在这个 DINO 案例中，由于 Transformer
                        接收不同大小的图像，我们为训练期间接收到的
                        最大尺寸学习一个位置编码，
                        并通过插值来
                        获得较小尺寸输入的位置编码。
                \end{itemize}
        \end{itemize}
\end{enumerate}
因此，最终我们有
\begin{equation}\label{eq:definition_of_embedding_module}
    f_{\theta}^{\emb}(\vX) \doteq \mat{\vz_{\cls}^{1},
    \vW^{\emb}f^{\patch}(\vX) + \vE^{\pos}}.
\end{equation}
所有参数 \(\vz_{\cls}^{1}, \vW^{\emb}, \vE^{\pos}\) 都
包含在参数集 \(\theta\) 中。

\begin{figure}
    \centering
    \includegraphics[width=\textwidth]{\toplevelprefix/chapters/chapter8/figs/transformer_backbone.pdf}
    \caption{\small\textbf{Transformer 主干网络的一层 \(f_{\theta}^{\ell}\)。}
        输入特征依次通过
        层归一化、多头自注意力机制和多层
        感知机块，形成该层的输出特征。}
    \label{fig:transformer_backbone}
\end{figure}

\paragraph{主干网络。} 给定一个嵌入序列
\(\vZ_{\theta}^{1}(\vX) \doteq f_{\theta}^{\emb}(\vX) \in
(\R^{d})^{*}\)，我们使用主干映射
\(f_{\theta}^{\backbone}\) 对其进行处理，如下所示且如
\Cref{fig:transformer_backbone} 所示。函数
\(f_{\theta}^{\backbone}\) 由 \(L\) 个 \textit{层}（layers）
\(f_{\theta}^{\ell}\) 组成，即，
\begin{equation}
    f_{\theta}^{\backbone} = f_{\theta}^{L} \circ \cdots \circ f_{\theta}^{1}.
\end{equation}
层 \(f_{\theta}^{\ell}\) 具有以下实现：
\begin{align}\label{eq:vit-res-block}
    \vZ_{\theta}^{\ell + 1/2}(\vX)
    &= \vZ_{\theta}^{\ell}(\vX) +
    \MHSA_{\theta}^{\ell}(\LN_{\theta}^{1, \ell}(\vZ_{\theta}^{\ell}(\vX))) \\
    \vZ_{\theta}^{\ell + 1}(\vX)
    &= \vZ_{\theta}^{\ell + 1/2}(\vX) +
    \MLP_{\theta}^{\ell}(\LN_{\theta}^{2, \ell}(\vZ_{\theta}^{\ell + 1/2}(\vX)))
\end{align}
并且 \(f_{\theta}^{\ell}\) 的定义使得
\(f_{\theta}^{\ell}(\vZ_{\theta}^{\ell}(\vX)) \doteq
\vZ_{\theta}^{\ell + 1}(\vX)\)。这里我们使用了一些算子，
例如 \(\MHSA_{\theta}^{\ell}, \MLP_{\theta}^{\ell}\) 和
\(\LN_{\theta}^{i, \ell}\)，它们的定义如下：
\begin{itemize}
    \item \(\MHSA_{\theta}^{\ell}\) 算子是
        多头自注意力（multi-head-self-attention），它是多头
        子空间自注意力（参见 \Cref{ch:deep-networks}）的前身。其
        公式如下：
        \begin{align}
            &\MHSA_{\theta}^{\ell}(\vZ) \label{eq:mhsa}
            \doteq \vU_{\out}^{\ell}\mat{\SA([\vU_{\query}^{1,
                    \ell}]^{\top}\vZ, [\vU_{\attnkey}^{1, \ell}]^{\top}\vZ,
                [\vU_{\val}^{1, \ell}]^{\top}\vZ) \\ \vdots
                \\ \SA([\vU_{\query}^{K, \ell}]^{\top}\vZ,
                    [\vU_{\attnkey}^{K, \ell}]^{\top}\vZ, [\vU_{\val}^{K,
            \ell}]^{\top}\vZ)} + \vb_{\out}^{\ell}\vone_{n}^{\top}, \\
            \label{eq:self_attention}
            &\text{其中} \qquad \SA(\vQ, \vK, \vV)
            \doteq \vV
            \underbrace{\softmax\rp{\frac{\vK^{\top}\vQ}{\sqrt{p}}}}_{\doteq
            \vA(\vQ, \vK)}
        \end{align}
        其中 \(p\) 是一个正整数，\(\vU_{\query}^{k, \ell},
            \vU_{\attnkey}^{k, \ell}, \vU_{\val}^{k, \ell} \in \R^{d
        \times p}\)、\(\vU_{\out}^{\ell} \in \R^{d \times Kp}\) 和
        \(\vb_{\out}^{\ell} \in \R^{d}\) 是包含在参数集
        \(\theta\) 中的可训练参数，并且
        \(\softmax\) 按列定义为
        \begin{align}
            \softmax(\vM)
            &\doteq \mat{\softmax(\vm_{1}) & \cdots & \softmax(\vm_{p})}, \\
            \forall \vM
            &= \mat{\vm_{1}, \dots, \vm_{p}} \in \R^{n \times p}.
        \end{align}
        在实践中，维度的选择通常使得 \(Kp
        = d\)。项
        \begin{eqnarray}
            \vA_{\theta}^{k, \ell}(\vZ) &\doteq&
            \vA([\vU_{\query}^{k, \ell}]^{\top}\vZ,
            [\vU_{\attnkey}^{k, \ell}]^{\top}\vZ), \nonumber \\
            \SA_{\theta}^{k, \ell}(\vZ) &\doteq&
            \SA([\vU_{\query}^{k, \ell}]^{\top}\vZ,
                [\vU_{\attnkey}^{k, \ell}]^{\top}\vZ, [\vU_{\val}^{k,
            \ell}]^{\top}\vZ)
            \label{eq:attention_map}
        \end{eqnarray}
        也分别被称为第 \(\ell\) 层的 \textit{第 \(k\) 个注意力图}（attention map）和
        \textit{第 \(k\) 个注意力头输出}（attention head output）。
        此外，操作 \(\SA(\vQ, \vK,
        \vV)\) 可以使用专门的软件（如 FlashAttention
        \citep{shah2025flashattention}）极其高效地
        计算。
    \item \(\MLP_{\theta}^{\ell}\) 是一个两层感知机，是
        深度网络中使用的常规非线性结构，其形式为
        \begin{equation}
            \MLP_{\theta}^{\ell}(\vZ) \doteq
            \vW_{\down}^{\ell}\ReLU(\vW_{\up}^{\ell}\vZ +
            \vb_{\up}^{\ell}\vone_{n}^{\top}) +
            \vb_{\down}^{\ell}\vone_{n}^{\top}
        \end{equation}
        其中 \(\vW_{\up}^{\ell} \in \R^{q \times d},
            \vW_{\down}^{\ell} \in \R^{d \times q}, \vb_{\up}^{\ell} \in
        \R^{q}, \vb_{\down}^{\ell} \in \R^{d}\) 是
        同样包含在参数集 \(\theta\) 中的可训练参数，
        而 \(\ReLU\) 是逐元素的 ReLU 非线性激活函数，即
        \(\ReLU(\vM)_{ij} = \max\{M_{ij}, 0\}\)。
    \item 对于 \(i \in \{1, 2\}\)，每个层归一化 \(\LN_{\theta}^{i, \ell}\)
        都是标准的归一化，它独立地按列应用于
        每个 token 特征：
        \begin{equation}
            \LN_{\theta}^{i, \ell}(\vZ) = \LN_{\theta}^{i,
            \ell}(\mat{\vz_{1}, \dots, \vz_{n}}) =
            \mat{\LN_{\theta}^{i, \ell}(\vz_{1}), \dots,
            \LN_{\theta}^{i, \ell}(\vz_{n})}
        \end{equation}
        其形式为
        \begin{equation}
            \LN_{\theta}^{i, \ell}(\vz) = \frac{\vz
            - \operatorname{mean}(\vz)\vone_{d}}{\norm{\vz
            - \operatorname{mean}(\vz)\vone_{d}}_{2}} \hada
            \valpha^{i, \ell} + \vbeta^{i, \ell} \qquad \text{其中}
            \qquad \operatorname{mean}(\vz) = \frac{1}{d}\vone_{d}^{\top}\vz
        \end{equation}
        其中 \(\hada\) 表示逐元素乘法，且
        \(\valpha^{i, \ell}, \vbeta^{i, \ell} \in \R^{d}\) 是
        包含在参数集 \(\theta\) 中的
        可训练参数。层归一化算子充当了对
        每个 token 的一种归一化，其中每个 token 之后的尺度
        是可学习的，并在所有 token 之间共享。
\end{itemize}

Transformer 是历史上最受欢迎的神经网络
架构之一，为几乎所有深度学习领域的
应用提供动力。

\paragraph{特征提取器。} 我们使用一个后处理步骤
\(f_{\theta}^{\ext}\)，它提取 \textit{类别 token
特征}（回想一下，这是旨在包含输入图像聚合
信息的特征），并对其应用 MLP 和
归一化。也就是说，我们有
\begin{equation}
    \vz_{\theta}(\vX) \doteq f_{\theta}^{\ext}(\vZ_{\theta}(\vX)) =
    f_{\theta}^{\ext}([\vz_{\theta}^{1}(\vX), \dots,
    \vz_{\theta}^{n}(\vX)]) \doteq
    \frac{\MLP_{\theta}^{\ext}(\vz_{\theta}^{1}(\vX))}{\norm{\MLP_{\theta}^{\ext}(\vz_{\theta}^{1}(\vX))}_{2}}.
\end{equation}

\paragraph{特定任务（“DINO”）头。} 对于 DINO，我们使用
特定任务的 DINO 头 \(h_{\vW, \vmu}\)。对于 SimDINO，如前所述，
我们\textit{完全不}使用特定任务头。

\subsection{优化策略}\label{sub:contrastive_learning_optimization}

\begin{figure}
    \centering
    \includegraphics[width=\textwidth]{\toplevelprefix/chapters/chapter8/figs/dino_pipeline.pdf}
    \caption{\small \textbf{DINO 流水线。} 为每个输入计算
        学生特征和教师特征。目标是
        尝试通过将两组特征投影到高维
        概率单纯形并计算交叉熵损失，
        使学生特征与教师特征对齐。值得注意的是，
        由于“停止梯度”（stop-grad），梯度仅
    相对于\textit{学生参数的输出}进行计算。}
    \label{fig:dino_pipeline}
\end{figure}

\paragraph{优化 DINO。} 我们已经有了损失函数和
架构，所以现在我们讨论优化策略。
DINO 的优化策略对\textit{同一架构使用两组权重}：
\textit{学生}（student）权重
\(\theta_{\student}\) 和 \textit{教师}（teacher）权重
\(\theta_{\teacher}\)。它们对应于两个具有相同架构的
不同神经网络，分别称为教师网络和学生网络。
教师网络对所有全局视图进行编码，
而学生网络对所有“其他”视图进行编码。
损失的目标是将教师的输出蒸馏到学生模型中。
也就是说，我们在损失 \(\cL_{\dino{}-\student\teacher}\) 上进行训练：
\begin{equation}\label{eq:dino_loss_teacherstudent}
    \cL_{\dino{}-\student\teacher}(\theta_{\student},
    \theta_{\teacher}, \vW_{\student}, \vW_{\teacher}, \vmu) \doteq
    \Ex[d_{\CE}(\vp_{\theta_{\teacher}, \vW_{\teacher},
    \vmu}(\vX_{g}), \vp_{\theta_{\student}, \vW_{\student}}(\vX_{c}))].
\end{equation}
现在，我们可以完整地描述 DINO 的整体流水线，如
\Cref{fig:dino_pipeline} 所示。

虽然很容易对 \eqref{eq:dino_loss_teacherstudent} 进行推理，
但在实践中，不可能使用由 \(\cL_{\dino{}-\student\teacher}\) 给出的损失
来实现诸如梯度下降之类的优化算法。
这是因为损失中的期望
是不可能评估的，更不用说对其求
梯度了。在这种极其常见的情况下，我们通过
有限样本来近似期望。也就是说，在每个时间步 \(k\)，我们：
\begin{itemize}
    \item 从我们的数据集中子采样 \(B\) 个数据点
        \(\{\vX_{1}^{(k)}, \dots, \vX_{B}^{(k)}\} \subset \cI\)。
    \item 对于每个数据点 \(\vX_{b}^{(k)}\)，采样 \(M_{\glo}\) 个
        全局视图 \(v_{b, g}^{(k), i}\) 和 \(M_{\loc}\) 个局部
        视图 \(v_{b, \ell}^{(k), i}\)。将这些视图应用于
        \(\vX_{b}^{(k)}\) 以获得 \(\vX_{b, g}^{(k), i} \doteq
        v_{b, g}^{(k), i}(\vX_{b}^{(k)})\) 和 \(\vX_{b, \ell}^{(k),
        i} \doteq v_{b, \ell}^{(k), i}(\vX_{b}^{(k)})\)。
    \item 对于每个 \textit{局部} 视图 \(\vX_{b, \ell}^{(k), i}\)，
        计算以下量：
        \begin{equation}
            \vz_{\theta_{\student}}(\vX_{b, \ell}^{(k), i}) \doteq
            (f_{\theta_{\student}}^{\ext} \circ
            f_{\theta_{\student}})(\vX_{b, \ell}^{(k), i}), \qquad
            \vp_{\theta_{\student}, \vW_{\student}}(\vX_{b,
            \ell}^{(k), i}) \doteq h_{\vW_{\student},
            \vzero_{m}}(\vz_{\theta_{\student}}(\vX_{b, \ell}^{(k), i}(\theta)))
        \end{equation}
        并且对于每个 \textit{全局} 视图 \(\vX_{b, g}^{(k), i}\)，
        计算以下量（滥用符号）：
        \begin{align}
            &\vz_{\theta_{\student}}(\vX_{b, g}^{(k), i}) \doteq
            (f_{\theta_{\student}}^{\ext} \circ
            f_{\theta_{\student}})(\vX_{b, g}^{(k), i}), \qquad
            \vp_{\theta_{\student}, \vW_{\student}}(\vX_{b, g}^{(k),
            i}) \doteq h_{\vW_{\student},
            \vzero_{m}}(\vz_{\theta_{\student}}(\vX_{b, g}^{(k), i})), \\
            &\vz_{\theta_{\teacher}}(\vX_{b, g}^{(k), i}) \doteq
            (f_{\theta_{\teacher}}^{\ext} \circ
            f_{\theta_{\teacher}})(\vX_{b, g}^{(k), i}), \qquad
            \vp_{\theta_{\teacher}, \vW_{\teacher}, \vmu}(\vX_{b,
            g}^{(k), i}) \doteq h_{\vW_{\teacher},
            \vmu}(\vZ_{\theta_{\teacher}}(\vX_{b, g}^{(k), i})).
        \end{align}
    \item 计算 \textit{代理的、近似的损失}
        \(\hat{\cL}_{\dino-\student\teacher}^{(k)}\)，定义如下：
        \begin{align}\label{eq:dino_loss_teacherstudent_empirical}
            &\hat{\cL}_{\dino{}-\student\teacher}^{(k)}(\theta_{\student},
            \theta_{\teacher}, \vW_{\student}, \vW_{\teacher}, \vmu) \doteq
            \frac{1}{BM_{\glo}(M_{\glo} + M_{\loc} - 1)}\sum_{b =
            1}^{B}\sum_{i = 1}^{M_{\glo}}\\
            &\Bigg[\sum_{j =
                1}^{M_{\loc}}d_{\CE}(\vp_{\theta_{\teacher},
                    \vW_{\teacher}, \vmu}(\vX_{b, g}^{(k), i}),
                    \vp_{\theta_{\student}, \vW_{\student}}(\vX_{b,
                \ell}^{(k), j})) + \sum_{\substack{j = 1 \\ j \neq
                i}}^{M_{\glo}}d_{\CE}(\vp_{\theta_{\teacher},
                    \vW_{\teacher}, \vmu}(\vX_{b, g}^{(k), i}),
                    \vp_{\theta_{\student}, \vW_{\student}}(\vX_{b, g}^{(k),
            j}))\Bigg]\nonumber
        \end{align}
        以及它相对于
        \(\theta_{\student}\) 和 \(\vW_{\student}\) 的梯度，这应该
        在假设 \(\theta_{\teacher}\)、
        \(\vW_{\teacher}\) 和 \(\vmu\) 是常数的情况下计算——即
        它们\textit{从计算图中分离出来}
        并且不依赖于 \(\theta_{\student}\) 和 \(\vW_{\student}\)。
    \item 通过迭代的基于梯度的优化算法更新学生参数 \(\theta_{\student}\) 和
        \(\vW_{\student}\)，并通过衰减参数分别为
        \(\nu^{(k)}\)、\(\nu^{(k)}\) 和 \(\rho^{(k)}\) 的指数移动平均
        更新 \(\theta_{\teacher}\)、
        \(\vW_{\teacher}\) 和 \(\vmu\)，即
        \begin{eqnarray}
            &&(\theta_{\student}^{(k + 1)}, \vW_{\student}^{(k + 1)})
            = \textsc{OptUpdate}^{(k)}(\theta_{\student}^{(k)},
                \vW_{\student}^{(k)}; \nabla_{(\theta_{\student},
            \vW_{\student})}\hat{\cL}_{\dino-\student\teacher}^{(k)}) \\
            &&\theta_{\teacher}^{(k + 1)}
            = \nu^{(k)}\theta_{\teacher}^{(k)} + (1 -
            \nu^{(k)})\theta_{\student}^{(k + 1)} \\
            &&\vW_{\teacher}^{(k + 1)}
            = \nu^{(k)}\vW_{\teacher}^{(k)} + (1 -
            \nu^{(k)})\vW_{\student}^{(k + 1)} \\
            &&\vmu^{(k + 1)}
            = \rho^{(k)}\vmu^{(k)} + (1 -
            \rho^{(k)})\cdot\frac{1}{BM_{\glo}}\sum_{b =
            1}^{B}\sum_{i =
            1}^{M_{\glo}}\vW^{(k)}\vz_{\theta_{\teacher}}(\vX_{b, g}^{(k), i}),
        \end{eqnarray}
        例如，如果选择的优化算法是
        随机梯度下降，我们将有更新
        \(\theta_{\student}^{(k + 1)} \doteq \theta_{\student}^{(k)}
            -
        \delta^{(k)}\nabla_{\theta_{\student}}\hat{\cL}_{\dino{}-\student\teacher}^{(k)}\)，
        等等。
\end{itemize}
请注意，优化过程相当不规则：尽管
所有四个参数在每次迭代中都会发生变化，但只有两个是
直接通过基于梯度的方法更新的。另外两个是
通过指数移动平均更新的，并且在计算任何梯度时
实际上被视为常数。训练结束后，我们丢弃
学生权重，并使用教师权重作为我们训练好的
网络 \(f\)，因为经验表明这种指数移动平均
可以稳定生成的模型（这种想法
被称为 Polyak 平均或迭代平均）。

\(\nu\) 和 \(\rho\) 在优化轨迹上的变化方式
（即函数 \(k \mapsto \nu^{(k)}\) 和 \(k
\mapsto \rho^{(k)}\)）是超参数或设计决策，通常
\(\nu^{(1)} < 1\) 且 \(\lim_{k \to \infty}\nu^{(k)} = 1\)，
\(\rho\) 也是如此。在 DINO 头 \(h_{\vW, \vmu}\) 中使用的
温度超参数 \(\tau\)，也会在
优化轨迹上发生变化（尽管这种依赖关系没有明确标出）。

使用代理（“经验”）损失将我们难以处理的
优化问题（如优化 \eqref{eq:dino_loss_teacherstudent} 中的损失）
转化为一个易于处理的随机
优化问题，世界上几乎每一个深度学习模型
都是通过运行这种优化来训练的。一旦你看到了一些例子，
这种转换就显得极其自然，我们希望在整个章节中
给出这些例子。

\begin{figure}
    \centering
    \includegraphics[width=0.7\textwidth]{\toplevelprefix/chapters/chapter8/figs/simdino_pipeline.pdf}
    \caption{\small\textbf{SimDINO 流水线。} 这里，与
        \Cref{fig:dino_pipeline} 中的 DINO 流水线相反，损失是
        直接在特征上计算的，不需要进一步的
        操作。这省去了几个大矩阵的
        参数并简化了流水线，同时使其
    训练更加稳定。}\label{fig:simdino_pipeline}
\end{figure}

\paragraph{优化 SimDINO。} 简化的 DINO 总体级别
目标在精神上非常相似，但在执行上要简单得多，即，
\begin{equation}\label{eq:simdino_loss_teacherstudent}
    \cL_{\simdino-\student\teacher}(\theta_{\student}, \theta_{\teacher}) \doteq
    \Ex\rs{d_{\ell^{2}}(\vz_{\theta_{\teacher}}(\vX_{g}),
    \vz_{\theta_{\student}}(\vX_{c}))} - \frac{\gamma}{2}\log\det\rp{\vI
+ \frac{d}{\eps^{2}}\Cov(\vz_{\theta_{\student}}(\vX_{g})))}.
\end{equation}
因此，如 \Cref{fig:simdino_pipeline} 中所详述，SimDINO
流水线严格来说比 DINO 流水线更简单。我们可以使用
DINO 训练流水线的更简单版本来优化 SimDINO。在
每个时间步 \(k\)，我们：
\begin{itemize}
\item 从我们的数据集中子采样 \(B\) 个数据点 \(\{\vX_{1}^{(k)},
    \dots, \vX_{B}^{(k)}\} \subset \cI\)。
\item 对于每个数据点 \(\vX_{b}^{(k)}\)，采样 \(M_{\glo}\) 个
    全局视图 \(v_{b, g}^{(k), i}\) 和 \(M_{\loc}\) 个局部视图
    \(v_{b, \ell}^{(k), i}\)。将这些视图应用于 \(\vX_{b}^{(k)}\) 以
    获得 \(\vX_{b, g}^{(k), i} \doteq v_{b, g}^{(k),
    i}(\vX_{b}^{(k)})\) 和 \(\vX_{b, \ell}^{(k), i} \doteq v_{b,
    \ell}^{(k), i}(\vX_{b}^{(k)})\)。
\item 对于每个 \textit{局部} 视图 \(\vX_{b, \ell}^{(k), i}\) 计算
    \(\vz_{\theta_{\student}}(\vX_{b, \ell}^{(k), i}) \doteq
        (f_{\theta_{\student}}^{\ext} \circ
    f_{\theta_{\student}})(\vX_{b, \ell}^{(k), i})\)。对于每个
    \textit{全局} 视图 \(\vX_{b, g}^{(k), i}\) 计算
    \(\vz_{\theta_{\student}}(\vX_{b, g}^{(k), i}) \doteq
        (f_{\theta_{\student}}^{\ext} \circ
    f_{\theta_{\student}})(\vX_{b, g}^{(k), i})\) 和
    \(\vz_{\theta_{\teacher}}(\vX_{b, g}^{(k), i}) \doteq
        (f_{\theta_{\teacher}}^{\ext} \circ
    f_{\theta_{\teacher}})(\vX_{b, g}^{(k), i})\)。
\item 计算 \textit{代理的、近似的损失}
    \(\hat{\cL}_{\simdino-\student\teacher}^{(k)}\)，定义如下：
    \begin{align}\label{eq:simdino_loss_teacherstudent_empirical}
        &\hat{\cL}_{\simdino{}-\student\teacher}^{(k)}(\theta_{\student},
        \theta_{\teacher}) \nonumber \\
        \doteq &
        \frac{1}{BM_{\glo}(M_{\glo} + M_{\loc} - 1)}\sum_{b =
        1}^{B}\sum_{i = 1}^{M_{\glo}}\Bigg[\sum_{j =
            1}^{M_{\loc}}d_{\ell^{2}}(\vz_{\theta_{\teacher}}(\vX_{b,
                g}^{(k), i}), \vz_{\theta_{\student}}(\vX_{b, \ell}^{(k),
            j})) \nonumber\\
            & + \sum_{j =
            1}^{M_{\glo}}d_{\ell^{2}}(\vz_{\theta_{\teacher}}(\vX_{b,
                g}^{(k), i}), \vz_{\theta_{\student}}(\vX_{b, g}^{(k),
        j}))\Bigg]\nonumber \\
        & - \frac{\gamma}{M_{\glo}}\sum_{i =
        1}^{M_{\glo}}R_{\eps}([\vz_{\theta_{\student}}(\vX_{1,
                g}^{(k), i}), \dots, \vz_{\theta_{\student}}(\vX_{B, g}^{(k),
        i})])\nonumber
    \end{align}
    其中 \(R_{\eps}\) 是在有限样本上估计的
    高斯编码率，在 \Cref{ch:deep-networks} 中有描述。
    \(\hat{\cL}_{\simdino-\student\teacher}^{(k)}\) 相对于
    \(\theta_{\student}\) 的梯度应该（再次）在
    假设 \(\theta_{\teacher}\) 为常数的情况下计算。
\item 通过迭代的基于梯度的优化算法更新学生参数 \(\theta_{\student}\)，
    并通过衰减参数为 \(\nu^{(k)}\) 的指数移动平均
    更新 \(\theta_{\teacher}\)，即
    \begin{align}
        \theta_{\student}^{(k + 1)}
        &= \textsc{OptUpdate}^{(k)}(\theta_{\student}^{(k)};
        \nabla_{\theta_{\student}}\hat{\cL}_{\simdino-\student\teacher}^{(k)})
        \\
        \theta_{\teacher}^{(k + 1)}
        &= \nu^{(k)}\theta_{\teacher}^{(k)} + (1 -
        \nu^{(k)})\theta_{\student}^{(k + 1)}.
    \end{align}
\end{itemize}
我们再次重申，梯度仅
相对于 \(\theta_{\student}\) 计算，将 \(\theta_{\teacher}\) 视为
常数。这里请注意，虽然 \(\nu\) 的选择仍然是一个
设计决策，但超参数 \(\rho\) 和 \(\tau\) 已被移除。

\subsection{评估方法}\label{sub:contrastive_learning_evals}
有几种方法可以评估训练好的 Transformer 模型。我们
在本节中重点介绍两种。让我们定义 \textit{中心裁剪}（center crop）
视图 \(v_{\cc} \colon \cI \to \cI\)，它是一个 \textit{确定性
调整大小的裁剪}：
\begin{itemize}
\item 它调整图像大小，使最短边的尺寸为
    \(S_{\rsz}\)（类似于面积百分比参数为 \(1\) 的
    随机调整大小的裁剪）；
\item 然后它进行 \textit{中心} \(S_{\cc} \times S_{\cc}\) 裁剪；
\end{itemize}
使得最终形状为 \((C, S_{\cc}, S_{\cc})\)。请注意，
给定输入时，视图 \(v_{\cc}\) 是完全确定性的。对于输入 \(\vX\)，我们记 \(\vX_{\cc} \doteq v_{\cc}(\vX)\)。这里 \(S_{\cc} \leq S_{\rsz}\)。

\paragraph{线性探测。}

评估编码器模型 \(\vX \mapsto \vz_{\theta}(\vX)\) 的第一种也是最与架构无关的方法是采用\textit{线性探测}（linear probing）。简而言之，线性探测就是在编码器计算出的聚合特征上运行逻辑回归。这告诉我们表示中存在多少语义信息，以及提取这些信息的难易程度。（也就是说：具有不同语义的图像特征在多大程度上存在于特征空间的不同子空间上？）

更正式地说，假设我们想在图像-标签数据 \((\vX, \vy)\) 上评估编码器特征的质量和保真度，其中有 \(N_{\cls}\) 个类别，且 \(\vy \in \{0, 1\}^{N_{\cls}}\) 是“独热编码”（one-hot encoding，即如果 \(\vX\) 属于第 \(i\) 类，则除了第 \(i\) 个位置为 \(1\) 外，其余所有位置均为零）。实现这一点的一种方法是求解逻辑回归问题
\begin{equation}\label{eq:linear_probing}
\min_{\vW \in \R^{N_{\cls} \times d}}\Ex[\CE(\vy, \vW \vz_{\theta}(\vX_{\cc}))].
\end{equation}
更实际地，如果我们有带标签的数据 \(\{(\vX_{b}, \vy_{b})\}_{b = 1}^{B}\)，我们可以求解由下式给出的\textit{经验}逻辑回归问题（类似于 \eqref{eq:dino_loss_teacherstudent} 与 \eqref{eq:dino_loss_teacherstudent_empirical} 的对比）：
\begin{equation}\label{eq:linear_probing_empirical}
\min_{\vW \in \R^{N_{\cls} \times d}}\frac{1}{B}\sum_{b =
1}^{B}\CE(\vy_{b}, \vW \vz_{\theta}(\vX_{b, \cc})).
\end{equation}
这个问题是关于 \(\vW\) 的凸优化问题，因此可以通过（随机）梯度下降或许多其他算法高效求解。这个线性探测器与编码器一起可以用作分类器，我们可以评估其分类准确率。通常的做法是首先在一个大型数据集（如 ImageNet-1K）上训练模型，然后在一个数据集（如 CIFAR-10 的训练集）上训练线性探测器，并在与第二个数据集同分布的第三个（“保留”）数据集（如 CIFAR-10 的评估集）上对其进行评估。

\paragraph{\(k\)近邻。} 我们还可以\textit{无需显式训练分类器}，通过使用 \(k\)近邻算法获得平均预测标签，来评估特征在分类任务上的性能。即，给定数据集 \(\{\vz_{b}\}_{b = 1}^{B} \subseteq \R^{d}\)，将另一个点 \(\vz \in \R^{d}\) 的 \(k\)个最近邻定义为 \(\operatorname{NN}_{k}(\vz, \{\vz_{b}\}_{b = 1}^{B})\)。使用此符号，我们可以计算预测标签 \(\hat{\vy}_{\theta}(\vX \mid \{(\vX_{b}, \vy_{b})\}_{b = 1}^{B})\) 如下：
\begin{eqnarray}
&&\hat{\vy}_{\theta}(\vX \mid \{(\vX_{b}, \vy_{b})\}_{b = 1}^{B}) =
\vone(i^{\star}) \quad \text{where} \nonumber \\
&& i^{\star} \doteq \argmax_{i \in [Q]}\sum_{b =
1}^{B}\vy_{b}\indvar[\vz_{\theta}(\vX_{\cc, b}) \in
\operatorname{NN}_{k}(\vz_{\theta}(\vX_{\cc}))].
\end{eqnarray}
这里，\(\vone(i) \in \Delta_{N_{\cls}}\) 是（滥用符号，参见指示变量）支撑在 \(i\) 处的独热概率向量，即第 \(i\) 个坐标为 \(1\)，其他地方为 \(0\)。也就是说，该过程在特征空间中取 \(k\) 个最近点中最常见的标签。\(k\)近邻分类准确率就是这个预测标签的准确率，即：
\begin{equation}
\Ex_{\vX, \vy}[\indvar(\hat{\vy}_{\theta}(\vX \mid \{(\vX_{b},
\vy_{b})\}_{b = 1}^{B}) = \vy)]
\end{equation}
或者更常见的是其对应的经验版本，其中 \((\vX, \vy)\) 在一个有限数据集上遍历（\textit{不是}用于 \(k\) 个邻居的现有样本 \((\vX_{b}, \vy_{b})\)）。

\paragraph{注意力图的保真度。}

对于基于 Transformer 的编码器，检查表示性能的另一种方法是检验注意力图 \(\vA^{L, k} \in \R^{n \times n}\) 的保真度，如 \Cref{eq:attention_map} 所定义，在最后一层 \(L\) 处，由以下流程给出：
\begin{align}
\vX \mapsto \cdots \mapsto \vZ^{L - 1} = [\underbrace{\vz_{1}^{L -
    1}}_{\text{class token}}, \underbrace{\vz_{2}^{L - 1} \dots,
\vz_{n}^{L - 1}}_{\text{patch tokens}}] \mapsto \vA^{k, L} =
\mat{\vA_{1, 1}^{k, L} & \vA_{1, 2:}^{k, L} \\ \vA_{2:, 1}^{k, L} &
\vA_{2:, 2:}^{k, L}}.
\end{align}
具体而言，我们检查给定输入的注意力图揭示了输入图像中哪些显著对象，即图像的哪些部分为类别标记（class token）提供了最具全局相关性的信息。实现这一点的一种特定方法是检查注意力图的某个分量，其中类别标记被提取为查询（query）并从值矩阵（value matrix）中移除，即 \(\vA_{2:, 1}^{k, L} \in \R^{1 \times (n - 1)} = \R^{1 \times N}\) 或其转置 \(\va^{k, L} = (\vA_{2:, 1}^{k, L})^{\top} \in \R^{N}\)。请注意，这个向量 \(\va^{k, L}\)（我们将其标记为“第 \(L\) 层第 \(k\) 个注意力头处的\textit{显著性向量}（saliency vector）”）对每个图块 \(1, \dots, N\) 都有一个值，我们使用该值来描述每个图块与全局信息的相关程度。特别是为了可视化的目的，我们创建了一幅新图像，其中每个图块被其在显著性向量中的对应值所替换，从而展示每个图块的贡献；我们将此图像称为“第 \(L\) 层第 \(k\) 个注意力头处的\textit{显著性图}（saliency map）”。为了可视化所有注意力头中每个图块对全局信息的总相关性，我们可以对显著性向量求平均，即 \(\tilde{\va}^{L} \doteq \frac{1}{K}\sum_{k =
1}^{K}\va^{k, L}\)，并将其展开为\textit{平均显著性图}。平均显著性图应能突出显示输入图像的相关部分。

\paragraph{目标检测与分割。}

我们可以通过将表示用于\textit{语义分割}（semantic segmentation），来评估它们如何捕获输入的细粒度（即更小或更详细的）属性。粗略地说，这意味着我们使用这些特征为输入中的所有对象构建边界框。有几种方法可以做到这一点，也有几种方法可以对生成的边界框与真实值（ground truth）进行评分。每种方法的组合对应于特定的分割指标。我们在此不正式描述它们，因为它们并没有提供特别深刻的见解，但 DINO 论文 \citep{caron2021emerging} 和 DINOv2 论文 \citep{oquab2023dinov2} 包含了实践中使用的所有指标的参考文献。

\subsection{实验设置与结果}
\label{sub:contrastive_learning_experiment_results}

由于 SimDINO 是直接建立在 DINO 之上的，为了公平比较，我们将原始论文 \citep{caron2021emerging} 给出的 DINO 最佳设置与应用于 SimDINO 的相同设置进行比较。

\paragraph{目标函数。} 在所有实验中，我们使用 \(10\) 个分辨率为 \(96 \times 96\) 的局部视图（即 \(M_{\loc} = 10\)，\(S_{\loc} = 96\)）和 \(2\) 个分辨率为 \(224 \times 224\) 的全局视图（即 \(M_{\glo} = 2\)，\(S_{\glo} = 224\)）。为局部和全局视图裁剪的原始图像的相应比例分别为 \(p_{\loc} \in [\frac{1}{20}, \frac{3}{10}]\) 和 \(p_{\glo} \in [\frac{3}{10}, 1]\)（每个视图随机选择）。调整大小后的裁剪图像中较小的边长为 \(S_{\rsz} = 256\)，中心裁剪（评估）视图的边长为 \(S_{\cc} = 224\)。所有这些设置均适用于 DINO 和 SimDINO。

\paragraph{模型架构。} 对于所有输入，我们将图块大小设置为 \(16 \times 16\)（即 \(P_{H} = P_{W} = 16\)）。我们使用 ViT \citep{dosovitskiy2020image} 架构的 small、base 和 large 模型作为嵌入和主干网络。特征提取器是一个三层 MLP，隐藏层大小为 \(2048\)，输出维度为 \(256\)，随后进行 \(\ell^{2}\) 归一化，如 \Cref{sub:contrastive_learning_architecture} 所述。对于 DINO 架构（即非 SimDINO 架构），DINO 头部 \(\vW\) 是一个 \(\R^{65536 \times 256}\) 中的矩阵，参数 \(\vmu\) 是一个 \(\R^{65536}\) 中的向量。

\paragraph{数据集与优化。} 对于预训练，我们的 DINO 复现和 SimDINO 在所有方法中均使用 ImageNet-1K 数据集。我们使用 AdamW \citep{Loshchilov2017DecoupledWD} 作为优化器，这是一个非常标准的选项。我们遵循以下超参数建议：
\begin{itemize}
\item 批量大小为 \(B = 1024\)。
\item 学习率（针对 AdamW 和学生模型）的“基础”值为 \(2 \times 10^{-3}\)。在前 \(10\) 个 epoch 中，学习率从 \(0\) 线性增加到基础值（即，在第 \(i\) 个 epoch，学习率为 \((i/10) \cdot 2
    \times 10^{-3}\)，其中 \(1 \leq i \leq 10\)）。然后在接下来的 \(90\) 个 epoch 中，学习率通过所谓的\textit{余弦调度}（cosine schedule）衰减回 \(0\)。余弦调度的定义在许多地方都有给出，包括 \href{https://pytorch.org/docs/stable/generated/torch.optim.lr_scheduler.CosineAnnealingLR.html}{PyTorch 文档}，它在训练深度视觉模型时非常常用。
\item 权重衰减（AdamW 中的 W）在训练过程中遵循从 0.04 到 \(0.4\) 的余弦调度。
\item EMA 衰减率 \(\nu\) 在训练过程中遵循从 \(0.996\) 到 \(1.0\) 的余弦调度。特别地，对于 DINO，中心化 EMA 衰减率 \(\rho\) 固定为 \(0.9\)。
\item 特别地，对于 DINO，教师温度 \(\tau_{\teacher}\) 固定为 \(0.1\)，而学生温度 \(\tau_{\student}\) 在前 \(30\) 个 epoch 期间从 \(0.04\) 线性增加到 \(0.07\)，此后固定为 \(0.07\)。
\end{itemize}
在训练期间，在提取局部和全局视图之前，我们对每张看到的图像使用一些（本质上保留信息的）数据增强，例如翻转、颜色抖动、高斯模糊和过度曝光（solarization）。控制这些增强的确切超参数未在此列出，但可参考 DINO 论文 \citep{caron2021emerging}。

对于线性探测，线性探测器通常使用 AdamW 优化器进行训练，学习率为 \(2 \times 10^{-4}\)，权重衰减为 \(0.01\)，批量大小为 \(512\)，但这些参数通常会根据具体情况进行修改以最小化损失。

\begin{table}
\caption{\small\textbf{分类性能}：DINO 和 SimDINO 在保留测试数据上的分类性能，使用 \(k\)近邻准确率（\(k = 20\)）和线性探测。在相同的迭代次数（\(100\)）下，SimDINO 在性能上明显更好，并且更稳定（使用提供的设置在 ViT-L 主干上运行的 DINO 训练优化非常不稳定，并在短时间内得到 NaN 损失）。我们还与 DINO 所基于的其他杰出方法（即 SwAV 和 MoCov3）进行了比较。}

\centering
\begin{tabular}{@{}lcccc@{}} %
    \toprule
    方法 & 模型 & Epochs & 20-NN & 线性探测
    \\
    \midrule
    DINO & ViT-B & 100 & 72.9 & 76.3 \\
    SimDINO & ViT-B & 100 & \bf 74.9 & \bf 77.3 \\
    DINO & ViT-L & 100 & -- & -- \\
    SimDINO & ViT-L & 100 & \bf 75.6 & \bf 77.4 \\
    \midrule
    \color{gray} SwAV & \color{gray} ViT-S & \color{gray} 800 &
    \color{gray} 66.3 & \color{gray} 73.5 \\
    \color{gray} MoCov3 & \color{gray} ViT-B & \color{gray} 300  &
    \color{gray} -- & \color{gray} 76.7 \\
    \bottomrule
\end{tabular}
\label{tab:dino_imagenet_linear_probing}
\end{table}

\begin{figure}
\centering
\includegraphics[width=\textwidth]{\toplevelprefix/chapters/chapter8/figs/dino_attention_maps.png}
\caption{\small\textbf{显著性图的定性比较}：由 DINO \textit{（中间行）} 和 SimDINO \textit{（底部行）} 生成的显著性图。对于每张图像，我们计算并显示最后一层 \(L\) 的平均显著性图。不同模型的显著性图相似，这意味着所有模型都收敛到关于哪些对象重要的相似概念。请注意，尽管 \(X_{\evaluation}\) 是正方形图像，但为了进行此可视化，它被插值回矩形形状。}
\label{fig:dino_attention_maps_saliency}
\end{figure}

\begin{table}
\caption{\small\textbf{分割性能}：预训练的 DINO 和 SimDINO 模型在 COCO val2017 \citep{lin2014microsoft} 上的分割性能，这是一个包含对象位置元数据的分割数据集。我们没有在 COCO 上进行训练，仅仅使用了预训练的嵌入和主干网络，边界框是通过一种称为 MaskCut \citep{wang2023cut} 的方法从特征中提取的。尽管如此，在公平比较下，SimDINO 在目标检测和分割方面超越了 DINO，甚至超越了使用更小图块大小（边长为 \(8\) 而不是 \(16\)）的 DINO。众所周知，较小的图块大小有助于提高性能，尤其是在检测和分割任务中，因此这一结果非常令人惊讶且令人鼓舞。}
\centering
\begin{tabular}{@{}llcccccccc@{}}
    \toprule
    &  & \multicolumn{3}{c}{检测 $\uparrow$} &
    \multicolumn{3}{c}{分割 $\uparrow$} \\
    方法 & 模型 & AP$_{50}$  & AP$_{75}$ & AP & AP$_{50}$ &
    AP$_{75}$ & AP  \\
    \midrule
    SimDINO &ViT-L/16 &\bf 5.4 &1.9 &2.4 &4.5 &1.4 &1.9 \\
    SimDINO &ViT-B/16 &5.2 & \bf 2.0 & \bf 2.5 & \bf4.7 & \bf 1.5 & \bf 2.0 \\
    DINO &ViT-B/16 &3.9 &1.5 &1.8 &3.1 &1.0 &1.4 \\
    \midrule
    \color{gray} DINO & \color{gray} ViT-B/8 & \color{gray}5.1 &
    \color{gray}2.3 & \color{gray}2.5 & \color{gray}4.1 &
    \color{gray}1.3 & \color{gray}1.8 \\
    \bottomrule
\end{tabular}
\label{tab:dino_segmentation}
\end{table}

\paragraph{评估结果。} 在下游分类性能方面，我们获得了 \Cref{tab:dino_imagenet_linear_probing} 中的性能。我们观察到，在公平比较下，SimDINO 的性能远高于 DINO。此外，它也稳定得多：DINO 的规定设置无法训练 ViT-L(arge) 模型。另一方面，\Cref{fig:dino_attention_maps_saliency} 展示了 DINO 和我们简化的 DINO 中平均显著性图的可视化，观察到不同模型之间的显著性图看起来非常相似，这表明模型学习到的特征在捕获细粒度细节方面至少同样出色。\Cref{tab:dino_segmentation} 中的分割和目标检测性能定量地证实了这一主张，其中 SimDINO 特征显示出比 DINO 特征实质性的改进。

\section{通过文本绑定的弱监督图像表示}\label{sec:CLIP}
%\yima{Robin, please provide a vignette about CLIP or its simplified version. Sam will give a general description of concepts and methods for joint image/text embedding and related work in Chapter 7. }

另一种有影响力的对比学习方法摆脱了纯粹的视觉比较，转而利用图像和语言之间的自然对齐。这类工作不再仅仅通过同一图像的多个视图来定义相似性，而是利用了现实世界中许多图像都伴随有提供弱语义监督的文本描述这一事实。这种范式的一个突出例子是 CLIP（对比语言-图像预训练，Contrastive Language–Image Pretraining）\citep{Radford2021-ir}。在 CLIP 中，图像及其对应的标题被视为正样本对，而不匹配的图像-文本对被视为负样本对。学习目标鼓励图像的表示接近其关联标题的表示，并远离其他图像的标题。通过将视觉表示建立在自然语言的基础上，CLIP 学习到了语义丰富的特征，这些特征捕获了超越单纯外观的高级概念。

\subsection{数据}

我们将用于探索 CLIP 方法的数据是{\em 文本-图像对}，而不仅仅是图像。最初的 CLIP 工作构建了一个大型网络规模的语料库，包含从互联网公开来源收集的约 4 亿个文本-图像对。每个样本由一张图像和一段相关的自然语言字符串（例如，标题、名称或简短描述）组成，该字符串在网络上与图像共同出现，提供了一种微弱且嘈杂的监督形式。为了鼓励对视觉概念的广泛覆盖，CLIP 论文描述了通过搜索文本与大约 500{,}000 个查询之一匹配的图像对来构建此数据集，并通过每个查询最多包含 20{,}000 个图像对来近似平衡结果。该数据集未公开发布，因此大多数社区对 CLIP 的复现和扩展都依赖于开放的网络规模文本-图像对数据集。其中最常见的是 LAION~\citep{schuhmann2022laion5bopenlargescaledataset} 系列：LAION-400M 提供了 4 亿个经过 CLIP 过滤的图像对，而 LAION-5B 将此方案扩展到了 58.5 亿个图像对。

由于原始的网络抓取图像对通常非常嘈杂（例如，图像与其关联文本不匹配、模板化的替代文本、垃圾信息、重复项或不安全内容），因此 CLIP 风格的数据集通常使用后处理管道构建，在训练之前对候选数据进行过滤和清理。典型的方法包括基本质量过滤（例如，丢弃文本极短或图像极小或损坏的图像对）、语言识别以保留目标语言的标题、去重以及删除潜在的恶意内容。

\subsection{CLIP 任务与目标}
我们的目标是学习数据的良好表示。然而，在 CLIP 设置中，“相似性”的概念不是由同一图像的两个增强视图引起的，而是由自然语言描述提供的弱语义监督引起的。具体而言，我们希望图像表示接近其关联标题的表示，并远离其他图像的标题。等价地，如果两张图像具有相似的文本描述，它们学习到的视觉特征应该是相似的，而由不相关文本描述的图像在潜在空间中应该被分开。

我们使用\textit{余弦相似度}（cosine similarity）来形式化表示之间的“相似性”概念。具体而言，设 \(f_{\theta}\) 为图像编码器，\(g_{\plainphi}\) 为文本编码器（其架构将在下一节中详细说明），两者都映射到一个公共潜在空间 \(\vZ \in \mathbb{R}^{d}\)。给定图像 \(\vX \in \cI\) 和文本字符串 \(\vT \in \cT\)，我们形成归一化嵌入：
\begin{equation}
    \vz^{I}_{\theta}(\vX) \doteq \frac{f_{\theta}(\vX)}{\|f_{\theta}(\vX)\|_2},
    \qquad
    \vz^{T}_{\plainphi}(\vT) \doteq \frac{g_{\plainphi}(\vT)}{\|g_{\plainphi}(\vT)\|_2}.
\end{equation}
那么嵌入之间的余弦相似度定义为：
\begin{equation}
    s(\vX,\vT)
    =
    \left\langle \vz^{I}_{\theta}(\vX), \vz^{T}_{\plainphi}(\vT)\right\rangle
    \in [-1,1].
\end{equation}
因此，我们的目标是最大化匹配对之间的余弦相似度，并最小化不匹配对之间的余弦相似度。为了实现这一点，我们可以使用如 \Cref{sub:text-cond} 中讨论的简单对称交叉熵损失。具体而言，给定一个包含 \(n\) 个文本-图像对的小批量 \((\vX_i, \vT_i)_{i=1}^{n}\)，我们可以将损失定义为：
\begin{equation}\label{eq:practical_clip_loss}
    \cL_{\text{CLIP}} = -\frac{1}{n}\left(
    \sum_{i=1}^{n}
    \frac{\exp(s(\vX_i,\vT_i)/\tau)}{\sum_{k=1}^{n}\exp(s(\vX_i,\vT_k)/\tau)}
    +
    \sum_{i=1}^{n}
    \frac{\exp(s(\vT_i,\vX_i)/\tau)}{\sum_{k=1}^{n}\exp(s(\vT_i,\vX_k)/\tau)}
    \right),
\end{equation}
其中 \(\tau > 0\) 是控制 softmax 函数锐度的温度参数。 

\subsection{架构：视觉塔}\label{sub:clip_vision_tower}
CLIP 的目标是在同一个潜在空间中学习图像和文本表示。因此，其架构训练了一个视觉编码器 \(f_{\theta}: \cI \to \mathbb{R}^{d} \) 和一个文本编码器 \(g_{\plainphi}: \cT \to \mathbb{R}^{d} \)，这也被称为“双塔架构”（dual-tower architecture）。

在架构方面，大多数视觉主干网络都可以使用：例如，卷积网络如 ResNet~\citep{He2016-lc}，以及 Vision Transformer~\citep{dosovitskiy2020image}（在 \Cref{sub:contrastive_learning_architecture} 中有详细说明）都是 CLIP 的有效选择。在实践中，由于更好的性能和多功能性，Vision Transformer 更为常用。与 DINO 类似，我们将类别标记特征 \(\vz_{\cls}\) 作为每张图像的表示：\(f_{\theta}(\vX) = \vz_{\cls}\)。这是因为类别标记旨在聚合关于整个图像的全局信息，并且通常用于下游预测任务。

% \paragraph{Feature Extractor.}
% Following the backbone, CLIP requires an additional feature extractor \(f^{\text{ext}}_{\theta}: \mathbb{R}^{t \times d^\prime} \to \mathbb{R}^{1 \times d^\prime} \), where \(t\) is the sequence length of the backbone's output and \(d^\prime\) is its hidden dimension. There are two mainstream designs of CLIP’s \(f^{\text{ext}}_{\theta}\), differing only in how they aggregate the backbone’s sequence output into a single \(d\)-dimensional image representation. The first is a \emph{mean extractor}: assume the backbone produces a sequence of latent features \([\vz_1,\dots,\vz_t] \in \mathbb{R}^{t\times d'}\), we take the pooled feature to be the average of token embeddings,
% \[
% \vz^{ext} =
% f^{\text{ext}}_{\theta}([\vz_1,\dots,\vz_t]) = \frac{1}{t}\sum_{i=1}^{t}\vz_i \in \mathbb{R}^{d'}.
% \]
% The second is a \emph{class-token extractor}, as is standard for Vision Transformers: assume the backbone output includes a designated class token along other latent features \([\vz_{\cls}, \vz_1,\dots,\vz_t] \in \mathbb{R}^{(t+1)\times d'}\), then we have
% \[
% \vz^{ext} = \vz_{\cls}.
% \]

% \paragraph{Task-specific Head.}
% The feature extractor is typically followed by a projection head \(f^{\text{emb}}_{\theta} : \mathbb{R}^{d'} \to \mathbb{R}^{d}\) (often a learned linear map) to project \( \vz^{ext} \) into the \(d\)-dimension latent space shared of by the final image and text representation
% \[
% \vz^{ext} = 
% \vW^{proj}\vz^{pool}
% \qquad
% \vW^{proj} \in \mathbb{R}^{d^\prime \times d}\,
% \]
% which is the final output of \(f^{I}_{\theta}\) and will later be used in the objective.

\subsection{架构：文本塔}
CLIP 文本塔 \(g_{\plainphi}\) 通常采用类似于前面描述的 Vision Transformer 的基于 Transformer 的架构。主要区别在于嵌入层，它用于将文本序列转换为潜在空间 \(\vZ \in \mathbb{R}^{d}\) 中的潜在表示序列。 

\paragraph{嵌入。} 给定文本序列 \(\vT \in \cT\)（例如，句子或文档），我们使用映射 \(g_{\plainphi}^{\text{emb}}\) 将其嵌入为 \(\R^d\) 中的标记序列，如下所示。
\begin{enumerate}
    \item 首先，我们将 \(\vT\) 标记化为来自大小为 \(|\cV|\) 的词汇表 \(\cV\) 的离散符号（子词/字符）序列。具体而言，分词器 \(t^{\tok} : \cT \to \cV^{N}\) 生成一个长度为 \(N\) 的标记序列\footnote{在实践中，人们通常会用特殊标记（例如，用于指示序列开始和结束的 \(\langle \text{bos}\rangle\) 和 \(\langle \text{eos}\rangle\) 标记）来扩充此序列，并可能对较短的序列进行填充，但我们保持符号的通用性，并将这些视为 \(\cV\) 的元素。} \((\tau_{1},\dots,\tau_{N})\)，其中 \(\tau_i \in \cV\)。
    \item 我们通过词汇映射 \(\cV\to\{1,\dots,|\cV|\}\) 将标记转换为整数索引，得到索引序列 \((s_{1},\dots,s_{N})\)。我们将这些索引放入矩阵 \(\vS \in \{0,1\}^{|\cV|\times N}\) 中，其中第 \(i\) 列是独热向量 \(\ve_{s_i}\)。
    \item 然后我们对 \(\vS\) 执行以下操作，将其投影到 \(\R^{d \times (N+1)}\)：
        \begin{equation}
            \vS \mapsto [\vz_{\cls}^{1}, \vW^{\tok}\vS] + \vE^{\pos}.
        \end{equation}
        这是 \(g_{\plainphi}^{\text{emb}}\) 的最终输出。这里我们有三个可训练参数 \(\vW^{\tok}\)、\(\vz_{\cls}^{1}\) 和 \(\vE^{\pos}\)，它们的用途如下：
        \begin{itemize}
            \item \(\vW^{\tok} \in \R^{d \times |\cV|}\) 是一个矩阵，其第 \(j\) 列是第 \(j\) 个词汇项的学习嵌入向量。因此 \(\vW^{\tok}\vS
                \in \R^{d\times N}\) 将每个离散标记映射为 \(\R^{d}\) 中的连续特征。
            \item \(\vz_{\cls}^{1} \in \R^{d}\) 是一个特殊的\textit{分类}（或\textit{摘要}）标记。启发式地，它旨在聚合关于整个序列的全局信息，并且通常用于下游预测任务。
            \item \(\vE^{\pos} \in \R^{d \times n}\) 是一个\textit{位置编码}，用于区分不同位置的标记。由于自注意力在标记维度上是排列等变的，\(\vE^{\pos}\) 注入了顺序信息，使得整体映射对标记的排列不再是不变的。
        \end{itemize}
\end{enumerate}
所有参数 \(\vz_{\cls}^{1}, \vW^{\tok}, \vE^{\pos}\) 都包含在参数集 \(\theta\) 中。词汇表 \(\cV\) 和分词器 \(t^{\tok}\) 是固定的。 

文本塔的其他组件，如主干网络和特定任务头部，遵循与图像塔相似的设计选择，为简洁起见予以省略。还值得注意的是，后续工作证明了在选择 CLIP 文本编码器方面具有极大的灵活性：大型语言模型中包含的固定文本表示，甚至简单的词袋特征，也可能是有效的。我们在 \Cref{sub:follow-up} 中讨论这些变体。

% \paragraph{Backbone.} Now, our text input (a string) has been transformed into a sequence of latent representations
% \([\vz_{\cls}, \vz_{1}, \ldots, \vz_{n}]\) in the latent space \(\R^{d}\). Therefore, its subsequent processing follows the same procedure as in \cref{sub:contrastive_learning_architecture} for image latent representations.

% \paragraph{Feature Extractor.} The text transformer adopts a class-token extractor similar to the one in \cref{sub:clip_vision_tower}.

% \paragraph{Task-specific Head.} The text transformer adopts a linear head similar to the one in \cref{sub:clip_vision_tower}.

\subsection{优化策略}
我们使用标准的端到端随机优化过程（例如带动量的 SGD、AdamW）来训练 CLIP。下面我们展示一个通用过程。

在每个时间步 \(k\)，我们：
\begin{itemize}
    \item \textbf{子采样成对数据。} 采样一个包含 \(n\) 个成对示例的小批量
    \[
        \{(\vX_i^{(k)}, \vT_i^{(k)})\}_{i=1}^{n} \subseteq \cX \times \cT,
    \]
    其中 \(\vX_i^{(k)}\) 是图像，\(\vT_i^{(k)}\) 是其关联文本。

    \item \textbf{计算潜在表示。} 对于每个数据对 \((\vX_i^{(k)}, \vT_i^{(k)})\)，使用视觉和文本编码器 \(f_{\theta}\) 和 \(g_{\plainphi}\) 计算归一化嵌入：
    \begin{equation}
        z_{\theta}^{I}(\vX_i^{(k)}) \doteq \frac{f_{\theta}(\vX_i^{(k)})}{\|f_{\theta}(\vX_i^{(k)})\|_2},
        \qquad
        z_{\plainphi}^{T}(\vT_i^{(k)}) \doteq \frac{g_{\plainphi}(\vT_i^{(k)})}{\|g_{\plainphi}(\vT_i^{(k)})\|_2}.
    \end{equation}

    \item \textbf{计算 CLIP 损失。} 计算如 \eqref{eq:practical_clip_loss} 中定义的 CLIP 损失：
    \begin{equation}
        \hat{\cL}_{\text{CLIP}}^{(k)}
        \doteq
        -\frac{1}{n}\left(    
        \textstyle
        \sum_{i=1}^{n}
        \frac{\exp(s(\vX_i^{(k)},\vT_i^{(k)})/\tau)}{\sum_{k=1}^{n}\exp(s(\vX_i^{(k)},\vT_k^{(k)})/\tau)}
        +
        \sum_{i=1}^{n}
        \frac{\exp(s(\vT_i^{(k)},\vX_i^{(k)})/\tau)}{\sum_{k=1}^{n}\exp(s(\vT_i^{(k)},\vX_k^{(k)})/\tau)}
        \right).
    \end{equation}

    \item \textbf{计算梯度。} 计算随机损失相对于两个参数集的梯度：
    \begin{equation}
        \nabla_{\theta}\hat{\cL}_{\text{CLIP}}^{(k)}(\theta,\plainphi),
        \qquad
        \nabla_{\plainphi}\hat{\cL}_{\text{CLIP}}^{(k)}(\theta,\plainphi).
    \end{equation}

    \item \textbf{更新参数。} 对 \((\theta,\plainphi)\) 应用一步优化算法，得到迭代：
    \begin{equation}
        (\theta^{(k+1)}, \plainphi^{(k+1)})
        \doteq
        \textsc{OptUpdate}^{(k)}\!\left(
        (\theta^{(k)}, \plainphi^{(k)});
        \nabla_{\theta}\hat{\cL}_{\text{CLIP}}^{(k)},
        \nabla_{\plainphi}\hat{\cL}_{\text{CLIP}}^{(k)}
        \right).
    \end{equation}
\end{itemize}


\subsection{评估方法}\label{sub:clip_eval}
由于 CLIP 是与架构无关的，因此可以应用大多数针对视觉和语言编码器的标准评估协议。例如，可以通过在视觉编码器之上训练一个轻量级分类头来执行线性探测，以评估在图像分类任务中学习到的表示的质量。此外，得益于 CLIP 的双塔架构，无需任何额外训练即可执行零样本分类（zero-shot classification）。

回顾在标准的 \(n\) 分类图像分类设置中，我们给定一个图像数据集
\(\{\vX_i\}_{i=1}^{N}\subseteq\cI\)，每张图像都标注有集合
\(\cL \doteq \{\ell_1,\ldots,\ell_n\}\) 中的 \(n\) 个类别标签之一（例如，\(\ell_j \in \{\text{cat},\text{dog},\ldots\}\)）。
我们的目标是利用预训练模型构建一个预测管道 \(P\)，使得
\(P:\cI\to\{1,\ldots,n\}\) 将每张图像 \(\vX_i\) 映射到预测的类别索引
\(\hat{y}_i \doteq P(\vX_i)\)。将真实类别标签写为其索引 \(y_i \in \{1,\ldots,n\}\)，
则 \(P\) 在该数据集上的（top-1）分类准确率计算为
\[
\frac{1}{N}\sum_{i=1}^{N}\boldsymbol{1}\{\hat{y}_i = y_i\}.
\]

在 CLIP 中，我们可以使用文本编码器 \(g_{\plainphi}\) 构建一个文本字典 \(\vD\in\R^{n\times d}\)，其第 \(j\) 行是第 \(j\) 个类别标签的归一化潜在表示。
具体而言，对于每个类别 \(\ell_j \in \cL\)，我们构建一个文本描述（提示，prompt）\(\vT_j\)（例如，
“一张 \(\ell_j\) 的照片”），并计算其归一化嵌入
\(\vz_j^T \doteq \frac{g_{\plainphi}(\vT_j)}{\|g_{\plainphi}(\vT_j)\|_2} \in \R^{d}\)。将这些按行堆叠得到
\[
\vD \doteq 
\mat{(\vz_1^T)^{\top}\\ \vdots \\ (\vz_n^T)^{\top}}
\in \R^{n\times d},
\]

给定图像 \(\vX_i \in \cI\)，我们通过视觉编码器计算其图像表示并将其归一化为
\[
\vz_i^I \doteq \frac{f_{\theta}(\vX_i)}{\|f_{\theta}(\vX_i)\|_2} \in \R^{d}.
\]
然后我们可以计算一个相似度向量
\[
\vd_i \doteq \vD\,\vz_i^I \in \R^{n},
\]
其第 \(j\) 个元素 \(\vd_i[j]\) 是第 \(i\) 张图像和第 \(j\) 个类别标签
在共享潜在空间中的余弦相似度。根据 CLIP 训练目标，鼓励图像与其匹配文本的相似度高于与不匹配文本的相似度；因此，较大的余弦相似度表明
\(\vX_i\) 与类别 \(\ell_j\) 之间的对齐更紧密。因此，我们通过下式预测类别索引：
\[
\hat{y}_i \doteq \arg\max_{j \in \{1,\ldots,n\}} \vd_i[j].
\]
最后，我们可以如上所述计算分类准确率。 

这种零样本决策协议不仅限于分类，还自然适用于\emph{检索}（retrieval）任务，
其中目标是给定查询从大型数据库中返回最相关的项目（例如，检索与文本查询最相关的图像，或检索与图像查询最相关的文本）。例如，在文本到图像的检索中，我们不是从类别名称提示构建字典，而是通过编码和归一化所有数据库图像 \(\{\vX_m\}_{m=1}^{M}\) 来构建\emph{图像字典}，
将它们的嵌入堆叠成
\(\vD\in\R^{M\times d}\)。给定查询文本 \(\vT\)，我们计算其归一化嵌入
\(\vz^T \doteq g_{\plainphi}(\vT)/\|g_{\plainphi}(\vT)\|_2\)，形成相似度得分
\(\vd \doteq \vD\vz^T\in\R^{M}\)，然后根据这些得分对图像进行排序（等价地，返回前 \(K\) 个
索引）。 
% More broadly, this simple embedding-and-dot-product pipeline is arguably the most distinctive and
% practically important feature of CLIP.



\subsection{实验设置与结果}

接下来，我们展示原始 CLIP 论文中的结果，将 CLIP 与监督 ResNet 基线在几个经典分类基准上进行比较。在所有基准测试中，CLIP 均使用 \Cref{sub:clip_eval} 中描述的零样本分类协议进行评估。ResNet 实验设置的完整细节在~\citep{He2016-lc} 中提供；在这里，我们简要介绍评估数据集，并总结与这些比较相关的 CLIP 训练和评估设置。

\paragraph{评估数据。}
\begin{itemize}
    \item \textbf{CIFAR-10~\citep{krizhevsky2009learning}。}
    CIFAR-10 是一个广泛使用的小规模图像分类基准，包含 60{,}000 张大小为 \(32\times 32\) 的 RGB 图像，分为 10 个类别（每个类别 6{,}000 张图像）。
    它通常被划分为 50{,}000 张训练图像和 10{,}000 张测试图像。

    \item \textbf{CIFAR-100~\citep{krizhevsky2009learning}。}
    CIFAR-100 具有与 CIFAR-10 相同的整体大小和图像分辨率（60{,}000 张 RGB \(32\times 32\) 图像，按 50{,}000/10{,}000 划分训练/测试集），但它更具挑战性，因为它包含 100 个细粒度类别（每个类别 600 张图像）。
    这 100 个类别还被额外分组为 20 个超类。

    \item \textbf{PASCAL VOC 2007 (VOC2007)~\citep{everingham2010pascal}。}
    PASCAL VOC 2007 是一个用于对象识别的经典计算机视觉基准，支持图像分类、目标检测和分割等任务。
    它侧重于 20 个对象类别（例如，人、车、狗）。

    \item \textbf{ImageNet-1K~\citep{imagenet_cvpr09}。}
    ImageNet-1K 是 ILSVRC 2012 的一个子集，包含 1{,}000 个对象类别。
    该子集包括 1{,}281{,}167 张训练图像、50{,}000 张验证图像和 100{,}000 张测试图像，并在大规模视觉识别和预训练中发挥了核心作用。
\end{itemize}

\paragraph{模型设置。}
在此处展示的 CLIP 模型中，Vision Transformer 被用作视觉编码器 \(f_{\theta}\)，并实例化为三种 ViT 变体之一：ViT-B/32、ViT-B/16 或 ViT-L/14。这里，\(\text{B}\)（“Base”）和 \(\text{L}\)（“Large”）指定了 Transformer 的宽度/深度（例如，ViT-B 通常使用 12 个 Transformer 块，嵌入维度 \(d=768\) 和 12 个注意力头，而 ViT-L 使用 24 个块，\(d=1024\) 和 16 个头），而后缀 \(/32\)、\(/16\) 和 \(/14\) 表示用于将输入图像分块的像素\emph{图块大小}（patch size）。较小的图块大小（例如 16 或 14）在固定分辨率下每张图像会产生更多的标记，这增加了计算量，但由于更细的空间粒度，通常会提高准确率。

\paragraph{优化设置。}
所有基于 ViT 的 CLIP 模型都使用带有解耦权重衰减的 Adam 优化器（即 AdamW 风格的正则化）训练 32 个 epoch，其中权重衰减应用于除增益和偏置之外的所有参数。学习率使用余弦调度进行衰减。可学习的温度参数 \(\tau\)（用于缩放对比对数几率）被初始化为等效于 \(0.07\)，并被裁剪以防止将对数几率缩放超过 100，这被发现对于训练稳定性是必要的。最后，训练使用非常大的小批量大小 \(32{,}768\)。

\paragraph{结果。}\Cref{tab:clip_vs_resnet} 报告了三种 CLIP ViT 变体（B/32、B/16、L/14）的零样本 top-1 准确率，并将它们与监督 ResNet-50 和 ResNet-101 基线进行了比较。
在所有四个数据集上，CLIP 始终优于监督 ResNet。
进一步扩展 CLIP 主干网络会产生额外的收益，ViT-L/14 在 CIFAR-10 上达到 \(98.0\)，在 CIFAR-100 上达到 \(87.5\)，在 VOC2007 上达到 \(89.6\)，在 ImageNet 上达到 \(83.9\)。
尽管如此，该表应被解释为两个\emph{系统}之间的比较，而不是受控的主干网络消融实验：CLIP 受益于大规模图像-文本预训练，并通过零样本提示进行评估，而 ResNet 基线依赖于监督预训练，并需要标记数据来在每个基准上拟合探测器。
因此，一致的优势主要证明了 CLIP 在零样本协议下可迁移表示的强度，而不是确立 ViT 在纯架构或数据匹配上优于 ResNet。

\begin{table}[t]
\centering
\begin{tabular}{@{}llcccc@{}}
\toprule
主干网络 & 变体 & CIFAR-10 & CIFAR-100 & VOC2007 & ImageNet \\
\midrule
CLIP-ViT & B/32 & 95.1 & 80.5 & 87.7 & 76.1 \\
CLIP-ViT & B/16 & 96.2 & 83.1 & 89.2 & 80.2 \\
CLIP-ViT & L/14 & 98.0 & 87.5 & 89.6 & 83.9 \\
\midrule
ResNet & 50  & 91.8 & 74.5 & 83.8 & 74.3 \\
ResNet & 101 & 93.0 & 77.2 & 84.4 & 75.8 \\
\bottomrule
\end{tabular}
\caption{各种 CLIP 主干网络和 ResNet 基线在下游数据集上的性能。}
\label{tab:clip_vs_resnet}
\end{table}

\subsection{简化扩展：LIFT}\label{sub:follow-up}

尽管有效，但 CLIP 表现出几个局限性。首先，从头开始联合训练文本编码器和图像编码器的计算成本很高，通常需要极大的批量和海量数据集才能达到强大的对齐和下游迁移。其次，以这种方式训练的模型可能在组合理解（compositional understanding）方面遇到困难——捕获文本中的词序、图像中的空间布局以及对象-属性或对象-对象关系——部分原因是检索式对比监督可能会奖励那些降低细粒度组合特征权重的捷径解决方案。LIFT（使用固定文本编码器的语言-图像对齐，Language–Image alignment with a Fixed Text encoder）~\citep{yang2025languageimagealignmentfixedtext} 重新审视了这些管道背后的一个核心假设：最佳对齐需要两个编码器的联合端到端训练。相反，LIFT 利用了现代大型语言模型（LLM）已经产生高度信息化的文本嵌入这一观察结果。具体而言，LIFT 固定了一个强大的预训练文本编码器（例如，源自 LLM 或在 LLM 上微调的编码器），离线计算文本嵌入，并仅训练图像编码器以对齐这些固定目标。

\paragraph{架构。}
LIFT 保留了与 CLIP 相同的双塔公式，即视觉编码器和文本编码器映射到共享的潜在空间，但关键区别在于文本塔是\emph{固定的}，并且可以由任何强大的基于 LLM 的文本编码器 \(g_{\plainphi}:\cT\to\R^{d}\) 实例化以生成文本嵌入，其中 \(\plainphi\) 在训练期间不更新。在我们的实现中，我们采用 NV-Embed-V2 文本编码器，其嵌入维度为 \(d=4096\)。由于文本塔 \(g_{\plainphi}\) 是固定的，我们可以离线预计算所有文本嵌入 \(\{g_{\plainphi}(\vT)\}\)，并在训练期间仅在线计算梯度并更新图像编码器 \(f_{\theta}\)。图像塔 \(f_{\theta}:\cI\to\R^{d}\) 遵循与 \Cref{sub:clip_vision_tower} 中相同的结构，投影头输出维度设置为 \(d\)，以匹配固定文本嵌入空间的维度。除了固定 \(g_{\plainphi}\) 和匹配投影维度外，其余组件（包括对比损失和优化过程）遵循与 CLIP 相同的设计原则。

\paragraph{评估方法。} 
组合理解是 CLIP 的一个已知局限性。我们在七个 SugarCrepe~\citep{hsieh2023sugarcrepe} 任务上评估 LIFT 和 CLIP，其中每张图像 \(\vX\) 都配有一个\emph{正样本}（正确）标题 \(\vT_{\text{pos}}\) 和一个\emph{负样本}标题 \(\vT_{\text{neg}}\)，后者是通过添加、替换或交换 \(\vT_{\text{pos}}\) 中的对象、属性或关系构建的。一些示例见图 \ref{fig:SugarCrepe-examples}。模型被要求通过比较标题-图像的余弦相似度来识别正确的标题。形式上，对于每个样本 \(i\in\{1,\dots,N\}\)，我们计算归一化嵌入：
\[
\vz_i^{I} \doteq \frac{f_{\theta}(\vX_i)}{\|f_{\theta}(\vX_i)\|_2},\qquad
\vz_{i,\text{pos}}^{T} \doteq \frac{g_{\plainphi}(\vT_{i,\text{pos}})}{\|g_{\plainphi}(\vT_{i,\text{pos}})\|_2},\qquad
\vz_{i,\text{neg}}^{T} \doteq \frac{g_{\plainphi}(\vT_{i,\text{neg}})}{\|g_{\plainphi}(\vT_{i,\text{neg}})\|_2}.
\]
然后我们通过余弦相似度计算相似度得分：
\[
s_{i,\text{pos}} \doteq \langle \vz_i^{I}, \vz_{i,\text{pos}}^{T} \rangle,
\qquad
s_{i,\text{neg}} \doteq \langle \vz_i^{I}, \vz_{i,\text{neg}}^{T} \rangle.
\]
如果 \(s_{i,\text{pos}} > s_{i,\text{neg}}\)，则认为模型在样本 \(i\) 上是正确的，整体组合准确率为
\[
\frac{1}{N}\sum_{i=1}^{N}\boldsymbol{1}\{s_{i,\text{pos}} > s_{i,\text{neg}}\}.
\]

\begin{figure}
\centering
\includegraphics[width=\textwidth]{\toplevelprefix/chapters/chapter8/figs/negative_caption_demo.pdf}
\caption{来自两个 SugarCrepe 任务的原始标题（顶部）及其负样本对应物（底部）：\texttt{replace relation}（左）和 \texttt{swap attribute}（右）。}
\label{fig:SugarCrepe-examples}
\end{figure}

\paragraph{实验结果。}
如 \Cref{tab:Compositionality} 所示，当在 DataComp-1B 上训练时，LIFT 在所有七个任务上均优于 CLIP，平均准确率提高了 \textbf{6.8\%}。LIFT 在 \texttt{add attribute}、\texttt{replace attribute} 和 \texttt{replace relation} 任务上取得了显著的收益。这些改进有力地证明了基于 LLM 的文本编码器 \(g_{\plainphi}\) 捕获了更具信息量的文本表示，从而能够更准确地建模对象-属性关联和对象-对象关系。另一方面，我们可以看到，与其他 SugarCrepe 任务相比，LIFT 在 \texttt{swap object} 和 \texttt{swap attribute} 上的准确率相对较低。我们将此局限性归因于对比学习目标，该目标主要侧重于对齐低阶统计量。应对这一挑战需要探索更精细的语言-图像对齐信息论度量，这是未来工作的一个关键方向。

\begin{table}[t]
\centering
\resizebox{\textwidth}{!}{%
\begin{tabular}{@{}lllccccccc@{}}
\toprule
& & & \multicolumn{2}{c}{添加} & \multicolumn{3}{c}{替换} & \multicolumn{2}{c}{交换} \\
方法 & 数据集 & 见过样本数 & 对象 & 属性 & 对象 & 属性 & 关系 & 对象 & 属性 \\
\midrule
OpenCLIP & DataComp & 1.28B & 82.3 & 73.7 & 91.7 & 79.4 & 61.2 & 59.6 & 56.9 \\
LIFT     & DataComp & 1.28B & \textbf{89.0} & \textbf{86.1} & \textbf{93.2} & \textbf{86.0} & \textbf{70.6} & \textbf{64.1} & \textbf{63.4} \\
\bottomrule
\end{tabular}%
}
\caption{LIFT 和 CLIP 在七个 SugarCrepe 任务上的性能。我们使用开源的 CLIP 实现 OpenCLIP~\citep{ilharco_gabriel_2021_5143773}，并在开源文本-图像对数据集 DataComp~\citep{gadre2023datacompsearchgenerationmultimodal} 上训练 OpenCLIP 和 LIFT。我们报告了每个任务的准确率，最佳结果以粗体显示。}
\label{tab:Compositionality}
\end{table}


\section{通过分类的监督图像表示}\label{sec:image_classification}

在前面的章节中，我们展示了如何基于本书介绍的压缩原理来简化一些过于复杂、通常是启发式设计的学习目标。然而，与一些最流行的学习任务相关的目标已经相当简单。在这些情况下，很难进一步简化目标。因此，在本节及以后的章节中，我们将重点关注以原则性方式修改各种任务的\textit{深度网络架构}。

让我们首先从机器学习中可以说是最经典的任务开始：\textit{图像分类}，它通常被用作评估模式识别算法或深度网络架构的标准任务。根据我们在 \Cref{ch:deep-networks} 中对白盒架构的讨论，我们只需要一个具有语义意义的任务就可以使用白盒架构学习到良好的表示。我们将在本节中验证这一想法。

首先，数据集在很大程度上与 \Cref{sub:contrastive_learning_data} 保持一致。训练和测试数据均由带标签的图像组成，即图像-标签对 \((\vX, \vy) \in \R^{C \times H \times W} \times \{0, 1\}^{N_{\cls}}\)。我们仍然对每个新批次中的每个样本应用各种数据增强（例如，翻转、高斯模糊、过度曝光等）。

\subsection{任务与目标} \label{sub:image_classification_objective}

与之前不同，我们的任务不仅是学习数据的良好表示，还要同时构建一个分类器。形式上，我们有带标签的数据对 \((\vX, \vy)\)，其中 \(\vy \in \{0, 1\}^{N_{\cls}}\) 是一个独热向量，表示 \(\vX\) 的类别归属。我们考虑输入数据 \(\vX\) 的确定性\textit{中心裁剪视图} \(v_{\cc}\)（参见 \Cref{sub:contrastive_learning_objective}）。我们希望联合训练特征映射 \((f_{\theta}, f_{\theta}^{\ext})\) 和\textit{分类头} \(h_{\theta}\)，定义如下：
\begin{equation}
h_{\theta}(\vz) \doteq \softmax(\vW^{\head}\vz + \vb^{\head}), \qquad
\forall \vz \in \R^{d}
\end{equation}
其中 \((\vW^{\head}, \vb^{\head}) \in \R^{N_{\cls} \times d} \times
\R^{N_{\cls}}\) 是参数集 \(\theta\) 中的可训练参数，使得映射 \(\vX_{\cc} \mapsto
\vp_{\theta}(\vX_{\cc}) \doteq h_{\theta}(\vz_{\theta}(\vX_{\cc}))\)
预测输入 \(\vX\) 的视图 \(\vX_{\cc} = v_{\cc}(\vX)\) 的平滑标签。学习问题试图最小化通过交叉熵测量的 \(\vp_{\theta}\) 和 \(\vy\) 之间的距离：
\begin{equation}\label{eq:classification_ce_loss}
\min_{\theta}\bc{\cL_{\CE}(\theta) \doteq \Ex[\CE(\vy,
\vp_{\theta}(\vX_{\cc}))]}.
\end{equation}

\subsection{CRATE 架构}\label{sub:image_classification_architecture}

我们使用的架构是 CRATE 架构，在 \Cref{ch:deep-networks} 中有详细描述。整体设置与 \Cref{sub:contrastive_learning_architecture} 中的常规 Transformer 相似，但有一些变化。虽然嵌入步骤与 \Cref{sub:contrastive_learning_architecture} 中的 DINO 和 SimDINO 相同，特征提取步骤与 \Cref{sub:contrastive_learning_architecture} 中的 SimDINO 相同（因为它只提取对应于类别标记的特征），并且分类头在 \Cref{sub:image_classification_objective} 中进行了描述，但主干架构是不同的。每一层的形式如下：
\begin{align}\label{eq:CARTE updates}
\vZ_{\theta}^{\ell + 1/2}(\vX)
&= \vZ_{\theta}^{\ell}(\vX) + \MSSA_{\theta}^{\ell}(\LN_{\theta}^{1,
\ell}(\vZ_{\theta}^{\ell}(\vX))), \\
\vZ_{\theta}^{\ell + 1}(\vX)
&= \ISTA_{\theta}^{\ell}(\LN_{\theta}^{2, \ell}(\vZ_{\theta}^{\ell +
1/2}(\vX))),
\end{align}
其中 \(\MSSA_{\theta}^{\ell}\) 和 \(\ISTA_{\theta}^{\ell}\)
块如 \Cref{ch:deep-networks} 所述，即：
\begin{itemize}
\item \(\MSSA\) 算子是多头子空间自注意力（multi-head-subspace-self-attention），定义如下：
    \begin{equation}
        \MSSA_{\theta}^{\ell}(\vZ) \doteq
        \vU_{\out}^{\ell}\mat{\SA([\vU^{1, \ell}]^{\top}\vZ, [\vU^{1,
            \ell}]^{\top}\vZ, [\vU^{1, \ell}]^{\top}\vZ)\\ \vdots
            \\ \SA([\vU^{K, \ell}]^{\top}\vZ, [\vU^{K, \ell}]^{\top}\vZ,
        [\vU^{1, \ell}]^{\top}\vZ)} + \vb_{\out}^{\ell}\vone_{n}^{\top}
    \end{equation}
    其中 \(\vU^{k, \ell} \in \R^{d \times p}\)、\(\vU_{\out}^{\ell}
    \in \R^{d \times Kp}\) 和 \(\vb_{\out}^{\ell} \in \R^{d}\) 是属于参数集 \(\theta\) 的可训练参数，并且（回想一下）自注意力算子 \(\SA\) 在 \eqref{eq:self_attention} 中定义。
\item \(\ISTA\) 算子是迭代收缩阈值算法（iterative-shrinkage-thresholding-algorithm）算子，定义如下：
    \begin{equation}
        \ISTA_{\theta}^{\ell}(\vZ) \doteq \ReLU(\vZ - \beta
            (\vD^{\ell})^{\top}(\vD^{\ell}\vZ - \vZ) + \beta\lambda
        \vone_{d}\vone_{n}^{\top}),
    \end{equation}
    之所以这样命名，是因为映射 \(\vX \mapsto \ReLU(\vX - \beta
            \vD^{\top}(\vD\vX - \vZ) + \beta  \lambda
    \vone_{d}\vone_{n}^{\top})\) 是成熟的 ISTA 算法的一步，用于寻找 \(\vZ\) 相对于完备字典 \(\vD\) 的逐元素非负稀疏表示（参见 \Cref{sec:dictionary_learning}）。
\end{itemize}

我们将此架构称为 CRATE，主干网络的一层如 \Cref{fig:crate_backbone} 所示。CRATE 模型除了具有可解释性外，通常还具有高性能和参数高效的特点。

\subsection{优化} \label{sub:image_classification_optimization}

我们使用简单的端到端随机优化过程来训练分类器，其中我们对数据和视图进行子采样，计算这些样本上的平均损失及其梯度，并使用优化算法来更改参数。在每个时间步 \(k\)，我们：
\begin{itemize}
\item 子采样 \(B\) 个不同的带标签样本 \(\{(\vX_{b}^{(k)},
    \vy_{b}^{(k)})\}_{b = 1}^{B} \subseteq \cI \times \{0, 1\}^{N_{\cls}}\)。
\item 对于每个带标签样本 \((\vX_{b}^{(k)}, \vy_{b}^{(k)})\)，计算中心裁剪视图 \(v_{b, \cc}^{(k)}\) 并将其应用于 \(\vX_{b}^{(k)}\) 以得到 \(\vX_{b, \cc}^{(k)} \doteq v_{b,
    \cc}^{(k)}(\vX_{b}^{(k)})\)。
\item 计算预测值 \(\vp_{\theta}(\vX_{b, \cc}^{(k)})
        \doteq (h_{\theta} \circ f_{\theta}^{\ext} \circ
    f_{\theta})(\vX_{b, \cc}^{(k)})\)。
\item 形成代理随机损失
    \begin{equation}
        \hat{\cL}_{\CE}^{(k)}(\theta) \doteq \frac{1}{B}\sum_{b =
        1}^{B}\CE(\vy_{b}^{(k)}, \vp_{\theta}(\vX_{b, \cc}^{(k)})).
    \end{equation}
\item 对 \(\theta\) 计算一步优化算法，得到以下迭代：
    \begin{equation}
        \theta^{(k + 1)} \doteq
        \textsc{OptUpdate}^{(k)}(\theta^{(k)};
        \nabla_{\theta}\hat{\cL}_{\CE}^{(k)}).
    \end{equation}
\end{itemize}

\subsection{评估方法} \label{sub:image_classification_evals}

我们使用与 \Cref{sub:contrastive_learning_evals} 相同的评估过程。总而言之，对于所有评估（以及训练），我们使用中心裁剪视图 \(v_{\cc}\)，它重塑输入图像并进行大小为 \((C, S_{\cc}, S_{\cc})\) 的大中心裁剪，其中 \(C\) 是输入图像的通道数。给定此视图的输出，我们随后可以进行线性探测、注意力图可视化以及检测/分割基准测试。

\subsection{实验设置与结果}\label{sub:image_classification_experiments}

由于 CRATE 直接基于 Transformer，为了公平比较，我们将 \cite{dosovitskiy2020image,touvron2020training} 给出的 ViT 最佳设置与应用于 CRATE 的相同设置进行比较。

\paragraph{模型架构。} 在评估和训练中，中心裁剪都会调整整个图像的大小，使较短的边长为 \(256\)（即 \(S_{\rsz} = 256\)），然后再进行大小为 \(224 \times 224\)（即 \(S_{\cc} = 224\)）的中心裁剪。我们采用图块大小 \(16\)（即 \(P_{H} = P_{W} = 16\)）。我们使用 ViT \cite{dosovitskiy2020image} 架构的 tiny、small、base 和 large 模型作为嵌入和主干网络，将 MHSA 和 MLP 组件分别替换为 MSSA 和 ISTA，在 MSSA 的情况下使用相同数量的头和头维度，从而大幅减少了训练参数的数量。对于 CRATE，我们设置 \((\beta, \lambda) = (1, 0.1)\)。

\paragraph{数据集与优化。} 对于预训练，我们使用 ImageNet-1K 数据集。我们使用 LION 优化器 \citep{chen2024symbolic} 来预训练我们的 ViT 复现以及 CRATE。我们将基础学习率设置为 \(2.4 \times
10^{-4}\)，权重衰减设置为 \(0.5\)，批量大小设置为 \(B = 2048\)。我们的学习率调度在前 \(5\) 个 epoch 内将学习率线性增加到基础学习率，并在接下来的 \(145\) 个 epoch 内使用余弦调度将其降低到 \(0\)（所有模型均训练 \(150\) 个 epoch）。对于预训练，我们对图像数据应用常规的数据增强方案（翻转、高斯模糊、过度曝光等），并向标签添加小噪声（这称为\textit{标签平滑}（label smoothing）\citep{muller2019does}）。

对于线性探测，我们使用几个评估数据集，如 CIFAR10、Oxford-Flowers 和 Oxford-IIT-Pets。我们使用 AdamW 优化器来训练线性探测器，学习率为 \(5 \times
10^{-5}\)，权重衰减为 \(0.01\)，批量大小为 \(B = 256\)。我们还对图像数据应用了上述数据增强。

\begin{table}
\caption{\small \textbf{CRATE 和 ViT 的线性探测分类准确率}：当主干网络在 ImageNet-1K 上进行分类预训练时，不同模型大小在各种数据集上的线性探测分类准确率。我们观察到，在给定相同模型配置的情况下，CRATE 具有相当的分类性能，且设计更简单、更有原则性、参数更高效。}
\centering
\resizebox{1\columnwidth}{!}
{
    \begin{tabular}{@{}lcccc|cc@{}}
        \toprule
        \textbf{模型} & CRATE-T  &  CRATE-S & CRATE-B & CRATE-L & {
        \color{gray} ViT-T} &  { \color{gray}ViT-S } \\
        \midrule
        \midrule
        参数量 & 6.09M & 13.12M & 22.80M & 77.64M & {
        \color{gray} 5.72M} & { \color{gray} 22.05M} \\
        \midrule
        ImageNet-1K & 66.7 & 69.2 & 70.8 & 71.3 & { \color{gray}
        71.5} & { \color{gray} 72.4} \\
        ImageNet-1K ReaL & 74.0 & 76.0 & 76.5 & 77.4 & { \color{gray}
        78.3 } & { \color{gray} 78.4} \\
        CIFAR10 & 95.5 & 96.0 & 96.8 & 97.2 & { \color{gray} 96.6} &
        { \color{gray} 97.2} \\
        CIFAR100 & 78.9 & 81.0 & 82.7 & 83.6 & { \color{gray} 81.8} &
        { \color{gray} 83.2}\\
        Oxford Flowers-102 & 84.6 & 87.1 & 88.7 & 88.3 & {
        \color{gray} 85.1} & { \color{gray} 88.5}\\
        Oxford-IIIT-Pets & 81.4 & 84.9 & 85.3 & 87.4 & { \color{gray}
        88.5} & { \color{gray} 88.6} \\
        \bottomrule
    \end{tabular}
}
\label{tab:crate_classification_linear_probing}
\end{table}

\begin{figure}[t]
\centering
\includegraphics[width=\textwidth]{\toplevelprefix/chapters/chapter8/figs/crate_semantic_heads.pdf}
\caption{\small\textbf{CRATE 中可解释的显著性图}（图块大小为 \(8\)）。当给定具有相似属性（可能但不一定来自同一类别）的图像时，最后一层中不同注意力头对应的显著性图各自突出显示了一个特定属性。可以观察到，平均显著性图（未包含）随后突出了图像中的所有相关对象，表明它使用了输入图像的所有细粒度细节进行分类。据作者所知，这是\textit{第一个}做到这一点的机器学习系统，更不用说在没有对任何分割数据进行训练的情况下自动完成的了。}
\label{fig:crate_semantic_heads}
\end{figure}

\begin{table}
\caption{\small \textbf{在 COCO {val2017}~\citep{lin2014microsoft} 上通过 MaskCut 进行的目标检测和细粒度分割}。这里所有模型都使用图块大小 \(8\) 而不是 \(16\) 进行训练。当两者都使用监督分类进行训练时，CRATE 在检测和分割指标上的表现最终优于 ViT。}
\centering
\begin{tabular}{@{}lcccccccc@{}}
    \toprule
    &  \multicolumn{3}{c}{检测 (\(\uparrow\))} &
    \multicolumn{3}{c}{分割 (\(\uparrow\))} \\
    模型 & AP$_{50}$ & AP$_{75}$ & AP & AP$_{50}$ & AP$_{75}$ & AP  \\
    \midrule
    CRATE-S/8 & \textbf{2.9} & \textbf{1.0} & 1.1 & 1.8 & \textbf{0.7} & 0.8 \\
    CRATE-B/8 & \textbf{2.9} & \textbf{1.0} & \textbf{1.3} &
    \textbf{2.2} & \textbf{0.7} & \textbf{1.0} \\
    ViT-S/8 & 0.1& 0.1 & 0.0 & 0.0 & 0.0 & 0.0 \\
    ViT-B/8 & 0.8 & 0.2 & 0.4 & 0.7 & 0.5 & 0.4 \\
    \bottomrule
\end{tabular}
\label{tab:crate_detection_segmentation}
\end{table}

\paragraph{实验结果。}

\Cref{tab:crate_classification_linear_probing} 表明，在相似的参数量下，CRATE 模型与流行的 Vision Transformer (ViT) 架构相比实现了持平或改进，至少在特征相对于不同类别的线性可分性方面是如此。在注意力图保真度方面，\Cref{fig:crate_semantic_heads} 展示了一个真正非凡的结果：无需在任何分割或目标检测数据上进行训练，显著性图\textit{不仅}有效地捕获了输入图像的所有相关部分，而且显著性图还\textit{自组织}地各自对应于一组离散的概念，甚至跨越了样本和类别！据作者所知，这是第一个做到这一点的系统，并且除了图像分类数据之外，它可以在不使用任何额外数据的情况下做到这一点。\Cref{tab:crate_detection_segmentation} 定量地证实了这些定性见解，显示出比在相同监督分类设置下训练的 ViT 有显著改进。

\section{通过掩码自编码的图像表示}\label{sec:image_completion}

在本节中，我们将展示如何将 CRATE 架构应用于\textit{图像补全}（image completion）问题，也称为\textit{掩码自编码}（masked autoencoding，MAE），它可以被视为 \Cref{ch:classic-models} 中讨论的低秩矩阵补全问题的推广。掩码自编码自 \cite{he2022masked} 将其引入深度学习背景以来，一直是一种主要且简单的自监督表征学习方法，其旨在赋予 \(\vZ_{\theta}\) 中的每个图块特征以聚合信息及其邻居的信息，使得图块特征和聚合特征都成为整个样本的丰富信息源。

我们用于此任务的数据与 \Cref{sub:contrastive_learning_data} 中讨论的图像数据集相同。像往常一样，我们仍然对每个新批次中的每个样本应用数据增强。

\subsection{MAE 任务与目标}\label{sub:image_completion_objective}

顾名思义，掩码自编码涉及一个视图 \(v_{m}\)，给定输入，它执行随机调整大小的裁剪（参见 \Cref{sub:contrastive_learning_objective}），将输入图像转换为大小为 \((C, S_{\mask}, S_{\mask})\) 的正方形图像，然后\textit{掩码}（即设置为零）输入中固定百分比 \(p_{\mask} \in [0, 1]\) 的像素。出于效率原因\footnote{\cite{he2022masked} 对 MAE 的原始实现是嵌入整个图像，\textit{移除}将被掩码的标记，将生成的标记集输入编码器，在掩码位置加回学习到的占位符标记并加回适当的位置编码，然后将生成的标记集输入解码器以获得自编码预测。这更有效率，因为编码器需要处理的标记更少，但在概念上与正文中讨论的方法相同，并且生成的模型在掩码自编码任务和下游评估中的性能非常相似。}，掩码是按图块进行的，即在嵌入整个图像后，将 \(p_{\mask}\) 百分比的\textit{图块}设置为零。MAE 的目标是训练一个编码器 \(f_{\theta} \colon \cI \to
(\R^{d})^{*}\) 和一个解码器 \(g_{\eta} \colon (\R^{d})^{*} \to \cI\)，它们可以从掩码中重建输入，即，记 \(\hat{\vX}_{\theta, \eta} \doteq g_{\eta} \circ f_{\theta}\)，我们有
\begin{equation}
\min_{\theta, \eta}\bc{\cL_{\mathrm{MAE}}(\theta, \eta) \doteq
\Ex\norm{\hat{\vX}_{\theta, \eta}(\vX_{m}) - \vX}_{F}^{2}}
\end{equation}
本质上，这意味着视图 \(\vX_{m} \doteq v_{m}(\vX)\) 的特征 \(\vZ_{\theta}(\vX_{m})\) 必须包含关于\textit{被掩码图块}以及现有图块的信息。从 \Cref{ch:deep-networks} 中基于压缩的白盒模型的角度来看，如果白盒自编码器 \((f_{\theta}, g_{\eta})\) 在此任务上取得成功，则意味着学习到的子空间和字典对数据执行了\textit{冗余}编码，从而可以从编码的数据的其他部分重建数据的缺失部分。这意味着关于每个图块的信息存储在其他图块中。因此，每个图块特征应同时包含关于该图块的信息和关于整个图像统计特征的信息。因此，我们再次期望表示应包含局部和全局的语义相关信息，因此具有相似局部和全局信息的不同图块的表示应该是相关的（即，在同一子空间上或由字典一起编码）。

\subsection{基于 CRATE 的 MAE 架构}\label{sub:image_completion_architecture}

\begin{figure}
\centering
\includegraphics[width=\textwidth]{\toplevelprefix/chapters/chapter8/figs/crate_ae_backbone.pdf}
\caption{\small\textbf{CRATE 自编码器主干网络中编码器和解码器的一层。} 编码器和解码器层都将其输入通过多头子空间自注意力和字典学习或字典编码步骤。请注意，编码器和解码器层是对称设计的；每个解码器层的概念目标是反转一个编码器层，因此这种对称性很大程度上是设计使然（例如，参见 \Cref{ch:consistent-representations}）。}
\end{figure}

我们使用 CRATE 编码器和解码器，如 \Cref{fig:transformer_backbone} 所示，尽管当然也可以使用常规的 Transformer 编码器和解码器。详情如下。

\paragraph{编码器。} 编码器与 \Cref{sub:image_classification_architecture} 中的 CRATE 编码器相同，需要注意的是
there is no feature extractor \(f_{\theta}^{\ext}\). However, both
the embedding \(f_{\theta}^{\emb}\) and the backbone
\(f_{\theta}^{\backbone}\) are the same.

\paragraph{解码器主干。} 解码器主干是 \Cref{ch:consistent-representations} 中描述的 CRATE 解码器。为了完整起见，我们现在对其进行描述。给定特征序列 \(\vZ_{\theta}(\vX) \doteq
f_{\theta}(\vX) \in (\R^{d})^{*}\)，我们可以使用解码器主干 \(g_{\eta}^{\backbone}\) 对其进行如下处理。函数 \(g_{\eta}^{\backbone}\) 由 \(L\) 层 \(g_{\eta}^{\ell}\) 组成，即，
\begin{equation}
g_{\eta}^{\backbone} = g_{\eta}^{L} \circ \cdots \circ g_{\eta}^{1}.
\end{equation}
层 \(g_{\eta}^{\ell}\) 具有以下实现。首先，定义 \(\tilde{\vZ}_{\theta, \eta}^{1}(\vX) \doteq
\vZ_{\theta}(\vX)\)。然后，我们得到
\begin{align}
\tilde{\vZ}_{\theta, \eta}^{\ell + 1/2}(\vX)
&= [\tilde{\vD}^{\ell}]^{\top}\LN_{\eta}^{1,
\ell}(\tilde{\vZ}_{\theta, \eta}^{\ell}(\vX)) \\
\tilde{\vZ}_{\theta, \eta}^{\ell + 1}(\vX)
&= \tilde{\vZ}_{\theta, \eta}^{\ell + 1/2}(\vX) -
\MSSA_{\eta}^{\ell}(\LN_{\eta}^{2, \ell}(\tilde{\vZ}_{\theta,
\eta}^{\ell + 1/2}))
\end{align}
并且 \(g_{\eta}^{\ell}\) 的定义使得 \(g_{\eta}^{\ell}(\tilde{\vZ}_{\theta, \eta}^{\ell}) \doteq
\tilde{\vZ}_{\theta, \eta}^{\ell + 1}(\vX)\)。这里的相关概念是，\(g_{\eta}^{\ell}\) 应该学习 \(f_{\theta}^{L + 1 - \ell}\) 的近似逆，分别作为正向和反向时间扩散过程的离散化。特别地，\(\tilde{\vD}^{\ell}\) 应该近似 \(\vD^{L + 1 -
\ell}\)，类似地，\(\MSSA_{\eta}^{\ell}\) 的参数应该与 \(\MSSA_{\theta}^{L + 1 - \ell}\) 的参数相似。输出为 \(\tilde{\vZ}_{\theta, \eta} \doteq
\tilde{\vZ}_{\theta, \eta}^{L + 1}\)。

\paragraph{逆嵌入模块（Un-embedding Module）。} 为了将 \(\tilde{\vZ}_{\theta, \eta}(\vX)\) 转换回对 \(\vX\) 的估计，我们需要使用逆嵌入模块 \(g_{\eta}^{\unemb}\) 来撤销嵌入模块 \(f_{\theta}^{\emb}\) 的影响。因此，回顾 \eqref{eq:definition_of_embedding_module} 中嵌入模块的函数形式，即，
\begin{equation}
f_{\theta}^{\emb}(\vX) \doteq \mat{\vz_{\cls}^{1},
\vW^{\emb}f^{\patch}(\vX) + \vE^{\pos}}
\end{equation}
这意味着我们的逆操作 \(g_{\eta}^{\unemb}\) 如下所示：
\begin{equation}
g_{\eta}^{\unemb}(\tilde{\vZ}) \doteq
g_{\eta}^{\unemb}(\mat{\tilde{\vz}^{1}, \dots, \tilde{\vz}^{n}}) =
g^{\unpatch}(\vW^{\unemb}([\tilde{\vz}^{2}, \dots, \tilde{\vz}^{n}] -
\tilde{\vE}^{\pos})),
\end{equation}
其中 \(g^{\unpatch}\) 执行与 \(f^{\patch}\) 的展开和展平操作相反的逆操作。\footnote{同样，“逆位置编码” \(\tilde{\vE}^{\pos}\) 是针对大输入学习的，对于较小的输入可以进行插值。甚至可以直接将 \(\tilde{\vE}^{\pos}\) 设置为等于位置编码 \(\vE^{\pos}\)，并在编码器和解码器中为每个输入使用相同的插值位置编码。}

该架构是一个白盒自编码器 \((f_{\theta},
g_{\eta})\)，其中（回想一下）\(f_{\theta} = f_{\theta}^{\backbone}
\circ f_{\theta}^{\emb}\) 且 \(g_{\eta} = g_{\eta}^{\unemb} \circ
g_{\eta}^{\backbone}\)。特别地，我们可以使用它来计算掩码视图的估计值 \(\hat{\vX}_{\theta, \eta}(\vX_{m}) =
(g_{\eta} \circ f_{\eta})(\vX_{m})\)，该估计值应近似等于 \(\vX\) 本身。

\subsection{优化}\label{sub:image_completion_optimization}

与 \Cref{sub:image_classification_optimization} 中一样，我们使用一个简单的优化设置：我们对图像和掩码进行采样，计算这些样本上的损失以及该损失的梯度，并使用通用的优化算法和上述梯度来更新参数。对于每个时间步 \(k\)，我们：
\begin{itemize}
\item 子采样 \(B\) 个不同的样本 \(\{\vX_{b}^{(k)}\}_{b =
    1}^{B} \subseteq \cI\)。
\item 对于每个样本 \(\vX_{b}^{(k)}\)，计算一个不同的随机调整大小的裁剪和掩码 \(v_{b, m}^{(k)}\)，并将其应用于 \(\vX_{b}^{(k)}\) 以得到 \(\vX_{b, m}^{(k)} \doteq v_{b,
    m}^{t}(\X_{b}^{(k)})\)。
\item 计算估计的自编码 \(\hat{\vX}_{\theta,
        \eta}(\vX_{b, r}^{(k)}) \doteq (g_{\eta} \circ
    f_{\theta})(\vX_{b, r}^{(k)})\)。
\item 构造代理随机损失
    \begin{equation}
        \hat{\cL}_{\mathrm{MAE}}^{(k)}(\theta, \eta) \doteq
        \frac{1}{B}\sum_{b = 1}^{B}\norm{\hat{\vX}_{\theta,
        \eta}(\vX_{b, r}^{(k)}) - \vX_{b}^{(k)}}_{F}^{2}.
    \end{equation}
\item 对 \((\theta,
    \eta)\) 计算优化算法的一个步骤，得到以下迭代：
    \begin{equation}
        (\theta^{(k + 1)}, \eta^{(k + 1)}) \doteq
        \textsc{OptUpdate}^{(k)}(\theta^{(k)}, \eta^{(k)};
        \nabla_{(\theta, \eta)}\hat{\cL}_{\mathrm{MAE}}^{(k)}).
    \end{equation}
\end{itemize}

\subsection{评估} \label{sub:image_completion_optimization_1}

这是我们在本章中讨论的第一个自编码器网络。我们使用与 \Cref{sub:contrastive_learning_evals,sub:image_classification_evals} 中相同的中心裁剪视图 \(v_{\cc}\)，将最终图像调整为边长为 \(S_{\cc} =
S_{\mask}\) 像素的正方形，以匹配训练期间看到的输入图像的形状。

除了评估掩码自编码损失本身之外，还可以直接评估数据 \(\vX\) 的视图 \(\vX_{\cc} \doteq v_{\cc}(\vX)\) 的特征 \(\vZ_{\theta}(\vX_{\cc})\)。对于注意力图保真度评估，获得 \(\vZ_{\theta}(\vX_{\cc})\) 就足够了，但对于线性探测（linear probing），我们需要从 \(\vZ_{\theta}\) 中提取一个总结性或聚合的特征。为此，我们可以使用一个（无参数的）特征提取映射，该映射返回对应于类别标记（class token）的特征，即，
\begin{equation}
f_{\theta}^{\ext}(\vZ) \doteq f_{\theta}^{\ext}([\vz^{1}, \dots,
\vz^{n}]) = \vz^{1},
\end{equation}
如（例如）\Cref{sub:image_classification_objective,sub:image_classification_architecture} 中所示。借此，我们有了一种获取聚合特征 \(\vz_{\theta}(\vX_{\cc}) \doteq (f_{\theta}^{\ext} \circ
f_{\theta})(\vX_{\cc})\) 的方法，此时我们可以执行线性探测、分割评估等。

\subsection{实验}\label{sub:image_completion_experiments}

由于 CRATE-MAE 直接基于 ViT-MAE，我们将 \citep{he2022masked} 给出的 ViT-MAE 最佳设置与应用于 CRATE-MAE 的相同设置进行比较，以进行公平的对比。

\paragraph{模型架构。} 在训练期间，掩码裁剪 \(v_{m}\) 调整整个图像的大小，使得较短的边尺寸为 \(256\)（即 \(S_{\rsz} = 256\)），然后进行 \(224 \times 224\) 大小的随机裁剪（即 \(S_{\mathrm{mask}} = 224\)），并掩蔽 \(p_{\mathrm{mask}} = \frac{3}{4}\) 的图块（patches）。我们采用图块大小 \(16\)（即 \(P_{H} = P_{W} = 16\)）。我们使用 ViT-MAE 架构的 small 和 base 变体作为编码器和解码器的嵌入和主干，分别将 MHSA 和 MLP 组件替换为 MSSA、ISTA 和线性层。在 MSSA 的情况下，我们使用相同数量的注意头和头部维度。然而，原始的 ViT-MAE 使用的编码器占据了几乎所有的总层数，而解码器仅使用少数几层；正如 \Cref{ch:consistent-representations} 中的概念和理论框架所建议的那样，我们将总层数（从 ViT-MAE 到 CRATE-MAE 保持不变）的一半分配给我们的编码器和解码器。对于 CRATE-MAE，我们设置 \((\beta, \lambda) = (1, 0.1)\)。

\paragraph{数据集与优化。} 对于预训练，我们使用 ImageNet-1K 数据集。我们使用 AdamW 优化器来预训练我们的 ViT-MAE 复现模型以及 CRATE-MAE。我们将基础学习率设置为 \(3 \times 10^{-5}\)，权重衰减设置为 \(0.1\)，批量大小设置为 \(B = 4096\)。我们的学习率调度在最初的 \(40\) 个 epoch 内将学习率线性增加到基础学习率，并在接下来的 \(760\) 个 epoch 内使用余弦调度将其降低到 \(0\)（所有模型各训练 \(800\) 个 epoch）。对于预训练，我们对图像数据应用常规的数据增强机制（翻转、高斯模糊、曝光等）。

对于线性探测，我们使用几个评估数据集，如 CIFAR10、CIFAR100、Oxford-Flowers 和 Oxford-IIT-Pets。对于线性探测，我们预先计算目标数据集中所有样本的特征，并应用快速线性回归求解器，例如来自 Scikit-Learn 等标准包的求解器。

\paragraph{实验结果。} \Cref{tab:crate_mae_linear_probing} 表明，粗略地说，CRATE-MAE 模型在相似的参数数量下实现了与流行的 ViT-MAE 架构相当的性能，并且特征学习性能（通过下游分类任务的性能来衡量）随着规模的增加而提高。同时，\Cref{fig:crate_mae_semantic_heads} 表明，编码器显著图（因此也包括编码器学习到的细粒度特征）确实隔离并突出了输入图像的关键部分。

\begin{table}
\caption{\small\textbf{CRATE-MAE 和 ViT-MAE 的线性探测分类准确率}。在主干网络于 ImageNet-1K 上进行掩码自编码预训练后，不同模型大小在各种数据集上的表现。在给定相同参数数量的情况下，CRATE-MAE 实现了大致相似的性能，同时享有更简单、更有原则的架构设计。}
\centering
\resizebox{1\columnwidth}{!}
{
    \begin{tabular}{@{}lcc|cc@{}}
        \toprule
        \textbf{Model} & CRATE-MAE-S(mall) & CRATE-MAE-B(ase) &
        {\color{gray} ViT-MAE-S} & {\color{gray} ViT-MAE-B} \\
        \midrule
        \midrule
        参数量 & 25.4M & 44.6M & 47.6M & {\color{gray}143.8M} \\
        \midrule
        CIFAR10 & 79.4 & 80.9 & {\color{gray} 79.9} & {\color{gray} 87.9} \\
        CIFAR100 & 56.6 & 60.1 & {\color{gray} 62.3} & {\color{gray} 68.0} \\
        Oxford Flowers-102 & 57.7 & 61.8 & {\color{gray} 66.8} &
        {\color{gray} 66.4} \\
        Oxford-IIIT-Pets & 40.6 & 46.2 & {\color{gray} 51.8} &
        {\color{gray} 80.1} \\
        \bottomrule
    \end{tabular}
}

\label{tab:crate_mae_linear_probing}
\end{table}

\begin{figure}
\centering
\includegraphics[width=\textwidth]{\toplevelprefix/chapters/chapter8/figs/crate_mae_semantic_heads.pdf}
\caption{\small\textbf{CRATE-MAE 的显著图。} 每对图像由原始图像（左）和对应于最后一层中某个注意力头的选定显著图（右）组成。正如 CRATE 模型的常见情况，但对于一般的类 Transformer 模型来说不常见的是，显著图对应于输入图像中的对象。}
\label{fig:crate_mae_semantic_heads}
\end{figure}

\section{通过自编码和采样进行图像生成}
\label{sec:auto-encode-img-gen}


在上一节中，我们介绍了掩码自编码（MAE）作为学习自然图像 $\x \sim p(\x)$ 自编码的一种方法：
\begin{equation}
\x\xrightarrow{\hspace{1mm} \mathcal{E} = f \hspace{1mm}} \z
\xrightarrow{\hspace{1mm} \mathcal{D} = g \hspace{1mm}} \hat{\x}.
\end{equation}
其中 $\z = f(\x)$ 是某个特征空间。这种自编码旨在学习自然图像图块之间的相关性或共现性。 

然而，MAE 并不强制学习到的特征 $\z$（的分布）具有任何显式的结构或属性，如果我们想将这样学习到的特征用于某些后续任务（如识别或受控生成），这一点可能非常重要。此外，用于图像补全的简单 MMSE 损失不一定尊重图像流形中可能的非线性结构，并且通常会为缺失的图像图块生成一个期望值，因此通常是模糊的版本，正如我们在上一章 \ref{ch:conditional-inference} 的图 \ref{fig:mae_autoencoding-small} 中所展示的那样。

因此，在本节中，我们将介绍一些其他通过自编码学习自然图像（分布）表示的方法，这些方法在很大程度上旨在弥补上述两个问题。更准确地说，我们希望编码特征的分布具有某些期望的结构，并且从这些特征解码出的图像更忠实于自然图像的分布——也就是说，具有更好的视觉质量。 


\subsection{图像数据的变分自编码}
% As we have introduced in Section \ref{sec:vae} of Chapter \ref{ch:consistent-representations}, variational auto-encoding is a method that aims to promote the distribution of learned features $\z$ and the conditional distributions between $\x$ and $\z$ to be more Gaussian like.\footnote{One may view such prior is to promote the learned the features to be diverse, in a similar spirit to maximizing the coding rate of all features.} 

% Here we introduce a popular implementation of a VAE for imagery data...




% \subsection{Variational Auto-Encoding of Imagery Data}
\label{sec:VAE-implement}


%refer back to Chapter 6 VAE here.
正如我们在第~\ref{ch:consistent-representations} 章第~\ref{sec:vae} 节中所介绍的，变分自编码是一种旨在促使学习到的特征 $\z$ 的分布以及 $\x$ 和 $\z$ 之间的条件分布更接近高斯分布的方法。\footnote{人们可以将这种先验视为促进学习到的特征具有多样性，这与最大化所有特征的编码率的精神相似。}

在这里，我们介绍一种流行的图像数据 VAE 实现，遵循 Stable Diffusion~\citep{rombach2022high} 中利用的感知压缩模型。其核心思想是训练一个自编码器，将高分辨率图像映射到低维潜在空间，同时对潜在分布施加 KL 正则化，以鼓励其接近标准高斯分布 $\mathcal{N}(\boldsymbol{0}, \boldsymbol{I})$。这个潜在空间随后可以作为下游生成模型（如潜在扩散模型）的高效操作空间。编码器和解码器都是端到端联合训练的。我们概述其细节如下。

\paragraph{编码器 ($\mathcal{E}$)。}
编码器 $\mathcal{E}$ 将输入的 RGB 图像 $\x \in \mathbb{R}^{H \times W \times 3}$ 映射到低维潜在空间。其架构采用自~\citep{esser2021taming}，是一个全卷积网络。它从一个输入卷积开始，将 3 通道 RGB 输入投影到基础通道数为 $128$ 的特征图中。随后是四个阶段，每个阶段包含两个 ResNet 块~\citep{he2016deep}。通道维度在各个阶段按乘数 $(1, 2, 4, 4)$ 进行缩放，分别产生具有 $128, 256, 512, 512$ 个通道的特征图。在连续的阶段之间，通过步幅卷积执行 $2\times$ 的空间下采样，导致三次下采样，总下采样因子为 $f = 2^3 = 8$。在这四个阶段之后，一个由两个额外 ResNet 块组成的中间块进一步细化特征。最后，一个输出卷积将特征投影到 $2c$ 个通道，这些通道被分成参数化潜在分布的两个量（如下所述）。

至关重要的是，与标准自编码器不同，编码器不输出单个确定性的潜在向量。相反，它通过输出两个量来参数化对角高斯后验 $q(\z | \x) = \mathcal{N}(\boldsymbol{\mu}, \boldsymbol{\sigma}^2 \boldsymbol{I})$：
\begin{itemize}
    \item 均值 $\boldsymbol{\mu} = \mathcal{E}_\mu(\x) \in \mathbb{R}^{h \times w \times c}$，以及
    \item 对数方差 $\log \boldsymbol{\sigma}^2 = \mathcal{E}_\sigma(\x) \in \mathbb{R}^{h \times w \times c}$。
\end{itemize}
然后通过\emph{重参数化技巧（reparameterization trick）}~\citep{kingma2013auto} 获得潜在表示：
\begin{equation}
    \z = \boldsymbol{\mu} + \boldsymbol{\sigma} \hada \boldsymbol{\epsilon}, \quad \boldsymbol{\epsilon} \sim \mathcal{N}(\boldsymbol{0}, \boldsymbol{I}),
\end{equation}
其中 $\hada$ 表示逐元素乘法。这种公式化使得能够通过随机采样步骤进行基于梯度的优化。

\paragraph{解码器 ($\mathcal{D}$)。}
解码器 $\mathcal{D}$ 学习从潜在空间回到像素空间的反向映射，生成重建图像 $\tilde{\x} = \mathcal{D}(\z) \in \mathbb{R}^{H \times W \times 3}$。其架构与编码器大致对称：输入卷积首先将 $c$ 通道潜在变量投影到特征空间，随后是一个包含两个 ResNet 块的中间块。然后，四个阶段——每个阶段包含两个 ResNet 块——逐步上采样空间分辨率，通道维度以相反的顺序镜像编码器的维度（$512, 512, 256, 128$）。最后的输出卷积将特征映射回 3 个 RGB 通道。编码器和解码器架构均遵循~\citep{esser2021taming} 中引入的卷积设计。

\subsubsection{数据}
我们假设可以访问来自原始数据分布的样本数据集 $\mathcal{D} = \{\x_i\}_{i=1}^N$。每个样本都是形式为 $\x \in \mathbb{R}^{H \times W \times 3}$ 的图像，其中 $H$ 和 $W$ 分别是高度和宽度。在~\citep{rombach2022high} 中，自编码器在 OpenImages 数据集上进行训练，图像被裁剪为 $256 \times 256$ 分辨率。

给定下采样因子 $f = 8$ 和潜在通道数 $c = 4$，编码器将每个输入图像映射到潜在张量 $\z \in \mathbb{R}^{H/f \times W/f \times c}$。对于默认的训练分辨率，这会产生大小为 $32 \times 32 \times 4$ 的潜在变量。在推理时，自编码器可以泛化到其他分辨率；例如，一个 $512 \times 512 \times 3$ 的输入被编码为一个 $64 \times 64 \times 4$ 的潜在变量。

\subsubsection{目标与评估指标}
为了训练 VAE，我们的目标是联合优化编码器 $\mathcal{E}$ 和解码器 $\mathcal{D}$，使得重建数据 $\tilde{\x} = \mathcal{D}(\mathcal{E}(\x))$ 尽可能接近原始数据 $\x$，同时正则化潜在分布使其接近标准高斯先验。重建损失的一个经验上有效的选择是 L1、LPIPS~\citep{zhang2018unreasonable} 和 GAN~\citep{goodfellow2014generative} 中的对抗损失的加权组合。此外，还包含一个 KL 散度项来正则化潜在空间。形式上，训练目标可以表示为：

\begin{equation}
\begin{split}
\min_{\theta_{\mathcal{E}}, \theta_{\mathcal{D}}} \; \mathcal{L}_{\mathrm{VAE}}(\theta_{\mathcal{E}}, \theta_{\mathcal{D}}) = \mathbb{E}_{\x \sim \mathcal{D}} \Big[
    &\|\tilde{\x} - \x\|_1 + \lambda_1 \, \mathrm{LPIPS}(\tilde{\x}, \x) + \lambda_2 \, \mathcal{L}_{\mathrm{GAN}}(\tilde{\x}, \x) \\
    &+ \; \lambda_3 \, D_{KL}\!\big(q(\z | \x) \,\big\|\, p(\z)\big)
\Big],
\end{split}
\label{eq:vae_loss}
\end{equation}
其中 $\theta_{\mathcal{E}}$ 和 $\theta_{\mathcal{D}}$ 分别是编码器和解码器的参数，$\lambda_1, \lambda_2, \lambda_3$ 是平衡每个损失项贡献的超参数。前三项以与第~\ref{sec:RAE-implement} 节中的重建目标相同的方式强制执行样本级一致性。第四项——编码器后验 $q(\z | \x) = \mathcal{N}(\boldsymbol{\mu}, \boldsymbol{\sigma}^2 \boldsymbol{I})$ 与标准高斯先验 $p(\z) = \mathcal{N}(\boldsymbol{0}, \boldsymbol{I})$ 之间的 KL 散度——对潜在空间进行正则化。对于对角高斯后验，这允许一个闭式表达式：
\begin{equation}
    D_{KL}\!\big(q(\z | \x) \,\big\|\, p(\z)\big) = \frac{1}{2} \sum_{j=1}^{d} \left( \mu_j^2 + \sigma_j^2 - \log \sigma_j^2 - 1 \right),
\end{equation}
其中 $d = h \times w \times c$ 是潜在空间的总维度。

来自~\citep{rombach2022high} 的实现细节是，KL 权重 $\lambda_3$ 被设置为一个非常小的值（$\approx 10^{-6}$）。这是因为这里 KL 项的主要目的不是为了实现从潜在空间的无条件采样（如在传统 VAE 中那样），而是为了防止潜在空间产生任意高的方差，这会降低在该空间中运行的下游模型的信噪比。

为了进行评估，我们通过计算 ImageNet-1k 验证集上重建图像的 FID 分数~\citep{heusel2017gans} 来衡量重建质量，记为 rFID。较低的 rFID 表示更好的重建质量。








\subsection{图像数据的表示自编码}
\label{sec:RAE-implement}
% \yima{Conceptual connections and necessity... Practical
% implementations details... (Robin and Peter)}

% rae stuff
\definecolor{deemph}{gray}{0.6}
\newcommand{\gc}[1]{\textcolor{deemph}{#1}}
\newcommand{\tablestyle}[2]{\setlength{\tabcolsep}{#1}\renewcommand{\arraystretch}{#2}\centering\scriptsize}
% \newcommand{\round}[1]{\num[round-mode=places, round-precision=2]{#1}}



正如在 \Cref{ch:lossy-compression} 中所讨论的，学习数据的结构化表示对于许多下游任务至关重要。在视觉表征学习中，这通常通过自监督学习方法来实现，例如 DINO \cite{caron2021emerging} 和 SimDINO（在第 \ref{sec:contrastive_learning} 节中研究）。在 \Cref{ch:consistent-representations} 中，我们假设一个好的数据表示应该能够有效地解码回原始数据空间。这需要一个一致的自编码框架，能够从其学习到的表示中准确地重建原始数据。那么很自然地会问：像 DINO 这样的方法学习到的表示是一致的吗？它能在自编码框架中被有效地解码吗？

为了找出答案，我们可以训练一个\textit{表示自编码器}（RAE），它由一个预训练的表示编码器（例如 DINO）和一个可训练的解码器组成。编码器 $f$ 将输入样本 $\x$ 映射到其学习到的表示 $\z = f(\x)$，而解码器 $g$ 将该表示映射回原始数据空间 $\x'=g(\z)$。正如在 \Cref{ch:consistent-representations} 中详细说明的，一个简单的目标是通过训练解码器以最小化原始数据 $\x$ 和重建数据 $\x'$ 之间的重建损失，从而施加样本级一致性。我们概述其细节如下。

\paragraph{编码器。} 我们使用一个预训练的表示编码器 $f \colon \cX \to \cZ$，它将输入数据样本 $\x \in \cX$ 映射到其学习到的表示 $\z = f(\x) \in \cZ$。在训练期间，编码器保持冻结状态。在这里，我们假设使用基于 ViT 的编码器，如 \Cref{sub:contrastive_learning_architecture} 中所介绍的。

\paragraph{解码器。} 为了架构的简单性，我们使用一个可训练的基于 ViT 的解码器 $g \colon \cZ \to \cX$，它将学习到的表示 $\z$ 映射回原始数据空间，生成重建数据 $\x' = g(\z)$。

\subsubsection{数据}
我们假设可以访问来自原始数据分布 $\cX$ 的样本数据集 $\cD = \{\x_i\}_{i=1}^N$。每个样本都是形式为 $\x \in \R^{3 \times H \times W}$ 的图像，其中 $H$ 和 $W$ 分别是高度和宽度。编码器 $f$ 的图块大小设置为 $p_e$，特征维度设置为 $d$。因此，对于每个输入图像 $\x$，我们将获得 $N = HW/p_e^2$ 个维度为 $d$ 的标记（tokens）。然后，图块大小为 $p_d$ 的解码器 $g$ 将这些视觉标记映射回去，以在原始像素空间中重建 $\x' \in \R^{3 \times H\frac{p_d}{p_e} \times W\frac{p_d}{p_e}}$。默认情况下，我们设置 $p_d = p_e$ 以确保重建图像具有与原始图像相同的分辨率。

\subsubsection{目标与评估指标} 为了训练 RAE，我们的目标是训练一个解码器 $g$，使得重建数据 $\x' = g(f(\x))$ 尽可能接近原始数据 $\x$。这可以通过最小化数据集 $\cD$ 上的样本级一致性损失来实现。损失函数的一个经验上有效的选择是 L1、LPIPS \cite{zhang2018unreasonable} \footnote{这是一种基于计算预训练网络不同层之间特征距离的感知损失。经验表明，它比像素级损失更好地捕捉了人类的感知相似性。详情请参阅原始论文。} 和 GAN \cite{goodfellow2014generative} 中的对抗损失的加权组合。形式上，训练目标可以表示为：

\begin{equation}
\min_{\theta} \cL_{\mathrm{RAE}}(\theta) = \Ex_{\x \sim \cD} \left[
    \norm{\x' - \x}_1 + \lambda_1 \mathrm{LPIPS}(\x', \x) + \lambda_2
\cL_{\mathrm{GAN}}(\x', \x) \right],
\end{equation}
其中 $\theta$ 是解码器 $g$ 的参数，$\lambda_1, \lambda_2$ 是平衡每个损失项贡献的超参数。

为了进行评估，我们通过计算 ImageNet-1k 验证集上重建图像的 FID 分数 \cite{heusel2017gans} 来衡量重建质量，记为 rFID。较低的 rFID 表示更好的重建质量。

\subsubsection{实验结果}
我们可以比较在不同预训练编码器（包括 DINOv2、MAE 和 SigLIP2）上训练的 RAE。作为基线，我们还考虑了直接在图像像素上训练而没有任何预训练编码器的普通变分自编码器（VAE）。一个流行的选择是 SD-VAE。RAE 解码器在 ImageNet-1K 训练集上进行训练，并在 ImageNet-1K 验证集上进行评估。结果总结在 \Cref{tab:rae_recon_results} 中。从 \Cref{tab:decoder_scaling_no_noise} 中，我们可以看到具有预训练编码器的 RAE 在重建质量上优于 SD-VAE。此外，如 \Cref{tab:decoder_scaling_no_noise} 所示，更大的解码器尺寸会带来重建质量的提高，同时仍然比 VAE 更高效。\Cref{tab:encoder_scaling_dinov2} 表明，在不同大小的 DINO 编码器中，重建质量是稳定的。最后，\Cref{tab:linear_probing} 证明 RAE 还保留了高质量的表示，与 VAE 相比，实现了高得多的线性探测准确率。

\begin{table*}[t]
\centering
{\captionsetup[subfloat]{labelfont=scriptsize}
    \caption{
    \small\textbf{RAE 的重建和线性探测结果。} }
    \begin{tabular}{@{}cc@{}}
        \subfloat[
            \scriptsize \textbf{编码器选择。} 所有编码器均优于 SD-VAE。
            \label{tab:recon_deterministic}
        ]
        {
            \begin{minipage}[t]{0.45\linewidth}{
                    \vspace{0pt}
                    \begin{center}
                        \tablestyle{4pt}{1.1}
                        \begin{tabular}{lc}
                            \toprule
                            Model & rFID \\
                            \midrule
                            \rowcolor{gray!20}
                            DINOv2-B & 0.49 \\
                            \rowcolor{white}
                            SigLIP2-B & 0.53 \\
                            MAE-B & 0.16 \\
                            \midrule
                            SD-VAE & \gc{0.62} \\
                            \bottomrule
                        \end{tabular}
                    \end{center}
                }
            \end{minipage}
        }
        &
        \subfloat[
            \scriptsize \textbf{更大的解码器} 改善了 rFID，同时仍然比 VAE 高效得多。
            \label{tab:decoder_scaling_no_noise}
        ]
        {
            \begin{minipage}[t]{0.45\linewidth}{
                    \vspace{0pt}
                    \begin{center}
                        \tablestyle{4pt}{1.1}
                        \begin{tabular}{lcc}
                            \toprule
                            Decoder & rFID & GFLOPs \\
                            \midrule
                            ViT-B & 0.58 & 22.2 \\
                            ViT-L & 0.50 & 78.1 \\
                            \rowcolor{gray!20}
                            ViT-XL & 0.49 & 106.7 \\
                            \midrule
                            \rowcolor{white}
                            SD-VAE & \gc{0.62} & \gc{310.4} \\
                            \bottomrule
                        \end{tabular}
                \end{center}}
            \end{minipage}
        }
        \\[1em]
        \subfloat[
            \scriptsize \textbf{编码器缩放。} rFID 在不同 RAE 尺寸下保持稳定。
            \label{tab:encoder_scaling_dinov2}
        ]{
            \begin{minipage}[t]{0.45\linewidth}{
                    \vspace{0pt}
                    \begin{center}
                        \tablestyle{4pt}{1.1}
                        \begin{tabular}{lc}
                            \toprule
                            Encoder & rFID \\
                            \midrule
                            DINOv2-S & 0.52 \\
                            \rowcolor{gray!20}
                            DINOv2-B & 0.49 \\
                            \rowcolor{white}
                            DINOv2-L & 0.52 \\
                            \bottomrule
                            \multicolumn{2}{c}{~} \\
                            \multicolumn{2}{c}{~} \\
                        \end{tabular}
                        \vspace{-0.1cm}
                    \end{center}
                }
            \end{minipage}
        }
        &
        \subfloat[
            \scriptsize \textbf{表示质量。} RAE 的线性探测准确率远高于 VAE。
            \label{tab:linear_probing}
        ]
        {
            \begin{minipage}[t]{0.45\linewidth}{
                    \vspace{0pt}
                    \begin{center}
                        \tablestyle{4pt}{1.1}
                        \begin{tabular}{lc}
                            \toprule
                            Model & Top-1 Acc. \\
                            \midrule

                            \rowcolor{gray!20} DINOv2-B & 84.5 \\
                            \rowcolor{white}
                            SigLIP2-B & 79.1 \\
                            MAE-B & 68.0 \\
                            \midrule
                            SD-VAE & \gc{8.0} \\
                            \bottomrule
                        \end{tabular}
                    \end{center}
                }
            \end{minipage}
        }
    \end{tabular}
}
% \vspace{-4pt}
\label{tab:rae_recon_results}
\end{table*}
为了补充 \Cref{tab:rae_recon_results} 中的定量指标，\Cref{fig:rae_recon_qualitative} 提供了重建的视觉比较。RAE 重建像 SD-VAE 基线一样忠实地保留了细粒度纹理、对象边界和语义细节，这与其更好的 rFID 分数相一致。

\begin{figure}[t]
    \centering
    \includegraphics[width=\textwidth]{chapters/chapter8/figs/recon_grid.pdf}
    \caption{\small\textbf{RAE 重建示例。} 从左到右：输入图像、RAE (DINOv2-B)、RAE (Siglp-B)、RAE (MAE-B)、SD-VAE。放大以查看细节。}
    \label{fig:rae_recon_qualitative}
\end{figure}

% \yima{Besides the quantitative results, please have some
% illustrations to show qualitative improvements... }

\subsection{通过去噪从学习到的表示中采样}

一旦我们学习到了自然图像分布的一致且结构化的表示（例如通过 VAE 或 RAE），我们通常希望从学习到的表示分布中高效地生成\textit{样本} $\z \in \cZ$，以便我们可以通过学习到的解码器 $\hat{\x} = g(\z) \in \cX$ 重新生成与这些样本相关的自然图像。这个自编码过程由第 \ref{ch:intro} 章中图 \ref{fig:mental-picture} 的上半部分说明。

为了学习并因此对潜在空间 $\cZ$ 中特征 $\z$ 的分布进行采样，我们可以利用第 \ref{ch:denoising} 章中介绍的去噪和扩散方法。具体来说，假设 $p(\z)$ 是特征空间中学习到的特征 $\z$ 的分布。我们可以学习一个去噪过程，该过程从初始的标准高斯分布 $\vg \sim \cN(0, \vI)$ 开始，收敛到该分布 $p(\z)$。按照第 \ref{ch:denoising} 章中介绍的方法，我们希望学习一个去噪器 $\bar{\z}$，使其迭代地将初始分布 $\cN(0, \vI)$ 中的样本去噪为学习到的表示的目标分布 $p(\z)$。这个过程由第 \ref{ch:intro} 章中图 \ref{fig:mental-picture} 的下半部分说明。

在这里，为了学习去噪，让我们以第 \ref{ch:denoising} 章中介绍的流匹配（flow matching）的简单情况为例。我们可以将去噪器定义为一个由 $\theta$ 参数化的随时间变化的向量场 $\bar{\z}_{\theta}(t,
\z_t)$，它随着时间 $t \in [0,
1]$ 将样本从初始分布转换为目标分布。这里，$\z_t$ 是在时间 $t$ 时 $\z_0$ 和 $\z_1$ 之间的插值：
\begin{equation}
\z_t = (1 - t)\z_0 + t\z_1,
\end{equation}
其中 $\z_0 \sim p(\z)$ 是来自学习到的表示分布的样本，且 $\z_1 \sim \cN(0, \vI)$。去噪器的训练目标可以表示为：
\begin{equation}
\min_{\theta} \cL_{\mathrm{FM}}(\theta) = \Ex_{\z_0 \sim p(\z), \z_1
\sim \cN(0, I), t \sim \mathrm{Unif}(0, 1)} \left[
\norm{\bar{\z}_{\theta}(t,\z_t) - (\z_1 - \z_0)}_2^2 \right].
\end{equation}

通常，去噪向量场 $\bar{\z}_{\theta}(t,
\z_t)$ 由深度网络建模。深度网络的流行选择可以是 U-Net 或扩散图像 Transformer（DiT）。下面我们提供有关如何在实践中实现它们的更多细节。




% %more detail about the UNet, and refer back to PCA. Add the references.
% \subsubsection{Realizing Denoising with a U-Net.}
% A common architecture for implementing the parametric denoiser
% $\bar{\z}_{\theta}(t, \z_t)$ is the \textit{U-Net}, an encoder-decoder 
% network with skip connections. The encoder progressively downsamples the 
% noisy input to extract multi-scale features, while the decoder upsamples 
% and refines predictions. Skip connections concatenate encoder features to 
% corresponding decoder layers, preserving spatial detail. To capture 
% long-range dependencies, self-attention layers are often inserted at 
% lower resolutions. The convolutional structure provides strong spatial 
% inductive biases, which can be beneficial when training data is limited.

% The timestep $t$ is typically embedded via 
% sinusoidal functions (analogous to positional encodings in transformers) 
% and injected through \textit{adaptive group normalization}: the scale and 
% shift parameters of normalization layers are regressed from the timestep 
% embedding. For class-conditional generation, a class embedding is simply 
% added to the timestep embedding before injection. For text conditioning, 
% a pretrained text encoder (e.g., CLIP) produces a sequence of token 
% embeddings, which are incorporated via \textit{cross-attention}: at 
% selected layers, the image features serve as queries while the text 
% embeddings serve as keys and values. This allows the network to attend to 
% relevant parts of the text prompt when denoising each spatial location.






\subsection{使用 U-Net 实现去噪}
\label{sec:UNet-implement}

在扩散模型中实现去噪网络 $\boldsymbol{\epsilon}_\theta$ 的主导架构一直是 \textit{U-Net}~\citep{ronneberger2015unet}，这是一种带有跳跃连接的编码器-解码器卷积网络，最初是为生物医学图像分割而开发的。Ho 等人~\citep{ho2020denoising} 首次采用 U-Net 作为去噪扩散概率模型（DDPM）的主干，确立了其作为该领域默认架构的地位。随后，Dhariwal 和 Nichol~\citep{dhariwal2021diffusion} 进行了一系列系统的架构消融实验，大幅改进了 U-Net 设计，产生了他们所谓的\emph{消融扩散模型}（ADM）。他们改进的模型是第一个在 ImageNet 上的样本质量超越 GAN 的扩散模型。在这些基础之上，Rombach 等人~\citep{rombach2022high} 将 U-Net 从像素空间移至预训练 VAE 的潜在空间（第~\ref{sec:VAE-implement} 节），并引入了交叉注意力层以实现灵活的多模态条件控制，从而催生了支撑 Stable Diffusion 的\emph{潜在扩散模型}（LDM）。在本节中，我们将追溯 U-Net 主干在这三个阶段的演变，然后讨论促使最近向基于 Transformer 的替代方案过渡的架构局限性（第~\ref{sec:DiT-implement} 节）。

\paragraph{带有跳跃连接的编码器-解码器结构。}
U-Net 遵循编码器-解码器结构。编码器通过一系列分辨率阶段逐步对输入进行下采样，在多个空间尺度上提取特征；然后解码器将这些特征上采样回原始分辨率以产生输出预测。U-Net 的显著特征是其\emph{跳跃连接（skip connections）}：在每个分辨率级别，编码器特征图在通道维度上与相应的解码器特征图拼接。这使得解码器能够将来自瓶颈层的粗糙、语义丰富的特征与来自早期编码器层的细粒度空间细节结合起来，这对于产生空间上精确的预测至关重要。

在 DDPM U-Net~\citep{ho2020denoising} 中，架构从一个输入卷积开始，将加噪输入投影到具有基础通道数（例如 128）的特征空间中。然后编码器经过多个分辨率阶段——对于 $32 \times 32$ 输入为四个阶段，对于 $256 \times 256$ 输入为六个阶段——其中每个阶段由两个卷积残差块组成。在连续的阶段之间，通过步幅卷积执行 $2\times$ 的空间下采样。自注意力层插入在卷积块之间的 $16 \times 16$ 分辨率级别，在低分辨率下提供有限的长程建模。在瓶颈处（最低分辨率），一个带有额外残差块和自注意力的中间块进一步细化特征。解码器镜像了编码器：每个分辨率级别的残差块，随后是通过转置卷积进行的上采样，并带有拼接相应编码器激活的跳跃连接。该架构自始至终使用组归一化（group normalization）~\citep{wu2018group}，取代了该设计最初起源的权重归一化。

当 U-Net 在像素空间中运行时，其输入和输出分别是加噪图像 $\x_t \in \mathbb{R}^{H \times W \times 3}$ 和预测噪声 $\hat{\boldsymbol{\epsilon}} \in \mathbb{R}^{H \times W \times 3}$。当在预训练 VAE 的潜在空间中运行（如在 LDM 中）时，输入变为 $\z_t \in \mathbb{R}^{h \times w \times c}$，输出为 $\hat{\boldsymbol{\epsilon}} \in \mathbb{R}^{h \times w \times c}$，其中 $h = H/f$，$w = W/f$，如第~\ref{sec:VAE-implement} 节中所定义。

\paragraph{通过自适应组归一化进行时间步条件控制。}
去噪网络必须以扩散时间步 $t$ 为条件，以便它能够适应不同噪声水平下的行为。遵循 Transformer~\citep{vaswani2017attention} 中的位置编码方案，时间步 $t$ 首先被映射为正弦频率嵌入：
\begin{equation}
    \mathrm{Embed}(t)_{2i} = \sin\!\left(\frac{t}{10000^{2i/d_e}}\right), \quad
    \mathrm{Embed}(t)_{2i+1} = \cos\!\left(\frac{t}{10000^{2i/d_e}}\right),
\end{equation}
其中 $d_e$ 是嵌入维度，$i$ 索引其条目。然后，该嵌入通过一个小型 MLP（通常是两个带有非线性的线性层）以产生时间步向量 $\boldsymbol{e}_t \in \mathbb{R}^{d_e}$。

U-Net 并非简单地将 $\boldsymbol{e}_t$ 拼接到输入中，而是通过\emph{自适应组归一化}（AdaGN）~\citep{dhariwal2021diffusion} 将时间步信息注入到每个残差块中。在每个残差块内的第一个卷积之后，应用组归一化操作，随后是仿射调制，其缩放和平移参数从时间步嵌入中回归得到：
\begin{equation}
    \mathrm{AdaGN}(\boldsymbol{h}, \boldsymbol{e}_t) = \boldsymbol{y}_s \hada \mathrm{GroupNorm}(\boldsymbol{h}) + \boldsymbol{y}_b,
    \label{eq:adagn}
\end{equation}
其中 $\boldsymbol{h}$ 表示中间特征图，$[\boldsymbol{y}_s, \boldsymbol{y}_b] = \mathrm{Linear}(\boldsymbol{e}_t)$ 是从时间步嵌入的线性投影中获得的缩放和平移向量，$\hada$ 表示逐元素乘法。这种机制类似于自适应实例归一化~\citep{dumoulin2017learned} 和 FiLM 条件控制~\citep{perez2018film}，并允许每一层根据当前的噪声水平调制其激活。~\citep{dhariwal2021diffusion} 中的消融实验表明，移除 AdaGN 会使 FID 降低超过 2 个点，证实了其重要性。

对于类别条件生成，在注入之前，只需将类别嵌入 $\mathrm{Embed}(c)$（一个学习到的查找表）添加到时间步嵌入中，使得等式~\eqref{eq:adagn} 中的 $\boldsymbol{e}_t$ 被替换为 $\boldsymbol{e}_t + \mathrm{Embed}(c)$。这种机制是轻量级的，但仅限于分类标签；更丰富的条件模态（如文本）需要第~\ref{sec:unet-ldm} 节中引入的交叉注意力机制。

\subsubsection{架构改进：从 DDPM 到 ADM}
\label{sec:unet-adm}

Dhariwal 和 Nichol~\citep{dhariwal2021diffusion} 进行了一项系统的消融研究，以确定哪些架构修改最能提高样本质量。从基线 DDPM U-Net~\citep{ho2020denoising} 开始，他们在 ImageNet $128 \times 128$ 上探索了以下变化：(a) 将自注意力从 $16 \times 16$ 分辨率扩展到三个分辨率（$32 \times 32$、$16 \times 16$、$8 \times 8$）；(b) 增加注意力头的数量，每个头使用 64 个通道以更好地匹配 Transformer 的惯例；(c) 在上采样和下采样中采用 BigGAN~\citep{brock2019large} 残差块，这改善了梯度流；以及 (d) 默认在每个残差块中包含 AdaGN 条件层（如上所述）。所有这些变化都改善了 FID，并且它们的效果是正向叠加的。由此产生的模型被他们命名为 ADM，作为后续工作的新默认 U-Net 设计。

除了架构改进之外，Dhariwal 和 Nichol 还引入了\emph{分类器引导（classifier guidance）}作为一种在样本多样性和保真度之间进行权衡的机制。在加噪图像上训练一个分类器 $p_\phi(c|\z_t, t)$，并使用其梯度将采样轨迹引向所需类别的高概率区域：
\begin{equation}
    \hat{\boldsymbol{\mu}}_\theta(\z_t, t) = \boldsymbol{\mu}_\theta(\z_t, t) + s \, \boldsymbol{\Sigma}_\theta(\z_t, t) \, \nabla_{\z_t} \log p_\phi(c|\z_t),
\end{equation}
其中 $s > 1$ 是一个引导尺度，它使有效类别分布变得更加尖锐。借助这项技术，ADM 在类别条件 ImageNet $256 \times 256$ 上实现了 4.59 的最先进 FID，首次超越了长期以来的 GAN 基准 BigGAN-deep（FID 6.95）。当与上采样堆栈（ADM-U）结合使用时，引导模型（ADM-G）进一步将 FID 提高到 3.94。

分类器引导的主要缺点是它需要在加噪图像上训练一个单独的分类器。这一局限性随后被\emph{无分类器引导（classifier-free guidance）}~\citep{ho2022classifier} 所解决，它通过在单个模型中联合学习条件和无条件去噪，完全避免了外部分类器——我们将在第~\ref{sec:DiT-implement} 节中详细描述这项技术。

\subsubsection{潜在扩散（Stable Diffusion）中的 U-Net}
\label{sec:unet-ldm}

虽然 ADM 直接在像素空间中运行，但 Rombach 等人~\citep{rombach2022high} 观察到，大部分计算负担都花在了对感知上无关的高频细节进行建模上。他们的核心思想是将\emph{感知压缩}与\emph{生成建模}分开：预训练的自编码器（第~\ref{sec:VAE-implement} 节）首先将图像压缩到低维潜在空间中，然后扩散模型完全在这个潜在空间内运行。对于下采样因子 $f = 8$ 的 $256 \times 256$ 图像，U-Net 处理的是 $32 \times 32 \times 4$ 的潜在张量，而不是 $256 \times 256 \times 3$ 的像素张量——空间元素减少了大约 $48\times$。这在训练和采样吞吐量上产生了至少 $2.7\times$ 的加速，同时改善了 FID。

除了节省计算量之外，潜在扩散框架还引入了一种通用的\emph{交叉注意力条件机制}，使 U-Net 能够以任意模态（如文本提示、语义图或边界框）为条件。简而言之，特定领域的编码器 $\tau_\theta$ 将条件输入 $\boldsymbol{y}$ 映射为中间表示 $\tau_\theta(\boldsymbol{y}) \in \mathbb{R}^{M \times d_\tau}$，然后通过交叉注意力层将其注入到 U-Net 的多个分辨率级别中：图像特征作为查询（queries），而条件嵌入作为键（keys）和值（values），允许每个空间位置关注条件信号的相关部分。我们将这种交叉注意力机制和条件训练目标的完整数学描述推迟到第~\ref{subsec:basic_realization} 节。

在 Stable Diffusion~\citep{rombach2022high} 中，文本编码器 $\tau_\theta$ 被实例化为预训练的 CLIP ViT-L/14 模型~\citep{radford2021learning}，它将文本提示映射为一系列标记嵌入。U-Net 本身保留了 ADM 的改进——BigGAN 残差块、三个分辨率级别的注意力、用于时间步注入的 AdaGN——并额外在这些相同的分辨率级别交错插入交叉注意力层以进行文本条件控制。完整模型在 U-Net 中包含大约 8.6 亿个参数，在文本编码器中包含 1.23 亿个参数。

\subsubsection{U-Net 架构的局限性}
\label{sec:unet-limitations}

尽管取得了成功，但 U-Net 存在一些局限性，这促使了替代主干的出现。首先，\textbf{扩展（scaling）} U-Net 是繁琐的：其编码器-解码器结构涉及许多耦合的设计选择（通道乘数、注意力位置、上/下采样块类型），并且不清楚如何最佳地分配额外的计算资源。相比之下，普通的 Transformer 沿着深度和宽度这两个简单的轴进行扩展，并表现出基于经验研究的特征明确的缩放定律~\citep{kaplan2020scaling}。其次，U-Net 的卷积层编码了很强的\textbf{空间归纳偏置}（局部性、平移等变性），这在低数据状态下有所帮助，但限制了大型数据集上的长程建模。自注意力部分解决了这个问题，但在标准 U-Net 设计中，由于其二次方成本，它仅应用于低分辨率阶段，使得高分辨率特征完全是局部的。第三，U-Net 是一种\textbf{特定于图像}的架构，不能很好地轻易泛化到其他模态。Transformer 已经被证明在语言、视觉、音频和多模态设置中有效~\citep{vaswani2017attention}，它提供了一个统一的主干，可以利用跨领域的扩展见解。

这些考虑促使了扩散 Transformer（DiT）~\citep{peebles2023scalable} 的产生，它完全用在潜在图块上运行的普通 Transformer 取代了 U-Net。正如我们在第~\ref{sec:DiT-implement} 节中所描述的，DiT 匹配并超越了最好的基于 U-Net 的模型，同时表现出更清晰的缩放行为（Gflops 与 FID 的相关性为 $-0.93$）和更简单的架构。









% % refer the transformer arch back to previous Chapters. Add the references.
% \subsubsection{Realizing Denoising with a DiT.}
% An alternative to the U-Net is the \textit{Diffusion Transformer (DiT)}, 
% which replaces convolutions entirely with transformer blocks. Following 
% the Vision Transformer, the input is divided into patches and linearly 
% embedded as token sequences. Unlike the U-Net, the DiT has no explicit 
% spatial inductive bias---all structure is learned from data. This 
% generality enables the same architecture to handle images, video, and 
% other modalities with minimal modification. Perhaps most importantly, DiT 
% follows favorable scaling laws: larger models and more compute yield 
% predictable improvements in sample quality.

% Timestep and class conditioning in DiT is 
% achieved through \textit{adaptive layer normalization} (adaLN-Zero): the 
% scale and shift parameters of each layer normalization are regressed from 
% the conditioning embedding (formed by summing timestep and class 
% embeddings), with these parameters initialized near zero for training 
% stability. For text conditioning, two approaches are common: (1) 
% cross-attention, as in the U-Net, where text token embeddings from a 
% frozen encoder are attended to by the image tokens; or (2) prepending the 
% text tokens directly to the image token sequence, allowing the 
% transformer's self-attention to jointly process both modalities.



% Once we have trained the denoiser, we can use it to generate new
% samples from the learned representation distribution via the denoising
% process. The generated feature samples can then be decoded back to
% the image space with the corresponding decoder $g$.



\subsection{使用扩散 Transformer 实现去噪}
\label{sec:DiT-implement}

实现去噪网络 $\boldsymbol{\epsilon}_\theta$ 的 U-Net 的一种替代方案是\textit{扩散 Transformer（DiT）}~\citep{peebles2023scalable}，它完全用 Transformer 块取代了卷积。关键的见解是，U-Net 的空间归纳偏置对扩散模型的性能\emph{并不至关重要}：在潜在图块序列上运行的标准 Transformer 可以匹配并超越 U-Net，同时继承 Transformer 架构类良好的缩放属性~\citep{vaswani2017attention}。与 U-Net 不同，DiT 没有显式的空间归纳偏置——所有的空间结构都是从数据中学习的。这种通用性使得相同的架构只需极少的修改就能处理图像、视频和其他模态。

DiT 在预训练 VAE（第~\ref{sec:VAE-implement} 节）的潜在空间内运行，使得完整的图像生成管道成为一种混合架构：用于感知压缩的卷积 VAE 和用于生成建模的基于 Transformer 的 DDPM。遵循视觉 Transformer（ViT）~\citep{dosovitskiy2020image}，输入潜在变量被划分为图块，并线性嵌入为标记序列。我们在下面详细描述该架构。

\paragraph{图块化与位置编码。}
DiT 的输入是加噪潜在变量 $\z_t \in \mathbb{R}^{h \times w \times c}$，其中 $h = H/f$ 和 $w = W/f$ 是由 VAE 编码器定义的潜在空间的空间维度（例如，对于 $f = 8, c = 4$ 的 $256 \times 256$ 图像，为 $32 \times 32 \times 4$）。第一层称为\emph{图块化（patchify）}，将 $\z_t$ 划分为大小为 $p \times p$ 的非重叠图块网格，并将每个图块线性嵌入为 $d$ 维标记表示。这产生了一个序列：
\begin{equation}
    T = \left(\frac{h}{p}\right) \times \left(\frac{w}{p}\right)
\end{equation}
包含 $\mathbb{R}^d$ 中的标记。标准的正弦位置嵌入（来自 ViT 的正弦-余弦变体）被添加到每个标记中以编码空间位置。图块大小 $p$ 是一个关键的设计参数：将 $p$ 减半会使 $T$ 增加四倍，从而至少使 Transformer 的总 Gflops 增加四倍，而参数数量基本保持不变。在~\citep{peebles2023scalable} 中，设计空间包括 $p \in \{2, 4, 8\}$；最精细的设置 $p = 2$ 产生最佳的生成质量。

\paragraph{带有 adaLN-Zero 的 DiT 块。}
图块化之后，标记序列由 $N$ 个 Transformer 块的堆栈处理。每个块遵循标准的 Transformer 结构~\citep{vaswani2017attention}：层归一化、多头自注意力（MHSA）、第二次层归一化以及逐点前馈网络（带有 GELU 激活的 MLP）。

核心的设计问题是如何将条件信息——扩散时间步 $t$ 和类别标签 $c$——注入到每个块中。Peebles 和 Xie~\citep{peebles2023scalable} 比较了四种策略：(1) \emph{上下文条件控制（in-context conditioning）}，将 $t$ 和 $c$ 作为额外的标记附加到输入序列中；(2) \emph{交叉注意力}，引入一个额外的多头交叉注意力层来关注条件嵌入；(3) \emph{自适应层归一化（adaLN）}，从条件信号中回归每个层归一化的缩放和平移参数；以及 (4) \emph{adaLN-Zero}，它用一个初始化为零的额外门控机制扩展了 adaLN。在这些策略中，adaLN-Zero 实现了最低的 FID，同时计算效率最高，因此被采用为默认策略。

在 adaLN-Zero 块中，条件信息首先通过对时间步和类别标签的嵌入求和来形成：
\begin{equation}
    \boldsymbol{e} = \mathrm{Embed}(t) + \mathrm{Embed}(c),
\end{equation}
其中 $\mathrm{Embed}(t)$ 是通过将 $t$ 的正弦频率嵌入传递给带有 SiLU 激活的两层 MLP 获得的，而 $\mathrm{Embed}(c)$ 是学习到的类别嵌入。从 $\boldsymbol{e}$ 中，一个线性层回归出六组逐维度的参数：用于自注意力子块的 $\boldsymbol{\gamma}_1, \boldsymbol{\beta}_1, \boldsymbol{\alpha}_1$ 和用于前馈子块的 $\boldsymbol{\gamma}_2, \boldsymbol{\beta}_2, \boldsymbol{\alpha}_2$。自适应层归一化应用如下：
\begin{equation}
    \mathrm{adaLN}(\boldsymbol{h}, \boldsymbol{\gamma}, \boldsymbol{\beta}) = \boldsymbol{\gamma} \hada \mathrm{LayerNorm}(\boldsymbol{h}) + \boldsymbol{\beta},
\end{equation}
其中 $\hada$ 表示逐元素乘法。参数 $\boldsymbol{\alpha}_1$ 和 $\boldsymbol{\alpha}_2$ 作为逐维度的\emph{门控因子}，在每个残差连接之前立即应用：
\begin{equation}
\begin{aligned}
    \boldsymbol{h}' &= \boldsymbol{h} + \boldsymbol{\alpha}_1 \hada \mathrm{MHSA}\!\big(\mathrm{adaLN}(\boldsymbol{h}, \boldsymbol{\gamma}_1, \boldsymbol{\beta}_1)\big), \\
    \boldsymbol{h}'' &= \boldsymbol{h}' + \boldsymbol{\alpha}_2 \hada \mathrm{MLP}\!\big(\mathrm{adaLN}(\boldsymbol{h}', \boldsymbol{\gamma}_2, \boldsymbol{\beta}_2)\big).
\end{aligned}
\end{equation}
adaLN-Zero 中的“Zero”指的是初始化：所有 $\boldsymbol{\alpha}$ 参数都由回归它们的 MLP 初始化为零向量，使得每个 DiT 块在训练开始时充当\emph{恒等函数}。这种零初始化策略受到 ResNets~\citep{goyal2017accurate} 和扩散 U-Nets 中类似实践的启发，对于训练稳定性很重要，并且产生比普通 adaLN 显著更低的 FID。

我们注意到，adaLN 将\emph{相同}的条件函数均匀地应用于所有标记。对于文本到图像的生成，其中条件信号是可变长度的文本标记序列，两种替代方法很常见：(1) 交叉注意力，其中来自冻结编码器的文本标记嵌入被图像标记所关注；或 (2) 将文本标记直接前置到图像标记序列中，允许 Transformer 的自注意力联合处理两种模态。

\paragraph{输出投影。}
在最后一个 DiT 块之后，模型必须将标记表示序列解码回空间预测。应用最后一次自适应层归一化（以 $\boldsymbol{e}$ 为条件），随后是一个线性层，将每个标记投影为形状为 $p \times p \times 2c$ 的张量。$2c$ 的因子解释了模型预测的两个量：噪声估计 $\hat{\boldsymbol{\epsilon}} \in \mathbb{R}^{h \times w \times c}$ 和对角协方差 $\boldsymbol{\Sigma}_\theta \in \mathbb{R}^{h \times w \times c}$。然后将解码后的标记重新排列（反图块化）回每个输出的原始空间布局 $\mathbb{R}^{h \times w \times c}$。


\subsubsection{训练}
DiT 作为类别条件潜在扩散模型在 ImageNet 数据集~\citep{deng2009imagenet} 上以 $256 \times 256$ 和 $512 \times 512$ 的图像分辨率进行训练。扩散过程在第~\ref{sec:VAE-implement} 节中预训练 VAE 产生的潜在表示上运行：对于 $256 \times 256$ 图像，DiT 的输入是 $\z_t \in \mathbb{R}^{32 \times 32 \times 4}$；对于 $512 \times 512$ 图像，它是 $\z_t \in \mathbb{R}^{64 \times 64 \times 4}$。

训练目标结合了两个项。噪声预测损失通过简单的均方误差训练 $\boldsymbol{\epsilon}_\theta$：
\begin{equation}
    \mathcal{L}_{\mathrm{simple}}(\theta) = \mathbb{E}_{\z_0, \boldsymbol{\epsilon}, t}\left[\|\boldsymbol{\epsilon}_\theta(\z_t, t) - \boldsymbol{\epsilon}\|_2^2\right],
\end{equation}
其中 $\z_0$ 是干净的潜在变量，$\boldsymbol{\epsilon} \sim \mathcal{N}(\boldsymbol{0}, \boldsymbol{I})$ 是采样的噪声，$\z_t = \sqrt{\bar{\alpha}_t}\,\z_0 + \sqrt{1 - \bar{\alpha}_t}\,\boldsymbol{\epsilon}$ 是时间步 $t$ 的加噪潜在变量。学习到的协方差 $\boldsymbol{\Sigma}_\theta$ 额外使用完整的变分下界 $D_{KL}$ 项进行训练，遵循 Nichol 和 Dhariwal~\citep{nichol2021improved} 的方法。使用线性方差调度，步数为 $t_{\max} = 1000$，范围从 $1 \times 10^{-4}$ 到 $2 \times 10^{-2}$。

所有模型均使用 AdamW~\citep{loshchilov2019decoupled} 进行训练，使用 $1 \times 10^{-4}$ 的恒定学习率，无权重衰减，批量大小为 $256$。数据增强仅包含随机水平翻转。模型权重保持指数移动平均（EMA），衰减率为 $0.9999$；所有报告的结果均使用 EMA 模型。值得注意的是，与监督 ViT 训练不同，DiT \emph{不需要}学习率预热，也\emph{不需要}正则化——在所有模型配置中，训练都高度稳定，没有观察到损失尖峰。

\subsubsection{无分类器引导}
为了产生类别控制的采样，DiT 可以轻松采用\emph{无分类器引导}~\citep{ho2022classifier}。在训练期间，类别标签 $c$ 以一定概率被随机丢弃，并替换为学习到的空嵌入 $\varnothing$，以便模型同时学习条件噪声预测 $\boldsymbol{\epsilon}_\theta(\z_t, c)$ 和无条件噪声预测 $\boldsymbol{\epsilon}_\theta(\z_t, \varnothing)$。在推理时，引导噪声预测计算如下：
\begin{equation}
    \hat{\boldsymbol{\epsilon}}_\theta(\z_t, c) = \boldsymbol{\epsilon}_\theta(\z_t, \varnothing) + s \cdot \big(\boldsymbol{\epsilon}_\theta(\z_t, c) - \boldsymbol{\epsilon}_\theta(\z_t, \varnothing)\big),
    \label{eq:cfg}
\end{equation}
其中 $s > 1$ 是引导尺度（$s = 1$ 恢复标准的无引导采样）。其动机源于贝叶斯法则：由于 $\nabla_{\z_t} \log p(c | \z_t) \propto \nabla_{\z_t} \log p(\z_t | c) - \nabla_{\z_t} \log p(\z_t)$，等式~\eqref{eq:cfg} 中的引导估计有效地将采样轨迹引向高条件似然 $p(c | \z_t)$ 的区域。这项技术对生成质量有显著影响：DiT-XL/2 在没有引导的情况下实现了 $9.62$ 的 FID，而在引导尺度为 $s = 1.5$ 时提高到了 $2.27$——在类别条件 ImageNet $256 \times 256$ 基准上超越了所有先前的扩散模型。

\subsubsection{生成}
一旦去噪网络 $\boldsymbol{\epsilon}_\theta$（即 DiT）训练完成，就可以通过运行学习到的反向过程来生成新图像。从纯噪声 $\z_{t_{\max}} \sim \mathcal{N}(\boldsymbol{0}, \boldsymbol{I})$ 开始，模型通过在每一步采样 $\z_{t-1} \sim p_\theta(\z_{t-1} | \z_t)$ 进行迭代去噪，使用预测的噪声 $\hat{\boldsymbol{\epsilon}}_\theta$ 和协方差 $\boldsymbol{\Sigma}_\theta$ 来参数化反向过渡。经过 $t_{\max}$ 个去噪步骤后，得到的干净潜在变量 $\z_0$ 通过 VAE 解码器解码回图像空间：$\tilde{\x} = \mathcal{D}(\z_0) \in \mathbb{R}^{H \times W \times 3}$。







\section{条件图像表示与生成}
\label{sec:cond_gen}
%\yima{From the above latent diffusion model for images, it is natural to transition to conditioned image generation. Text condition via CLIP or SimCLIP... and DPS implementation, classifier-free, or cross-attention etc. (Robin, Jingfeng, Daniel, and Peter, decide what to feature...)}

% --- draft by Daniel Yeh ---
在前面的小节中，我们建立了生成建模的基本机制：一个自编码器架构，将数据 $\x$ 编码为（压缩且更结构化的）潜在变量 $\z$：
\begin{equation}
\x\xrightarrow{\hspace{1mm} \mathcal{E} \hspace{1mm}} \z
\xrightarrow{\hspace{1mm} \mathcal{D} \hspace{1mm}} \hat{\x}.
\end{equation}

然后，学习一个生成过程来对这些潜在变量的分布 $p(\z)$ 进行建模。在本节中，我们将重点从\textit{无条件生成}转移到\textit{控制}。我们假设表示 $\z$ 已经被学习到（通过第 \ref{ch:conditional-inference} 章中详述的方法）；我们这里的具体任务是使用用户提供的“控制”信号 $\vc$ 来引导采样过程。也就是说，我们的目标是从条件分布 $p(\z | \vc)$ 中采样。在实践中，$\vc$ 可以表示所需样本图像的部分属性：掩码图像、语义掩码、类别信息或描述所需内容的文本标题。

\subsection{任务公式化}
\label{subsec:cond_problem}

形式上，条件生成的目标是学习一个映射，将简单的源分布（通常是高斯噪声）转换为目标条件分布。生成的特征 $\hat{\z}$ 随后被解码为相应的图像 $\hat{\x}$。如下所示，该系统由其输入和输出定义：
\begin{enumerate}
    \item \textbf{随机输入 ($\epsilon$)：} 从标准正态先验 $\epsilon \sim \mathcal{N}(\bm{0}, \I)$ 中采样的噪声变量。这为多样性提供了“种子”，确保我们可以生成同一概念的多个变体。
    \item \textbf{控制信号 ($\vc$)：} 定义所需语义或空间内容的结构化输入。这可以是文本标题（例如，“一辆红色的车”）、语义分割掩码或布局。
    \item \textbf{输出 ($\hat{\x}$)：} 一个高保真数据样本 $\hat{\x} = \mathcal{D}(\hat{\z})$，它与用户的意图 $\vc$ 保持一致，同时保留在自然图像流形上。
\end{enumerate}
我们采用一个特定领域的\textbf{条件编码器}，记为 $\tau_\phi$，将控制信号 $\vc$ 投影到一个有用的嵌入空间中。因此，整个受控生成过程可以形式化地描述为生成器 $G_\theta$ 和解码器 $\mathcal{D}$ 的组合：
\begin{equation}
    \underbrace{\epsilon \sim \mathcal{N}(\bm{0}, \I)}_{\text{Source}}, \vc 
    \xrightarrow{\quad G_\theta(\epsilon, \tau_\phi(\vc)) \quad} 
    \hat{\z} \sim p(\z\mid \vc) 
    \xrightarrow{\quad \mathcal{D}(\hat{\z}) \quad} 
    \hat{\x} \approx \x.
\end{equation}

\paragraph{模态鸿沟。}
虽然该任务在概念上很简单——映射 $(\epsilon, \vc) \to \x$——但它需要桥接不相交的拓扑空间，特别是在文本到图像的生成中。视觉数据 $\x$ 位于连续的高维像素空间 $\mathbb{R}^{H \times W \times 3}$ 中，而条件 $\vc$ 通常位于离散的符号空间（语言）中。没有自然的欧几里得度量可以直接测量距离 $\| \x - \vc \|$。
因此，上面提到的关键组件 $\tau_\phi(\vc)$ 必须对齐这些模态。正如在第 \ref{ch:conditional-inference} 章中所讨论的，现代方法依赖于预训练的对比编码器（如 CLIP）或大型语言模型（如 T5~\citep{raffel2020exploring}）将离散的语言标记映射到连续的语义嵌入空间中。然后，生成主干 $G_\theta$（即去噪 U-Net）必须学习将噪声分布映射到数据分布，并通过特定的融合机制（如交叉注意力）在这些语义嵌入的引导下进行。

\paragraph{一对多逆问题。}
条件控制本质上是适定性不佳的（ill-posed），因为控制信号通常是\textit{欠定的（under-determined）}。文本提示 $\vc=$ “一只狗”提供了语义类别信息，但省略了狗的姿势、光照、背景和毛发图案。因此，逆映射 $\vc \to \x$ 不是一个单点，而是一个复杂的多模态分布。如果我们训练一个简单的确定性回归模型来最小化欧几里得误差 $\|\x - G_\theta(\vc)\|^2$，该模型将收敛到条件均值 $\mathbb{E}[\x|\vc]$——所有可能的狗的平均值——导致模糊、不真实的输出。为了生成高保真结果，我们可以利用生成框架（如第~\ref{sec:auto-encode-img-gen} 节中提到的 U-Net 或 DiT）来对完整的条件分布 $p(\z|\vc)$ 进行建模，从而允许我们采样出满足条件的多样化、清晰的实例。

% To generate high-fidelity results, we
% must utilize the Flow Matching framework (Eq. 8.7 \yima{what is this?}) to model the full
% conditional distribution $p(\z|\vc)$, allowing us to sample diverse,
% sharp instances that satisfy the condition.




% \subsection{The Basic Realization: Text-to-Image Generation}
% \label{subsec:basic_realization}

% To realize the conditional framework defined in Section
% \ref{subsec:cond_problem}, we require a concrete architecture capable
% of bridging the modality gap between discrete text tokens and
% continuous pixel arrays. As illustrated in Figure
% \ref{fig:sd_pipeline}, modern state-of-the-art models (such as Stable
% Diffusion~\citep{rombach2022high}, DALL-E
% 2~\citep{ramesh2022hierarchical}, and
% Imagen~\citep{saharia2022photorealistic}) universally follow a
% pipeline consisting of three decoupled components: (1) a Text Encoder
% to interpret the prompt, (2) a Generation Model to synthesize a
% compressed ``latent representation'' of the image, and (3) a Decoder to map
% this representation back to a high-fidelity image in the pixel domain. We detail the
% implementation of each component below, using the Latent Diffusion
% Model (LDM)~\citep{rombach2022high} as the primary case study.

% \begin{figure}[ht]
% \centering
% \includegraphics[width=\textwidth]{\toplevelprefix/chapters/chapter8/figs/sd_pipeline.pdf}
% \caption{\textbf{The three-stage inference pipeline of Stable Diffusion.}
%     (1) \textbf{Text Encoder} (e.g., CLIP) converts the input prompt
%     $\vc$ into semantic embeddings.
%     (2) \textbf{Generation Model} (a time-conditional U-Net)
%     iteratively denoises a random latent tensor $\z_T$ to produce a
%     clean intermediate representation $\hat{\z}$, conditioned on the
%     text embeddings via cross-attention.
%     (3) \textbf{Decoder} (from a VAE) maps the denoised latent
%     $\hat{\z}_0$ back to the pixel space to synthesize the final
% high-resolution image $\hat{\x}$.}
% \label{fig:sd_pipeline}
% \end{figure}

% \paragraph{The Text Encoder.}
% The first component converts the user's text prompt $\vc$ (e.g., ``A
% fluffy ginger cat on the couch'') into a machine-readable vector
% sequence. Its role is semantic understanding: because the generative
% model's knowledge is limited to its training data, a powerful
% pre-trained encoder is crucial for generalizing to novel concepts or
% rare words. In practice, Stable Diffusion typically utilizes a frozen
% CLIP text encoder~\citep{radford2021learning}, producing a sequence
% of semantically rich embeddings $\tau_\theta(\vc) \in \mathbb{R}^{d_\tau \times M}$ (where $d_\tau$ is typically 768 or 1024). \yima{Please refer to Section \ref{sec:CLIP}, and make sure notation and dimension match!} Recent
% works like Imagen have demonstrated that scaling up this encoder
% (e.g., using T5-XXL \yima{citation?}) yields significantly larger improvements in
% image fidelity than scaling the generative backbone itself.







% \paragraph{The Generation Model (Latent Representation).}
% The core generative process does not produce the final image
% directly. Instead, to avoid the prohibitive computational cost of
% modeling high-resolution pixels (e.g., $1024 \times 1024$), it
% generates a compressed ``intermediate product''—a latent representation $\z$.
% \begin{itemize}
% \item \textbf{Data Format:} The input image $\x \in \mathbb{R}^{H
%     \times W \times 3}$ is compressed by a factor of $f=8$ into a
%     latent tensor $\z_0 \in \mathbb{R}^{h \times w \times 4}$.\yima{How is it ``compressed''? Via a VAE? Do you follow some work here? citation?} For a
%     standard $512^2$ image, this results in a manageable $64 \times
%     64 \times 4$ latent space.
% \item \textbf{Conditioning Mechanism:} The generation model is
%     implemented as a Time-Conditional U-Net. It predicts the noise
%     $\epsilon$ added to the latent $\z_t$, conditioned on the text
%     embeddings via {\em Cross-Attention}. To enable this interaction,
%     the spatial latent features $\phi_i(\z_t)$ \yima{what ``spatial'' mean here? it was not introduced. what is $\phi_i(\cdot)$? how it is specified?} are flattened into a
%     sequence $\mathbb{R}^{N \times d}$ to act as the Query ($Q$),
%     while the text embeddings serve as Key ($K$) and Value ($V$):
%     \begin{equation}
%         Q = W_Q^{(i)} \cdot \phi_i(\z_t), \quad K = W_K^{(i)} \cdot
%         \tau_\theta(\vc), \quad V = W_V^{(i)} \cdot \tau_\theta(\vc).
%     \end{equation}
%     \yima{is $i$ here the index for layer?}
%     This mechanism allows the model to spatially align specific
%     visual regions (e.g., ``cat's eye'') with specific text tokens.
% \end{itemize}

% Training involves a noise prediction task. We sample a random
% timestep $t$, add Gaussian noise $\epsilon$ to the clean latent
% $\z_0$ to create a noisy latent $\z_t$, and train the U-Net backbone
% $\epsilon_\theta$ to predict the added noise. The model minimizes the
% standard reweighted diffusion objective:
% \begin{equation}
% \mathcal{L}_{\text{LDM}} = \mathbb{E}_{\z_0, \vc, \epsilon \sim
% \mathcal{N}(0,1), t} \left[ \| \epsilon - \epsilon_\theta(\z_t, t,
% \tau_\theta(\vc)) \|^2 \right].
% \end{equation}
% During inference (sampling), we start with pure Gaussian noise $\z_T
% \sim \mathcal{N}(\bm{0}, \I)$. In each step, the U-Net predicts the
% noise component, which is then subtracted from the current state to
% gradually refine the latent representation $\z_t \to \z_{t-1}$ until
% a clean latent $\hat{\z}_0$ is formed.



% \paragraph{The Generation Model (Latent Representation).}
% The core generative process does not produce the final image
% directly. Instead, to avoid the prohibitive computational cost of
% modeling high-resolution pixels (e.g., $1024 \times 1024$), it
% generates a compressed "intermediate product"—a latent representation $\z$.
% \begin{itemize}
% \item \textbf{Data Format:} The input image $\x \in \mathbb{R}^{H
%     \times W \times 3}$ is compressed by a factor of $f=8$ into a
%     latent tensor $\z_0 \in \mathbb{R}^{h \times w \times 4}$.
%     This compression is performed by a pre-trained Variational Autoencoder (VAE)~\citep{rombach2022high} that learns a perceptually equivalent but computationally efficient lower-dimensional space.
%     For a standard RGB input image of size $512 \times 512 \times 3$ (i.e., $H=W=512$), this results in a manageable $64 \times
%     64 \times 4$ latent space.

% \item \textbf{Conditioning Mechanism:} The generation model is
%     implemented as a Time-Conditional U-Net. It predicts the noise
%     $\epsilon$ added to the latent $\z_t$, conditioned on the text
%     embeddings via \textbf{Cross-Attention}.
%     To enable this interaction, we treat the intermediate visual features from the U-Net layer $i$ as the Query source, and the text embeddings as the Key/Value source.
%     Let $\phi_i(\z_t) \in \mathbb{R}^{d_\epsilon^{(i)} \times N}$ denote the flattened spatial features at layer $i$ (where $d_\epsilon^{(i)}$ varies by layer, e.g., 320, 640, or 1280 channels), and let $\tau_\theta(\vc) \in \mathbb{R}^{d_\tau \times M}$ denote the sequence of $M$ text embeddings (typically $d_\tau=768$).
%     Since these dimensions differ, we project both into a common subspace of dimension $d$ (head dimension) using learnable weights.

%     We first define the Query ($\boldsymbol{Q}$), Key ($\boldsymbol{K}$), and Value ($\boldsymbol{V}$) matrices using projection weights $\boldsymbol{U}^{(i)}_{Q} \in \mathbb{R}^{d \times d_\epsilon^{(i)}}$ and $\boldsymbol{U}^{(i)}_{K}, \boldsymbol{U}^{(i)}_{V} \in \mathbb{R}^{d \times d_\tau}$:
%     \begin{equation}
%         \boldsymbol{Q} = \boldsymbol{U}^{(i)}_{Q} \phi_i(\z_t), \quad
%         \boldsymbol{K} = \boldsymbol{U}^{(i)}_{K} \tau_\theta(\vc), \quad
%         \boldsymbol{V} = \boldsymbol{U}^{(i)}_{V} \tau_\theta(\vc).
%     \end{equation}
%     Adapting the column-wise formulation from [placeholder], the cross-attention operation is then computed as:
%     \begin{equation}
%         \text{CrossAttn}(\boldsymbol{Q}, \boldsymbol{K}, \boldsymbol{V}) \doteq \boldsymbol{V} \cdot \operatorname{softmax}\left( \frac{\boldsymbol{K}^\top \boldsymbol{Q}}{\sqrt{d}} \right),
%     \end{equation}
%     where the $\operatorname{softmax}$ is defined column-wise as
%     \begin{align}
%         \operatorname{softmax}(\boldsymbol{A})
%         &\doteq \begin{bmatrix} \operatorname{softmax}(\boldsymbol{a}_{1}) & \cdots & \operatorname{softmax}(\boldsymbol{a}_{N}) \end{bmatrix}, \\
%         \forall \boldsymbol{A}
%         &= \begin{bmatrix} \boldsymbol{a}_{1}, \dots, \boldsymbol{a}_{N} \end{bmatrix} \in \mathbb{R}^{M \times N}.
%     \end{align}
%     In this context, the matrix $\boldsymbol{A} = \frac{\boldsymbol{K}^\top \boldsymbol{Q}}{\sqrt{d}}$ represents the alignment scores, and the column-wise softmax normalizes the attention weights over the $M$ text tokens for each of the $N$ visual regions. This mechanism allows the model to spatially align specific visual regions (e.g., ``cat's eye'') with specific text tokens.
% \end{itemize}

% Training involves a noise prediction task. We uniformly sample a random
% timestep $t \in \{1, \dots, T\}$, add Gaussian noise $\epsilon$ to the clean latent
% $\z_0$ to create a noisy latent $\z_t$. To inform the model of the noise level, the scalar timestep $t$ is first mapped to a sinusoidal embedding (encoding low and high frequency components), followed by a Multi-Layer Perceptron (MLP) composed of `Linear` $\to$ `SiLU` $\to$ `Linear` layers. This produces a rich time embedding vector (e.g., 1280-dim) which is added to the feature maps in each ResNet block of the U-Net backbone $\epsilon_\theta$. The model then minimizes the standard reweighted diffusion objective:
% \begin{equation}
% \mathcal{L}_{\text{LDM}} = \mathbb{E}_{\z_0, \vc, \epsilon \sim
% \mathcal{N}(0,1), t} \left[ \| \epsilon - \epsilon_\theta(\z_t, t,
% \tau_\theta(\vc)) \|^2 \right].
% \end{equation}
% During inference (sampling), we start with pure Gaussian noise $\z_T
% \sim \mathcal{N}(\bm{0}, \I)$. In each step, the U-Net predicts the
% noise component, which is then subtracted from the current state to
% gradually refine the latent representation $\z_t \to \z_{t-1}$ until
% a clean latent $\hat{\z}_0$ is formed.






% \paragraph{The Decoder (Image Reconstruction).}
% The final component is a neural network trained to translate the sample latent vector $\hat{\z}$ back into a natural image
% $\hat{\x}$. Unlike the generation model which requires image-text
% pairs, this Decoder is part of a VAE (Variational Autoencoder) that
% can be trained separately on large-scale unlabeled image data.
% During the training phase, the VAE minimizes a reconstruction error
% (combining L1 loss and Perceptual Loss) to ensure the decoder
% $\mathcal{D}$ can accurately map compressed latents back to pixels
% ($\mathcal{D}(\mathcal{E}(\x)) \approx \x$). Once trained, the
% Encoder is discarded from the inference pipeline, and the Decoder is
% frozen. It simply takes the denoised latent $\hat{\z}_0$ produced by
% the Generation Model and upsamples it to the final pixel result $\hat{\x}$.

% \textit{Note on Visualization:} When visualizing the generation for every time-step using Stable Diffusion, users often observe the image evolving from a
% "blurry blob" to a sharp result. This distinctive "blur-to-sharp"
% transition occurs because, at every step of the sampling process, the
% system passes the current noisy latent $\z_t$ through the Decoder to
% visualize progress. Since the Decoder is trained on clean data, it
% interprets the noise in the latent space as high-frequency texture
% information, rendering the noisy latent as a blurry, low-detail image
% 而不是静态噪声。






\subsection{基本实现：文本到图像生成}
\label{subsec:basic_realization}

为了实现\ref{subsec:cond_problem}中定义的条件生成框架，我们需要一个具体的架构来弥合离散文本标记（tokens）和连续像素阵列之间的模态鸿沟。如图~\ref{fig:sd_pipeline}所示，现代最先进的模型（如Stable Diffusion~\citep{rombach2022high}）将这一复杂任务解耦为三个独立的组件：（1）一个\textbf{感知压缩}模块（VAE），用于在像素和潜在变量（latents）之间进行转换；（2）一个\textbf{语义条件}接口（CLIP文本编码器），用于解释提示词（prompts）；以及（3）一个\textbf{潜在生成}主干网络（U-Net），用于合成内容。我们在下文详细介绍每个组件的实现。


\begin{figure}[ht]
\centering
\includegraphics[width=\textwidth]{\toplevelprefix/chapters/chapter8/figs/sd_pipeline.pdf}
\caption{\textbf{Stable Diffusion的三阶段推理流水线。} (1) \textbf{文本编码器}（例如CLIP）将输入提示词$\boldsymbol{c}$转换为语义嵌入。(2) \textbf{生成模型}（时间条件U-Net $\epsilon_\theta$）迭代地对随机潜在张量$\boldsymbol{z}_T$进行去噪，以生成干净的中间表示$\hat{\boldsymbol{z}}$，该过程通过交叉注意力机制以文本嵌入为条件。(3) \textbf{解码器}（来自VAE）将去噪后的潜在变量$\hat{\boldsymbol{z}}$映射回像素空间，以合成最终的高分辨率图像$\hat{\boldsymbol{x}}$。\textit{注：VAE编码器仅在训练期间使用，在此推理流水线中被省略。}}
\label{fig:sd_pipeline}
\end{figure}


\paragraph{组件1：感知压缩（VAE）。}
该架构的第一个支柱是变分自编码器（VAE），其经过训练以抽象掉那些在感知上不重要但在计算上建模成本高昂的高频细节。正如我们在第~\ref{sec:VAE-implement}节中所介绍的，该组件定义了发生生成过程的\textbf{潜在空间}。
\begin{itemize}
    \item \textbf{编码器（$\mathcal{E}$）：} 编码器将输入的RGB图像$\boldsymbol{x} \in \mathbb{R}^{H \times W \times 3}$映射为低维潜在表示$\boldsymbol{z} = \mathcal{E}(\boldsymbol{x}) \in \mathbb{R}^{h \times w \times c}$。遵循\cite{rombach2022high}的方法，它使用了一个经过KL正则化的编码器，下采样因子为$f=H/h=W/w=8$。对于$512 \times 512 \times 3$的标准输入，这将产生一个尺寸为$64 \times 64 \times 4$的紧凑潜在张量。在扩散过程（组件3）的背景下，这种编码表示作为“干净”的真实潜在变量$\boldsymbol{z}_0$。
    
    \item \textbf{解码器（$\mathcal{D}$）：} 解码器学习从潜在空间重建图像的逆映射，记为$\tilde{\boldsymbol{x}} = \mathcal{D}(\boldsymbol{z}) = \mathcal{D}(\mathcal{E}(\boldsymbol{x}))$。其训练目标是最小化$\boldsymbol{x}$和$\tilde{\boldsymbol{x}}$之间的感知差异。在扩散模型的推理阶段，该解码器被冻结，并用于将生成的（去噪后的）潜在向量映射回高分辨率像素空间。
\end{itemize}
请注意，VAE是结合感知损失和基于图块（patch-based）的对抗损失进行预训练的，以确保重建质量。一旦训练完成，其权重将被冻结，并作为扩散模型的静态接口。

\paragraph{组件2：语义条件（文本编码器）。}
为了引导生成过程，我们需要对用户的提示词$\vc$（例如，“一只毛茸茸的橘猫”）进行鲁棒的语义表示。潜在扩散模型通常不从头开始训练文本编码器，而是利用来自CLIP~\citep{radford2021learning}的冻结文本编码器。这个预训练的Transformer将输入文本标记映射为一系列语义嵌入$\tau_\theta(\vc) \in \mathbb{R}^{d_\tau \times M}$，其中$M$是序列长度（例如，77个标记），$d_\tau$是嵌入维度（通常为768或1024）。这些嵌入作为生成主干网络的语义“键（Key）”和“值（Value）”来源。

\paragraph{组件3：潜在生成（时间条件U-Net）。}
核心生成引擎被实现为一个时间条件U-Net $\epsilon_\theta$。与像素空间模型不同，该网络完全在压缩的潜在空间内运行。其架构由一系列用于特征提取的ResNet块和用于条件注入的交叉注意力（Cross-Attention）层组成。
\begin{itemize}
    \item \textbf{输入流：} 模型处理三种不同的信号：
    \begin{enumerate}
        \item \textbf{视觉流：} 噪声潜在变量$\z_t \in \mathbb{R}^{h \times w \times 4}$，代表当前扩散步骤的图像。
        \item \textbf{时间步（$t$）：} 为了告知模型当前的噪声水平，标量时间步$t \in \{1, \dots, T\}$被映射为正弦嵌入（编码低频和高频分量），随后通过一个由`Linear` $\to$ `SiLU` $\to$ `Linear`层组成的多层感知机（MLP）。这产生了一个全局时间嵌入向量（例如，1280维），该向量被注入到U-Net中的每个ResNet块中。
        \item \textbf{上下文流：} 来自组件2的文本嵌入$\tau_\theta(\vc)$。
    \end{enumerate}

    \item \textbf{条件机制（交叉注意力）：} 视觉内容和文本提示词之间的交互是通过交错在整个U-Net中的交叉注意力层来处理的。令$\phi_i(\z_t) \in \mathbb{R}^{d_\epsilon^{(i)} \times N}$表示U-Net第$i$层的展平空间特征（其中$N = h_i \times w_i$是视觉标记的数量，$d_\epsilon^{(i)}$是通道深度，例如320、640或1280），并令$\tau_\theta(\vc) \in \mathbb{R}^{d_\tau \times M}$表示文本嵌入。
    
    由于视觉维度$d_\epsilon^{(i)}$和文本维度$d_\tau$不同，我们使用可学习的权重将两者投影到一个维度为$d$（注意力头维度）的公共子空间中。我们首先定义查询（$\boldsymbol{Q}$）、键（$\boldsymbol{K}$）和值（$\boldsymbol{V}$）矩阵：
    \begin{equation}
        \boldsymbol{Q} = \boldsymbol{U}^{(i)}_{Q} \phi_i(\z_t), \quad
        \boldsymbol{K} = \boldsymbol{U}^{(i)}_{K} \tau_\theta(\vc), \quad
        \boldsymbol{V} = \boldsymbol{U}^{(i)}_{V} \tau_\theta(\vc),
    \end{equation}
    
    其中$\boldsymbol{U}^{(i)}_{Q} \in \mathbb{R}^{d \times d_\epsilon^{(i)}}$且$\boldsymbol{U}^{(i)}_{K}, \boldsymbol{U}^{(i)}_{V} \in \mathbb{R}^{d \times d_\tau}$。借鉴\eqref{sub:image_classification_architecture}中的按列公式，交叉注意力操作计算如下：
    \begin{equation}
        \text{CrossAttn}(\boldsymbol{Q}, \boldsymbol{K}, \boldsymbol{V}) \doteq \boldsymbol{V} \cdot \operatorname{softmax}\left( \frac{\boldsymbol{K}^\top \boldsymbol{Q}}{\sqrt{d}} \right),
    \end{equation}
    其中$\operatorname{softmax}$按列定义为
    \begin{align}
        \operatorname{softmax}(\boldsymbol{A})
        &\doteq \begin{bmatrix} \operatorname{softmax}(\boldsymbol{a}_{1}) & \cdots & \operatorname{softmax}(\boldsymbol{a}_{N}) \end{bmatrix}, \\
        \forall \boldsymbol{A}
        &= \begin{bmatrix} \boldsymbol{a}_{1}, \dots, \boldsymbol{a}_{N} \end{bmatrix} \in \mathbb{R}^{M \times N}.
    \end{align}
    在此背景下，矩阵$\boldsymbol{A} = \frac{\boldsymbol{K}^\top \boldsymbol{Q}}{\sqrt{d}}$表示对齐分数，按列的softmax对每个$N$视觉区域在$M$个文本标记上的注意力权重进行归一化。这种机制允许模型在空间上将特定的视觉区域（例如，“猫的眼睛”）与特定的文本标记对齐。
\end{itemize}

\paragraph{训练目标。}
训练过程将这些组件统一起来。我们采样一张随机图像$\x$，使用冻结的编码器将其压缩为$\z_0 = \mathcal{E}(\x)$，并添加高斯噪声$\epsilon$以产生$\z_t$。然后训练U-Net $\epsilon_\theta$以预测该噪声，条件是文本$\vc$和时间步$t$：
\begin{equation}
\mathcal{L}_{\text{LDM}} = \mathbb{E}_{\mathcal{E}(\x), \vc, \epsilon \sim \mathcal{N}(0,1), t} \left[ \| \epsilon - \epsilon_\theta(\z_t, t, \tau_\theta(\vc)) \|^2 \right].
\end{equation}
在推理期间，编码器被丢弃。该过程从纯噪声$\z_T \sim \mathcal{N}(\bm{0}, \I)$开始，使用U-Net对其进行迭代去噪，最后将生成的$\hat{\z}_0$通过解码器$\mathcal{D}$以合成图像。




\paragraph{评估指标。}
由于条件生成缺乏单一的“正确”真实值（ground truth），我们依赖统计指标来评估性能。我们将真实图像的分布记为$p_r$，将生成图像的分布记为$p_g$。

\begin{itemize}
    \item \textbf{FID（Fréchet Inception Distance）：} 该指标通过测量$p_r$和$p_g$在预训练的Inception-v3网络特征空间中的距离，来量化生成样本的真实性和多样性。我们将这些特征分布近似为多维高斯分布$\mathcal{N}(\bm{\mu}_r, \bm{\Sigma}_r)$和$\mathcal{N}(\bm{\mu}_g, \bm{\Sigma}_g)$。然后将FID分数定义为它们之间的Fréchet距离：
    \begin{equation}
        \text{FID}(p_r, p_g) = \| \bm{\mu}_r - \bm{\mu}_g \|^2 + \text{Tr}\left( \bm{\Sigma}_r + \bm{\Sigma}_g - 2(\bm{\Sigma}_r \bm{\Sigma}_g)^{1/2} \right).
    \end{equation}
    较低的FID表明生成的分布在统计上更接近真实分布，这意味着更高的图像质量和多样性。

    \item \textbf{CLIP分数（CLIP Score）：} 该指标衡量生成图像$\hat{\x}$与其条件提示词$\vc$之间的语义对齐程度。令$E_{\text{img}}$和$E_{\text{txt}}$为第\ref{ch:conditional-inference}章中讨论的预训练CLIP模型的图像和文本编码器。该分数计算为它们嵌入之间的余弦相似度：
    \begin{equation}
        \text{CLIP-Score}(\hat{\x}, \vc) = \max \left( 0, \frac{E_{\text{img}}(\hat{\x}) \cdot E_{\text{txt}}(\vc)}{\|E_{\text{img}}(\hat{\x})\| \|E_{\text{txt}}(\vc)\|} \right) \cdot w,
    \end{equation}
    其中$w$是缩放因子（通常为2.5）。较高的分数表明视觉内容忠实地反映了文本描述。
\end{itemize}

\paragraph{高级文本条件架构（SDXL与Flux）。}
为了解决分辨率和提示词遵循度方面的局限性，高级架构已经超越了基本的LDM。SDXL~\citep{podell2023sdxl}引入了双编码器策略以捕捉更丰富的文本细微差别，它同时从CLIP ViT-L和OpenCLIP ViT-bigG计算嵌入，并将它们沿通道轴拼接。在数据处理方面，SDXL采用了\textit{宽高比分桶（Aspect Ratio Bucketing）}，将具有相似宽高比的训练图像分组到批次中，以防止模型学习到任意裁剪主体。此外，它还实现了\textit{微条件（Micro-Conditioning）}，将图像分辨率和裁剪坐标等元数据嵌入并与文本条件$\vc$拼接，使模型能够明确区分居中和被裁剪的主体。

同时，Flux~\citep{flux2024}代表了从U-Net架构向多模态扩散Transformer（MMDiT）的范式转变。它使用联合注意力（Joint Attention）块在统一的流中处理文本和视觉标记，这得益于旋转位置嵌入（RoPE）和RMSNorm等现代Transformer组件。在数学上，Flux用\textbf{修正流匹配（Rectified Flow Matching）}取代了标准扩散。与预测噪声$\epsilon$的标准扩散不同，流匹配将训练目标简化为预测噪声和数据之间的直线速度场$v_\theta$。我们在连续时间$t \in [0, 1]$上定义该过程，其中$t=0$表示噪声分布（$\z_0$），$t=1$表示数据分布（$\z_1$）。\footnote{请注意，这种时间约定（$0 \to \text{噪声}, 1 \to \text{数据}$）与前面章节中使用的标准扩散符号相反。} 其目标是最小化：
\begin{equation}
    \mathcal{L}_{\text{RFM}} = \mathbb{E}_{t, \z_1, \z_0} \left[ \| v_\theta(\z_t, t, \vc) - (\z_1 - \z_0) \|^2 \right],
\end{equation}
其中$\z_t = t \z_1 + (1-t) \z_0$是插值的潜在状态。模型$v_\theta$学习预测将噪声$\z_0$沿直线轨迹直接传输到数据$\z_1$所需的恒定速度$(\z_1 - \z_0)$。与传统扩散模型的弯曲随机路径相比，这种线性显著提高了采样效率和提示词遵循度。



\subsection{高级条件机制}
\label{subsec:advanced_conditioning}

\begin{figure}[ht]
    \centering
    \includegraphics[width=\textwidth]{\toplevelprefix/chapters/chapter8/figs/cnet_pipeline.pdf}
    \caption{\textbf{ControlNet架构流水线。} Stable Diffusion (SD) U-Net编码器和中间块的一个可训练副本，通过提示编码器$\mathcal{E}_{\text{hint}}$以提示图像为条件，通过零卷积将特征图注入到冻结的SD U-Net解码器中，以实现空间控制。（图像示例来自\url{https://github.com/huggingface/diffusers/blob/main/docs/source/en/using-diffusers/controlnet.md}）。}
    \label{fig:controlnet_pipeline}
\end{figure}


虽然文本条件提供了语义控制，但它难以处理细粒度的空间约束（例如，特定的姿势或边缘图）。为了解决这个问题，ControlNet~\citep{zhang2023adding}引入了一种神经架构，为大型预训练文本到图像扩散模型添加了空间控制，如图~\ref{fig:controlnet_pipeline}所示。

\textbf{架构与连接。}
ControlNet通过克隆预训练扩散模型的可训练参数来运行。具体来说，它创建了Stable Diffusion U-Net的12个编码器块和1个中间块的可训练副本，同时保持原始参数$\Theta$完全锁定（在\texttt{torch.no\_grad()}下执行）。这保留了从数十亿张图像中学习到的生成质量，同时允许模型在小数据集上学习条件任务。

为了匹配潜在特征图$\z \in \mathbb{R}^{h \times w \times c}$的空间分辨率（通常为第~\ref{subsec:basic_realization}节中定义的$64 \times 64$），空间条件$\vc_s$（例如，$512 \times 512$的Canny边缘图）首先由一个轻量级的\textbf{提示编码器（Hint Encoder）}$\mathcal{E}_{\text{hint}}(\cdot)$处理。与VAE不同，该编码器更深但更窄，由\textbf{八个}具有$3 \times 3$卷积核和SiLU激活函数的卷积层组成。它使用三个步长为2的层来下采样分辨率，逐渐增加通道数（$16 \to 32 \to 96 \to 256$），最后投影到模型维度（320个通道）。至关重要的是，该编码器的最后一层是零初始化的。

允许这两个流合并的关键创新是\textbf{零卷积（Zero Convolution）}——一个$1 \times 1$的卷积层$\mathcal{Z}(\cdot; \theta_z)$，其权重$W$和偏置$B$均被严格初始化为零。其连接遵循特定的模式：
\begin{enumerate}
    \item \textbf{输入注入：} 编码后的提示$\mathcal{E}_{\text{hint}}(\vc_s)$\textbf{仅}被添加到ControlNet副本的第一个输入块中。
    \item \textbf{输出注入：} ControlNet副本中每个可训练块的输出通过一个零卷积，然后添加到冻结的Stable Diffusion解码器相应的\textbf{跳跃连接（skip connection）}中。
\end{enumerate}



形式上，考虑U-Net的第$i$个块。令$\boldsymbol{h}_i$表示该层的输入特征图。我们将这个冻结块的操作抽象为函数$\mathcal{F}_i(\cdot; \Theta_i)$。相应地，令$\mathcal{F}_{\text{trainable}}^{(i)}(\cdot; \Theta_{c,i})$表示其在ControlNet分支中的可训练副本。输出$\boldsymbol{y}_i$计算如下：
\begin{equation}
    \boldsymbol{y}_i = \mathcal{F}_i(\boldsymbol{h}_i; \Theta_i) + \mathcal{Z}\left( \mathcal{F}_{\text{trainable}}^{(i)}(\boldsymbol{h}_i + \mathcal{E}_{\text{hint}}(\vc_s); \Theta_{c,i}); \theta_z \right).
    \label{eq:controlnet}
\end{equation}





由于零卷积被初始化为零，控制项在训练开始时消失（$\mathcal{Z}(\cdot) = 0$），确保模型的行为与原始的Stable Diffusion完全一致。然而，权重$W$的梯度是非零的：
\begin{equation}
    \frac{\partial \mathcal{Z}(I; \{W, B\})}{\partial W_{i,j}} = I_{p,i} \neq 0,
\end{equation}
其中$I$是输入特征图。这确保了尽管初始影响为零，但学习在第一次反向传播时立即开始。

\textbf{训练动态与实验。}
ControlNet的训练旨在最小化标准扩散目标，但将额外的空间条件$\vc_s$注入到噪声预测器中。具体的学习目标是：
\begin{equation}
    \mathcal{L} = \mathbb{E}_{\z_0, t, \vc_t, \vc_s, \epsilon \sim \mathcal{N}(0,1)} \left[ \| \epsilon - \epsilon_\theta(\z_t, t, \vc_t, \vc_s) \|^2 \right],
\end{equation}
其中$\vc_t$表示文本提示条件，$\z_t$是噪声潜在变量。观察到的一个独特现象是“突然收敛（sudden convergence）”：模型最初忽略空间条件，然后随着误差急剧下降，突然对齐。为了迫使模型依赖结构信号而不仅仅是文本（文本可能与边缘图相矛盾），一个关键策略涉及\textbf{CFG-Dropout}：在训练期间随机将50\%的文本提示$\vc_t$替换为空字符串。这迫使网络直接从空间条件$\vc_s$中提取语义。经验表明，ControlNet非常高效，与标准微调相比仅需要增加约23\%的GPU内存，却能从仅1,000个样本的小数据集中有效地学习到鲁棒的控制。










\paragraph{主题驱动的个性化。}
虽然基础模型拥有庞大的语义先验，但它们缺乏在不同上下文中一致地合成特定主题实例（例如，用户的某只特定的猫）的能力。为了解决这个问题，DreamBooth \citep{ruiz2023dreambooth}引入了一种“个性化”微调协议。核心挑战是在不发生“灾难性遗忘（catastrophic forgetting）”的情况下，将新主题植入模型的输出域中——即模型失去对一般类别的理解（例如，忘记了普通的猫长什么样），或者遭受“语言漂移（language drift）”（用特定猫的特征覆盖了“猫”这个词）。


\begin{figure}[ht]
    \centering
    \includegraphics[width=\textwidth]{\toplevelprefix/chapters/chapter8/figs/db_pipeline.pdf}
    \caption{\textbf{DreamBooth的训练和推理流水线。} 在\textit{主题驱动的微调}期间，文本到图像模型使用组合损失进行更新：用于学习特定主题（例如，“[V] cat”）的重建损失，以及用于保持通用类别知识（例如，“a cat”）的先验保留损失。在\textit{推理}时，个性化模型可以使用唯一标识符在新的上下文中（例如，“laying on a wooden table”）生成特定主题。}
    \label{fig:dreambooth_pipeline}
\end{figure}



\textbf{标识符与数据策略。}
与需要数百万个数据对的标准训练不同，DreamBooth在用户提供的“少样本（few-shot）”数据集上运行，通常仅包含该主题的3到5张图像。为了将该主题绑定到模型的文本空间，需要一个唯一的标识符$[V]$。作者提出了一种罕见标记（rare-token）选择策略：该方法不是使用随机字符（这可能会被分词为具有强先验的常见字母），而是在分词器的词汇表中搜索具有弱语义关联的罕见标记序列（例如，在T5-XXL分词器范围$\{5000, \dots, 10000\}$内）。该标识符在提示词中始终与粗略的类别名词（Class Noun）配对（例如，“a $[V]$ cat”）。这种“绑定（tethering）”至关重要：它利用了模型对该类别的强大先验（知道如何渲染各种姿势的“猫”），同时将特定的身份特征绑定到标识符$[V]$上。

\textbf{先验保留训练。}
为了防止模型对少样本图像过拟合并失去多样性，DreamBooth采用了一种自生成的特定类别先验保留损失。如图\ref{fig:dreambooth_pipeline}所示，训练过程最小化一个组合目标：
\begin{equation}
    \mathcal{L}_{\text{DB}} = \mathbb{E}_{\z_t, \vc, \epsilon, t} \left[ \underbrace{\| \epsilon - \epsilon_\theta(\z_t, t, \vc_{\text{subj}}) \|^2}_{\text{Reconstruction Loss}} + \lambda \underbrace{\| \epsilon - \epsilon_\theta(\z'_t, t, \vc_{\text{prior}}) \|^2}_{\text{Prior Preservation Loss}} \right].
    \label{eq:dreambooth}
\end{equation}
第一项微调模型权重（通常是整个U-Net和文本编码器），以在唯一提示词$\vc_{\text{subj}}$（例如，“a $[V]$ cat”）的条件下重建主题图像。第二项作为正则化器：它使用模型自身生成的通用类别$\vc_{\text{prior}}$（例如，“a cat”）的样本来监督模型。这迫使模型保留其对类别分布的原始理解，确保标识符$[V]$学习到特定主题的细微差别，而类别名词“cat”保持通用。






\textbf{多概念个性化。}
虽然DreamBooth很有效，但微调整个模型在计算上非常昂贵且占用大量存储空间。Custom Diffusion \cite{kumari2023multi}通过仅更新交叉注意力层中的键（$W_K$）和值（$W_V$）投影矩阵来提高效率，证明了这些特定权重足以编码新概念。


\begin{figure}[ht]
    \centering
    \includegraphics[width=\textwidth]{\toplevelprefix/chapters/chapter8/figs/gen4gen_teaser.pdf}
    \caption{\textbf{用于多概念个性化的Gen4Gen流水线。} 为了克服现有数据集在多主题场景中的局限性，Gen4Gen \cite{yeh2024gen4gen}自动创建了一个高质量的基准测试（MyCanvas）。在这些数据上进行训练显著提升了模型生成包含多个特定主题的复杂场景的能力（例如，“a $[V1]$ cat and a $[V2]$ dog”）。}
    \label{fig:gen4gen_pipeline}
\end{figure}

然而，由于在Stable Diffusion的预训练阶段缺乏用于复杂构图的高质量、文本和图像对齐的训练数据，在单张图像中准确合成“多个”个性化概念（例如，“a $[V1]$ cat and a $[V2]$ dog”）仍然很困难。为了解决这个问题，Gen4Gen \cite{yeh2024gen4gen}提出了一种以数据为中心的解决方案，如图\ref{fig:gen4gen_pipeline}所示。它没有修改模型架构，而是引入了一个生成式数据流水线，利用基础模型构建“MyCanvas”——一个包含复杂多对象场景并配有详细、密集字幕的基准数据集。通过在这些合成数据上进行微调（其中对象身份和空间关系经过明确策划），模型可以显著提高其同时解耦和保留多个个性化主题的能力（例如，特定的猫和狗在复杂场景中互动），而不会遭受身份混合或遗漏的问题。


最近，该领域已转向免微调（Tuning-Free）架构，如PhotoMaker \cite{li2024photomaker}和InstantID \cite{wang2024instantid}。这些方法没有采用迭代优化，而是采用了一种“堆叠（stacking）”策略：一个外部的ID编码器（通常基于人脸识别主干网络）从参考图像中提取高保真身份嵌入，然后通过解耦的交叉注意力或提示词堆叠将其直接注入到扩散过程中。这使得在单次前向传播中能够对任意主题进行零样本（zero-shot）个性化，而无需任何梯度更新。



\subsection{扩展到视频生成}
\label{subsec:video_gen}

通过将视频视为时间上相关的帧序列$\boldsymbol{x} \in \mathbb{R}^{T \times H \times W \times 3}$，条件2D生成的原理自然地扩展到了视频生成。然而，直接对这个高维空间进行建模在计算上是令人望而却步的。因此，现代最先进的方法通常在由VAE（第~\ref{subsec:basic_realization}节）定义的压缩潜在空间内运行，其中视频被表示为张量$\boldsymbol{z} \in \mathbb{R}^{T \times h \times w \times c}$。

\paragraph{时空膨胀（Spatiotemporal Inflation）。}
为了利用2D图像生成器建立的鲁棒视觉先验，诸如VideoLDM~\citep{blattmann2023align}和AnimateDiff~\citep{guo2023animatediff}等架构引入了将2D层“膨胀（inflating）”为伪3D层的概念。这是通过重塑输入张量$\boldsymbol{z}$以在空间和时间处理之间交替来实现的：
\begin{itemize}
    \item \textbf{空间层（冻结）：} 2D主干网络（ResNet和空间注意力）通过将输入重塑为$(B \cdot T, h, w, c)$来独立处理帧。这保留了预训练的纹理和对象生成能力。
    \item \textbf{时间层（可训练）：} 为了对运动进行建模，张量被重塑以隔离每个空间位置的时间轴。我们将空间网格点（$B \cdot h \cdot w$）视为批次维度，将时间（$T$）视为序列长度。然后应用\textbf{时间自注意力（Temporal Self-Attention）}层在帧之间交换信息。
\end{itemize}

在数学上，令$\boldsymbol{u} \in \mathbb{R}^{c \times T}$表示特定空间像素位置随时间变化的特征向量（类似于第~\ref{subsec:basic_realization}节中的空间特征$\phi_i(\z_t)$）。我们使用可学习的投影权重$\boldsymbol{U}^{(temp)}$定义时间查询（$\boldsymbol{Q}$）、键（$\boldsymbol{K}$）和值（$\boldsymbol{V}$）矩阵：
\begin{equation}
    \boldsymbol{Q} = \boldsymbol{U}^{(temp)}_Q \boldsymbol{u}, \quad \boldsymbol{K} = \boldsymbol{U}^{(temp)}_K \boldsymbol{u}, \quad \boldsymbol{V} = \boldsymbol{U}^{(temp)}_V \boldsymbol{u},
\end{equation}
其中$\boldsymbol{U}^{(temp)}_{Q}, \boldsymbol{U}^{(temp)}_{K}, \boldsymbol{U}^{(temp)}_{V} \in \mathbb{R}^{c \times c}$。遵循第~\ref{subsec:basic_realization}节中建立的按列公式，时间自注意力计算如下：
\begin{equation}
    \text{TempAttn}(\boldsymbol{Q}, \boldsymbol{K}, \boldsymbol{V}) \doteq \boldsymbol{V} \cdot \operatorname{softmax}\left( \frac{\boldsymbol{K}^\top \boldsymbol{Q}}{\sqrt{c}} \right) + \boldsymbol{u}.
\end{equation}
这里，项$\boldsymbol{K}^\top \boldsymbol{Q} \in \mathbb{R}^{T \times T}$表示时间注意力图，将每一帧与所有其他帧关联起来。至关重要的是，该注意力块的输出投影层通常是零初始化的。这确保了在训练开始时，时间层的贡献消失，允许模型的行为与预训练的2D生成器完全一致，并逐渐学习时间动态。

\paragraph{交互式世界模型。}
最近的进展已经超越了被动的文本到视频合成，转向交互式世界模拟，有效地将视频生成器视为物理环境的预测模型。形式上，目标是对以智能体动作$\boldsymbol{a}_t$为条件的潜在转移概率进行建模：
\begin{equation}
    p(\boldsymbol{z}_{t+1} \mid \boldsymbol{z}_{1:t}, \boldsymbol{a}_t),
\end{equation}
其中$\boldsymbol{a}_t$表示时间$t$的控制信号（例如，键盘输入或电机命令）。在Genie~\citep{bruce2024genie}等框架中，动作$\boldsymbol{a}_t$被离散化为一个标记，并使用与文本相同的条件机制注入到网络中。例如，$\boldsymbol{a}_t$可以映射为嵌入向量并添加到时间步嵌入中，或者在交叉注意力层中用作键/值对。这使得生成模型能够作为数据驱动的物理引擎运行，通过最小化重建或流匹配目标来预测动作的视觉后果（例如，角色跳跃）。像Sora~\citep{openai2024sora}这样的大规模实现进一步表明，扩展这种基于自回归或扩散的预测会导致涌现出3D一致性和对象持久性。










\section{图像和文本条件的3D对象生成}
\label{sec:Michelangelo}\label{sec:object-generation}
%\yima{Zibo和Sam，请将图像-文本条件的3D形状生成材料（Michelangelo）移至本节。尽量保持自包含，并与本书其余部分的符号和术语保持一致，特别是与下一节基于生成的3D重建保持一致。}

到目前为止，在前面的章节中，我们主要展示了如何学习{\em 二维}（2D）图像分布的表示，以及这种表示所促进的许多基于图像的应用或任务：例如基于图像的对象识别、图像分割或补全，以及图像生成。\footnote{基于前面介绍的技术：视觉Transformer（ViT）~\cite{NIPS2017_3f5ee243}、语言-图像对比学习（CLIP）~\cite{pmlr-v139-radford21a}和潜在扩散模型（VAE-SD）~\cite{rombach2022high}。}

然而，我们的世界是三维（3D）的。\footnote{如果我们考虑到世界是动态的，将时间作为一个额外的维度，那么就是四维（4D）的。} 尽管我们的眼睛只感知世界的2D投影——通过双眼获得的立体对，但我们能够发展出一种视觉记忆，为我们的生活环境和其中的对象表示一个完整的3D世界模型。

本节证明，使用本书中介绍的方法，可以学习自然对象3D形状的分布，然后从该分布中生成新的3D对象形状，条件是输入的2D图像或对象的文本描述。特别是，这里介绍的工作基于Michelangelo项目~\cite{NEURIPS2023_ea1a7f7b}，该项目极大地将我们为2D图像表示和生成引入的技术推广到了3D形状数据\footnote{一般来说，多边形网格和视频数据都可以被称为“3D数据”；然而，在本节中，术语3D数据指的是多边形网格或其等效表示（显式或隐式）。}。

我们首先简要介绍表示3D形状的常用数据形式和相应的数据集。然后，遵循名为\textbf{Michelangelo}的开源项目\cite{NEURIPS2023_ea1a7f7b}：\\
\centerline{\url{https://github.com/NeuralCarver/Michelangelo},}\\
我们提出了一个计算框架、目标和相关的网络架构，试图解决图像和文本条件的3D对象生成任务。最后，我们展示了实验结果和分析以验证所提出的解决方案，并讨论了进一步的改进。

% Introduce basic 3D representations and common used 3D dataset.
\subsection{3D形状表示、数据集与标记化}

\paragraph{3D形状表示}
\label{sec:3D representation}
与3D图像不同，3D形状和数据在本书前面尚未正式介绍。因此，本小节简要介绍3D数据。
希望了解更多细节的读者建议查阅引用文献中的来源。

我们这里所指的3D数据主要意味着连续的真实物理3D形状，这些形状通常在特定的表示空间中被数字化为离散的、可计算的形式。寻找3D形状的良好表示是计算机视觉和计算机图形学领域的核心问题之一。四种最常用且在数学上基本等效的3D形状表示类型是：点云（point clouds）、体素（voxels）、多边形网格（polygon meshes）和隐函数（implicit functions）\footnote{例如距离符号函数（DSF）。}，如图~\ref{fig:3d_representations}所示\footnote{前三列的图像来自\cite{shi2022deep}，最后一列的图像来自https://medium.com/@sunil7545/implicit-vs-parametric-3d-shape-representation-9d4c01c8c60c}。

\begin{figure}[ht]
\centering
\includegraphics[width=\textwidth]{\toplevelprefix/chapters/chapter8/figs/3drep.pdf}
\caption{\textbf{斯坦福兔子（Stanford Bunny）不同3D表示的图示。} (a) 点云；(b) 体素；(c) 多边形网格；(d) 隐函数。}
\label{fig:3d_representations}
\end{figure}

\textbf{点云}通过在对象表面离散采样的点精确地表征对象的3D形状。它们直接记录表面上点的3D空间坐标$(x,y,z)$（可能还有这些点处的一些额外光度信息），如图~\ref{fig:3d_representations}(a) 点云所示。
点云被定义为3D欧几里得空间中的一组离散坐标：
\begin{equation}
\label{eq:pointcloud}
P = \{ \vp_i \in \mathbb{R}^d\ |\ i=1,2,\dots,N \}
\end{equation}
其中$d \geq 3$表示特征维度，可能包括坐标$(x_i,y_i,z_i)$、颜色$(r_i,g_i,b_i)$和表面法线$(n_{xi},n_{yi},n_{zi})$的（拼接）。

\textbf{体素}
表示3D空间中的体积单元，构成了2D像素概念向三维的自然延伸。这种表示是通过将3D空间离散化为结构化网格来获得的。
典型的体素模型将3D空间划分为均匀的立方体单元，其中每个体素在空间上被索引，并存储特定区域的属性，如占用率、密度、材质或颜色，如图~\ref{fig:3d_representations}(b) 体素所示。

\textbf{多边形网格}
可以说是计算机图形学和物理模拟领域中最常用的3D形状表示。它通过离散的顶点、边和多边形面来近似连续表面，其中三角形网格是主要的变体，如图~\ref{fig:3d_representations}(c) 多边形网格所示。
在几何上，多边形网格由$\mathbb{R}^3$中的顶点坐标及其拓扑连接性组成。形式上，网格$M$是一个元组$(V,F)$，其中顶点集$V=\{\vv_i\in\mathbb{R}^3\}_{i=1}^n$表示顶点坐标，面集$F=\{f_j\}_{j=1}^m$定义了拓扑结构。

\textbf{隐函数}
通过标量函数的水平集来表示表面。这种方法为复杂表面的数学描述提供了一个统一的框架（图~\ref{fig:3d_representations}(d) 隐函数）。
标准的隐式公式为：
\begin{equation}
\label{eq:implicit_function}
F: \mathbb{R}^3 \rightarrow \mathbb{R}, \quad F(x,y,z) = d,
\end{equation}
其中$(x,y,z)$表示查询点的坐标，$d$表示查询点与感兴趣表面之间的距离。\footnote{请注意，我们在\Cref{sec:manifold-distance}中使用了相同的概念来描述高维空间中分布的低维支撑。} 因此，感兴趣的表面被隐式定义为：
\begin{equation}
\label{eq:implicit_function_surface}
\mathcal{S} = \{ (x,y,z) \in \mathbb{R}^3 \mid F(x,y,z)=0. \}
\end{equation}
从3D形状的隐函数表示到多边形网格的转换通常可以通过等值面提取方法（如移动立方体算法（marching cubes）~\cite{10.1145/280811.281026}）来实现。

从数学上讲，所有四种形状表示在表示能力方面基本等效。
然而，每种数据形式可能具有不同的计算优势或局限性。因此，对于不同的任务，通常会偏好某种表示。在本节中，我们将采用传统的点云表示，并展示如何基于这种形式的数据学习3D形状的分布。在下一节中，我们将展示如何使用（广义的）体素表示来学习3D形状和姿态的分布。

\paragraph{3D数据集。}
对于上述提到的每种表示，都有许多公开可用的数据集。它们具有不同的数量和质量。在本节中，我们将使用ShapeNet~\cite{chang2015shapenet}来研究联合的形状-图像-文本表示及其生成性能。ShapeNet提供了55个类别中约55,000个制造网格。
每个网格都有一个类别标签和相应的文本，例如细粒度类别或创建者给出的简短描述。
为了准备三元组数据（3D形状、图像、文本），提供的文本通过两种方式进行增强。
首先，形状标签和相应的描述以“a 3D model of (shape tag), in the style of (description)”或“a 3D model of (shape tag), (description)”的格式组合。
然后，受ULIP~\cite{10203465}的启发，利用包含65个预定义短语的多个模板在训练期间提供更多文本信息。
至于图像数据，每个网格在四个相机姿态下进行渲染，并通过深度条件的ControlNet~\cite{10377881}增强和改善渲染多样性。

\paragraph{3D形状数据标记化。}
\label{sec:3d_tokenization}
如上所述，3D数据以各种形式表现，包括点云、体素、多边形网格和隐式神经表示。从根本上说，点云、体素和网格作为连续3D流形的离散近似。在计算机图形学和工程中，多边形网格尤为普遍，因为它们显式地保留了表面拓扑，在视觉保真度和几何结构之间提供了有效的平衡。

然而，将现代深度学习架构——特别是Transformer——直接应用于这些原始表示面临着重大挑战。与具有规则网格结构、易于分割成图块（patches）的2D图像不同，3D数据通常是稀疏、非结构化和不规则的。为了在3D中利用序列建模和自注意力机制的强大功能，我们必须弥合原始几何数据与Transformer所需的离散、有序序列之间的鸿沟。这个过程被称为\textit{3D标记化（3D Tokenization）}。

\begin{itemize}
\item \textbf{点云标记化。} 
原始点云通常被标记化为局部组以捕获几何上下文。这个过程通常从最远点采样（FPS）~\cite{eldar1997farthest}开始，这是一种迭代算法，从原始集合$\mathcal{P}$中选择一个点子集$\{x_1, \dots, x_m\}$，使得每个选定的点到先前选定点集合的距离最大化：
$x_{i} = \arg\max_{x \in \mathcal{P}} \min_{j=1}^{i-1} \|x - x_j\|$。
在FPS之后，通过k最近邻（k-NN）对点进行分组以形成局部簇。
每个簇作为一个标记，为了处理簇内点的排列不变性，应用了如PointNet~\cite{qi2017pointnet}这样的轻量级网络。
在数学上，对于一个点簇$\{p_1, \dots, p_k\}$，PointNet通过应用共享的多层感知机（MLP）$h$，然后进行对称聚合操作（通常是最大池化）来近似对称函数$f$：
\begin{equation}
    \boldsymbol{t} = \gamma \left( \max_{i=1}^k \{ h(p_i) \} \right)
\end{equation}
其中$\boldsymbol{t}$是生成的特征标记，$\gamma$是后处理MLP。这将非结构化点云转换为特征向量序列。



% {\color{blue} Update: We have added the implement details of Farthest Point
%     Sampling, and use math formula to describe the process of PointNet, and the output token format. The following is the original version, and is replaced by the above paragraph.}
    
% Point clouds are typically
%     tokenized using grouping techniques such as Farthest Point
%     Sampling (FPS) \yima{what is this?} followed by k-Nearest Neighbors (k-NN). Each
%     cluster of points is treated as a token, often processed by a
%     permutation-invariant networkwork (such as PointNet \yima{reference? what does it do exactly, in mathematical terms?} ) to extract a
%     feature vector.
    
\item \textbf{体素标记化。} 
类似于2D图像中视觉Transformer（ViT）的图块嵌入策略（见第\ref{sub:contrastive_learning_architecture}章），体素网格可以被划分为非重叠的3D子网格。给定一个体素网格$V \in \mathbb{R}^{D \times H \times W \times C}$，我们提取大小为$P \times P \times P$的图块。每个图块被展平为向量$v_p \in \mathbb{R}^{P^3 \cdot C}$，并通过可学习的线性投影映射到标记维度$d_{model}$。结果是一个标记序列$Z \in \mathbb{R}^{L \times d_{model}}$，其中$L = \frac{DHW}{P^3}$。虽然简单直接，但这种方法继承了体素分辨率的立方复杂度（$O(N^3)$），通常会为高分辨率形状造成内存瓶颈。


% Similar to 2D image patches, voxel
%     grids can be partitioned into 3D sub-grids or ``cublets." While
%     straightforward, this method often suffers from the cubic
%     complexity of voxel resolution (the $O(N^3)$ memory bottleneck). 
    
%     Extending the 2D patch embedding strategy to 3D, 
% we partition the input voxel grid $V \in \mathbb{R}^{D \times H \times W \times C}$ into 
% non-overlapping 3D patches or ``cubelets'' of size $p \times p \times p$. 
% These cubelets are flattened into 1D vectors and mapped via a learnable linear projection 
% to the token dimension $d_{model}$. 
% The resulting output is a sequence of tokens with length $L = \frac{DHW}{p^3}$. 
% However, this approach suffers from the cubic complexity of voxel resolution, 
% as the sequence length $L$ scales with $O(N^3)$, causing significant memory bottlenecks.
% {\color{blue} Update: We have added the implement details of the output tokens and the reference section about the patch embedding strategy in Vision Transformers for 2D images.}
    
%     \yima{Not clear at all, what are the output tokens? how are they generated in what format? If this process is similar to tokenizing 2D image, please refer to an earlier section that describes tokenizing 2D image, if that is what you mean here.}


\item \textbf{网格标记化。} 
由于网格涉及顶点和面的类图结构，对其进行标记化具有挑战性。
与图像不同，网格缺乏规范的顺序。为了实现序列建模，常见的方法涉及显式序列化，其中原始几何和拓扑元素被直接量化并排序为序列。

% Current methods typically serialize the mesh elements into a sequence of discrete tokens to enable autoregressive modeling (where the model predicts the next vertex/face token based on previous ones to generate valid topology).

% Prominent approaches like MeshGPT~\cite{siddiqui2024meshgpt} serialize the mesh into two interleaved sequences: (1) a **Vertex Sequence**, where vertex coordinates are quantized and ordered (often via z-order curve sorting), and (2) a **Face Sequence**, where faces are represented as indices pointing to the vertex list. Alternatively, graph-based methods may treat local mesh patches (geodesic neighborhoods) as tokens, encoding them via Graph Convolutional Networks (GCNs) into feature vectors.

以PolyGen~\cite{nash2020polygen}为例，网格被表示为一组顶点$V = \{v_1, \dots, v_n\}$（其中每个$v_i \in \mathbb{R}^3$）和面$F = \{f_1, \dots, f_m\}$（其中每个$f_j$是顶点索引的元组）。
模型通过将联合分布$P(V, F)$分解为两个连续阶段来学习它：预测量化坐标的顶点模型（Vertex Model），以及预测顶点索引的面模型（Face Model）
\begin{equation}
    P(V, F) = P(V) \cdot P(F \mid V) = \prod_{i=1}^{n} P(v_i \mid v_{<i}) \cdot \prod_{j=1}^{m} P(f_j \mid f_{<j}, V).
\end{equation}
至关重要的是，这种分解决定了对应于每个阶段的两种不同的标记格式。
对于顶点模型（$P(V)$），标记是量化的坐标值，连续坐标被离散化（例如，转换为8位整数），从而从大小为256的固定词汇表中创建一个标记序列。
相反，对于面模型（$P(F \mid V)$），标记是顶点索引（指针）。
虽然有效，但对这些原始图元进行操作通常会导致极长的序列长度，从而限制了可扩展性。

% {\color{blue} Update: We take the PolyGen as a representative work and add the implement details of the mesh tokenization process. The following is the original version, and is replaced by the above paragraph.}
% Tokenizing meshes presents unique
%     challenges due to their irregular connectivity and graph-like
%     structure. Common approaches involve traversing the surface to
%     serialize faces or vertices into sequences, or decomposing the
%     mesh into discrete primitives. These methods are particularly
%     prevalent in autoregressive mesh generation \yima{what is this, how it is related to the tokenization we are discussing here}, where the objective
%     is to generate not just the 3D geometry, but also a valid, clean topology.\yima{This paragraph does NOT describe exactly how to tokenize meshes at all, not even mention any existing method! }

\end{itemize}

虽然这些显式标记化策略有效地离散化了连续流形，但它们直接在高维观察空间中操作，通常导致过长的序列长度和计算冗余。为了解决这些可扩展性限制并实现更高的压缩比，该领域已倾向于基于学习的离散表示。以下方法没有直接对原始几何图元进行标记化，而是首先将3D数据投影到紧凑的低维潜在空间中，在其中聚合语义信息，从而有效地将几何重建与结构建模解耦。

\textbf{标准潜在向量量化。} 该方法通过向量量化（VQ）学习一个离散密码本（codebook）来表示形状。一个代表性的方法是AutoSDF~\cite{mittal2022autosdf}。在这个框架中，高维输入（例如，体素网格$X$或隐函数场$F_\theta$）被编码器$E$压缩为潜在向量的低维空间网格$Z \in \mathbb{R}^{h \times w \times d \times c}$。
然后，每个向量$z_i \in Z$被替换为来自可学习密码本$\mathcal{C}=\{c_k\}_{k=1}^K$的最近邻：
\begin{equation}
q(z_i) = \arg\min_{c_k \in \mathcal{C}} | z_i - c_k |_2^2.
\end{equation}
这将形状简化为一系列离散的密码本索引，在保留空间结构的同时大幅降低了维度。这些离散索引作为后续Transformer处理的标记。然而，这种基于规则网格的方法通常统一处理占用空间和空白空间，导致计算冗余。

% {\color{blue} Update: We add more details of latent vector quantization process and wew specify that the 3D shape is converted to a sequence of tokens.}
% \yima{Then each 3D shape is converted to a single vector, instead of a sequence of tokens?}

\textbf{从结构优化到语义对齐。}
为了解决规则网格的冗余问题，诸如3DILG~\cite{zhang20223dilg}等方法提出了不规则潜在网格（Irregular Latent Grids）。
3DILG没有使用将标记浪费地分配给空白空间的密集体素网格，而是将形状定义为一组漂浮在连续空间中任意自适应位置的潜在代码，通常集中在形状的表面上。
具体来说，给定表示为点云的输入形状，该方法首先将其编码为一组特征-位置对$\mathcal{P} = \{(x_i, f_i)\}_{i=1}^M$，其中$x_i \in \mathbb{R}^3$是通过最远点采样（FPS）选择的关键点，$f_i \in \mathbb{R}^C$是由基于PointNet~\cite{qi2017pointnet}的编码器提取的其关联的连续特征向量。将这个非结构化集合转换为标记序列涉及三个关键步骤：
\begin{itemize}\item \textbf{坐标量化。} 首先，对于通过最远点采样（见点云标记化）选择的一组$M$个关键点，连续空间坐标$x_i \in \mathbb{R}^3$被离散化为8位整数：
\begin{equation}
x_i \rightarrow (x_{i,1}, x_{i,2}, x_{i,3}), \quad \text{where } x_{i,k} \in {0, 1, \dots, 255}.
\end{equation}
这一步有效地将连续位置转换为离散的、可标记化的单元。

\item \textbf{特征量化。} 同时，与每个点关联的局部特征向量通过向量量化（VQ）被压缩为离散的密码本索引$z_i$：
\begin{equation}
    z_i \in \{0, 1, \dots, D-1\},
\end{equation}
其中$D$表示密码本的大小。

\item \textbf{序列构建。} 为了促进使用Transformer进行自回归生成，必须对非结构化的元组集合进行序列化。3DILG采用了基于量化坐标的字典序排序。最终形状被表示为一个结构化序列$\mathcal{S}$，它显式地耦合了位置坐标和局部几何特征信息：
\begin{equation}
    \mathcal{S} = \big\{ (x_{0,1}, x_{0,2}, x_{0,3}, z_0), \dots, (x_{M-1,1}, x_{M-1,2}, x_{M-1,3}, z_{M-1}) \big\}.
\end{equation}
\end{itemize}
与密集网格相比，这种混合公式在保留高保真表面细节的同时，大幅减少了序列长度。

虽然3DILG优化了标记的空间结构，但它在很大程度上仍然是一种几何压缩方法。在下文中，我们详细阐述了Michelangelo，这项工作通过解决语义鸿沟推进了这一范式。它将3D形状编码为一系列潜在标记，这些标记不仅在几何上紧凑，而且与2D视觉-语言模型（例如CLIP~\cite{pmlr-v139-radford21a}）的特征空间显式对齐，从而促进了高质量的条件生成。

% Introduce task and challenges
\subsection{任务}
我们这里的目标是学习一个条件3D形状生成模型，该模型能够从学习到的分布中采样3D形状，并受给定条件的约束，例如给定图像或某些文本，如图\ref{fig:michelangelo_network}所示。

遵循潜在扩散模型~\cite{rombach2022high}，为了充分利用生成模型的能力，首先训练一个VAE将3D数据样本压缩到潜在空间中。
通常，一个好的潜在表示应该具备两个关键特征：对于3D数据样本，它必须既高度可压缩又能忠实重建，并且有利于生成模型的训练。
然而，与2D图像数据相比，3D数据仍然稀缺\footnote{即使在Objaverse-XL~\cite{10.5555/3666122.3667676}发布之后，这种情况依然存在。}。
因此，为3D数据学习有效的潜在表示极具挑战性。图\ref{fig:michelangelo_network}展示了Michelangelo项目提出的框架的总体图。我们将在下面解释每个组件的实现细节。
\begin{figure}[t]
\centering
\includegraphics[trim=0cm 0cm 0cm 0cm,
clip=true,width=.99\linewidth]{\toplevelprefix/chapters/chapter8/figs/michelangelo_network.pdf}
\caption{
    \textbf{先对齐后生成（Alignment-before-generation）流水线}。该方法包含两个模型：形状-图像-文本对齐变分自编码器（SITA-VAE）和对齐形状潜在扩散模型（ASLDM）。SITA-VAE由四个模块组成：图像编码器、文本编码器、3D形状编码器和3D形状解码器。编码器将输入对编码到对齐空间中，3D形状解码器在给定来自对齐空间的嵌入的情况下重建3D形状。ASLDM将图像或文本条件映射到对齐的形状潜在空间，以采样高质量的3D形状嵌入，随后由3D形状解码器将其重建为高保真3D形状。
}
\label{fig:michelangelo_network}
\end{figure}



% Introduce network and loss func
\subsection{架构与目标}

\paragraph{形状-图像-文本对齐变分自编码器。}

受CLIP~\cite{pmlr-v139-radford21a}的启发，Michelangelo项目采用的方法通过对比损失将语义信息注入潜在空间，以缓解仅从3D数据训练VAE时观察到的问题，这构成了形状-图像-文本对齐变分自编码器（SITA-VAE）。

SITA-VAE包含四个组件：预训练且固定的CLIP图像编码器$\mathcal{E}_{\mathrm{i}}$和CLIP文本编码器$\mathcal{E}_{\mathrm{t}}$，可训练的3D形状编码器$\mathcal{E}_{\mathrm{s}}$和神经场解码器$\mathcal{D}_{\mathrm{s}}$。
CLIP图像编码器和文本编码器将2D图像$\vI \in \R^{H \times
W \times 3}$和标记化文本$\vT \in \R^{d_{{\mathrm{t}}} \times
L_{\mathrm{t}}}$作为输入，并生成图像标记
$\vE_{\mathrm{i}} \in \R^{d \times (1 + L_{\mathrm{i}})}$和文本标记$\vE_{\mathrm{t}} \in \R^{d \times L_{\mathrm{t}}}$，其中
$(1+L_{\mathrm{i}})$和$L_{\mathrm{t}}$分别是图像标记$\vE_{\mathrm{i}}$和文本标记$\vE_{\mathrm{t}}$的序列长度。
% Advantage
该方法利用了来自CLIP的预训练图像编码器和文本编码器。这两个编码器在大型图像-文本对上进行训练，具有足够的鲁棒性以捕获良好对齐的视觉-语言空间，这通过对比学习进行多模态对齐后丰富了3D形状表示的语义。

% 3D shape encoder
\textbf{3D形状编码器}旨在提取强大的特征表示以有效地表征每个3D形状。为了实现这一点，首先从3D形状的表面采样点云$\vP \in \R^{(3 + C) \times N}$，其中$N$表示点的数量，$C$表示额外的点特征，如法线或颜色。

为了更好地捕获高频几何细节，我们利用傅里叶位置编码将每个点的低维空间坐标$\boldsymbol{v} \in \mathbb{R}^3$映射到高维频率空间。该编码$\gamma(\boldsymbol{v})$被定义为一系列频率呈指数增加的正弦函数：$$\gamma(\boldsymbol{v}) = [\sin(2^0\pi \boldsymbol{v}), \cos(2^0\pi \boldsymbol{v}), \dots, \sin(2^{L-1}\pi \boldsymbol{v}), \cos(2^{L-1}\pi \boldsymbol{v})],$$
其中$L$是频带的数量。
为了构建最终的输入表示，我们显式地将输入点云$\vP \in \mathbb{R}^{(3+C) 
\times N}$分解为空间坐标$\boldsymbol{v}$和辅助特征$\boldsymbol{f} \in \mathbb{R}^{C \times 
N}$。然后将高维位置嵌入$\gamma(\boldsymbol{v})$与保留的特征$\boldsymbol{f}$拼接，并通过可学习的线性层进行投影。形式上，3D形状编码器输入$\boldsymbol{X} \in \mathbb{R}^{d \times N}$的获取方式如下：$$\boldsymbol{X} = \text{Linear}([\gamma(\boldsymbol{v}); \boldsymbol{f}])$$其中$[\cdot ; \cdot]$表示沿特征维度的拼接操作。


% Next, a linear layer is used to project the concatenation of
% the Fourier positional encoded point clouds $\vP$ to the 3D shape
% encoder input $\vX \in \R^{d \times N}$. 
% \yima{What does ``Fourier positional encoded'' mean? How is $\X$ related to $\vP$, exactly?} {\color{blue} Update: We have added the formula for Fourier positional encoding and provided a mathematical expression describing the process of transforming from $\vP$ to $\X$.}

为了高效处理这个密集的特征序列$\vX$，
% \yima{Which data, $\vP$ or $\X$?} {\color{blue} Update: We have changed the expression to avoid the ambiguity.} 
编码器采用了基于Perceiver的架构~\cite{pmlr-v139-jaegle21a}。与随输入长度呈二次缩放的标准Transformer不同，Perceiver采用潜在瓶颈机制将处理复杂度与输入大小解耦。模型没有直接对输入点进行标记化以进行自注意力计算，而是利用了一组固定的、小型的可学习潜在查询$\vQ \in \R^{(1+L_{\mathrm{s}}) \times 
d}$，其中$1 + L_{\mathrm{s}}$是查询标记的数量。在机制上，交叉注意力层将高维输入空间投影到这个紧凑的潜在空间中。形式上，遵循第~\ref{subsec:basic_realization}章中的公式，令$\boldsymbol{U}^{(temp)}_{Q}, \boldsymbol{U}^{(temp)}_{K}, \boldsymbol{U}^{(temp)}_{V}$为可学习的投影矩阵。初始潜在表示$\vZ_0$计算如下：
$$\vZ_0 = \text{Cross-Attention}
(\vQ, \vX) = (\boldsymbol{U}^{(temp)}_{V}\vX) \cdot \text{softmax}\left(\frac{(\boldsymbol{U}^{(temp)}_{K}\vX)^\top(\boldsymbol{U}^{(temp)}_{Q}\vQ)}
{\sqrt{d_k}}\right).$$
这里，输入点云$\vX$作为键和值，而可学习的查询$\vQ$作为查询。该操作将可变长度的输入$N$压缩为固定长度的潜在序列$\vZ_0 \in \R^{(1+L_{\mathrm{s}}) \times 
d}$。

在Michelangelo的背景下，这些查询被结构化以捕获分层语义：它们由一个具有高级语义的全局头标记$\vQ_g \in \R^{1 \times d}$和$L_{\mathrm{s}}$个包含低级几何结构信息的局部标记$\vQ_l \in \R^{L_{\mathrm{s}} \times d}$组成。然后，使用几个自注意力块来迭代改进特征表示并获得最终的形状嵌入$\vE_{\mathrm{s}} 
\in \R^{(1+L_{\mathrm{s}}) \times d}$。

% Drawing inspiration from
% previous transformer-based architectures for point cloud
% understanding~\cite{pmlr-v139-jaegle21a}, the 3D shape encoder is
% built on a perceiver-based transformer. Specifically, a
% cross-attention layer is used to inject the 3D shape information from
% the input $\vX$ into a series of learnable query tokens $\vQ \in
% \R^{d \times (1+L_{\mathrm{s}})}$, where $1 + L_{\mathrm{s}}$ is the
% length of query tokens $\vQ$, consisting of one global head token
% $\vQ_g \in \R^{d}$ with high-level semantics and $L_{\mathrm{s}}$
% local tokens $\vQ_l \in \R^{d \times L_{\mathrm{s}}}$ containing
% low-level geometric structure information. Then, several
% self-attention blocks are used to iteratively improve the feature
% representation and obtain the final shape embeddings,
% $\vE_{\mathrm{s}} \in \R^{d \times (1+L_{\mathrm{s}})}$.

% \yima{Probably more detail or illustration for the above description about how to tokenize 3D shape data, say point clouds, through a perceiver. How is it trained, separately or jointly, with what architecture and data etc.? How is it related to the ``3D tokenization'' methods introduced earlier?} {\color{blue} Update: We have supplemented the description of how the input is tokenized through the perceiver-based architecture, and the entire process is trained jointly. }

% Learn aligned latent
\textbf{3D形状、图像和文本之间的对齐}在SITA-VAE和条件生成模型中起着至关重要的作用。由于3D数据比图像和文本数据小一个数量级，为了在3D形状、图像和文本之间学习到更好对齐的表示，3D形状编码器被约束为接近预对齐的视觉-语言空间，该空间利用对比学习策略在具有丰富图像和文本表示的大规模图像-文本对上进行了预训练。考虑3D形状$\vX$、图像$\vI$和标记化文本$\vT$的输入三元组。三元组编码器通过将提取的形状标记$\vE_{\mathrm{s}}$、图像标记$\vE_{\mathrm{i}}$和文本标记$\vE_{\mathrm{t}}$投影为三个具有相同维度的向量，生成相应的形状嵌入$\ve_{\mathrm{s}}$、图像嵌入$\ve_{\mathrm{i}}$和文本嵌入$\ve_{\mathrm{t}}$，表示为：$\ve_{\mathrm{s}} = \mathcal{F}_{\mathrm{s}}(\vE_{\mathrm{s}}),
\ve_{\mathrm{i}} = \mathcal{F}_{\mathrm{i}}(\vE_{\mathrm{i}})$，以及$\ve_{\mathrm{t}} = \mathcal{F}_{\mathrm{t}}(\vE_{\mathrm{t}})$，其中$\mathcal{F}_{\mathrm{s}}$是可学习的形状嵌入投影器，图像嵌入投影器$\mathcal{F}_{\mathrm{i}}$和文本嵌入投影器$\mathcal{F}_{\mathrm{t}}$在训练和推理期间是预训练且冻结的。对比损失为：

\begin{equation}
\begin{aligned}
    \mathcal{L}_{(\mathrm{s},\mathrm{i})} &= - \frac{1}{2} \sum\limits_{(j,k)}
    \left(
        \log \frac{\exp (\ve_{\mathrm{s}}^j \ve_{\mathrm{i}}^k) }{
        \sum\limits_l \exp (\ve_{\mathrm{s}}^j \ve_{\mathrm{i}}^l) } +
        \log \frac{\exp (\ve_{\mathrm{s}}^j \ve_{\mathrm{i}}^k) }{
        \sum\limits_l \exp (\ve_{\mathrm{s}}^l \ve_{\mathrm{i}}^k) }
    \right), \\
    \mathcal{L}_{(\mathrm{s},\mathrm{t})}  &= - \frac{1}{2} \sum\limits_{(j,k)}
    \left(
        \log \frac{\exp (\ve_{\mathrm{s}}^j \ve_{\mathrm{t}}^k) }{
        \sum\limits_l \exp (\ve_{\mathrm{s}}^j \ve_{\mathrm{t}}^l) } +
        \log \frac{\exp (\ve_{\mathrm{s}}^j \ve_{\mathrm{t}}^k) }{
        \sum\limits_l \exp (\ve_{\mathrm{s}}^l \ve_{\mathrm{t}}^k) }
    \right), \\
\end{aligned}
\end{equation}
其中$(j,k)$表示训练批次中的正样本对，并且由于利用了来自CLIP的预训练编码器，模型不受$\mathcal{L}_{(\mathrm{i}, \mathrm{t})}$约束的限制。

\textbf{3D形状解码器}$\mathcal{D}_{\mathrm{s}}$将形状嵌入$\vE_{\mathrm{s}}$作为输入，以高质量重建3D神经场。KL散度损失$\mathcal{L}_{KL}$用于促进生成过程，以保持潜在空间为连续分布。
% pre-fc, post-fc
此外，利用投影层将潜在变量从维度$d$压缩到较低维度$d_0$以获得紧凑表示。然后，使用另一个投影层将采样的潜在变量从维度$d_0$转换回高维度$d$，用于重建3D形状的神经场。
% decode neural field
与编码器一样，解码器模型也建立在具有交叉注意力机制的Transformer之上。给定场中的查询3D点$\vx \in \R^{3}$及其相应的形状潜在嵌入$\vE_{\mathrm{s}}$，解码器迭代计算交叉注意力以预测查询点的占用率$\mathcal{O} (\vx)$。训练损失表示为：

\begin{equation}
\mathcal{L}_r = \mathbb{E}_{\vx \in \R^3} [\text{BCE}(
\mathcal{D}(\vx \mid \vE_{\mathrm{s}} ), \mathcal{O}(\vx) )],
\end{equation}
其中BCE是二元交叉熵损失，训练形状-图像-文本对齐变分自编码器（SITA）的总损失写为：

\begin{equation}
\mathcal{L}_{SITA} = \lambda_c (\mathcal{L}_{(\mathrm{s},\mathrm{i})}
+ \mathcal{L}_{(\mathrm{s},\mathrm{t})}) + \mathcal{L}_r +
\lambda_{KL} \mathcal{L}_{KL}.
\label{eq:sita_vae}
\end{equation}

\paragraph{对齐形状潜在扩散模型。}
在训练SITA-VAE之后，获得了3D形状、图像和文本之间的对齐空间，以及一个3D形状编码器和解码器，它们将3D形状压缩为低维形状潜在嵌入，并将形状潜在嵌入高质量地重建为神经场。基于潜在扩散模型（LDM）~\cite{rombach2022high}在文本到图像生成中的成功（它在计算开销和生成质量之间取得了平衡），在对齐空间上提出了一种形状潜在扩散模型，以学习从2D图像或文本到3D形状潜在嵌入的更好概率映射。通过利用对齐空间和形状潜在扩散模型，可以生成更符合视觉或文本条件输入的高质量3D形状。

对齐形状潜在扩散模型（ASLDM）建立在类似U-Net的Transformer~\cite{10203178}之上，旨在拟合形状潜在嵌入的分布，并伴随一个自编码器用于将数据样本编码到潜在空间并在给定采样潜在变量的情况下重建数据样本。通过在潜在空间中学习，潜在扩散模型在计算上是高效的，并且利用这种紧凑的表示使模型能够更快地拟合目标分布。具体来说，模型$\epsilon_{\theta}$专注于在条件$\vC$下生成形状潜在嵌入$\vE_{\mathrm{s}}$，条件$\vC$由CLIP图像或文本编码器表示。遵循LDM~\cite{rombach2022high}，目标是
\begin{equation}
\mathcal{L} = \mathbb{E}_{\vE_{\mathrm{s}}, \boldsymbol{\epsilon} \sim
\mathcal{N}(\Zero,\vI),t} \left[ \Vert \boldsymbol{\epsilon}
- \epsilon_\theta(\vE_{\mathrm{s}}^{(t)}, \vC, t) \Vert^2_2 \right],
\end{equation}
其中$t$从$\{1, ..., T\}$中均匀采样，$\vE_{\mathrm{s}}^{(t)}$是$\vE_{\mathrm{s}}^{(0)}$的噪声版本。在推理期间，在采样高斯噪声后，模型逐渐对信号进行去噪，直到达到$\vE_{\mathrm{s}}^{(0)}$。遵循无分类器引导（CFG）~\cite{ho2021classifierfree}，条件潜在扩散模型使用无分类器引导进行训练。在训练阶段，条件$\vC$以10\%的固定概率随机转换为空集$\emptyset$。然后，使用条件和无条件样本的线性组合执行采样：
\begin{equation}
\epsilon_\theta ( \vE_{\mathrm{s}}^{(t)}, \vC, t ) = \epsilon_\theta(
\vE_{\mathrm{s}}^{(t)}, \emptyset, t ) + \lambda
( \epsilon_\theta( \vE_{\mathrm{s}}^{(t)}, \vC, t ) -\epsilon_\theta(
\vE_{\mathrm{s}}^{(t)}, \emptyset, t ) ),
\label{eq:cfg-michelangelo}
\end{equation}
其中$\lambda$是用于权衡采样保真度和多样性的引导尺度（guidance scale）。

\subsection{实验与分析}

\textbf{指标。}
交并比（IoU）用于反映重建的准确性。提出了两个新指标用于评估3D形状生成方法。第一个是形状-图像分数（SI-S），它使用3D形状编码器和图像编码器提取相应的形状和图像嵌入，并计算这两种模态的余弦相似度。另一个是形状-文本分数（ST-S），它在对齐的形状嵌入和文本嵌入空间中计算生成的3D形状与条件文本输入之间的相似度。这两个指标都评估了结果与其相应条件之间的相似度。此外，预训练的ULIP~\cite{10203465}和SITA都用于计算SI-S和ST-S，分别记为SI-S (ULIP)、ST-S (ULIP)、SI-S (SITA)和ST-S (SITA)。此外，遵循Point-E~\cite{nichol2022point}中引入的P-IS和P-FID指标，使用预训练的PointNet++~\cite{NIPS2017_d8bf84be}计算点云的类似Inception Score~\cite{NIPS2016_8a3363ab}和FID~\cite{NIPS2017_8a1d6947}，以评估生成的3D形状的多样性和质量。

\textbf{基线（Baselines）。}
我们将Michelangelo方法与其他方法进行比较，包括OccNet~\cite{8953655}、ConvOcc~\cite{10.1007/978-3-030-58580-8_31}、IF-Net~\cite{9156292}、3DILG~\cite{NEURIPS2022_894ca1c4}以及3DS2V的可学习查询版本~\cite{10.1145/3592442}，在重建任务上验证模型在ShapeNet数据集~\cite{chang2015shapenet}上给定形状嵌入恢复神经场的能力。对于条件生成阶段，选择了两种最近强大的3D生成方法3DILG和3DS2V作为基线。它们的形状表示模块首先在ShapeNet和3D Cartoon Monster的混合数据集上进行微调。然后，使用相同的协议重新训练3DILG和3DS2V的文本和图像条件生成模型。

\textbf{数值比较}。
数值结果报告在表~\ref{tab:michelangelo_recon}和表~\ref{tab:michelangelo_gen}中。表~\ref{tab:michelangelo_recon}显示，Michelangelo在55个整体类别上具有最佳的重建性能，超越了其他方法。所选类别的结果进一步证明，该模型可以忠实地重建55个类别中每个类别的3D形状。表~\ref{tab:michelangelo_gen}报告了条件3D形状生成的数值结果。Michelangelo在所有SI-S和ST-S指标上均取得了最佳性能，表明它可以将来自图像或文本的信息映射到其相应的3D形状信息，以生成高保真结果。此外，P-FID证明该模型可以产生高质量的形状标记以生成逼真的3D形状，而P-IS表明了样本的多样性。具体而言，左侧四列显示Michelangelo在图像条件生成上超越了基线，证明它可以更好地将视觉信息映射到3D形状。右侧四列验证了文本条件生成的生成质量。由于与2D图像相比，自然语言通常提供有限且抽象的信息，因此学习一个将文本信息映射到3D形状的模型具有挑战性。然而，得益于在对齐潜在空间上的训练，Michelangelo显著改善了文本条件生成，如表~\ref{tab:michelangelo_gen}的右侧列所示，这反映了该模型有效地将自然语言信息映射到3D形状，并生成了多样化且高质量的结果。

\begin{table}[th]
\centering
\resizebox{1\columnwidth}{!}
{
    \begin{tabular}{rcccccccc}
        \hline
        & Overall & Selected & Table & Chair & Airplane & Car & Rifle
        & Lamp   \\ \hline
        OccNet~\cite{8953655}
        & 0.825 & 0.81 & 0.823 & 0.803 & 0.835 & 0.911 & 0.755 & 0.735  \\
        ConvOccNet~\cite{10.1007/978-3-030-58580-8_31}
        & 0.888 & 0.873 & 0.847 & 0.856 & 0.881 & 0.921 & 0.871 & 0.859  \\
        IF-Net~\cite{9156292}
        & 0.934 & 0.924 & 0.901 & 0.927 & 0.937 & 0.952 & 0.914 & 0.914 \\
        3DILG~\cite{NEURIPS2022_894ca1c4}
        & 0.950 & 0.948 & 0.963 & 0.95 & 0.952 & 0.961 & 0.938 & 0.926 \\
        3DS2V(LQ)~\cite{10.1145/3592442}
        & 0.955 & 0.955 & \textbf{0.965} & 0.957 & 0.962 & 0.966 &
        0.947 & 0.931 \\
        Ours
        & \textbf{0.966} & \textbf{0.964} & \textbf{0.965} &
        \textbf{0.966} & \textbf{0.966} & \textbf{0.969} &
        \textbf{0.967} & \textbf{0.95} \\ \hline
    \end{tabular}
}
\caption{\textbf{IoU重建比较的数值结果（$\uparrow$，值越大越好）}。结果表明，Michelangelo在55个类别中表现最佳。所选类别的结果进一步证明该模型可以忠实地重建每个类别。}
\label{tab:michelangelo_recon}
\end{table}

\begin{table}[th]
\centering
\resizebox{1\columnwidth}{!}{%
    \begin{tabular}{rcccccccc}
        \hline
        & \multicolumn{4}{c}{Image-Conditioned} &
        \multicolumn{4}{c}{Text-Conditioned} \\
        \cmidrule(lr){2-5}\cmidrule(lr){6-9}
        & SI-S (ULIP)$\uparrow$ & SI-S (SITA)$\uparrow$   &
        P-FID$\downarrow$ & P-IS$\uparrow$ & ST-S (ULIP)$\uparrow$ &
        ST-S (SITA)$\uparrow$  & P-FID$\downarrow$ & P-IS$\uparrow$ \\ \hline
        3DILG
        & 9.134  & 11.703 & 4.592 & 12.247 & 10.293 & 6.878  & 10.283 & 12.921\\
        3DS2V
        & 13.289 & 15.156  & 2.921 & 12.92 & 12.934 & 9.833  & 5.704 & 13.149\\
        Ours
        & \textbf{13.818} & \textbf{15.206}  & \textbf{1.586} &
        \textbf{13.233} & \textbf{16.647} & \textbf{13.128}  &
        \textbf{2.075} & \textbf{13.558}\\ \hline
    \end{tabular}%
}
\caption{\textbf{条件生成比较的数值结果}。结果表明，Michelangelo取得了最佳的生成性能。SI-S和ST-S表明，该模型通过有效地将条件信息映射到其相关的3D形状，生成了高保真结果。此外，P-FID反映了该模型生成了最逼真的3D形状，而P-IS表明生成的样本是多样化的。$\uparrow$表示较大值更好，$\downarrow$则相反。}
\label{tab:michelangelo_gen}
\end{table}

\textbf{视觉比较}。
图像/文本条件下的3D形状生成的视觉比较如图~\ref{fig:michelangelo_img_cond_gen} 和图~\ref{fig:michelangelo_text_cond_gen} 所示。
图~\ref{fig:michelangelo_img_cond_gen} 表明，3DILG~\cite{NEURIPS2022_894ca1c4} 在自回归生成过程中更关注全局形状，其结果缺乏对3D形状细节的刻画。虽然 3DS2V~\cite{10.1145/3592442} 生成了更多3D形状的细节，但其结果表面存在不连续或噪点。此外，当给定条件映射到复杂物体、精密机械或罕见怪物时，这两种方法都难以生成完整的形状。
图~\ref{fig:michelangelo_text_cond_gen} 展示了文本条件生成的视觉比较。在上半部分行中，展示了给定简单和抽象概念的结果，而在下半部分行中，展示了给定详细文本（如目标形状确定性部分的描述）的结果。
与上述观察结果类似，3DILG~\cite{NEURIPS2022_894ca1c4} 生成的形状表面过于平滑且细节较少，而 3DS2V~\cite{10.1145/3592442} 在不连续的物体表面上产生的细节较少。相比之下，Michelangelo 生成了符合给定概念或详细描述的正确形状，并在平滑表面上呈现出精致的细节。

\begin{figure}[th]
\centering
\includegraphics[width=.98\linewidth]{\toplevelprefix/chapters/chapter8/figs/michelangelo_img_cond_gen.pdf}
\caption{
    \textbf{图像条件生成比较的视觉结果}。该图表明，3DILG~\cite{NEURIPS2022_894ca1c4} 生成了过于平滑的表面并缺乏形状细节，而 3DS2V~\cite{10.1145/3592442} 生成的细节很少，且形状表面存在噪点和不连续性。与基线方法相比，Michelangelo 生成了平滑的表面并刻画了形状细节。请放大以查看更多视觉细节。
}
\label{fig:michelangelo_img_cond_gen}
\end{figure}

\begin{figure}[th]
\centering
\includegraphics[width=.98\linewidth]{\toplevelprefix/chapters/chapter8/figs/michelangelo_text_cond_gen.pdf}
\caption{
    \textbf{文本条件生成比较的视觉结果}。
    在前两行中，模型使用抽象文本进行了测试，结果表明只有 Michelangelo 生成了符合目标文本的3D形状，具有平滑的表面和精细的细节。
    最后两行展示了给定包含详细描述的文本的结果，这进一步证明了 Michelangelo 能够捕获全局条件信息和局部信息，以生成高保真度的3D形状。关键词以红色突出显示；请放大以查看更多视觉细节。}
\label{fig:michelangelo_text_cond_gen}
\end{figure}

\textbf{在对齐空间中训练生成模型的有效性}。我们进行了视觉比较，以消融在对齐空间中训练生成模型的有效性，如图~\ref{fig:michelangelo_ablate_align} 所示。上半部分的样本是由在对齐空间中训练的生成模型生成的，而下半部分的样本是由未在对齐空间中训练的生成模型生成的。结果表明，上半部分的样本符合给定的文本，而下半部分的样本则不符合，这表明在对齐空间中训练生成模型能够产生高保真度的样本。

\begin{figure}[ht]
\centering
\includegraphics[width=.98\linewidth]{\toplevelprefix/chapters/chapter8/figs/michelangelo_1aligned_space_for_generation.pdf}
\caption{
    \textbf{消融研究：在对齐空间中训练生成模型的有效性}。该图展示了关于在对齐空间中训练生成模型有效性的消融研究的视觉比较。与基于条件文本的下半部分样本相比，上半部分样本在语义上更接近条件，这表明了在对齐空间中训练生成模型的有效性。
}
\label{fig:michelangelo_ablate_align}
\end{figure}

\textbf{视觉-语言模型的有效性}。除了著名的视觉-语言模型（VLM）CLIP~\cite{pmlr-v139-radford21a} 之外，我们还引入了另一种视觉-语言模型（VLM）SLIP~\cite{10.1007/978-3-031-19809-0_30} 来训练 SITA-VAE，以进行全面比较。首先，评估了视觉-语言模型对 SITA-VAE 重建能力的影响，结果如图~\ref{fig:michelangelo_ablate_vlm} 所示。结果表明，与 CLIP 结合的模型取得了最佳性能。然后，评估了视觉-语言模型对齐多模态空间能力的影响。选择了标准和零样本分类任务来反映视觉-语言模型的影响。请注意，分类是通过特征匹配操作执行的，其中向 SITA-VAE 提供多个3D形状和短语；它返回3D形状与每个短语之间的相似度作为分类结果，这表明多模态空间越对齐，分类准确率就越高。结果表明，与 CLIP 结合的模型取得了最佳性能。

\textbf{可学习查询嵌入的影响}。使用与上述相同的实验对可学习查询嵌入进行了消融研究，结果表明，使用 512 个可学习查询嵌入在重建和分类上取得了最佳性能。

\begin{figure}
\centering
\includegraphics[width=1.\linewidth]{\toplevelprefix/chapters/chapter8/figs/michelangelo_2aligned_space_for_generation.pdf}
\caption{
    \textbf{消融研究：视觉-语言模型的有效性以及可学习查询嵌入的影响}。该图展示了关于视觉-语言模型有效性以及可学习查询嵌入影响的消融研究。根据该表，具有 CLIP 和 512 个可学习查询嵌入的模型实现了最佳的重建和分类性能，表明其恢复3D形状和对齐多模态空间的能力。
}
\label{fig:michelangelo_ablate_vlm}
\end{figure}

\paragraph{最近邻分析。}
为了验证模型是学会了基于给定条件生成结果，还是仅仅记住了数据集的训练集，我们遍历了训练集以检索生成样本的最近邻，如图~\ref{fig:michelangelo_nearestneighbor} 所示。具体而言，展示了每个样本的三个最近邻以供比较。根据可视化结果，生成的3D形状与检索到的每个最近邻都不同，这证明了模型是在“幻觉”出新的3D形状，而不是对训练集过拟合。

\begin{figure}[th]
\centering
\includegraphics[width=1.\linewidth]{\toplevelprefix/chapters/chapter8/figs/michelangelo_nearestneighbor.pdf}
\caption{
    \textbf{最近邻分析}。我们遍历了整个训练集，为生成的3D形状寻找三个最近邻，结果表明我们的模型能够基于给定图像生成新颖的3D形状，而不是记住特定的形状。
}
\label{fig:michelangelo_nearestneighbor}
\end{figure}

\paragraph{讨论。}
尽管 Michelangelo 项目表明我们在生成3D物体方面可以取得优异的结果，但它仍然存在一些局限性。首先，它需要大量真实3D形状样本来训练和学习分布，而此类3D形状数据的获取成本要高得多，且可用数据集通常比2D图像数据集小几个数量级。正如我们在上一章末尾的 \Cref{sec:3D-visual-model} 中所提到的，我们的大脑能够仅通过2D图像学习完整的3D（或4D）视觉模型，因为此类数据获取起来更加自然和经济。因此，最终，为了真正扩展生成式3D模型，我们必须开发更先进的方法，使我们能够仅从（多视图）2D图像中（例如通过可微渲染）学习形状表示（在3D形状-图像-文本对齐空间内）。
此外，由于我们将每个3D形状表示为占用场（occupancy field），因此需要将3D网格转换为水密（watertight）网格，这不可避免地会降低3D网格的原始质量。

\section{基于生成的3D形状与姿态重建}\label{sec:CUPID}\label{sec:shape-and-pose}

% refactor
\newcommand{\pose}{\boldsymbol{\theta}}
\newcommand{\object}{\mathcal{O}}
\newcommand{\Image}{\boldsymbol{I}}
\newcommand{\image}{\vI}
% \newcommand{\mask}{\boldsymbol{M}^{\textrm{occ}}}
\newcommand{\render}{\boldsymbol{\mathcal{P}}}
\newcommand{\latent}{\boldsymbol{z}}
\newcommand{\timestep}{t}
\newcommand{\latentt}{\boldsymbol{z}_{\timestep}}
\newcommand{\latentzero}{\boldsymbol{z}_{0}}
\newcommand{\noise}{\boldsymbol{\epsilon}}
\newcommand{\velocity}{\boldsymbol{v}}
\newcommand{\parameter}{\phi}
\newcommand{\encoder}{\varphi}
\newcommand{\point}{\boldsymbol{x}}
\newcommand{\pixcoord}{\boldsymbol{u}}
\newcommand{\feat}{\boldsymbol{f}}
\newcommand{\sflow}{\boldsymbol{\mathcal{G}_{S}}}
\newcommand{\lflow}{\boldsymbol{\mathcal{G}_{L}}}
\newcommand{\transpose}{\operatorname{T}}
\newcommand{\Grid}{\mG}
\newcommand{\sgrid}{\mS}
\newcommand{\numpts}{L}
% \newcommand{\res}{N}
\newcommand{\sres}{D}
\newcommand{\vaeencoder}{\psi}
\newcommand{\trellis}{{\textsc{TRELLIS}}\xspace}
% \newcommand{\dino}{\textsc{DINO}\xspace}
\newcommand{\interp}{\operatorname{BilinearInterp}}
\newcommand{\slatenc}{\operatorname{SlatEncoder}}
\newcommand{\twotothree}{$2$D $\to$ $3$D\xspace}
\newcommand{\threetotwo}{$2$D $\to$ $3$D\xspace}
\newcommand{\low}{\mathrm{l}}
\newcommand{\high}{\mathrm{h}}



% \paragraph{基于生成的3D重建}
我们发展视觉的根本原因是为了学习感知我们的三维（3D）环境。我们的大脑投入了很大一部分区域——{\em 视觉皮层}（visual cortex）——来执行这项任务。它使我们能够从眼睛感知到的2D图像中准确高效地重建3D物体的形状及其空间关系\footnote{例如相对于我们视角的相对姿态，或物体之间的相对姿态。}。在数学上理解原则上如何从（多张）2D图像中完全重建3D几何结构，有着丰富而悠久的历史。感兴趣的读者可以参考教科书 \cite{MaY2003} 以获取关于这一重要主题的详细论述。 

\begin{figure}[th]
\centering
\includegraphics[width=0.8\linewidth]{chapters/chapter8/figs/cupid_definition.pdf}
\caption{\textbf{任务：从单张图像推断具有相对于相机姿态的规范姿态的物体3D形状。}
    左图：输入图像。中图：输出恢复了物体的规范3D模型（具有形状和纹理等内在属性）以及生成输入的精确相机姿态（外在属性）。
    右图：从估计的姿态重新渲染3D模型，忠实地再现了输入。}
\label{fig:cupid_teaser}
\end{figure}

然而，正如我们在上一章 \ref{ch:conditional-inference} 的 \ref{sec:3D-visual-model} 节中所提到的，理想情况下，3D重建应该以贝叶斯方式进行，就像我们的大脑所做的那样。我们对3D环境的感知似乎毫不费力，主要是因为我们基于丰富的“视觉记忆”来重建3D场景并识别3D物体。更准确地说，我们通过多年积累的经验学习了3D环境及其内部物体的分布。这种分布是相当低维且结构化的，因此即使我们只看到了它的一小部分（例如从单一视角），我们也能正确且高效地补全场景及其内部物体的3D几何结构，如图 \ref{fig:cupid_teaser} 中的示例所示。从概念上讲，这种“补全”缺失信息的能力与自然图像如何从一小部分图块中补全非常相似，正如我们在 \ref{sec:image_completion} 节中所看到的那样。

作为一个分布，我们的3D视觉记忆不仅是低维的，而且是高度结构化的。例如，我们的大脑似乎将场景的语义（物体的身份）存储在{\em 下颞叶}皮层（IT cortex）中，而将场景的3D空间信息存储在{\em 海马体}（hippocampus）中。我们的大脑不仅重建场景的3D几何结构，\footnote{例如以不同但在很大程度上等效的形式表示的3D几何结构：点云、网格、符号距离函数、神经辐射场或高斯泼溅（Gaussian splatters）。}而且将其内容解析为单个物体，以及它们与观察者之间和物体相互之间的空间关系。这种结构化的“什么在哪里（what-is-where）”表示使我们能够从{\em 自我中心（egocentric）、物体中心（object-centric）或异中心（allocentric）}的视角在场景内进行空间推理。这些能力对于在场景中导航、交互和操作内部物体等任务是必要且至关重要的。

因此，如果学习3D场景模型的目标是支持与环境及其内部物体的物理交互，\footnote{这也是智能体或机器人的所谓“具身智能（Embodied AI）”的目标。}我们必须学习一种与上述视觉记忆具有相同属性和特征的3D环境表示。




\subsection{任务与目标}
\label{sub:cupid_objective}


\paragraph{从单一2D视图获取物体3D形状与姿态}
在本节中，我们迈出了第一步，通过说明一个具有挑战性且看似适定性不佳（ill-posed）的问题来展示如何实现这一点：{\em 从单张2D图像中推断物体的完整3D形状和姿态，包括自我中心和物体中心的姿态。} 这个问题本身就处于计算机视觉的核心，旨在逆转图像形成的物理过程。图像 $\image$ 是渲染过程的结果，
\[
\image = \render(\object, \pose),
\]
其中 $\object$ 表示在规范的、与视角无关的坐标系中的3D物体，$\pose$ 表示相机视角。仅给定 $\image$，将物体的内在属性（例如形状、纹理）与外在的相机姿态解耦的任务从根本上来说是具有歧义的。


以前的方法通常过度简化了这种歧义。生成模型通常优先考虑规范物体 $\object$ 而牺牲相机姿态 $\pose$，导致重建结果无法忠实地与输入视图对齐。相反，传统的以视图为中心的方法固定了姿态，这限制了它们生成或“修复（in-paint）”被遮挡区域并将重建的物体无缝粘贴回原始图像的能力。


我们通过 \textbf{Cupid}~\citep{huang2025cupid} (\url{https://cupid3d.github.io}) 解决了这个问题，这是一个开源框架，联合建模了3D物体和相机姿态的分布（图~\ref{fig:cupid_teaser}）。通过同时生成3D模型并推理视角，Cupid 确保输出可以准确地\textit{粘贴}到原始图像中。这种方法将单视图重建重新表述为条件采样问题，通过将观测数据与学习到的生成先验相结合来实现几何保真度，从而实现高效的物体级和场景级重建（图~\ref{fig:cupid_results}）。

\begin{figure}
    \centering
    \includegraphics[width=0.95\textwidth]{chapters/chapter8/figs/cupid/figures/results.pdf}
    \caption{\textbf{从单张测试图像进行\textit{生成式3D重建}的结果。}给定输入图像（左上），Cupid 估计相机姿态（左下）并重建3D模型（右下），重新渲染输入（右上）。它对尺度、位置和光照的变化具有鲁棒性，同时保留了精细的细节，并支持组件对齐的场景重建（底行）。所有结果都是通过学习模型的全连接采样在几秒钟内生成的。 }
    \label{fig:cupid_results}
    \vspace{-5mm}
\end{figure}

\paragraph{任务公式化与目标。}
我们将生成式3D重建公式化为估计联合后验分布的任务
\[
p(\object, \pose \mid \image)
\]
在观测值 $\image$ 是物体 $\object$ 从姿态 $\pose$ 渲染而来的约束下。为了对这种条件分布进行建模，我们采用了基于流（flow-based）的生成模型。具体来说，我们首先将3D物体和相机姿态映射到统一的体积潜在表示中
\(\latent = \encoder(\object, \pose)\)。
然后，我们使用修正流（Rectified Flow）~\citep{lipmanflow} 模型来学习该潜在变量的条件生成。该模型定义了数据分布样本 $\latentzero$ 和噪声分布样本 $\noise$ 在归一化时间区间 $\timestep \in [0,1]$ 上的线性插值：
\begin{equation}
\latentt = (1 - \timestep) \latentzero + \timestep \noise.
\end{equation}
生成过程涉及通过学习一个随时间变化的速度场来逆转这一轨迹
\[
\velocity(\latentt, \image, \timestep) = \nabla_{\timestep} \latentt
\]
该速度场在输入图像 $\image$ 的条件下，将噪声样本传输到数据流形。我们使用神经网络 $\velocity_{\parameter}$ 对该速度场进行参数化，并使用条件流匹配（CFM）目标~\cite{lipmanflow} 对其进行训练：
\begin{equation}
\mathcal{L}_{\mathrm{CFM}}(\parameter) \doteq \mathbb{E}_{\timestep,
\latentzero, \noise} \left\| \velocity_{\parameter}(\latentt, \image,
\timestep) - (\noise - \latentzero) \right\|_{2}^{2}.
\end{equation}
通过从这个学习到的模型中采样，我们可以生成一个潜在代码 $\latent$，然后可以将其解码为最终的3D物体和相机姿态。

\subsection{3D形状表示与标记化}


3D形状生成的一个核心挑战是找到一种既紧凑又富有表现力、能够实现高保真重建的标记化（tokenization）方法。虽然 \textbf{Michelangelo}~\citep{zhao2023michelangelo} 使用了学习到的潜在向量表示——将形状编码为 1024 个标记的序列——但这些向量缺乏显式的3D空间坐标。在这个项目中，我们寻求一种保留显式3D位置信息的潜在表示，使我们能够将形状标记直接与图像像素对齐。通过建立这种空间对应关系，我们可以利用局部像素信息来增强生成细节。因此，\textbf{Cupid} 采用了如 \ref{sec:3d_tokenization} 节所述的基于3D体素的潜在表示和标记化方法。在这里，我们用带有特征的稀疏体素来表示形状，这些特征通常由 RGB 颜色或来自2D基础模型的更高级的 DINO 特征~\citep{DINOv2} 定义。然后，我们训练一个解码器，将这种表示映射回由 SDF 参数化的形状，以及由高斯泼溅~\citep{kerbl20233d} 参数化的外观。 

具体来说，给定一组位于坐标 $\boldsymbol{x}_i$ 的稀疏体素特征 $\boldsymbol{F} = \{\boldsymbol{f}_i\}_{i=1}^N$，我们采用两个不同的解码器 $\mathcal{D}_{geo}$ 和 $\mathcal{D}_{app}$，将几何结构与外观解耦。解码过程公式化为：
\begin{equation}
s_i = \mathcal{D}_{geo}(\boldsymbol{f}_i), \quad \mathcal{G}i = \mathcal{D}_{app}(\boldsymbol{f}_i),
\end{equation}
其中 $s_i \in \mathbb{R}$ 表示用于表面重建的符号距离场（SDF）值，$\mathcal{G}_i = \{\boldsymbol{c}_i, \alpha_i, \boldsymbol{s}_i, \boldsymbol{q}_i\}$ 表示体积渲染所需的3D高斯参数集——包括颜色、不透明度、缩放和旋转。然后，我们将使用稀疏体素特征作为潜在流匹配的潜在表示~\citep{rombach2022high}。

\subsection{架构与实现}
\label{sub:3d_reconstruction_architecture}
由于我们使用3D体素进行物体表示，如 \ref{sec:3d_tokenization} 节所述，网格的很大一部分由不必要的空白空间组成，因此简单地应用扩散会导致过多的内存消耗。为了缓解这个问题，\textbf{Cupid} 架构被设计为一个级联的两阶段流模型，逐步细化3D重建~\citep{xiang2025structured}。第一阶段生成粗略的几何形状并估计相机姿态，而第二阶段生成高保真度的几何结构和外观。该框架的一个关键创新在于相机姿态的显式表示及其在生成过程中的集成。整体流程如图 \Cref{fig:cupid_architecture} 所示。

\paragraph{解耦的物体与姿态表示。}

正如在 \ref{sec:3D representation} 节中所讨论的，3D物体 $\object$ 被张量化为基于稀疏体素的表示：
\begin{equation}
    \object \triangleq \{\point_i, \feat_i\}_{i=1}^{\numpts},
\end{equation}
其中 $\point_i \in \mathbb{R}^3$ 表示活跃体素的坐标，$\feat_i$ 是编码局部几何结构和外观的特征向量。这些特征是通过3D VAE压缩多视图 DINO~\citep{oquab2023dinov2} 特征的密集3D网格得出的。随后训练 VAE 解码器将 $\object$ 重建为符号距离场（SDF）以进行网格提取。这种3D潜在表示如图~\ref{fig:cupid_architecture} 中部蓝色块的上半部分所示。

相机姿态 $\pose$ 由其投影矩阵
$\boldsymbol{P} = \boldsymbol{K}[\boldsymbol{R}|\boldsymbol{t}]$ 表示，为了促进联合生成，我们以一种新颖的方式对其进行了重新参数化。我们没有使用紧凑的向量，而是将姿态表示为3D到2D对应关系的密集场，我们将其称为\textit{立方体内像素分布（in-cube pixel distribution）}：
\[
\pose \triangleq \{\point_i, \pixcoord_i\}_{i=1}^{\numpts},
\]
其中 $\pixcoord_i = (u_i, v_i) \in [0,1]^2$ 是对应于3D体素中心 $\point_i$ 的归一化2D像素坐标。这些对应关系是通过透视投影获得的：
\[
\pixcoord_i = \pi(\boldsymbol{P}, \point_i).
\]
这组对应关系可以可视化为3D UV立方体，其中 $(u,v)$ 值就像3D网格上依赖于视角的颜色。这在图 \ref{fig:cupid_architecture} 中部蓝色块的底部进行了说明。

给定生成的一组对应关系，通过使用最小二乘求解器~\citep{abdel2015direct} 求解{\em Perspective-n-Point}（PnP）问题来恢复全局相机矩阵 $\boldsymbol{P}^{*}$：
\begin{equation}
\boldsymbol{P}^{*} = \argmin_{\boldsymbol{P}} \sum_{i=1}^{\numpts}
\big\|\pi(\boldsymbol{P},\point_i) - \pixcoord_i\big\|^2.
\label{eq:pnp}
\end{equation}
这种公式化优雅地将联合物体-姿态生成问题转化为生成以物体为中心的3D形状/外观，以及相关的以观察者/相机为中心的物体姿态（通过 UV 坐标）。这在图 \ref{fig:cupid_architecture} 右侧红色块的左下角进行了说明。


\begin{remark}[结构化3D建模]
请注意，此处为形状和姿态选择的体素类型表示与前面 \Cref{sec:Michelangelo} 中研究的条件形状生成任务中采用的点云类型表示不同。在那个任务中，目标是生成一个只需要与给定图像（或文本）强相关的物体。它不需要物体（在其规范姿态下）与给定图像（视角）之间进行精确的姿态对齐。然而，对于这里考虑的任务，我们希望“重建”一个形状、外观和姿态都忠实地与给定图像一致的物体。因此，内部表示必须同时显式地建模物体的规范形状和相对姿态。请注意，这种将物体（规范）形状/外观与其被观察的姿态“解耦”的做法，密切模仿了前面提到的大脑中 IT 皮层和海马体各自的作用。正如我们稍后将看到的，这种解耦的表示将使我们能够获得一个结构化且一致的场景3D模型，该场景可能包含来自多个视角的多个物体。我们认为，这样结构化的3D模型对于我们在场景中进行以视角为中心或以物体为中心的推理、\footnote{例如，回答诸如“哪个物体在哪个物体的前面、后面、左边或右边，或者在哪个物体的上面”等问题。}生成以及与场景的交互是必要的。
\end{remark}

\begin{figure*}[t]
\centering
\includegraphics[width=\textwidth]{chapters/chapter8/figs/cupid_method.pdf}
\caption{\textbf{Cupid 两阶段生成式重建流程。}给定输入图像，第一阶段（$\sflow$）生成粗略的占用立方体和 UV 立方体，后者编码了3D到2D的对应关系。PnP 求解器从这些对应关系中恢复相机姿态 $\boldsymbol{P}^{*}$。第二阶段（$\lflow$）以恢复的姿态为条件。它将像素对齐的特征（从 DINOv2 和低级特征图中采样）注入生成过程，以合成最终3D物体的高保真几何结构和外观。}
\label{fig:cupid_architecture}
\end{figure*}



\subsection{级联流建模}
\label{sec:flow_model}

我们采用两阶段级联流模型来联合采样3D物体及其对应的相机姿态，记为 $\latent = \{ (\point_i, \feat_i), (\point_i,\pixcoord_i) \}_{i=1}^{\numpts}$。在第一阶段，占用和姿态生成模型（$\sflow$）产生两个关键输出：（i）识别活跃体素的占用立方体，以及（ii）编码以物体为中心的相机姿态的 UV 立方体。在第二阶段，以这些预测为条件，姿态对齐的几何和外观模型（$\lflow$）为每个活跃体素合成 DINO~\citep{oquab2023dinov2} 特征 $\feat_i$，以产生最终的潜在变量 $\latent$。

\paragraph{占用与姿态生成。}
给定条件图像 \(\image\)，此阶段生成分辨率均为 $r$ 的占用立方体和 UV 立方体。占用立方体 \(\boldsymbol{G}_o \in \{0,1\}^{r \times r \times r \times 1}\) 包含区分活跃体素和非活跃体素的二进制值。UV 立方体 \(\boldsymbol{G}_{uv} \in [0,1]^{r \times r \times r \times 2}\) 存储每个体素的归一化像素坐标 $(u,v)$；值得注意的是，共享相同像素坐标的体素位于同一条相机光线上。\footnote{我们没有显式地建模遮挡，因为最近的研究表明 Transformer 可以有效地隐式建模光线传输~\citep{jin2024lvsm, zeng2025renderformer}。}

为了提高计算效率，我们训练了一个3D VAE将 UV 立方体压缩为低分辨率特征网格 \(\boldsymbol{S}_{uv} \in \mathbb{R}^{s\times s\times s \times C}\)，实现了近乎无损的姿态恢复（平均 RRE/RTE $< 0.5^\circ$）。为了微调 TRELLIS，我们将原始特征网格 \(\boldsymbol{S}_o\) 与 \(\boldsymbol{S}_{uv}\) 拼接，并在流网络 $\sflow$ 的输入和输出处引入线性投影层。一旦生成了占用和 UV 立方体，我们通过收集活跃体素并通过公式~\ref{eq:pnp} 求解相机姿态，提取集合 $\{ (\point_i, \pixcoord_i(\pose)) \}_{i=1}^{\numpts}$。

% {\color{blue} Updated: We have included 3D tokenization in section \ref{sec:3d_tokenization}.}
% \yima{How these voxel-type 3D data are tokenized. Discuss
% relationship or difference from the Perceiver for tokenizing point
% clouds in the Michelangelo project. Also, describe basic deep
% architectures used for the two flow-matching based generative models learned. }


\paragraph{姿态对齐的几何与外观生成。}
在第二阶段，我们仅在活跃体素位置生成详细的3D潜在特征 $\{ \feat_i \}_{i=1}^{\numpts}$。实验表明，TRELLIS~\citep{xiang2025structured} 中标准的 $\lflow$ 模型依赖于全局关注的图像信息，容易出现颜色漂移和细粒度细节的丢失。我们通过利用第一阶段得出的相机姿态将局部关注的像素信息注入每个体素来缓解这一问题。

具体来说，我们使用估计的姿态计算第 $i$ 个体素的特征如下：
\begin{equation}
\begin{aligned}
\feat^\dino_{i} &= \interp(\pixcoord_i, \dino(\image)) \in \mathbb{R}^{1024}, \\
\{\feat^{\high}_{i}\}_{i=1}^{\numpts} &= \slatenc\big(\{\point_i, \feat^\dino_i\}_{i=1}^{\numpts}\big), \quad \feat^{\high}_{i} \in \mathbb{R}^{8},
\end{aligned}
\label{eq:latent}
\end{equation}
其中 $\pixcoord_i$ 表示第 $i$ 个3D体素中心在图像平面上的投影，$\interp$ 表示双线性插值，$\slatenc$ 是3D VAE编码器。尽管 DINO~\citep{oquab2023dinov2} 捕获了高级语义，但它通常缺乏精确重建所需的低级线索。为了弥补这一点，我们使用轻量级卷积头从 $\image$ 中提取互补的低级特征 $\feat^{\low}$，并通过 $\interp$ 在 $\pixcoord_i$ 处对其进行采样。最后，在每个时间步 $t$，我们通过线性层将噪声体素特征 $\feat_i^{t}$ 与像素对齐的特征融合，然后再将它们输入到流 Transformer 中：
\begin{equation}
l_t = \mathrm{Linear}\big([\feat_i^{t}\,\oplus\, \feat^{\high}_{i} \,\oplus\, \feat^{\low}_{i}]\big).
\end{equation}
这种姿态对齐的融合策略显著提高了相对于输入图像的3D几何精度和外观保真度。

\begin{remark}[解耦的生成过程]
请注意，上述两阶段框架将3D物体的分布解耦为几个相关的组件：形状、姿态和外观，并分别对它们进行学习。虽然这意味着我们必须一起学习两个生成去噪模型，但由此产生的学习表示结构更好，并能够实现更精确、更灵活的空间推理和生成。
\end{remark}

\subsection{实验}
\label{sec:experiments}
\label{sec:Evaluation}

为了验证该架构，我们在三个关键维度上评估了系统：单目几何预测、输入视图一致性以及完整的单图像到3D重建。基准测试过程侧重于基于学习的系统，排除了逐场景优化方法，以确保对推理能力进行公平比较。我们对姿态条件进行了消融实验。

\paragraph{比较概览。}
我们以目前主导该领域的三个不同范式为基准进行了比较：
\begin{enumerate}[label=(\roman*)]
    \item \textit{点图回归（Point-map Regression）}，例如 VGGT~\citep{wang2025vggt} 和 MoGe~\citep{wang2025moge}，它们同时预测逐像素的3D点和相机姿态，但通常缺乏针对遮挡几何结构的鲁棒先验；
    \item \textit{以视图为中心的重建（View-centric Reconstruction）}，包括 LRM~\citep{hong2023lrm} 和 LaRa~\citep{LaRa}，它们直接在视图空间中生成3D模型，以绕过测试时的姿态估计；以及
    \item \textit{解耦生成（Decoupled Generation）}，以 OnePoseGen~\citep{geng2025one} 为代表，它将规范的3D生成与姿态估计分开，依赖于事后对齐。
\end{enumerate}

\newcommand{\cmark}{\checkmark} % Checkmark
\newcommand{\xmark}{\times} % Cross

\begin{table*}[t]
\renewcommand{\arraystretch}{1.0}
\centering
\caption{
\textbf{Monocular geometry accuracy.} Cupid outperforms all 3D reconstruction and generation baselines and matches point-map regression methods that predict only partial geometry. Note that VGGT uses a ground-truth object mask, which may overestimate accuracy.
}
\label{tab:geometry}
\resizebox{0.99\textwidth}{!}{
\begin{tabular}{l c c c c c c | c c c c c}
\toprule
\multirow{2}{*}{\textbf{Method}} & \multirow{2}{*}{\textbf{3D}} &
\multicolumn{5}{c|}{\textbf{Toys4k}} &
\multicolumn{5}{c}{\textbf{GSO}} \\
& &
\textbf{mIOU} & \textbf{CD} & \textbf{CD} & \textbf{F-score} & \textbf{F-score} &
\textbf{mIOU} & \textbf{CD} & \textbf{CD} & \textbf{F-score} & \textbf{F-score} \\
& &
(avg)$\uparrow$ & (avg)$\downarrow$ & (med)$\downarrow$ & (0.01)$\uparrow$ & (0.05)$\uparrow$ &
(avg)$\uparrow$ & (avg)$\downarrow$ & (med)$\downarrow$ & (0.01)$\uparrow$ & (0.05)$\uparrow$ \\
\midrule
\textcolor{gray}{VGGT} & $\xmark$ &
\textcolor{gray}{--} & \textcolor{gray}{1.144} & \textcolor{gray}{0.498} & \textcolor{gray}{61.85} & \textcolor{gray}{95.90} &
\textcolor{gray}{--} & \textcolor{gray}{1.396} & \textcolor{gray}{0.388} & \textcolor{gray}{65.98} & \textcolor{gray}{95.95} \\
\textcolor{gray}{MoGe} & $\xmark$ &
\textcolor{gray}{92.80} & \textcolor{gray}{1.284} & \textcolor{gray}{0.581} & \textcolor{gray}{58.54} & \textcolor{gray}{95.31} &
\textcolor{gray}{96.18} & \textcolor{gray}{1.743} & \textcolor{gray}{0.575} & \textcolor{gray}{58.99} & \textcolor{gray}{94.68} \\ 
OnePoseGen & $\cmark$ &
9.34 & 153.2 & 59.92 & 6.11 & 24.10 &
12.16 & 116.2 & 60.56 & 7.28 & 25.77 \\
LaRa & $\cmark$ &
68.11 & 32.15 & 16.59 & 18.57 & 57.67 &
70.63 & 34.23 & 19.36 & 13.48 & 49.95 \\
OpenLRM & $\cmark$ &
86.26 & 2.726 & 1.291 & 40.42 & 90.60 &
91.35 & 3.741 & 1.858 & 34.14 & 87.20 \\
Ours & $\cmark$ &
\textbf{92.43} & \textbf{2.534} & \textbf{0.236} & \textbf{69.82} & \textbf{97.76} &
\textbf{95.27} & \textbf{1.823} & \textbf{0.434} & \textbf{61.01} & \textbf{95.59} \\
\bottomrule
\end{tabular}
}
\end{table*}

\begin{table}[t]
\vspace{6pt}
\centering
\caption{\textbf{Comparison on full 3D quality. }We report CLIP image scores of novel views following \citep{gao2024cat3d}.}
\label{tab:clip}

\renewcommand{\arraystretch}{1.0}
\resizebox{1.0\linewidth}{!}{
\begin{tabular}{l|ccccc}
\toprule
 & OnePoseGen & LaRa & OpenLRM & TRELLIS & Cupid \\
\midrule
\textbf{ViT-B/16} & 0.7933 & 0.8334 & 0.8939 & 0.9465 & \textbf{0.9501} \\
\textbf{ViT-L/14} & 0.7193 & 0.7682 & 0.8410 & 0.9210 & \textbf{0.9291} \\
\bottomrule
\end{tabular}
}
\end{table}




\begin{figure}[ht]
    \centering
    \includegraphics[width=1.0\linewidth]{chapters/chapter8/figs/cupid/figures/nvs.pdf}
    \caption{
\textbf{
输入视图一致性的定性比较。
}我们使用生成的相机姿态渲染输入视图。
对于以视图为中心的方法（LRM、LaRa），我们使用真实的内参进行渲染，因为它们不建模内参。 
我们的方法产生了最高保真度的几何结构和外观；LRM 幻觉出了不正确的细节，LaRa 由于2D扩散的不一致性而过于模糊，而3D生成方法 OnePoseGen 经常无法可靠地配准姿态。
}
\label{fig:view_consistency}
\end{figure}
\renewcommand{\arraystretch}{1.0}
\begin{table}[t]
\centering
\caption{
\textbf{Input-view consistency.} Cupid achieves superior input view consistency, producing accurate appearance alignment.
}
\label{tab:appearance}
\centering
\resizebox{\linewidth}{!}{
\begin{tabular}{lc|ccc|ccc}
\toprule
\textbf{Dataset} & &
\multicolumn{3}{c|}{\textbf{Toys4K}} &
\multicolumn{3}{c}{\textbf{GSO}} \\
\textbf{Method} & \textbf{Pose} & \textbf{PSNR}$\uparrow$ & \textbf{SSIM}$\uparrow$ & \textbf{LPIPS}$\downarrow$ & \textbf{PSNR}$\uparrow$ & \textbf{SSIM}$\uparrow$ & \textbf{LPIPS}$\downarrow$ \\
\midrule
LaRa & $\xmark$ & 22.00 & 93.42 & 0.0884 & 19.81 & 91.61 & 0.1119 \\
OpenLRM & $\xmark$ & 26.41 & 80.17 & 0.1156 & 25.79 & 78.80 & 0.1268 \\
OnePoseGen & $\cmark$ & 17.43 & 89.37 & 0.1174 & 14.87 & 86.46 & 0.1386 \\
Cupid & $\cmark$ & \textbf{30.05} & \textbf{96.81} & \textbf{0.0251} & \textbf{28.68} & \textbf{95.49} & \textbf{0.0354} \\
\bottomrule
\end{tabular}
}
\end{table}

\paragraph{几何精度与流程简化。}
在单目几何任务中，所提出的模型表现出明显优于以视图为中心的方法（表~\ref{tab:geometry}）。在 GSO 数据集~\citep{downs2022google} 上，相对于 LRM~\citep{hong2023lrm}，它将平均倒角距离（Chamfer Distance, CD）降低了 50\%。这种改进归功于模型的双重能力：它通过生成先验准确地建模了物体的内在属性，同时预测了精确的相机外参。

关键的是，该模型实现了与专门的点图回归器（如 MoGe~\citep{wang2025moge}）相媲美的几何精度。这一发现对流程设计具有重要意义：它表明，传统上将部分几何估计作为先决步骤的复杂多阶段工作流~\citep{yao2025castcomponentaligned3dscene} 可以得到简化。我们的方法在单个步骤中实现了相当的精度，使得那些中间表示变得不必要。

\paragraph{视觉保真度与一致性。}
我们评估了输入视图一致性，以衡量3D重建与源图像重新对齐的程度（表~\ref{tab:appearance}）。结果突出了基线方法中的特定弱点：LaRa 的纹理质量由于不一致的新视图生成而下降，而 OnePoseGen 在配准过程中遭受了严重的纹理错位和颜色偏移。

相比之下，我们的方法在不牺牲纹理细节的情况下保持了较高的姿态精度（图~\ref{fig:view_consistency}）。它成功地从合成数据集（Toys4K）和具有挑战性的“野外（in-the-wild）”图像（图~\ref{fig:cupid_results}）中重建了高保真物体，在 CLIP 相似度~\citep{radford2021learning} 等语义对齐指标上始终优于基线方法（表~\ref{tab:clip}）。


\paragraph{消融研究}
\label{sec:ablation}

\begin{figure}
    \centering
    \includegraphics[width=1\linewidth]{chapters/chapter8/figs/cupid/figures/ablation.pdf}
    \caption{
\textbf{各种姿态对齐条件的定性比较。}我们的方法（e）在颜色保真度和细节方面实现了最佳的视觉质量。
}\label{fig:abl_appearance}
\end{figure}
% \begin{table}[t]
% \centering
% \renewcommand{\arraystretch}{1.0}
% \begin{minipage}[t]{\linewidth}
% \vspace{0pt} % Ensures baseline alignment
% \captionof{table}{Comparison on CLIP image scores of novel views.}
% \label{tab:clip}
% \centering
% \resizebox{0.55\linewidth}{!}{
% \begin{tabular}{l|cc}
% \toprule
% \textbf{Method} & \textbf{ViT-B/16} & \textbf{ViT-L/14} \\
% \midrule
% OnePoseGen   & 0.7933 & 0.7193 \\
% LaRa         & 0.8334 & 0.7682 \\
% OpenLRM      & 0.8939 & 0.8410 \\
% TRELLIS      & 0.9465 & 0.9210 \\
% Ours         & \textbf{0.9501} & \textbf{0.9291} \\
% \bottomrule
% \end{tabular}
% }
% \end{minipage}

% \hfill

% \renewcommand{\arraystretch}{1.065}

% \begin{minipage}[t]{\linewidth}
% \vspace{0pt} % Ensures baseline alignment
% \captionof{table}{Ablation studies of pose-aligned conditioning.}
% \label{tab:abl_appearance}
% \centering
% \resizebox{\linewidth}{!}{
% \begin{tabular}{l|ccc|ccc}
% \toprule
% \multirow{2}{*}{\textbf{Method}} & \multicolumn{3}{c|}{\textbf{GT Geo \& Pose}} & \multicolumn{3}{c}{\textbf{Sampled Geo \& Pose}} \\
%  & \textbf{PSNR} & \textbf{SSIM} & \textbf{LPIPS} & \textbf{PSNR} & \textbf{SSIM} & \textbf{LPIPS} \\
% \midrule
% (a) Baseline (TRELLIS)   & 31.84 & 97.50 & 0.0219 & 27.47 & 95.64 & 0.0327 \\
% (b) Position Embedding   & 32.07 & 97.58 & 0.0211 & 27.56 & 95.67 & 0.0323 \\
% (c) Latent (w/o Occ.)    & 32.37 & 97.72 & 0.0201 & 27.85 & 95.87 & 0.0309 \\
% (d) Latent (Occ.)        & 32.39 & 97.77 & 0.0199 & 27.74 & 95.80 & 0.0313 \\
% (e) Latent (Visual Feat.)& \textbf{34.86} & \textbf{98.24} & \textbf{0.0168} & \textbf{30.05} & \textbf{96.81} & \textbf{0.0251} \\
% \bottomrule
% \end{tabular}
% }
% \end{minipage}
% \end{table}


\begin{table}[t]
\centering
\caption{\textbf{Ablation studies of pose-aligned conditioning.}}
\label{tab:abl_appearance}
\renewcommand{\arraystretch}{1.065}
\resizebox{1.0\linewidth}{!}{
\begin{tabular}{l|ccc|ccc}
\toprule
\multirow{2}{*}{\textbf{Method}} & \multicolumn{3}{c|}{\textbf{GT Geo \& Pose}} & \multicolumn{3}{c}{\textbf{Sampled Geo \& Pose}} \\
 & \textbf{PSNR} & \textbf{SSIM} & \textbf{LPIPS} & \textbf{PSNR} & \textbf{SSIM} & \textbf{LPIPS} \\
\midrule
(a) Baseline (w/o PAC)   & 31.84 & 97.50 & 0.0219 & 27.47 & 95.64 & 0.0327 \\
(b) Position Embedding   & 32.07 & 97.58 & 0.0211 & 27.56 & 95.67 & 0.0323 \\
(c) Latent (w/o Occ.)    & 32.37 & 97.72 & 0.0201 & 27.85 & 95.87 & 0.0309 \\
(d) Latent (Occ.)        & 32.39 & 97.77 & 0.0199 & 27.74 & 95.80 & 0.0313 \\
(e) Latent (Visual Feat.)& \textbf{34.86} & \textbf{98.24} & \textbf{0.0168} & \textbf{30.05} & \textbf{96.81} & \textbf{0.0251} \\
\bottomrule
\end{tabular}
}
\end{table}

对第二阶段细化的消融研究证实了姿态对齐条件（Pose-Aligned Conditioning, PAC）的关键作用。我们测试了从标准潜在基线到结合 DINOv2 特征体的各种变体。

如表~\ref{tab:abl_appearance} 所示，最有效的配置被证明是一种混合方法：将 DINOv2 特征与通过卷积层提取的低级视觉线索拼接起来。这种组合产生了最高的视觉质量，有效地将语义理解与像素级细节联系起来（图~\ref{fig:abl_appearance}）。此外，这种条件机制表现出显著的鲁棒性，即使输入几何结构包含来自第一生成阶段的随机扰动，也能保持高性能。



\paragraph{应用一：组合式场景重建。}

\begin{figure*}[t]
\centering
\includegraphics[width=\textwidth]{chapters/chapter8/figs/cupid_scene.pdf}
\caption{\small\textbf{组件对齐的场景重建。}每个物体都使用 \textbf{Cupid} 独立重建；显式的3D-2D对应关系能够实现将其精确放置到共享坐标系中。}
\label{fig:cupid_scene_reconstruction}
\end{figure*}


\label{sub:cupid_application}
为了将 Cupid 扩展到场景级重建，我们利用基础模型将场景分解为 $K$ 个带有掩码 $\{ \boldsymbol{M}^{(k)} \}_{k=1}^K$ 的物体，并估计全局度量点图 $\mathcal{P}$~\citep{wang2025moge}。

在第一阶段，我们独立重建每个物体以恢复其规范几何结构，通过随机掩码微调使生成器适应处理相互遮挡~\citep{wu2025amodal3r}。对于第 $k$ 个物体，这会产生一组对应于像素 $\boldsymbol{u}_i \in \boldsymbol{M}^{(k)}$ 的规范3D点 $\boldsymbol{x}_i^{(k)}$。为了解决独立生成中固有的尺度歧义，我们将对齐公式化为生成的规范空间与估计的相机空间之间的对应问题：
\begin{equation}
    \boldsymbol{p}_i = \mathcal{P}(\boldsymbol{u}_i), \quad \forall \boldsymbol{u}_i \in \boldsymbol{M}^{(k)},
\end{equation}
其中 $\boldsymbol{p}_i$ 表示从深度基础模型得出的相机坐标系中的目标度量坐标。

在第二阶段，我们计算每个物体的相似变换 $\mathcal{T}_{k} = \{s_{k}, \boldsymbol{R}_{k}, \boldsymbol{t}_{k}\}$，以将所有组件放置到统一的坐标系中。我们使用 Umeyama 方法~\citep{umeyama2002least} 通过最小化每个物体可见像素上的对齐误差来求解这些参数：
\begin{equation}
    \mathcal{T}_{k}^* = \argmin_{s_k, \boldsymbol{R}_k, \boldsymbol{t}_k} \sum_{i} \left\| \boldsymbol{p}_i - \left(s_k \boldsymbol{R}_k \boldsymbol{x}_i^{(k)} + \boldsymbol{t}_k\right) \right\|^2.
\end{equation}
应用这些最优变换会产生度量一致的场景组合，其中所有生成的组件都相对于彼此正确定位和缩放。请参见图~\ref{fig:cupid_scene} 获取更多场景组合示例。

\begin{figure}
    \centering
    \includegraphics[width=0.95\linewidth]{chapters/chapter8/figs/cupid/figures/scene_result.pdf}
    %
    \caption{\textbf{组件对齐的场景重建的更多示例。}对于显示的每个示例，面板显示：（左上）输入图像，（右上或左下）最终渲染输出，以及（底部）重建的各个组件，为了清晰起见进行了颜色编码。
    }
    \label{fig:cupid_scene}
    \vspace{-5mm}
\end{figure}


\paragraph{应用二：多视图一致性重建。}

\begin{figure}
    \centering
    \includegraphics[width=\textwidth]{chapters/chapter8/figs/cupid/figures/mv.pdf}
    \caption{
    \textbf{多视图条件。}我们解耦的联合建模自然支持多视图条件。在有多个输入视图可用的情况下，我们跨流路径融合共享的\textit{与视图无关的}物体潜在变量（类似于 MultiDiffusion~\citep{bar2023multidiffusion}），从而实现跨所有视图的物体和相机细化。上图：输入；中图：重建的3D物体和相机姿态；下图：渲染的图像和几何结构。}
    \label{fig:mv_images}
\end{figure}

\label{sub:cupid_application_mv}
% 为了实现 \ref{sec:flow_model} 节中描述的多视图功能，我们推广了采样过程以适应 $K$ 个输入视图 $\{ \image^{(k)} \}_{k=1}^K$。 
得益于我们的贝叶斯公式化，即使仅在单张图像上进行训练，我们的模型也能泛化到多视图场景。为了实现 \ref{sec:flow_model} 节中描述的多视图功能，我们扩展了采样过程以适应 $K$ 个输入视图 ${ \image^{(k)} }_{k=1}^K$。

在第一阶段，我们的目标是恢复单一的规范几何结构，同时估计 $K$ 个不同的相机姿态。我们初始化一个共享的几何潜在变量 $\boldsymbol{S}_{o}$ 和 $K$ 个独立的姿态潜在变量 $\{ \boldsymbol{S}_{uv}^{(k)} \}_{k=1}^K$。在每个去噪步骤 $t$，我们为每个视图并行预测几何和姿态的流场。令第 $k$ 个视图的流网络输出表示为几何和姿态更新的拼接向量：
\begin{equation}
    [d_{o, t}^{(k)}, d_{uv, t}^{(k)}] = \sflow(\text{concat}(\boldsymbol{S}_{o}^t, \boldsymbol{S}_{uv}^{(k), t}), t, \image^{(k)}).
\end{equation}
为了确保所有观测结果的几何一致性，我们聚合几何更新，同时保持特定于视图的姿态轨迹。步骤 $t$ 的更新规则变为：
\begin{equation}
\begin{aligned}
    \boldsymbol{S}_{o}^{t-1} &\leftarrow \text{Step}\left(\boldsymbol{S}_{o}^{t}, \frac{1}{K} \sum_{k=1}^K d_{o, t}^{(k)}\right), \\
    \boldsymbol{S}_{uv}^{(k), t-1} &\leftarrow \text{Step}\left(\boldsymbol{S}_{uv}^{(k), t}, d_{uv, t}^{(k)}\right),
\end{aligned}
\end{equation}
其中 $\text{Step}(\cdot)$ 表示流采样器的更新规则。生成的特征网格将被解码为共享的粗略占用网格 $\boldsymbol{G}_o$ 和姿态 $\{\pose^{k}\}_{k=1}^{K}$。

在第二阶段，在固定几何结构的情况下，我们将 Multi-Diffusion 策略~\citep{bar2023multidiffusion} 应用于外观潜在变量 $\boldsymbol{F}_{o}$。我们维护一个单一的共享潜在变量，并使用来自所有 $K$ 个视图的平均流预测对其进行更新，条件是第一阶段估计的粗略几何结构 $\boldsymbol{G}_o$ 和姿态 $\{\pose^{k}\}_{k=1}^{K}$：
\begin{equation}
    d_{\text{feat}, t}^{(k)} = \lflow(\boldsymbol{F}_{o}^t, \boldsymbol{G}_o, \pose^{k}, t, \image^{(k)}),
\end{equation}
\begin{equation}
    \boldsymbol{F}_{o}^{t-1} \leftarrow \text{Step}\left(\boldsymbol{F}_{o}^{t}, \frac{1}{K} \sum_{k=1}^K d_{\text{feat}, t}^{(k)}\right).
\end{equation}
这种平均操作确保生成的3D表示同时满足所有输入视图施加的视觉约束，从而产生如图~\ref{fig:mv_images} 所示的类似 SfM 的一致性。


\subsection{结论}
本节提出了一个用于3D生成和重建的统一框架，将物体身份（“是什么”）和空间上下文（“在哪里”）之间的生物学区别付诸实践。通过联合建模规范的3D物体 $\object$ 和相机姿态 $\pose$，我们的方法有效地将内在形状与外在视角解耦。

利用姿态先验和对齐条件，该模型执行类似贝叶斯的推理，从有限的单视图输入中解析出完整的3D几何结构。这确保了对视角变化的几何一致性和鲁棒性。最终，这种结构化表示提升了具身智能（Embodied AI）的能力，为智能体配备了必要的空间理解能力，使其能够有效地在3D环境中导航和交互。

\section{条件人体运动生成}
\label{sec:egoallo}
% \yima{EgoAllo (Brent, Annie and Chen)，将第6章的内容移到这里，并进一步详细说明方法和实现。}
% \subsection{以头部和手部为条件的身体姿态生成}\label{sub:ego-allo}

%\input{chapters/chapter8/applications_egoallo}
% 这种类型的条件估计或生成问题在许多实际应用中非常自然地出现。
% 这类问题的一个典型例子是如何在给定头部姿态和自我中心图像的条件下估计和生成身体姿态和手势，如图 \ref{fig:pose-teaser} 所示。
% 这通常是当人们佩戴头戴式设备（如苹果的 Vision Pro 或 Meta 的 Project Aria）时我们需要解决的问题。需要推断全身的姿态和手部的姿势，以便我们可以利用这些信息来控制人与之交互的虚拟物体。
% \begin{figure*}[t]
%   \centering
%   \includegraphics[width=\linewidth]{\toplevelprefix/chapters/chapter7/figs/pose-teaser.pdf}
%     \caption{
% 一个系统，在给定自我中心 SLAM 姿态和图像（左）的条件下，估计人体身高、姿态和手部参数（中）。输出在场景的异中心参考系中捕获佩戴者的动作，我们在此用3D重建（右）进行可视化。
%   }
%   \label{fig:pose-teaser}

% \end{figure*}

% 请注意，在这种情况下，人们只有设备提供的头部姿态，以及手部和上肢部分非常有限的视野。身体其余部分的姿态需要基于这些部分信息进行“推断”或“补全”。
% 随时间估计身体姿态的唯一方法是预先学习头部和身体姿态序列的联合分布，然后在实时部分输入的条件下对该先验分布进行采样。图 \ref{fig:pose-method} 概述了一个名为 EgoAllo \cite{yi2024egoallo} 的系统，该系统基于学习到的条件扩散去噪模型来解决这个问题。
% \begin{figure*}[t]
%   \centering
%   \includegraphics[width=\linewidth]{\toplevelprefix/chapters/chapter7/figs/pose-method.pdf}
% \caption{
%     \textbf{EgoAllo \cite{yi2024egoallo} 技术组件概述。}
%     预训练了一个扩散模型，该模型可以基于局部身体参数生成身体姿态序列（中）。
%     使用 SLAM 姿态（左）的不变参数化 $g(\cdot)$ 来作为扩散模型的条件。这些可以通过与输入姿态的全局对齐放置到全局坐标系中。
%     当可用时，自我中心视频通过 HaMeR~\cite{pavlakos2023reconstructing} 用于手部检测（左），这可以通过生成手势的引导合并到样本中。
%   }
%   \label{fig:pose-method}
% \end{figure*}

% 图 \ref{fig:comparison} 比较了一些真实运动序列与 EgoAllo 生成的采样结果。虽然该图为每个输入头部姿态序列显示了一个结果，但不同的运行可以生成与给定头部姿态一致的不同身体姿态序列，所有这些都从自然全身运动序列的分布中提取。
% \begin{figure}[t]
% \centering
% \begin{subfigure}{0.45\textwidth}
%   \centering
%   \includegraphics[width=\linewidth]{\toplevelprefix/chapters/chapter7/figs/qual0_gt.png}
%   \caption{\centering Ground-truth}
% \end{subfigure}
% \hfill
% \begin{subfigure}{0.45\textwidth}
%   \centering
%   \includegraphics[width=\linewidth]{\toplevelprefix/chapters/chapter7/figs/qual0_ours.png}
%   \caption{EgoAllo}
% \end{subfigure}

% \begin{subfigure}{0.45\textwidth}
%   \centering
%   \includegraphics[width=\linewidth]{\toplevelprefix/chapters/chapter7/figs/qual1_gt.png}
%   \caption{\centering Ground-truth}
% \end{subfigure}
% \hfill
% \begin{subfigure}{0.45\textwidth}
%   \centering
%   \includegraphics[width=\linewidth]{\toplevelprefix/chapters/chapter7/figs/qual1_ours.png}
%   \caption{EgoAllo}
% \end{subfigure}

% \caption{
% \textbf{跑步（上）和蹲下（下）序列的自我中心人体运动估计。}
% 将真实运动与 EgoAllo 输出的一个与给定头部姿态序列一致的结果进行比较。
% }
% \label{fig:comparison}
% \end{figure}

% 严格来说，EgoAllo \cite{yi2024egoallo} 中提出的解决方案并没有使用上面介绍的技术来强制执行测量匹配。相反，它通过利用 Transformer 架构中的交叉注意力机制启发式地强制执行条件。正如我们将在 \Cref{sub:text-cond} 中的配对数据设置中更精确地描述的那样，有理由相信交叉注意力机制在某种程度上近似地实现了去噪后验的条件采样。我们相信，这里介绍的更有原则的技术如果得到正确实施，可以带来更好的方法，进一步改善身体姿态和手势估计。

%EgoAllo：自我中心人体运动估计
% 第7章小节 - 学习真实世界数据的表示
%
% 遵循 Cupid（3D物体生成）小节的结构

% \section{自我中心人体运动估计}
% 我们必须知道，我们的大脑感知外部3D世界以发展视觉记忆的主要原因，是为了让我们用身体与物理环境进行协商，并用手与其中的物体进行交互。因此，我们的身体运动和手部运动可以说是我们想要了解的另一个最重要的数据。事实上，我们的大脑投入了很大一部分区域——运动皮层——来记录和协调我们自己的身体运动。在本节中，我们将展示如何使用本书介绍的方法预先学习人体运动的分布，以及如何利用它来促进实时感知和估计我们自身运动的任务。

到目前为止，我们在前面的章节中已经展示了如何学习包括2D图像和3D物体形状在内的视觉数据的分布。请注意，我们感知3D世界并建立其模型的主要目标是为了在其中导航、交互和操作物体。基于感知到的环境信息，我们的大脑会根据手头的任务制定计划并采取行动。通过“手眼协调”所能达到的灵巧程度可以说是自然界中最高形式的物理智能。因此，准确估计和控制我们的身体和手部运动对于我们与环境中的物体进行交互至关重要。我们的大脑为此投入了大量资源：位于大脑额叶的{\em 运动皮层}（motor cortex）记录并协调我们的身体运动，而来自我们身体神经系统的本体感觉信号则提供关于我们四肢位置的持续反馈。因此，如果我们希望机器人（尤其是类人机器人）模仿人类的动作，它们就需要学习我们身体运动（的分布）。

在本节中，我们将展示如何应用第~\ref{ch:conditional-inference} 章中的条件推理框架来学习人体运动的分布，然后将其用于身体姿态估计任务。为此，我们考虑一个简单但自然的设置：一个人佩戴头戴式设备（如 AR 头显或智能眼镜），该设备捕获前向视频并通过运动恢复结构（SFM）或同时定位与建图（SLAM）跟踪其3D姿态。这些方法通过跟踪和三角测量场景中的视觉特征来估计设备的位置，以恢复3D结构和姿态。\footnote{感兴趣的读者可以参考教科书 \cite{MaY2003} 以了解这一领域的工作。} 我们的目标是从这些自我中心的观测中估计佩戴者的全身姿态、身高和手部配置，尽管在捕获的视频中很少能看到身体。关键的观察结果是\emph{头部运动编码了关于身体运动的丰富信息}。当我们走路时，我们的头部会上下摆动；当我们伸手时，它会倾斜；当我们坐下时，它会遵循一条独特的弧线。这些从动作捕捉数据中学习到的相关性形成了一个条件先验 $p(\text{body} \mid \text{head})$。我们在本节中介绍的方法基于一个名为 \textbf{EgoAllo}（代表从自我中心到异中心估计）的开源项目：\\
\centerline{\url{https://egoallo.github.io}},\\
该项目通过引导采样结合了手部检测和地面接触等额外测量。图~\ref{fig:egoallo_teaser} 说明了整体设置和任务。

\begin{figure}[t]
    \centering
    \includegraphics[width=\textwidth]{\toplevelprefix/chapters/chapter8/figs/egoallo/teaser_oct9.pdf}
    \caption{\textbf{自我中心人体运动估计。}给定来自头戴式设备的 SLAM 头部姿态和自我中心图像（左），EgoAllo 估计世界（异中心）参考系中的全身姿态、身高和手部参数（右）。估计的身体在场景中是接地的，脚接触地板，手处于物理上合理的位置。}
    \label{fig:egoallo_teaser}
\end{figure}

% ==============================================================================
\subsection{表示人体}
\label{sub:egoallo_body_models}
% ==============================================================================

在公式化估计问题之前，我们需要一种人体的数学表示。我们身体的骨骼是一棵运动学树：刚性骨骼通过关节连接，以骨盆为根，链条延伸到头部、手部和脚部。每个关节都有旋转自由度：肩膀在三个轴上旋转，肘部主要在一个轴上旋转。有关像人体这样的链接刚体的运动学表示的更正式和详细的介绍，感兴趣的读者可以参考经典教科书 \cite{MurrayR1994}。

\paragraph{SMPL 身体模型。}
SMPL（Skinned Multi-Person Linear model）\citep{loper2015smpl} 为人体的（运动学）姿态和形状提供了一种紧凑的、可微的表示。它将身体配置分为两个组件：
\begin{itemize}
    \item \textbf{姿态参数} $\theta$：指定身体关节活动的关节旋转。SMPL 使用 23 个关节（加上根节点），每个关节由一个3D旋转参数化。这些参数共同决定了一个人是站立、坐着、伸手、走路等。
    \item \textbf{形状参数} $\beta$：一个低维向量，通常为 10 个系数，控制身体比例，如身高、体重、四肢长度和躯干宽度。这些是通过对大量3D身体扫描语料库进行主成分分析得出的。
\end{itemize}

给定姿态 $\theta$ 和形状 $\beta$，SMPL 通过学习到的混合蒙皮（blend skinning）生成3D网格。它还通过正向运动学提供关节位置：从根（骨盆）位置开始，沿着运动学链依次应用每个关节的旋转，我们可以计算出每个关节在3D空间中的最终位置。例如，为了找到手腕的位置，我们组合骨盆、脊柱、肩膀和肘部的旋转，累加由骨骼长度定义的平移。

\paragraph{添加手部：SMPL-H 和 MANO。}
原始的 SMPL 模型将手表示为简单的端点。对于需要详细手部姿势（如抓取、做手势或打字）的应用来说，这是不够的。SMPL-H \citep{romero2017embodied} 基于 MANO 手部模型扩展了 SMPL，增加了铰接的手部，每只手增加 15 个关节（每个手指三个），总共 52 个关节。

手部旋转 $\phi$ 遵循与身体关节相同的参数化。完整的状态包括身体旋转 $\theta$、手部旋转 $\phi$ 和形状参数 $\beta$。

\paragraph{为什么身体模型很重要。}
为什么要使用 SMPL-H，而不是直接回归关节位置或网格顶点？参数化模型提供了关键优势。输出始终是有效的人体，没有解剖学上不可能的配置。该模型是可微的，支持基于梯度的优化。参数化是紧凑的，只需要几百个数字，而不是完整的网格。而且我们获得了语义结构，可以推理“左手腕”而不是顶点索引。

% ==============================================================================
\subsection{任务与设置}
\label{sub:egoallo_objective}
% ==============================================================================

\paragraph{输入。}
输入由来自头戴式设备的两个数据流组成：
\begin{itemize}
    \item \textbf{头部姿态} $\boldsymbol{h}_{1:T} \in \mathsf{SE}(3)^T$\footnote{这里 $\mathsf{SE}(3)$ 表示 $\mathbb{R}^3$ 中的特殊欧几里得群，本质上是由旋转和平移组成的变换群。}：来自视觉惯性 SLAM 的设备在每个时间步在世界坐标系中的位置和方向。这里 $\mathsf{SE}(3)$ 表示3D中的刚体变换群，包括旋转和平移。\footnote{感兴趣的读者可以参考 \cite{MaY2003} 以获取关于刚体运动及其表示和属性的正式介绍。} 这些姿态具有度量尺度，这意味着我们知道以米为单位的实际距离，而不仅仅是相对几何结构，并且它们是重力对齐的，因此我们知道哪个方向是上方。
    \item \textbf{自我中心图像} $I_{1:T}$：前向视频。佩戴者的身体大部分是不可见的，但在伸手或操作物体时，手部会频繁出现。
\end{itemize}

从图像中，现成的手部估计器（如 HaMeR \citep{pavlakos2023hamer}）在手部可见时产生逐帧测量值：
\begin{equation}
    \boldsymbol{y}^{\text{hand},t} = \big(\hat{\phi}^{\text{MANO},t},\,
    \hat{\boldsymbol{q}}^{3D,t},\, \hat{\boldsymbol{q}}^{2D,t}\big),
\end{equation}
包括估计的手部旋转、相机坐标系中的3D关键点和2D关键点。

\paragraph{输出。}
EgoAllo 估计 $T$ 个时间步的运动序列：
\begin{equation}
    \boldsymbol{x}_0 = \{\boldsymbol{x}_0^{1}, \dots, \boldsymbol{x}_0^{T}\},
    \quad\text{其中}\quad
    \boldsymbol{x}_0^{t} = \big(\theta^t,\, \phi^t,\, \beta,\, s\big).
\end{equation}
这里 $\theta^t$ 是身体关节旋转，$\phi^t$ 是手部旋转，$\beta$ 是跨时间共享的形状参数，$s$ 是同样共享的身高参数。下标 $0$ 表示干净的运动，遵循扩散符号，其中 $\boldsymbol{x}_n$ 是第 $n$ 步的噪声版本。

除了姿态之外，EgoAllo 还估计身高，这对于接地（grounding）至关重要。没有它，身体会漂浮在任意的垂直位置；有了它，脚就能接触地板，头部也能与相机对齐。

\paragraph{后验公式化。}
遵循第~\ref{ch:conditional-inference} 章的贝叶斯框架，估计被公式化为从后验中采样。令 $\boldsymbol{y}$ 表示所有观测值，$\boldsymbol{c} = g(\boldsymbol{h}_{1:T})$ 为头部姿态的处理后表示（将在 \ref{sub:egoallo_representation} 节中很快解释）。后验分解为：\footnote{为了使这种分解有效，我们本质上假设给定 $\boldsymbol{x}_0$，$\boldsymbol{y}^{\text{hand}}, \boldsymbol{y}^{\text{phys}}$ 在条件上独立于头部姿态 $\boldsymbol{c}$。}
\begin{equation}
    p(\boldsymbol{x}_0 \mid \boldsymbol{y}) \;\propto\; p_\theta(\boldsymbol{x}_0 \mid
    \boldsymbol{c}) \cdot p(\boldsymbol{y}^{\text{hand}},
    \boldsymbol{y}^{\text{phys}} \mid \boldsymbol{x}_0),
\end{equation}
其中：
\begin{itemize}
    \item $p_\theta(\boldsymbol{x}_0 \mid \boldsymbol{c})$ 是学习到的条件先验，具体来说是一个扩散去噪模型，它在给定头部运动的情况下生成合理的身体运动。
    \item $p(\boldsymbol{y}^{\text{hand}}, \boldsymbol{y}^{\text{phys}} \mid
        \boldsymbol{x}_0)$ 是在给定身体运动的情况下，手部观测和物理约束（如地面接触和无脚部滑动）的似然。
\end{itemize}

似然通过引导能量 $\mathcal{L}_{\text{guide}}$ 表示：
\begin{equation}
    p(\boldsymbol{y}^{\text{hand}}, \boldsymbol{y}^{\text{phys}} \mid
    \boldsymbol{x}_0) \propto \exp\big(-\mathcal{L}_{\text{guide}}(\boldsymbol{x}_0)\big),
\end{equation}
因此后验变为：
\begin{equation}
    \label{eq:egoallo_posterior}
    p(\boldsymbol{x}_0 \mid \boldsymbol{y}) \;\propto\; p_\theta(\boldsymbol{x}_0 \mid
    \boldsymbol{c}) \exp\big(-\mathcal{L}_{\text{guide}}(\boldsymbol{x}_0)\big).
\end{equation}
这是第~\ref{ch:conditional-inference} 章中的形式：学习到的先验通过基于能量的似然与测量约束相结合。

% ==============================================================================
\subsection{设计条件表示}
\label{sub:egoallo_representation}
% ==============================================================================

条件先验 $p_\theta(\boldsymbol{x}_0 \mid \boldsymbol{c})$ 在给定头部运动表示 $\boldsymbol{c}$ 的情况下生成身体运动。这种表示应该是什么样的？正如本书自始至终所强调的，表示的选择从根本上影响学习和泛化。

一种朴素的方法是直接以原始头部姿态 $\boldsymbol{h}_{1:T} \in \mathsf{SE}(3)^T$（即每个时间步的绝对位置和方向）为条件。这种方法效果很差。要理解原因，请考虑该表示应该满足哪些属性。

\paragraph{空间不变性。}
从头部运动到身体运动的映射不应依赖于运动发生在世界上的哪个位置。在厨房里走路应该产生与在客厅里走路相同的身体姿态估计。只有局部运动模式才重要，绝对位置并不重要。

以绝对姿态为条件违反了这一点。如图~\ref{fig:egoallo_conditioning} 所示，相同的身体运动对应于不同的绝对头部轨迹，具体取决于世界位置。以绝对姿态为条件的模型必须分别学习“头部在 $(3,2,1)$ 向前移动”和“头部在 $(10,5,2)$ 向前移动”都对应于走路——这是一种不必要的负担，会损害泛化能力。

\paragraph{时间不变性。}
该映射也不应依赖于运动发生在观测窗口中的哪个时间。如果包含相同的底层运动，处理第 $0$--$100$ 帧的行为应该与处理第 $50$--$150$ 帧的行为相同。这排除了累积量（如从序列开始的总位移）以及绝对时间索引。

人们可能会尝试通过将第一帧放置在原点来规范化到第一个时间步。但这违反了时间不变性：同一运动的不同时间切片，规范化到它们各自的第一帧后，会产生不同的轨迹。对于相同的底层身体运动，模型会看到不同的输入。

\paragraph{解决方案：逐时间步规范化。}
为了同时实现空间和时间不变性，EgoAllo 在\emph{每个}时间步（而不仅仅是在开始时）定义了一个规范坐标系。表示 $\boldsymbol{c} = \{\boldsymbol{c}^{1}, \dots,
\boldsymbol{c}^{T}\}$ 结合了：
\begin{itemize}
    \item 连续帧之间的相对头部运动 $\Delta\boldsymbol{T}_{\text{head}}^{t}$，捕获头部从 $t-1$ 到 $t$ 是如何移动的。
    \item 逐时间步规范化姿态
        $\boldsymbol{T}_{\text{canonical}}^{t}$：头部相对于重力的方向和离地面的高度，去除了地平面位置和偏航角（yaw）。
\end{itemize}

这些量对于运动发生在世界上的哪个位置以及我们提取哪个时间窗口是不变的。具有相同局部运动的两个序列会产生相同的条件标记 $\boldsymbol{c}$，而不管绝对世界位置或时间偏移如何。

\begin{figure}[t]
    \centering
    \includegraphics[width=2in]{\toplevelprefix/chapters/chapter8/figs/egoallo/conditioning_oct9.pdf}
    \caption{\textbf{不变的头部运动条件化。} 条件化表示必须对空间变换（在不同世界位置的相同运动）和时间偏移（在不同绝对时间的相同运动）都具有不变性。每时间步的规范化实现了这两种属性。}
    \label{fig:egoallo_conditioning}
\end{figure}

% ==============================================================================
\subsection{条件先验}
\label{sub:egoallo_architecture}
% ==============================================================================

条件先验 $p_\theta(\boldsymbol{x}_0 \mid \boldsymbol{c})$ 是一个扩散模型，它在给定头部运动条件的情况下生成身体运动序列。其架构遵循第~\ref{sec:image_classification} 节中讨论的基于 Transformer 的设计，并针对序列运动数据进行了调整。

\paragraph{架构。}
该模型接收三个输入：
\begin{itemize}
    \item 扩散步 $n$ 处的含噪运动序列 $\boldsymbol{x}_n$，按时间步进行分词（tokenized）。
    \item 扩散时间步嵌入，指示噪声水平。
    \item 不变的头部条件化 token $\boldsymbol{c}$，每个时间步一个。
\end{itemize}

头部 token 首先由一个包含六个带有自注意力的 Transformer 块的编码器进行处理，以生成上下文表示。这些表示作为解码器（另外六个 Transformer 块）的条件，用于对运动序列进行去噪。条件化通过交叉注意力注入：在每个解码器层中，运动 token 形成查询（queries），对作为键（keys）和值（values）的头部 token 进行注意力计算，从而允许模型动态选择头部运动的哪些方面与预测每个时间步的每个身体关节相关。模型直接输出对干净运动 $\boldsymbol{x}_0$ 的估计，而不是预测噪声，其参数化为头部相对坐标系中的关节旋转。

\paragraph{训练数据。}
该模型在 AMASS \citep{mahmood2019amass} 上进行训练，这是一个大规模的动作捕捉数据集，将多个数据源聚合为统一的 SMPL-H 格式。AMASS 包含各种人类运动——行走、跑步、跳舞、坐姿、伸手、锻炼——但没有头戴式设备的记录。头部轨迹是通过从动作捕捉数据中提取头部关节的位置和方向来模拟的，就好像设备刚性地附着在头部一样。

\paragraph{训练目标。}
训练使用来自 \Cref{ch:denoising} 的去噪分数匹配损失：
\begin{equation}
    \mathcal{L}_{\text{diff}} = \mathbb{E}_{\boldsymbol{x}_0, \boldsymbol{n}} \left[ w_n
    \left\| D_\theta(\boldsymbol{x}_n, \boldsymbol{n}, \boldsymbol{c}) - \boldsymbol{x}_0 \right\|^2 \right],
\end{equation}
其中 $D_\theta$ 是去噪器网络，$w_n$ 是依赖于噪声的权重，$\boldsymbol{x}_n$ 是根据前向扩散过程向干净运动 $\boldsymbol{x}_0$ 添加噪声获得的。模型学习在以头部轨迹为条件的情况下，从含噪运动中预测干净运动。

% ==============================================================================
\subsection{通过引导结合测量}
\label{sub:egoallo_guidance}
% ==============================================================================

条件先验捕获了头部运动和身体运动之间的关系，但没有使用图像观测或物理约束。为了结合这些信息，EgoAllo 定义了一个引导能量 $\mathcal{L}_{\text{guide}}$，用于衡量运动假设与观测结果之间的一致性。

\paragraph{引导能量。}
总引导能量结合了几个项：
\begin{equation}
    \mathcal{L}_{\text{guide}}(\boldsymbol{x}_0) =
    \lambda_{3D}\mathcal{L}_{3D} + \lambda_{2D}\mathcal{L}_{2D} +
    \lambda_{\text{skate}}\mathcal{L}_{\text{skate}} +
    \lambda_{\text{prior}}\mathcal{L}_{\text{prior}},
\end{equation}
其中：
\begin{itemize}
    \item $\mathcal{L}_{3D}$：在世界坐标系中，估计的手部关键点与来自 HaMeR 的 3D 手部检测之间的对齐。
    \item $\mathcal{L}_{2D}$：重投影误差。投影到图像中的估计手部关键点应与 2D 检测相匹配。
    \item $\mathcal{L}_{\text{skate}}$：脚部滑动惩罚。脚在接触地面时不应滑动。
    \item $\mathcal{L}_{\text{prior}}$：对扩散先验预测的弱正则化，防止引导偏离合理运动太远。
\end{itemize}

遵循第~\ref{ch:conditional-inference} 章中基于能量的解释，该引导能量对应于负对数似然：$p(\boldsymbol{y} \mid
\boldsymbol{x}_0) \propto \exp(-\mathcal{L}_{\text{guide}}(\boldsymbol{x}_0))$。低能量意味着与观测结果一致；高能量意味着不一致。

\paragraph{为什么重投影很重要。}
像 HaMeR 这样的单帧手部估计器从单目图像中生成 3D 手部姿态，但存在尺度和深度模糊的问题。靠近相机的小手与远处的大手看起来是一样的。因此，即使手指关节等局部手部姿态是准确的，3D 位置在绝对意义上也是不可靠的。

2D 重投影损失通过在实际观察到手部的图像空间中进行监督来解决这个问题。估计的 3D 手部位置被投影回图像中，并与检测到的 2D 关键点进行比较。这避开了深度模糊，同时仍然提供了关于手部位置的有用信号。

% ==============================================================================
\subsection{条件采样算法}
\label{sub:egoallo_sampling}
% ==============================================================================

第~\ref{ch:conditional-inference} 章中的测量匹配框架告诉我们如何从后验中采样：修改去噪器以包含朝向测量一致性的梯度校正。回想一下，条件去噪器分解为：
\begin{equation}
    \mathbb{E}[\boldsymbol{x}_0 \mid \boldsymbol{x}_n, \boldsymbol{y}]
    =
    \mathbb{E}[\boldsymbol{x}_0 \mid \boldsymbol{x}_n]
    +
    \sigma_n^2 \nabla_{\boldsymbol{x}_n} \log p(\boldsymbol{y} \mid \boldsymbol{x}_n),
\end{equation}
其中第一项是先验去噪器，第二项是测量匹配校正。

在我们的设置中，先验去噪器以头部运动为条件：$D_n^{\text{head}}(\cdot; \boldsymbol{c})$。测量校正使用引导能量。完整的条件去噪器变为：
\begin{equation}
    D_n^{\text{guided}}(\boldsymbol{z})
    =
    D_n^{\text{head}}(\boldsymbol{z};\boldsymbol{c})
    -
    \sigma_n^2
    \nabla_{\boldsymbol{z}}
    \mathcal{L}_{\text{guide}}(\boldsymbol{x}_0(\boldsymbol{z})),
\end{equation}
这里使用了关系 $\nabla \log p(\boldsymbol{y} \mid \boldsymbol{x}) =
-\nabla \mathcal{L}_{\text{guide}}(\boldsymbol{x})$。

\paragraph{交替优化。}
在实践中，EgoAllo 并不是在每一步都通过完整的引导能量计算梯度，而是在两个操作之间交替进行：

\begin{enumerate}
    \item \textbf{先验去噪步骤。} 使用以头部为条件的去噪器
        $D_n^{\text{head}}(\cdot;\boldsymbol{c})$ 应用 DDIM 风格的更新，以获得提议
        $\tilde{\boldsymbol{x}}_{n-1}$。DDIM（去噪扩散隐式模型）提供了一种确定性的采样过程，将含噪样本映射到更干净的样本；与随机采样器不同，每个去噪步骤都是当前状态的确定性函数，这允许以更少的步骤进行更快的采样。

    \item \textbf{测量匹配步骤。} 通过求解近端优化问题来细化提议：
        \begin{equation}
            \boldsymbol{x}_{n-1}
            =
            \operatorname*{arg\,min}_{\boldsymbol{x}}
            \left\{
                \frac{1}{2}\|\boldsymbol{x}-\tilde{\boldsymbol{x}}_{n-1}\|^2
                +
                \gamma_n\,\mathcal{L}_{\text{guide}}(\boldsymbol{x})
            \right\},
        \end{equation}
        其中 $\gamma_n$ 是步长参数。这会找到接近先验预测但与观测结果更一致的运动。
\end{enumerate}

近端优化使用 Levenberg-Marquardt 算法求解，这是非线性最小二乘问题的经典算法。它在梯度下降和高斯-牛顿优化之间进行插值，根据线性化对真实目标的近似程度来调整其步长。这非常适合 SMPL-H 参数估计，因为其目标涉及可微的正向运动学。在每个去噪步骤中，只需几次迭代即可。

\begin{figure}[t]
    \centering
    \includegraphics[width=\linewidth]{\toplevelprefix/chapters/chapter8/figs/egoallo/egoallo-overview.png}
    \caption{\textbf{EgoAllo 框架组件概述。} 扩散模型仅限于局部身体参数。SLAM（头部）姿态的不变参数化 g(·) 用于对扩散模型进行条件化。通过与输入姿态的全局对齐，可以将这些参数放置到全局坐标系中。当可用时，第一人称视频用于手部检测（例如通过 HaMeR
        \citep{pavlakos2023hamer}），这可以通过引导结合到样本中。}
    \label{fig:egoallo_framework}
\end{figure}

\paragraph{全局对齐。}
扩散模型在头部相对坐标系中运行，预测相对于头部的身体姿态，而不是在绝对世界坐标系中。在最后的去噪步骤之后，EgoAllo 通过与观测到的 SLAM 头部姿态进行组合，将身体放置在世界中：
\begin{equation}
    \boldsymbol{T}_{\text{world},\text{root}}^t
    =
    \boldsymbol{T}_{\text{world},\text{head}}^t
    \cdot
    \boldsymbol{T}_{\text{head},\text{root}}^{(\theta^t, \beta)},
\end{equation}
其中 $\boldsymbol{T}_{\text{head},\text{root}}$ 是通过正向运动学从估计的姿态和形状计算得出的。这确保了身体正确接地：双脚踩在地板上，头部与相机对齐。

\paragraph{处理长序列。}
扩散模型在固定长度的序列（32--128 帧）上进行训练。对于更长的记录，EgoAllo 使用 MultiDiffusion 方法 \citep{bar2023multidiffusion}：将序列划分为重叠的窗口并并行去噪。在每个去噪步骤之后，重叠区域的结果在进入下一步之前通过平均进行融合。这种逐步骤的融合确保了窗口在整个扩散过程中保持彼此一致，从而产生时间上连贯的输出，而不会在窗口边界处出现不连续性。

% ==============================================================================
\subsection{实验见解}
\label{sub:egoallo_experiments}
% ==============================================================================

EgoAllo 在三个数据集上进行了评估：AMASS
\citep{mahmood2019amass}（带有合成头部轨迹的动作捕捉）、RICH \citep{huang2022rich}（具有挑战性的人机交互的动作捕捉）和 Aria Digital Twin
\citep{pan2023aria}（来自 Project Aria 眼镜的真实记录，带有真实身体姿态）。对于手部估计，还额外使用了 EgoExo4D
\citep{grauman2023egoexo4d}。

\paragraph{表示很重要。}
表~\ref{tab:egoallo_cond_ablation} 显示了不同条件化表示对身体估计准确性的影响，比较了五种在不变性属性上有所不同的参数化方法：
\begin{itemize}
    \item \textbf{每时间步不变（Per-timestep invariant）}：第~\ref{sub:egoallo_representation} 节中描述的表示，具有空间和时间不变性。
    \item \textbf{绝对 + 相对（Absolute + relative）}：绝对姿态加上局部坐标系相对运动。
    \item \textbf{绝对 + 全局增量（Absolute + global deltas）}：绝对姿态加上世界坐标系位置增量。
    \item \textbf{序列规范化（Sequence canonicalization）}：与第一帧对齐，提供空间不变性但不提供时间不变性。
    \item \textbf{绝对（Absolute）}：原始的 $\mathsf{SE}(3)$ 姿态，没有任何不变性属性。
\end{itemize}

\begin{table}[t]
    \centering
    \caption{\textbf{条件化表示对 AMASS 的影响。}
        头部运动条件化的不同参数化会导致截然不同的估计准确性。满足空间和时间不变性的每时间步不变表示表现最好。MPJPE 是平均每关节位置误差，单位为毫米。}
    \label{tab:egoallo_cond_ablation}
    \small
    \begin{tabular}{lcccc}
        \toprule
        & \multicolumn{2}{c}{\textbf{序列长度 32}} &
        \multicolumn{2}{c}{\textbf{序列长度 128}} \\
        \cmidrule(lr){2-3} \cmidrule(lr){4-5}
        \textbf{条件化} & MPJPE & 影响 & MPJPE & 影响 \\
        \midrule
        每时间步不变 & \textbf{129.8} & --- & \textbf{119.7} & --- \\
        绝对 + 相对 & 133.0 & +2.4\% & 124.5 & +4.0\% \\
        绝对 + 全局增量 & 136.2 & +4.9\% & 127.4 & +6.4\% \\
        序列规范化 & 153.1 & +17.9\% & 134.0 & +11.9\% \\
        绝对 & 159.9 & +23.2\% & 148.3 & +23.9\% \\
        \bottomrule
    \end{tabular}
\end{table}

结果表明，条件化表示对准确性有显著影响，最佳和最差选择之间的 MPJPE 差异高达 24\%。这强化了本书的一个中心主题：表示的选择从根本上影响学习和泛化。

\paragraph{高度估计实现接地。}
表~\ref{tab:egoallo_body_results} 比较了完整估计（包括高度）与不包含高度估计的变体。估计高度将 MPJPE 提高了 6--7\%，并且对于接地至关重要：如果没有它，估计的身体通常会“漂浮”在地面上方，因为系统无法确定正确的垂直位置。接地指标（GND）衡量估计的脚是否接触地板，高度估计显著改善了这一点。

\begin{table}[t]
    \centering
    \caption{\textbf{跨数据集的身体估计结果。} MPJPE
        和 PA-MPJPE 单位为毫米；GND 是脚接触地面的序列比例。高度估计对于正确接地至关重要。}
    \label{tab:egoallo_body_results}
    \small
    \begin{tabular}{llccc}
        \toprule
        \textbf{数据集} & \textbf{变体} &
        \textbf{MPJPE}$\downarrow$ & \textbf{PA-MPJPE}$\downarrow$ &
        \textbf{GND}$\uparrow$ \\
        \midrule
        AMASS & 完整（含高度） & \textbf{119.7} & \textbf{101.1}
        & \textbf{1.00} \\
        & 不含高度 & 128.1 & 110.3 & 0.98 \\
        \midrule
        RICH  & 完整（含高度） & \textbf{176.2} & \textbf{160.1}
        & \textbf{0.96} \\
        & 不含高度 & 185.7 & 169.9 & 0.82 \\
        \midrule
        ADT   & 完整（含高度） & \textbf{155.1} & \textbf{129.3} & 0.94 \\
        & 不含高度 & 163.7 & 140.0 & \textbf{0.96} \\
        \bottomrule
    \end{tabular}
\end{table}

\paragraph{身体上下文改善手部估计。}
一个值得注意的发现是身体和手部估计之间的共生关系。像 HaMeR 这样的单帧单目手部估计器可以生成准确的局部手部姿态（包括手指关节），但由于尺度和深度模糊，无法确定手部的绝对 3D 位置。通过首先估计全身，我们获得了运动学约束：手连接到手腕，手腕连接到前臂、肘部、上臂和躯干。这条运动链为全局手部定位提供了上下文。图 \ref{fig:egoallo_hands} 显示了一些定性结果。

表~\ref{tab:egoallo_hands} 量化了这种对 EgoExo4D 的影响。原始的 HaMeR 估计具有 237.9mm 的世界坐标系误差。添加身体上下文将其减少到 131.5mm，提高了 45\%，即使仅使用单目手部观测也是如此。当可以从立体相机获得额外的 3D 手腕观测时，误差进一步降至 60.1mm。

\begin{table}[t]
    \centering
    \caption{\textbf{EgoExo4D 上的手部估计。} MPJPE 衡量世界坐标系准确性；PA-MPJPE 衡量局部姿态准确性。身体上下文显著改善了世界坐标系定位，同时保留了局部姿态质量。}
    \label{tab:egoallo_hands}
    \small
    \begin{tabular}{lcc}
        \toprule
        \textbf{方法} & \textbf{MPJPE}$\downarrow$ &
        \textbf{PA-MPJPE}$\downarrow$ \\
        \midrule
        HaMeR（单目，无身体） & 237.9 & \textbf{13.0} \\
        \midrule
        带身体（2D 重投影） & 131.5 & 14.7 \\
        带身体（来自立体的 3D 手腕） & \textbf{60.1} & 14.4 \\
        \bottomrule
    \end{tabular}
\end{table}

这种关系是双向的：手部观测也改善了身体估计。当在采样过程中加入手部引导时，AMASS 上的身体 MPJPE 从 119.7mm 提高到 78.8mm，减少了 34\%。额外的测量约束收紧了后验，并在整个运动链中产生了更准确的估计。

\begin{figure}[t]
    \centering
    \includegraphics[width=\textwidth]{\toplevelprefix/chapters/chapter8/figs/egoallo/hands_oct1_v4.png}
    \caption{\textbf{身体上下文改善手部估计。} 蓝色：来自 HaMeR 的单目手部估计，由于尺度/深度模糊，其具有准确的局部姿态但世界坐标系位置不正确。紫色：带有身体上下文的估计，在世界坐标系中正确接地。}
    \label{fig:egoallo_hands}
\end{figure}

\paragraph{定性结果。}
图~\ref{fig:egoallo_qualitative} 显示了来自 Project Aria 眼镜真实记录的估计身体姿态。估计的身体正确接地：脚适当地接触地板，身体高度与环境匹配，并且在包括行走、坐姿和物体操作在内的各种活动中，姿态与观察到的头部运动一致。

\begin{figure}[t]
    \centering
    \includegraphics[width=\textwidth]{\toplevelprefix/chapters/chapter8/figs/egoallo/qual_grounded_outputs_v3.png}
    \caption{\textbf{定性结果。} 来自真实世界第一人称记录的估计身体姿态，在 3D 场景重建中可视化。该方法在各种活动中生成了物理上合理的姿态并具有适当的接地。}
    \label{fig:egoallo_qualitative}
\end{figure}

% ==============================================================================
\subsection{讨论}
\label{sub:egoallo_discussion}
% ==============================================================================

此应用展示了第~\ref{ch:conditional-inference} 章中的条件推理框架如何从图像扩展到结构化、关节化的输出（如人体运动）。

\paragraph{表示设计至关重要。}
贯穿本书的一个反复出现的主题是，表示的选择从根本上影响学习和泛化。这一原则在这里尤为明显。通过分析条件化表示应满足哪些不变性，我们推导出了一种性能大幅优于其他替代方案的参数化方法。空间不变性确保了独立于世界位置；时间不变性确保了独立于时间偏移。经验教训：在训练模型之前，仔细考虑表示应该编码什么结构。

\paragraph{先验加引导模式的泛化。}
采样过程直接应用了测量匹配框架：学习到的先验（这里是以头部为条件的扩散模型）通过引导采样与手部观测和物理等测量约束相结合。这种先验加引导的相同模式在从图像生成到运动估计的各个领域中反复出现。这些组件是模块化的：我们可以在不改变整体框架的情况下换入不同的先验或不同的引导项。

\paragraph{整体估计优于模块化流水线。}
身体和手部估计之间的共生关系表明，对整个运动链的联合推理优于独立估计各个组件。头部运动为身体姿态提供信息；身体姿态为手部位置提供上下文；手部观测细化了身体和手部的估计。这种互惠互利自然地从后验公式中产生，其中所有变量都通过似然和先验耦合在一起。

\paragraph{局限性。}
仍然存在一些局限性。首先，该方法需要具有度量尺度和重力对齐的准确 SLAM；头部姿态估计中的误差会直接传播到身体估计中。其次，泛化受到训练数据多样性的限制，因此 AMASS 中未表示的异常姿态或活动可能估计不佳。第三，只有在第一人称视角中可见时才能细化手部；被遮挡的手完全依赖于身体先验。最后，该方法假设地板是平坦的；楼梯或不平坦的地形将需要额外的建模。

\paragraph{开放问题。}
有前景的方向包括结合场景几何以实现更好的物理接地，其中手应接触表面，身体应避开障碍物。其他方向包括为实时 AR/VR 应用启用实时推理，以及扩展到佩戴者与他人互动的多人场景。该框架足够灵活，可以通过额外的引导项或修改后的先验来适应这些扩展。




\section{自然语言表示与生成}\label{sec:clm_text}
到目前为止，我们已经展示了如何为与视觉感知和身体运动相关的数据分布学习深度结构化表示。在很大程度上，这样学到的结构化表示可以被视为存储在大脑视觉皮层和运动皮层中的记忆。这些记忆加在一起可以被宽泛地称为“世界模型（model of the world）”，或者简称为“世界模型（world model）”——使我们能够感知世界并采取行动。请注意，人类和动物都有能力学习和发展这种关于世界及其身体的记忆，这些记忆记录了从感知中可以预测的内容以及对其在环境中生存至关重要的运动技能。

在物理环境和身体运动的这种世界模型之上，人类发展了（口头和书面）语言，以便我们能够将学到的关于世界的知识传达给他人——我们大脑中的布洛卡区（Broca's area）和韦尼克区（Wernicke's area）分别以语音处理和语言理解而闻名。在很大程度上，语言记录了我们世界模型中值得在我们之间分享的一小部分但很重要的内容。因此，自然语言作为一种非常特殊的数据类型，必然具有编码非常丰富信息的非常特殊的结构。如果机器能够学习自然语言的结构，它们也许能够更好地与人类交流。

因此，在本节中，我们将展示如何将本书中介绍的方法应用于\textit{语言建模}任务。或者更具体地说，如何训练大型语言模型（LLM）来学习自然语言的分布并随后生成它们。该任务及其设置与用于训练 GPT-2 和许多其他语言模型的任务非常相似。

\subsection{数据} \label{sub:clm_text_data}

我们将用于研究 CRATE 在语言任务上性能的数据是 OpenWebText (OWT)~\cite{Gokaslan2019OpenWeb}，这是 OpenAI 用于训练 GPT2 的未发布 WebText 数据集的开源复现。OWT 中的每个样本都是一个 Web 文档，通常来源于高质量的网页、博客、文章或在线讨论，使用格式良好的自然语言编写。OpenWebText 数据集包含约 801 万个不同长度的文档，总计约 41.70GB 的文本。为了进行评估，我们将使用几个数据集，例如
WikiText~\cite{merity2016pointer}\footnote{对于 WikiText2 和
WikiText103~\cite{merity2016pointer}，测试集划分是相同的，因此我们将它们合并为一个称为 WikiText 的单一数据集。}、LAMBADA~\cite{paperno2016lambadadatasetwordprediction}\footnote{为了获得 LAMBADA 数据集上的准确率，我们使用贪婪解码。} 和 PTB~\cite{marcus-etal-1993-building}。与其他数据集相比，PTB 和 OWT 通常更容易。PTB 侧重于较简单的新闻文本，非常适合传统的语言建模，而 OWT 则具有多样性和非正式性，涵盖各种主题，但在语言结构或长程依赖性方面复杂性较低。
WikiText 具有正式的结构和特定领域的内容，需要比 OWT 更复杂的理解，但仍然是可控的。LAMBADA 是最具挑战性的，因为它涉及长程依赖，要求模型掌握更广泛的上下文信息以准确完成句子。

在更正式的层面上，我们的数据 \(\vX\) 将是文本或字符串；我们令 \(\cT\) 为所有字符串的集合。

\subsection{任务与目标} \label{sub:clm_text_objective}

对于因果语言建模预训练，其思想是我们希望\textit{训练模型输出类似人类的文本}。到目前为止，最流行的方法是使用两阶段训练过程：\footnote{现代语言模型训练有几个额外的训练步骤，需要不同的数据分布和算法方法。然而，训练一个仅仅模仿人类写作的模型只需要这里介绍的这几个步骤。}
\begin{itemize}
\item \textit{首先}，我们希望\textit{学习}一种方法，将文档最佳地编码为基本（“构建块”）字符串序列，称为 \textit{token}。这被称为\textit{分词（tokenization）}，我们构建一个\textit{分词器（tokenizer）}。
\item \textit{其次}，我们希望\textit{学习}一种方法，\textit{在给定所有先前 token 的情况下预测下一个 token 的分布}。这被称为\textit{下一个 token 预测（next-token prediction）}，我们构建一个\textit{语言模型}。
\end{itemize}
这个过程实际上可以追溯到马尔可夫（Markov），他首先注意到在给定适当分词的情况下，自然语言可以通过同名的马尔可夫链结构进行建模 \citep{markov2006example}，然后是香农（Shannon），他提出使用字符级分词器（即每个字符是一个 token）和所谓的“\(n\)-gram”（即从训练数据中计算出的显式查找表，用于在给定前 \(n\) 个 token 的情况下获取下一个 token 的分布）来代替语言模型，从而实现这种精确的语言建模设置 \citep{Shannon-1948}。\footnote{最近一项扩展 \(n\)-gram 模型的研究 \citep{liu2024infini} 表明，对于较大的 \(n\)，它们能够相当好地对文本进行建模，但当然，存储这样一个查找表所需的内存是 \(V^{n}\) 数量级的，因此是完全难以处理的。}

\subsubsection{训练分词器}

构建分词器相当于构建一个词表 \(\cV\)，它是一个 token 集合，并具有某个预先指定的大小 \(V\)。有几种方法可以做到这一点。一种流行的算法被称为字节对编码（Byte Pair Encoding, BPE），可以描述如下：
\begin{itemize}
\item 从训练数据中所有唯一字符及其频率的列表开始。确保此类字符少于 \(V\) 个，并将每个字符作为单独的字符串（“token”）连同其频率一起添加到词表中。
\item 直到词表中有 \(V\) 个 token 为止：
    \begin{itemize}
        \item 通过获取两个最频繁的现有 token 并将它们合并来构造一个 token。
        \item 计算该 token 在数据集中的频率。
        \item 将其添加到词表中（连同其频率）。
    \end{itemize}
\item 此时，不再需要频率信息，可以将其丢弃。
\end{itemize}
BPE 的整个过程如 \Cref{fig:BPE} 所示。请注意，此过程是经典信息论压缩过程的修改版，用于\textit{学习字节流数据（如文本）的无损编码}，因此，可以将其解释为寻找数据的最佳无损压缩。请注意，这是可能的，因为（与图像不同）这里的数据从根本上说是离散且无噪声的。
\begin{figure}
\centering
\includegraphics[width=0.45\textwidth]{\toplevelprefix/chapters/chapter8/figs/BPE1.png}\hspace{0.3in}
\includegraphics[width=0.45\textwidth]{\toplevelprefix/chapters/chapter8/figs/BPE2.png}
\caption{\small {\bf 使用 BPE 对文本数据进行分词的过程。}
    （图片来源：\url{https://huggingface.co/learn/nlp-course/chapter6/5}）。（左）我们首先分析给定的文本语料库，并构建一个由单个字符（或在字节级 BPE 的情况下为字节）组成的初始词表。然后，我们计算语料库中相邻字符对的频率。这涉及扫描整个文本并计算每个双字符序列（bigram）出现的频率。（右）在计算相邻字符对的频率后，我们识别语料库中最频繁的字符对。然后将这对字符合并为一个新的子词单元，作为单个 token 添加到词表中。迭代重复此过程，直到达到预定义的词表大小。 }
\label{fig:BPE}
\end{figure}
在构建了这样的词表之后，分词器就可以将文档分解为 token（即对其进行“分词”）。BPE 使用与训练中类似的过程对数据进行分词：
\begin{itemize}
\item 将文档分离为一个由单字符长度 token 组成的长列表。也就是说，如果文档是“Hello”，那么初始列表就是 `H'、`e'、`l'、`l'、`o'。
\item 只要任何两个相邻的 token 可以连接并且它们的连接是另一个 token，我们就这样做，即我们用合并后的 token 替换这对 token。也就是说，如果 `He' 是词表中的一个 token，`H'、`e'、`l'、`l'、`o' 将变成 `He'、`l'、`l'、`o'。
\item 重复上述过程，直到无法进行更多合并。此时，文档被划分为最终的 token 列表（序列）。
\end{itemize}

在分词过程中需要考虑许多基于实用性和效率的因素。例如，如果天真地实现，上面介绍的算法\textit{远非最佳}。我们不会非常详细地讨论这个主题；网上有许多资源可以了解更多信息，例如
\href{https://huggingface.co/learn/nlp-course/en/chapter6/5}{HuggingFace
教程}。

例如，每个 token 都有一个对应的\textit{索引}，这只是它在词表中的索引（词表毕竟只是一个长度为 \(V\) 的列表）。因此，大多数分词器的输出是一个索引列表，比如 \([V]^{*}\) 的一个元素。请记住，它们对应于原始文档的子字符串，如上所示。

一旦学习了分词器，任何语言模型都可以将其用作黑盒。例如，许多模型具有基于 \texttt{tiktoken} 库的相同（基于 OpenAI 的）分词器。在本节的其余部分，我们将对所有内容使用这样一个固定的、预先构建的分词器，从而将每个文本文档 \(\vX \in \cT\) 与其在 \([V]^{*}\) 中的分词版本等同起来。因此，我们不妨认为文本空间 \(\cT\) \textit{等于} token 序列空间 \([V]^{*}\)（并且不会丢失任何本质内容）。

\subsubsection{训练语言模型}

一旦我们将每个文档表示为 token 序列 \(\vX \in [V]^{N} \subseteq [V]^{*} = \cT\)，我们希望执行下一个 token 预测。也就是说，给定一个\textit{上下文} \(\vX_{:n} \in [V]^{n}\)（即文档中的前 \(n\) 个 token \(\vx_{1}, \dots, \vx_{n} \in [V]\)）\footnote{注意与 Python 符号的不一致：这里的符号\textit{包含}索引 \(n\)。}，我们希望预测位置 \(n + 1\) 处的 token \(\vx_{n + 1} \in [V]\)。为此，我们通过 \(\vz_{\theta}(\vX_{:n}) \doteq (f_{\theta}^{\ext} \circ
f_{\theta})(\vX_{:n}) \in \R^{d}\) 计算 \(\vX_{:n}\) 的聚合特征，并使用分类头 \(h_{\theta} \colon \R^{d} \to \Delta_{V}\)（实现为线性层、MLP 或稍微复杂一点的结构）将此特征投影到 \(V\) 维概率单纯形 \(\Delta_{V}\) 中。这个投影 \(\vp_{\theta}(\vX_{:n}) \doteq h_{\theta}(\vz_{\theta}(\vX_{:n}))\) 用作下一个 token 的估计概率分布。然后，使用符号 \(\vone(\vx_{n + 1}) \in \Delta_{V}\) 表示在第 \(\vx_{n + 1}\) 个分量中为 \(1\)，在其他地方为 \(0\)，因果语言建模损失为
\begin{equation}\label{eq:clm_loss}
\min_{\theta}\bc{\cL_{\mathrm{CLM}}(\theta) \doteq
    \Ex_{\vX}\rs{\frac{1}{N - 1}\sum_{n = 1}^{N - 1}\CE(\vone(\vx_{n +
1}), \vp_{\theta}(\vX_{:n}))}}
\end{equation}
请注意这与分类损失（例如图像分类）有多么相似；人们使用交叉熵并尝试将预测的概率向量与真实值对齐。这两个损失之间的主要区别在于，在这个损失中，我们在整个序列上计算损失，其中每个预测都与其他预测相关（不像独立同分布的分类情况）。

优化这种损失通常被称为“预训练（pre-training）”（与“后训练（post-training）”以及最近的“中训练（mid-training）”形成对比，后者是修改下一个 token 预测器以执行有用任务的方法）。

\textit{旁注：} 为什么 \eqref{eq:clm_loss} 的第一项预测 \(\vone(\vx_{2})\)，而没有衡量预测第一个 token 损失的项？这是因为如果我们想预测第一个 token，我们将以\textit{空序列}作为上下文，因此使用与适用于其他 token 性质不同的机制来进行这第一个 token 的预测。所以实际上，这个模型并没有被训练来预测任何文档的\textit{最开始的} token。这之所以没问题，是因为分词器的一个实现细节：通常，在构建分词器之后，我们在其词表中插入一个\textit{特殊} token，称为\textit{字符串（或文档）开始} token，标记为 \texttt{<|bos|>}。\footnote{通常有几个用于不同目的的特殊 token。包含特殊 token 的现有文本会经过特殊处理。} 然后，在处理每个文档时，我们将 \texttt{<|bos|>} token 添加到文档 token 序列的开头，使分词序列的长度增加 \(1\)。因此，上述因果语言建模目标包含一项，该项涉及在仅给定 \texttt{<|bos|>} token 作为上下文的情况下尝试预测文档的第一个 token，因此它在概念上是一个正确的损失。

\subsection{架构：因果 CRATE}

对于架构，我们使用标准的 GPT-2 风格 Transformer，用 CRATE 层替换 Transformer 层。\footnote{与本书和许多其他社区的惯例直接相悖，NLP 社区将这种 GPT-2 风格的 Transformer（涵盖几乎所有当前的 LLM）称为“仅解码器（decoder-only）”Transformer。“仅编码器（encoder-only）”Transformer 具有不同的架构，而“编码器-解码器（encoder-decoder）”Transformer 将“仅编码器”Transformer 与“仅解码器”Transformer 连接起来。尽管事实上“仅解码器”Transformer \textit{也}计算数据的编码！} 为了完整起见，我们在此详细说明该架构。

\paragraph{嵌入。} 我们首先将 token 序列 \(\vX \in
[V]^{N}\) 嵌入到欧几里得空间。这通常是通过使用一个\textit{庞大的}\footnote{这里的“庞大”是指这种结构通常占语言模型总大小的很大一部分。}数组 \(\vE \in \R^{V \times d}\) 将 \([V]\) 中的每个索引与 \(\R^{d}\) 中的一个向量相关联，并直接形成序列 \([\vE_{\vx_{1}}, \dots, \vE_{\vx_{N}}] \in \R^{d \times
N}\) 来完成的。完整的嵌入映射 \(f_{\theta}^{\emb}\) 还应用了位置编码 \(\vE^{\pos} \in \R^{d  \times N_{\max}}\)，其中 \(N_{\max}\) 是可能处理的最大 token 数量，\footnote{现代位置编码方法已经解决了这个问题，并允许（在理论上）无限外推，但这些方法的开发更为复杂，为了简单起见，我们这里只介绍绝对加性位置编码。} 这产生了嵌入映射
\begin{equation}
f_{\theta}^{\emb}(\vX) \doteq [\vE_{\vx_{1}}, \dots, \vE_{\vx_{N}}] +
\vE_{:N}^{\pos}
\end{equation}
参数 \(\vE\) 和 \(\vE^{\pos}\) 是直接可训练的。由于 \(\vE\) 非常大（并且关于它的梯度更新非常稀疏，因为每个样本中只使用了词表的一小部分），因此使用专门的软件来确保内存更新不会太繁重。还要注意，我们没有像其他小节那样使用类 token（class token）；稍后会详细介绍。

\paragraph{主干网络。} 我们使用类似 CRATE 的主干网络处理嵌入，该网络使用因果掩码（causal masking）。为了说明因果掩码的动机，考虑 \eqref{eq:clm_loss} 中定义的因果语言建模损失 \(\cL_{\mathrm{CLM}}\)。最天真的实现将要求我们计算 \(N\) 次前向传播才能进行一次反向传播。显然这是极其低效的，因为 \(N\) 通常可能达到数千。为了有效地扩展使用这种损失的训练，我们施加了一个\textit{因果}约束，即
\begin{equation}\label{eq:causal_backbone_def}
\vZ_{\theta}(\vX_{:n}) = \vZ_{\theta}(\vX)_{:n}
\end{equation}
即，无论 \(n\) 和 \(N\) 的正值如何（只要 \(N \geq n\)），token 特征 \(\vZ_{\theta}(\vX_{:n}) \in \R^{d \times n}\) 的 \(n\) 列应该与 token 特征 \(\vZ_{\theta}(\vX) \in
\R^{d \times N}\) 的前 \(n\) 列相同。实际上，这意味着我们可以对整个序列应用\textit{一次}主干网络并计算 \(\vZ_{\theta}(\vX)\)，然后随着 \(n\) 增长到序列长度 \(N\)，将 \(f_{\theta}^{\ext}\) 应用于每个递增的子集 \(\vZ_{\theta}(\vX_{:n}) = \vZ_{\theta}(\vX)_{:n}\)。然后我们可以使用所有这些来计算损失。

那么，既然我们想要一个用于主干网络的因果架构，我们该如何获得它呢？由于每个 Transformer 层内的 MLP 和层归一化单独影响每个 token，因此对因果关系唯一重要的是注意力块（或在 CRATE 的情况下的 \(\MSSA\)）。为了使 \(\MSSA\) 具有因果性，我们将 \(\mathrm{CausalMSSA}\) 块定义为
\begin{align}
&\operatorname{CausalMSSA}_{\theta}^{\ell}(\vZ) \doteq
\vU_{\out}^{\ell}\mat{\operatorname{CausalSA}([\vU^{1,
        \ell}]^{\top}\vZ, [\vU^{1, \ell}]^{\top}\vZ, [\vU^{1,
    \ell}]^{\top}\vZ) \\ \vdots \\ \operatorname{CausalSA}([\vU^{K,
        \ell}]^{\top}\vZ, [\vU^{K, \ell}]^{\top}\vZ, [\vU^{1,
\ell}]^{\top}\vZ)} + \vb_{\out}^{\ell}\vone_{N}^{\top} \\
&\text{where} \quad  \operatorname{CausalSA}(\vQ, \vK, \vV) \doteq
\vV\softmax\rp{\frac{\operatorname{CausalMask}(\vK^{\top}\vQ)}{\sqrt{p}}}, \\
&\text{where} \quad  \operatorname{CausalMask}(\vM)_{ij} =
\casework{M_{ij}, & \text{if}\ i \geq j, \\ -\infty, & \text{if}\ i < j}.
\end{align}
在这里，从业者说因果掩码\textit{允许未来的 token \(i\) 关注过去的 token \(j\)，反之则不然}。为了理解原因，让我们写出 \(\operatorname{CausalSA}(\vQ, \vK, \vV)\) 第 \(t\) 列的表达式：
\begin{equation}
\operatorname{CausalSA}(\vQ, \vK, \vV)_{t} = \sum_{i =
1}^{t}\vV_{i}\softmax\rp{[\vK_{:t}]^{\top}\vQ_{t}}_{i}
\end{equation}
（这里非冒号下标表示列）。这个第 \(t\) 个 token 的表达式没有使用任何关于索引 \(t\) 之后的 token 的信息。因此 \(\operatorname{CausalSA}\)，进而 \(\operatorname{CausalMSSA}\)，进而整个因果 CRATE 主干网络在 \eqref{eq:causal_backbone_def} 的定义下都是因果的，并且我们解锁了所承诺的巨大效率提升。

\paragraph{特征提取器。} 我们使用一个后处理步骤 \(f_{\theta}^{\ext}\)，它提取\textit{最后一个已知 token 的特征向量}，以便预测下一个 token。理论上，这意味着每个 token \(\vZ_{\theta}(\vX)_{n}\) 都应该包含关于索引 \(n\) 之前或在索引 \(n\) 处的所有 token（即 \(\vx_{1}, \dots, \vx_{n}\)）的丰富信息，因为所有这些信息都应该可用于预测索引 \(n + 1\) 处的下一个 token。在实践中，每个预测任务实际上只需要这些 token 中的一小部分。无论如何，\(f_{\theta}^{\ext}\) 的方程是
\begin{equation}
f_{\theta}^{\ext}(\vZ_{\theta}(\vX_{:n})) \doteq (\vZ_{\theta}(\vX))_{n}
\end{equation}
其中（再次）非冒号下标是列。在这种情况下，如前所述，我们只是直接提取序列中最后一个 token 的特征向量。

\paragraph{特定任务头。} 对于我们的分类头 \(h_{\theta}\)，GPT-2 架构使用一个简单的线性层和一个 softmax 来获得所需的概率向量：
\begin{equation}
h_{\theta}(\vz) \doteq \softmax(\vW^{\out}\vz + \vb^{\out}),
\end{equation}
其中 \(\vW^{\out} \in \R^{V \times d}, \vb^{\out} \in \R^{V}\)。其他一些更现代的架构使用小型 MLP 和层归一化，但思想非常相似。请注意，这个线性层也有很大的内存使用量（因为 \(V\) 非常大），并构成了训练中的瓶颈；人们已经付出了巨大的努力试图规避它。

所有这些架构选择意味着因果训练相对于非因果训练极其高效：
\begin{itemize}
\item 我们只需要通过主干网络进行\textit{一次前向传播}即可计算整个序列的损失。
\item 特征提取基本上是\textit{免费的}。
\item 所有 token 都可以\textit{并行}地通过特定任务头。
\end{itemize}

\subsection{优化策略}

我们使用端到端随机优化来训练我们的语言模型。剩下的一个问题是，在实践中，一个批次中的不同文档具有不同的长度（就每个序列所需的 token 数量而言），但在撰写本书时，主要的深度学习框架大体上只允许“矩形”张量，这不能适应这种行为。为了尝试解决这个问题，我们只需为批次中所有较短的样本插入一个特殊的填充 token \texttt{<|pad|>}，这样我们就可以使用矩形张量对所有内容进行批处理。在每个时间步 \(k\)，我们：
\begin{itemize}
\item 子采样 \(B\) 个不同的分词文档 \(\{\vX_{b}^{(k)}\}_{b = 1}^{B} \subseteq \cT = [V]^{*}\)，每个长度为 \(N_{b}^{(k)}\)。
\item 计算 \(N_{\max}^{(k)} \doteq \max_{b \in [B]}N_{b}^{(k)}\) 并使用特殊的填充 token 将每个 \(\vX_{b}^{(k)}\) 填充到长度 \(N_{\max}^{(k)}\)。
\item 计算特征 \(\vZ_{\theta}(\vX_{b}^{(k)})\)。
\item 计算预测分布 \(\vp_{\theta}(\vX_{b,
        :n}^{(k)}) \doteq (h_{\theta} \circ
    f_{\theta}^{\ext})(\vZ_{\theta}(\vX_{b}^{(k)})_{:n})\)。
\item 形成代理随机损失
    \begin{equation}
        \hat{\cL}_{\mathrm{CLM}}^{(k)}(\theta) \doteq
        \frac{1}{B(N_{\max}^{(k)} - 1)}\sum_{b = 1}^{B}\sum_{n =
        1}^{N_{\max}^{(k)}-1}\CE(\vone(\vx_{b, n + 1}^{(k)}),
    \vp_{\theta}(\vX_{b, :n}^{(k)}))).
\end{equation}
\item 对 \(\theta\) 计算优化算法的一步，给出以下迭代：
\begin{equation}
    \theta^{(k + 1)} \doteq \textsc{OptUpdate}^{(k)}(\theta^{(k)};
    \nabla_{\theta}\hat{\cL}_{\mathrm{CLM}}^{(k)}).
\end{equation}
\end{itemize}

\subsection{评估方法} \label{sub:clm_text_evals}

有几种方法可以评估训练好的 Transformer 语言模型。
\begin{itemize}
\item 在任意文本的保留数据集上，我们可以在其上评估 \(\cL_{\mathrm{CLM}}\)；损失越低越好，因为这意味着模型的采样产生更好的性能。
\item 在多项选择题数据集上，对于每个问题，我们可以将其作为上下文，并检查生成正确答案的估计概率。
\item 我们还可以测试\textit{文本生成}能力。也就是说，在给定上下文的情况下，我们可以从模型关于下一个 token 的概率分布中重复采样。每次采样我们都会生成一个新的 token，我们将其打印出来并添加到上下文中。这允许我们从 LLM 中采样，并随心所欲地判断生成的样本。\footnote{每次都必须在每个 token 上重新运行模型可能会变得极其昂贵。巧妙地存储语言模型的不同内部特征（例如所谓的\textit{\(K\)-\(V\) 缓存}），以及架构的因果性，可以显著降低采样成本。}
\end{itemize}
在本节中，我们专门执行第一种评估。

\subsection{实验设置与结果}

由于我们的因果 CRATE 架构直接建立在 GPT-2 之上，我们将 NanoGPT 仓库 \citep{nanogpt} 给出的 GPT-2 最佳设置与应用于 CRATE 的相同设置进行比较，以进行公平的比较。

\paragraph{模型架构。} 我们使用 GPT-2 分词器，其词表大小为 \(V = 50257\)，包括一个用于 \texttt{<|pad|>} 的特殊 token。\footnote{此设置中不包含 \texttt{<|bos|>} token，尽管它在现代语言模型中非常常见。} 上下文长度为 \(N_{\max} = 1024\)。主干模型遵循 GPT2-Base 架构 \citep{radford2019language}，并进行了适当的修改以具有因果 CRATE 层，我们将其与 GPT2-Small 和 GPT2-Base 进行比较。

\paragraph{数据集与优化。} 为了训练因果 CRATE，我们遵循 NanoGPT 仓库 \citep{nanogpt} 中的实现。具体来说，我们使用 384 的批量大小，并使用 Adam 优化器~\citep{kingma2014adam} 训练 600,000 步。对于 Adam 优化器，我们使用 $(\beta_1, \beta_2)=(0.9, 0.95)$ 和 $0.1$ 的权重衰减。对于学习率调度，我们应用线性预热和余弦衰减，在第 $2,000$ 次迭代时达到峰值 $\eta=6\times 10^{-4}$，最小值为 $6\times 10^{-5}$。随迭代次数变化的训练和验证损失如 \Cref{fig:crate-text-evals} 所示。在使用 $384$ 的批量大小和 $600,000$ 次迭代进行训练后，训练/验证损失收敛在 $3.37$ 左右。相比之下，开源的 GPT-2 实现使用 $512$ 的批量大小和 $600,000$ 步在 OpenWebText 上进行预训练，并收敛到 $2.85$ 的验证损失 \citep{nanogpt}。

\begin{figure}
\centering
\includegraphics[width=0.5\textwidth]{\toplevelprefix/chapters/chapter8/figs/gpt-loss.pdf}
\caption{\bf 在 OpenWebText 数据集上训练的 CRATE-GPT-Base 的损失曲线。}
\label{fig:crate-text-evals}
\end{figure}

\paragraph{实验结果。}

\Cref{tab:gpt-eval} 表明，与具有相似参数数量和相似架构的 GPT-2 模型相比，CRATE 模型在各种数据集上的因果语言建模损失方面取得了合理的性能。

\begin{table}
\def\arraystretch{1.1}
\small
\caption{\small CRATE-GPT2-Base 模型以及 GPT2-Small、GPT2-Base 模型在数据集测试集上评估的零样本交叉熵损失（$\downarrow$ 越低越好）。
}
\centering
\resizebox{1\columnwidth}{!}
{
\begin{tabular}{ccccccc}
    \hline
    & 参数量 & \textbf{OWT} & \textbf{LAMBADA} &
    \textbf{WikiText} & \textbf{PTB} & \textbf{平均} \\
    \hline
    GPT2-Base  & {124M} & 2.85$\downarrow$ & 4.12$\downarrow$ &
    3.89$\downarrow$ & 4.63$\downarrow$ & 3.87$\downarrow$ \\
    {GPT2-Small } &  {64M} & {3.04} & {4.49} & {4.31} & {5.15} & {4.25} \\
    Causal-CRATE-Base & {60M} & 3.37 & 4.91 & 4.61 & 5.53 & 4.61 \\
    \hline
\end{tabular}
}
\label{tab:gpt-eval}
\end{table}

\section{扩展与改进白盒 Transformer }\label{sec:scalable}\label{sec:improvements}

在最后一节中，我们将讨论几种方法，通过这些方法可以扩展 CRATE 类型模型的各个部分，或者使其在某些特殊任务中更高效，同时仍然保持完全可解释的白盒特性。这些发展融合了概念和经验上的见解，可以被视为关于如何在实践中使用白盒理解来改进深度学习模型的案例研究。我们用来评估这些方法的任务将是图像分类和下一个 token 预测，数据将分别是 ImageNet 和 OpenWebText，优化过程将是相同的反向传播，唯一改变的是架构。

\subsection{增加网络宽度：
CRATE-\texorpdfstring{\(\alpha\)}{alpha}}\label{sub:crate_alpha_experiments}

\begin{figure}[t]
\centering
\includegraphics[width=\textwidth]{\toplevelprefix/chapters/chapter8/figs/crate_alpha_backbone.pdf}
\caption{\small\textbf{CRATE-\(\alpha\) 主干网络的一层。}
与 CRATE 的区别在于 \(\ISTA_{\theta}^{\ell}\) 块被 \(\operatorname{ODL}_{\theta}^{\ell}\) 块取代，后者使用过完备字典执行几个 \(\ISTA\) 步骤。}
\label{fig:crate_alpha_backbone}
\end{figure}

CRATE 框架强制执行的一个设计决策是网络中非线性的\textit{宽度}。在常规的 Transformer 中，宽度通常设置为特征维度的 \(4\)、\(8\) 或 \(\frac{11}{3}\) 倍。然而，CRATE 强制要求宽度完全等于特征维度，即字典 \(\vD^{\ell}\) 是方阵，这可能会导致减少
performance. CRATE框架限制我们做出这一选择的根本原因如下：
\begin{itemize}
\item ISTA模块在字典学习中只进行\textit{单}步优化。
\item 通常，任何迭代优化算法的单步操作都无法有效优化目标。那么为什么这种方法有效呢？
\item 当且仅当优化算法具有良好的初始化（即\textit{热启动}（warm starts））时，它们通常会非常快地收敛。ISTA模块具有热启动——它将输入特征作为最终稀疏编码的初始化。
\item 这强制要求输入特征和稀疏编码具有相同的维度。也就是说，ISTA学习了一个完备的稀疏化字典（参见\Cref{ch:classic-models}）。
\end{itemize}
因此，如果我们想使用宽字典，我们需要ISTA执行\textit{过完备}（overcomplete）字典学习。这意味着我们无法拥有相同的热启动（因为我们的稀疏编码的维度大于特征的维度），并且需要更多的迭代才能收敛到稀疏编码。因此，从特征 \(\vZ_{\theta}^{\ell + 1/2}\) 到稀疏编码 \(\vZ_{\theta}^{\ell + 1}\) 的步骤将不再是
\begin{equation}
\vZ_{\theta}^{\ell + 1} = \ISTA_{\theta}^{\ell}(\vZ_{\theta}^{\ell +
1/2} \mid \vZ_{\theta}^{\ell + 1/2})
\end{equation}
其中 \(\ISTA_{\theta}^{\ell}\) 函数定义为（滥用前几节的符号）
\begin{equation}
\ISTA_{\theta}^{\ell}(\vZ \mid \vY) \doteq \ReLU(\vZ - \beta
(\vD^{\ell})^{\top}(\vD^{\ell}\vZ - \vY) + \beta \lambda
\vone_{s}\vone_{n}^{\top})
\end{equation}
而是以下迭代：
\begin{equation}
\vZ_{\theta}^{\ell + 1} = \vA_{\theta}^{\ell, T}; \qquad
\vA_{\theta}^{\ell, t + 1} =
\ISTA_{\theta}^{\ell}(\vA_{\theta}^{\ell, t} \mid \vZ_{\theta}^{\ell
+ 1/2}) \quad \forall 0 \leq t < T; \qquad \vA_{\theta}^{\ell, 0} =
\vzero_{s \times n},
\end{equation}
即在每一层的前向传播中，对LASSO目标运行 \(T \geq
1\) 步近端梯度下降，初始化为
\(\vzero_{s \times n}\)。在这种情况下，字典可以根据需要足够宽，即 \(\vD^{\ell} \in \R^{s \times d}\)，其中 \(s
\geq d\)（在实践中通常取 \(s = 4d\)）。

\begin{table}[b]
\caption{\small \textbf{在COCO val2017~\citep{lin2014microsoft}上通过MaskCut进行目标检测和细粒度分割}。
这里所有模型都使用大小为 \(8\) 的图像块（patch size）而不是 \(16\) 进行训练。
与CRATE和ViT等现有模型相比，
CRATE-\(\alpha\) 模型系列在性能和可扩展性上都有显著提升。}
\centering
\begin{tabular}{@{}lcccccccc@{}}
\toprule
&  & \multicolumn{3}{c}{检测} &  \multicolumn{3}{c}{分割} \\
模型 & AP$_{50} \uparrow $ & AP$_{75} \uparrow $ & AP $\uparrow$ &
AP$_{50} \uparrow$ & AP$_{75} \uparrow $ & AP $\uparrow$ \\
\midrule
\midrule
CRATE-\(\alpha\)-B/8 & 3.5 & 1.1 & 1.5 & 2.2 & 1.0 & 1.1 \\
CRATE-\(\alpha\)-L/8 & \textbf{4.0} & \textbf{1.7} & \textbf{2.0} &
\textbf{2.7} & \textbf{1.1} & \textbf{1.4} \\
\midrule
\color{gray}CRATE-B/8 & \color{gray}2.9 & \color{gray}1.0 &
\color{gray}1.3 & \color{gray}2.2 & \color{gray}0.7 & \color{gray}1.0 \\
\color{gray}ViT-B/8 & \color{gray}0.8 & \color{gray}0.2 &
\color{gray}0.4 & \color{gray}0.7 & \color{gray}0.5 & \color{gray}0.4 \\
\bottomrule
\end{tabular}
\label{tab:crate_alpha_detection_segmentation}
\end{table}

然而，这带来了一个经验上的问题。使用上述配置，如果 \(\vZ^{\ell + 1/2} \in \R^{d \times n}\)，那么
\(\vZ^{\ell + 1} \in \R^{s \times n}\)，它可以具有任意大的特征维度。在实践中，我们希望每一层的特征维度相同。因此，这为设计网络设定了一个实际的三难困境（trichotomy），即我们\textit{不能}同时满足以下\textit{所有}要求：
\begin{enumerate}
\item 每一层的特征维度相同。
\item 字典是宽的，即过完备的。
\item 非线性的输出是输入相对于字典的稀疏编码。
\end{enumerate}
在实践中，出于效率原因，放弃（1）较难处理。放弃（2）会导致通常的CRATE框架。放弃（3）会导致CRATE的宽版本，即CRATE-\(\alpha\)，它具有以下非线性操作，以从 \(\vZ^{\ell + 1/2}\) 得到 \(\vZ^{\ell + 1}\)：
\begin{equation}
\vZ_{\theta}^{\ell + 1} = \vD^{\ell}\vA_{\theta}^{\ell, T}; \qquad
\vA_{\theta}^{\ell, t + 1} =
\ISTA_{\theta}^{\ell}(\vA_{\theta}^{\ell, t} \mid \vZ_{\theta}^{\ell
+ 1/2}); \qquad \vA_{\theta}^{\ell, 0} = \vzero,
\end{equation}
即，它获取通过近端梯度下降得到的稀疏编码，并乘以字典以获得输入的去噪版本。因此，CRATE-\(\alpha\) 的非线性计算的是易于进行稀疏编码的输入的去噪版本，而不是实际的稀疏编码本身。这里从 \(\vZ_{\theta}^{\ell
+ 1/2}\) 到 \(\vZ_{\theta}^{\ell + 1}\) 的映射被称为过完备字典学习（Overcomplete Dictionary Learning, ODL）模块，记为
\(\operatorname{ODL}_{\theta}^{\ell}\)，即，
\begin{equation}
\vZ_{\theta}^{\ell + 1}(\vX) \doteq
\operatorname{ODL}_{\theta}^{\ell}(\vZ_{\theta}^{\ell + 1/2}(\vX)).
\end{equation}

\begin{figure}[t]
\centering
\includegraphics[width=\textwidth]{\toplevelprefix/chapters/chapter8/figs/crate_alpha_semantic_heads.pdf}
\caption{\small\textbf{使用大小为 \(8\) 的图像块的CRATE-\(\alpha\) 的显著性图（Saliency maps）。} 每一行是不同的图像，每一列对应于最后一层中不同的注意力头。我们观察到显著性图与输入图像中的对象强烈对应。}
\label{fig:crate_alpha_saliency_maps}
\end{figure}

\begin{table}
\caption{\small\textbf{语言建模中的验证损失。} 这里所有模型都在OpenWebText的大部分数据上进行了预训练，并在OpenWebText的一个保留子集上测量了验证交叉熵损失。
CRATE-\(\alpha\) 相比CRATE设计显示出显著的改进，尽管与GPT-2等传统Transformer相比仍存在差距。}
\centering
\begin{tabular}{@{}lcccc@{}}
\toprule
模型 & GPT-2-B(ase) & CRATE-B & CRATE-\(\alpha\)-S(mall) &
CRATE-\(\alpha\)-B \\
\midrule
\midrule
参数量 & 124M & 60M & 57M & 120M \\
OWT验证损失 & 2.85 & 3.37 & 3.28 & 3.14 \\
\bottomrule
\end{tabular}
\label{tab:crate_alpha_lm}
\end{table}

CRATE-\(\alpha\) 层如\Cref{fig:crate_alpha_backbone}所示。在实践中，这种对CRATE的修改在更大规模上表现非常出色。例如，当我们在ImageNet-21K上预训练CRATE-\(\alpha\) 模型时，与CRATE相比，分割等无监督任务（见\Cref{fig:crate_alpha_saliency_maps}和\Cref{tab:crate_alpha_detection_segmentation}）的性能通常有显著提升。在使用因果自注意力（causal self-attention）的语言模型训练中也存在类似的趋势（见\Cref{tab:crate_alpha_lm}）。总的来说，这是扩大规模以匹配Transformer等黑盒模型性能的一条有前途的途径。\footnote{请注意，本节中的实验结果使用了一种略有不同的模型架构，这增加了一些微小的经验收益。这些变化是：（1）在ODL模块上增加了一个残差连接，（2）修改 \(\ISTA\) 以使用两个独立的字典，而不是 \(\vD^{\ell}\) 和 \((\vD^{\ell})^{\top}\)。}

\subsection{线性时间复杂度Transformer}\label{sub:tost_experiments}

在实践中，深度学习模型受到空间和时间复杂度瓶颈的困扰，这代表了在给定固定资源的情况下它们无法扩展的问题规模。当处理每个样本本身都是高维且丰富的数据（例如长文本流或视频）时，一个特别有意义的瓶颈是处理长序列数据的\textit{时间复杂度}。为了缓解使用Transformer处理数据的时间复杂度，在\Cref{sub:tost}中，我们提出了一种\textit{词元统计自注意力}（token statistics self-attention）算子 \(\TSSA_{\theta}^{\ell}\)。现在我们围绕它构建一个\textit{词元统计Transformer}，称为ToST，我们可以将其用于长上下文任务。特别是，我们可以使用以下层（如\Cref{fig:tost_backbone}所示）作为CRATE中主干层的直接替换（drop-in replacement）：
\begin{align}
\vZ_{\theta}^{\ell + 1/2}(\vX)
&= \vZ_{\theta}^{\ell}(\vX) + \TSSA_{\theta}^{\ell}(\LN_{\theta}^{1,
\ell}(\vZ_{\theta}^{\ell}(\vX))) \\
\vZ_{\theta}^{\ell + 1}(\vX)
&= \vZ_{\theta}^{\ell + 1/2}(\vX) +
\MLP_{\theta}^{\ell}(\LN_{\theta}^{2, \ell}(\vZ_{\theta}^{\ell + 1/2}(\vX)))
\end{align}
其中 \(\TSSA\) 模块的定义如\Cref{sub:tost}所述。请注意，这与\Cref{sub:contrastive_learning_architecture}中讨论的视觉Transformer架构完全相同，只是 \(\TSSA\) 替换了传统的多头自注意力模块 \(\MHSA\)。无论如何，该层前向传播的计算复杂度在所有问题变量（序列长度、特征维度、头数和头维度）上都是线性的。

\begin{figure}[!htbp]
\centering
\includegraphics[width=\textwidth]{\toplevelprefix/chapters/chapter8/figs/tost_backbone.pdf}
\caption{\small\textbf{ToST主干网络的一层}。词元表示经过层归一化（layer-norms）、词元统计自注意力（TSSA）算子和一个MLP，以形成该层的输出。}
\label{fig:tost_backbone}
\end{figure}

\begin{table}[t]
\centering
\resizebox{\textwidth}{!}{
\begin{tabular}{@{}lccc|cc|cc@{}}
    \toprule
    数据集 & ToST-T(iny) & ToST-S(mall) & ToST-M(edium) &  {
    \color{gray} XCiT-S} &  { \color{gray}XCiT-M} &
    {\color{gray}ViT-S} & {\color{gray}ViT-B(ase)}\\
    \toprule
    参数量 & 5.8M & 22.6M & 68.1M &  { \color{gray}24.9M }& {
    \color{gray}80.2M } & { \color{gray} 22.1M } & { \color{gray}86.6 M } \\
    \midrule
    ImageNet & 67.3 & 77.9 & 80.3 &  { \color{gray} 80.5} & {
    \color{gray} 81.5 } & { \color{gray} 79.8} & { \color{gray} 81.8} \\
    ImageNet ReaL & 72.2  & 84.1 & 85.6 &  { \color{gray} 85.6 } & {
    \color{gray} 85.9} & { \color{gray} 85.6 } & { \color{gray} 86.7 }\\
    \midrule
    CIFAR10 & 95.5 & 96.5 & 97.5 &  { \color{gray} 98.1} & {
    \color{gray} 98.3} & { \color{gray} 98.6} & { \color{gray} 98.8}\\
    CIFAR100 & 78.3 & 82.7 & 84.5 &  { \color{gray} 86.1} & {
    \color{gray} 87.6} & { \color{gray} 88.8} & { \color{gray} 89.3}\\
    Oxford Flowers-102 & 88.6 & 92.8 & 94.2 &  { \color{gray} 93.9} &
    { \color{gray} 94.0} & { \color{gray} 94.0}& { \color{gray} 95.7}\\
    Oxford-IIIT-Pets & 85.6 & 91.1 & 92.8 &  { \color{gray} 92.9} & {
    \color{gray} 94.0} & { \color{gray} 92.8} & { \color{gray} 94.1} \\
    \bottomrule
\end{tabular}%
}
\caption{\small \textbf{ToST的线性探测分类准确率}。当主干网络在ImageNet-1K分类任务上进行预训练时，不同模型大小的ToST在各种数据集上的表现。我们观察到，与XCiT（一种专门用于高效处理长序列的经验设计的类Transformer架构）和ViT相比，ToST保持了相对相似的性能，同时还享有更快的运行时间和白盒设计等优势。}
\label{tab:tost_linear_probing}
\end{table}

\begin{table}[!htbp]
\centering
\begin{tabular}{@{}lcccccc@{}}
\toprule
模型 & 参数量 & OWT & Lambada & Wikitext & PTB & 平均 $\downarrow$
\\ \midrule
GPT-2-Base & 124M & 2.84 & 4.32 & 4.13 & 5.75 & 4.26 \\
ToST-Base & 110M & 3.20 & 4.98 & 4.77 & 6.39 & 4.84 \\
ToST-Medium & 304M & 2.88 & 4.45 & 4.30 & 5.64 & 4.32 \\
ToST-Large & 655M & 2.72 & 4.32 & 3.99 & 5.03 & 4.02 \\ \bottomrule
\end{tabular}%
\caption{\small\textbf{语言建模验证损失}。在各种自然语言数据集（的保留集）上预训练模型后计算得出。我们观察到ToST具有良好的可扩展性，因此ToST-Large在因果语言建模中超越了基线GPT-2-Base，同时在长上下文中享有卓越的效率。}
\label{tab:tost_lm}
\end{table}

\begin{table}
\centering

\begin{tabular}{@{}lccccccc@{}}
\toprule
模型        & ListOps  & 文本     & 检索 & 图像    &
Pathfinder & 平均      \\
\midrule
\midrule
Reformer              & \textbf{37.27} & 56.10             & 53.40
& 38.07             & 68.50               & 50.56             \\
BigBird               & 36.05             & 64.02             & 59.29
& 40.83             & 74.87                & 54.17             \\
LinFormer         & 16.13             & \underline{65.90} & 53.09
& 42.34             & \underline{75.30}               & 50.46             \\
Performer             & 18.01             & 65.40             & 53.82
& 42.77             & \textbf{77.05}                & 51.18             \\
Transformer           & 37.11             & 65.21             &
\underline{79.14}              & \underline{42.94}             &
71.83              & \underline{59.24}            \\
ToST & \underline{37.25}    & \textbf{66.75}    & \textbf{79.46}
& \textbf{46.62}    &    69.41      &     \textbf{59.90}\\

\bottomrule
\end{tabular}%
\caption{\small \textbf{ToST(-B)与针对长上下文优化的顶级Transformer变体的Long-Range Arena (LRA)性能比较。} Long-Range Arena是一系列基准测试，通过固定数据集和评估机制来测试算法和架构的长序列建模能力。与所有已知的Transformer变体（包括XCiT和常规（ViT）Transformer，参见\Cref{tab:tost_linear_probing}）相比，ToST在排行榜上名列前茅。此外，ToST具有最低的时间和空间复杂度推理。（在此表中，特定基准测试的最佳分数以粗体显示，第二好的分数带有下划线。）}
\label{tab:tost_lra_results}
\end{table}

此外，所提出的名为ToST（代表“词元统计Transformer”，Token
Statistics Transformer）的架构在视觉任务（即\Cref{tab:tost_linear_probing}）和语言任务（即\Cref{tab:tost_lm}）上表现良好。对于长序列长度任务（参见\Cref{tab:tost_lra_results}）尤其如此，在这些任务中，它比传统的Transformer和所有其他类Transformer架构性能更好且效率高得多。

\subsection{纯注意力Transformer（Attention-Only Transformers）} \label{sub:aot_experiments}

深度学习模型（特别是类Transformer架构）中需要消除的另一个瓶颈是MLP中大规模矩阵乘法带来的内存瓶颈，其中内部维度远大于特征维度 \(d\)。因此，提出一个有趣且重要的问题：我们\textit{真的}需要Transformer内部的MLP吗？如果没有它，性能能达到多好？为了探索这个问题，我们使用了纯注意力Transformer（attention-only-transformer, AoT）架构（见\Cref{sub:aot}），如\Cref{fig:aot_backbone}所示。也就是说，每一层的形式简单如下：
\begin{equation}
\vZ_{\theta}^{\ell + 1}(\vX) = \vZ_{\theta}^{\ell}(\vX) +
\MSSA_{\theta}^{\ell}(\LN_{\theta}^{\ell}(\vZ_{\theta}^{\ell}(\vX))).
\end{equation}
在我们的实现中，我们还尝试了使用多头自注意力（MHSA）代替MSSA。事实证明，这种架构\textit{也是}可行的，尽管网络的深度需要大得多才能达到与通常的CRATE或Transformer架构相当的性能。 %

\begin{figure}[!htbp]
\centering
\includegraphics[width=0.7\textwidth]{\toplevelprefix/chapters/chapter8/figs/aot_backbone.pdf}
\caption{\small\textbf{AoT主干网络的一层。} 词元表示仅经过层归一化和多头（子空间）自注意力算子以形成该层的输出。请注意，这里没有诸如MLP、ISTA或ODL之类的逐词元非线性操作。}
\label{fig:aot_backbone}
\end{figure}

我们使用提出的AoT架构进行了实验，并展示了其潜力。我们在OpenWebText~\citep{Gokaslan2019OpenWeb}上预训练了不同大小的AoT-MSSA和AoT-MHSA模型，以及GPT-2。我们在\Cref{fig:loss}(a)和(b)中分别绘制了训练损失和验证损失随训练迭代次数的变化。可以观察到，中型和大型的基于AoT的模型实现了与GPT-2基础模型相当的训练和验证损失。此外，与GPT-2基础模型相比，AoT-MHSA模型除了架构中没有MLP层之外，与GPT-2基础模型完全相同。如\Cref{fig:loss}所示，加入MLP层可以加速训练过程。使用上述预训练模型，我们在\Cref{tab:zeroshot}中计算了在不同数据集上未经训练的交叉熵验证损失。可以观察到，具有中等和较大参数规模的AoT模型可以实现与GPT-2基础模型相当的性能。此外，我们发现向AoT添加MLP层并不能提高零样本（zero-shot）性能。这些结果突显了纯注意力模型在保持可解释性的同时取得竞争性结果的潜力。

\begin{table*}[!htbp]
\caption{在几个语言基准数据集和任务上的零样本结果：评估具有MSSA和MHSA算子的不同大小的AoT，并与GPT2模型进行比较。}\vskip 0.1in
\centering
\resizebox{1\columnwidth}{!}
{
\begin{tabular}{l|cccccc}
    \toprule
    模型  & {\bf LAMBADA} & {\bf PTB} & {\bf WikiText} & {\bf
    LAMBADA} & {\bf CBT CN} & {\bf CBT NE} \\
    参数量  & (验证损失) $\downarrow$ &  (验证损失)
    $\downarrow$ &(验证损失) $\downarrow$ & (准确率) $\uparrow$ &(准确率)
    $\uparrow$ &(准确率) $\uparrow$ \\
    \midrule
    AoT-MSSA Base (102M) & 4.70 & 6.03 & 4.65 & 0.25 & 0.80 & 0.74\\
    AoT-MSSA Medium (182M) & 4.47 & 5.08 & 4.22 & 0.29 & 0.84 & 0.77 \\
    AoT-MHSA Base (122M) & 4.42 & 5.52 & 4.19 & 0.38 & 0.86 & 0.82\\
    GPT-2 Base (124M) & 4.32 & 5.75 & 4.13 &  0.40 &  0.87 &  0.84 \\
    \bottomrule
\end{tabular}
}
\label{tab:zeroshot}
\end{table*}
\vspace{-0.05in}

\begin{figure*}[t]
\begin{center}
\includegraphics[width=0.4\linewidth]{\toplevelprefix/chapters/chapter8/figs/training_loss.png}
\hspace{0.4in}
\includegraphics[width =
0.4\linewidth]{\toplevelprefix/chapters/chapter8/figs/val_loss.png}
\vspace{-0.15in}
\caption{\centering \textbf{在语言任务上评估模型。} 我们绘制了在OpenWebText上预训练的AoT和GPT-2模型的训练损失（左）和验证损失（右）。}  \label{fig:loss}
\end{center}
\vspace{-0.15in}
\end{figure*}

\section{总结与注释}

本章中的所有工作都是Transformer架构的下游工作，该架构由 \citet{vaswani2017attention} 引入。Transformer架构在\Cref{sec:contrastive_learning}中进行了正式描述。近年来，在Transformer架构的普及和性能的推动下，一个主要的经验创新是将给定的学习问题表述为序列到序列（sequence-to-sequence）问题，并应用Transformer架构。这使得Transformer架构在（几乎）所有深度学习应用中变得无处不在。因此，对Transformer的直接改进可以传播开来，成为许多问题的解决方案，并产生相当大的影响；同样，我们可以将我们对类Transformer架构的白盒理解应用于许多现代问题。本章涵盖的材料仅仅是已经完成的工作的一个子集；其他工作包括文本数据的掩码补全（即BERT）\citep{devlin2019bert,yu2024white}、语言和视觉模型的（机制）可解释性 \citep{bai2024improving}，以及纠错编码 \citep{zheng2025white}。还有更多的工作要做。

还有更多专门关于扩展神经网络实践的理论，这在实践中是非常可行的，我们至少在这里提及一下。这一系列工作由“张量程序”（Tensor Programs）系列工作 \citep{yang2022tensor} 推广。基本的方案是，我们希望Transformer中的初始梯度更新具有恒定的大小，并且通过仔细推导反向传播方程（\Cref{app:optimization}），我们可以确定实现这一点所需的初始化和学习率（逐层选择）的尺度。在实践中，这种方案极大地提高了大规模训练的稳定性和收敛性；它们还规定了一种仅使用小规模训练来寻找大规模训练的“最优”\footnote{“最优”一词加了引号，因为关于这方面的工作仅仅使用了一些关于初始化时的权重大小、特征大小和梯度大小的期望条件来确定“最优性”，而不是例如收敛时的测试损失。}超参数的方法。这项工作的后续研究试图适应特征几何 \citep{bernstein2024oldoptimizernewnorm}，这可以从本书关于表征学习的工作中获得启发。其他后续研究将这种逐权重的信息整合到优化器本身中，以自动获得这些扩展优势，从而获得了诸如Muon \citep{jordan6muon} 之类的优化器，这些优化器最近被用于非常稳定地训练万亿参数模型 \citep{moonshot2025kimi}。总的来说，这两种深度学习理论方法是正交或互补的。

\section{练习与扩展}

\begin{exercise}
阅读DINO论文 \cite{caron2021emerging}。
\end{exercise}

\begin{exercise}
DINO v2 \cite{oquab2023dinov2} 使用了DINO v1的所有内容，但在数据增强阶段，还会随机掩码每个视图中的图像块。这种增强应该强制具有相似局部信息的图像特征相似。构建一个在编码器中促进这一点的优化问题，并实现它。
\end{exercise}

\begin{exercise}
本练习考虑实现随机优化算法以最小化涉及期望的损失。
\begin{enumerate}[(a)]
\item 提出一个替代 \eqref{eq:simdino_loss_teacherstudent_empirical} 中涉及 \(R_{\eps}\) 的项的方法，用于近似 \eqref{eq:simdino_loss_teacherstudent} 中的协方差正则化项。评估计算你提出的项及其梯度所需的时间复杂度。包括在单个计算节点与多个节点上计算它的分析。
\item 评估计算 \eqref{eq:simdino_loss_teacherstudent_empirical} 中现有项及其梯度所需的时间复杂度。
\end{enumerate}
\end{exercise}

\begin{exercise}
证明 \eqref{eq:linear_probing} 和 \eqref{eq:linear_probing_empirical} 是凸优化问题。
\end{exercise}

\begin{exercise}
\phantom{}
\begin{enumerate}[(a)]
\item 实现CRATE和CRATE-\(\alpha\)模型。
\item 比较它们在CIFAR-10数据集上的性能和效率。
\item 从两个方面比较它们的可解释性：
    \begin{itemize}
        \item 表示 \(\vZ\) 的稀疏度 \(\norm{\vZ}_{0}\)
        \item 注意力图 \(\va_{\theta}^{k, \ell}\)
    \end{itemize}
\end{enumerate}
\end{exercise}

\end{document}